% Generated by build_chapter.py from the GY1976 files. Do not edit by hand; edit the files and rebuild.
\chapter{The academic years 1976--77 and 1977--78}\label{ch:1976}
\textbf{Scope and chronology.} This chapter audits all 23 requested catalogue
records and all 55 supplied scan pages. The mathematical core belongs to
1976--77. One requested record does not: P000753 dates record~437 to
13 February 1978, not the catalogue's 13 February 1977. It is retained as
a chronological appendix, with corrected date and \texttt{1977-437}
occurrence identifiers. The five adjacent pages P000705--P000706,
P000752 and P000757--P000758 are accounted for without assigning their
new content to 1976--77. Plans and year headings remain visible records.
The archive list is a finding aid, not the authority for a conflicting
manuscript date \citep{mccme2026seminararchive}.

\textbf{The year's lines.} The opening talk connects real integral geometry
on \(SL(2,\mathbb R)\) with holomorphic domains and separation of discrete
series. Its strongest exact match is the 1977 Gelfand--Gindikin paper;
the notebook's projector kernel omits a nonconstant trace factor
\citep[Eq.~22 and Proposition~9]{gelfandgindikin1977complex}. Fuchs and
Kirillov revisit the Weil algebra and transgression. Arnold's discussion
of second-order equations is principally classical projective geometry:
point linearization requires the full Lie--Tresse conditions, not a
vanishing coefficient along one solution
\citep[\S\S VII--VIII]{cartan1924projective}.

Superalgebras, the Wess--Zumino multiplet, Berezin integration and Stokes
boundaries continue the preceding year's supergeometry. The distinction
between an ordinary differential form and an integrable superdensity is
essential. The Bernstein--Leites boundary theorem is a precise later
publication, while the Popov--Daikhin Einstein--Yang--Mills construction
had already appeared in 1975
\citep{bernsteinleites1977integral,popovdaikhin1975einstein}.
The latter match exposes a speaker-initial conflict that the notes
alone cannot resolve. Migdal's confinement discussion remains a physical
programme; its class E does not assert a theorem of continuum confinement.

The winter material links Euler products and affine Lie algebras with
integrable equations. Manin's filtered bimodules, spectral curves and
nodal gluing have exact counterparts in his 1978 survey, translated in
1979. Comparing formulas, rather than titles, also identifies an erroneous
flatness sign shared by the notebook and the printed survey
\citep[Ch.~III, Theorem~2.6]{manin1978algebraic}.
Duistermaat's variational index formula had circulated in May 1974;
his spectral discussion leads to the later joint paper with Kolk and
Varadarajan. Gabrielov's hypersimplicial chains and Graev's admissible
line complexes likewise belong to established programmes with earlier
publications, not newly discovered orphan problems.

The spring notes connect Poncelet closure, almost Euclidean sections,
Young diagrams and representation packets. Vershik's limit-shape
announcement is distinguished from the later Gaussian fluctuation theorem:
the latter is a generalized random field with a specified covariance,
not the notebook's unqualified white noise
\citep[\S7]{ivanovolshanski2002kerov}. Tyurin's universal assertion that
varying the complex structure can merge any two irreducible representations
fails already for distinct unitary characters; his own later survey records
the relevant uniqueness principle
\citep[Ch.~II, \S2, Exercise~10]{tyurin1978periods}. The 1978 appendix
records the BGG correspondence and distinguishes it from Beilinson's
related construction \citep{bernsteingelfandgelfand1978bundles,beilinson1978coherent}.

\begin{center}
\small
\begin{tabular}{lrr}
\hline
Historical class & Entire supplied package & 1976--77 core\\
\hline
A: published soon after & 55 & 45\\
B: later development by others & 2 & 2\\
C: partly developed & 0 & 0\\
D: not developed & 0 & 0\\
E: already known & 176 & 163\\
F: not identifiable precisely & 17 & 17\\
\hline
Total & 250 & 227\\
\hline
\end{tabular}
\end{center}

\textbf{Questions and limits.} The problem index contains 15
definite questions, including resolved questions and explicitly linked
earlier occurrences. No item meets the C/D development gate. Unspecified
range spaces, continuous-order matrices, deleted global conditions and
unqualified behaviour questions remain F; they supply no evidence of an
undeveloped theorem. The development section contains complete critical
checks and counterexamples, marked as our verifications, with no claim
of research novelty. Historical fate, literal correctness and depth of
verification are separate assessments throughout.

\textbf{Continuities.} The supplied earlier index provides particularly
useful routes through 1975 records~366 and~370 (Euler products), 374--385
and~391 (spectral curves and inverse constructions), 377 and~386 (gauge
theory), and~389 (supergeometry). Older Morse/Maslov and integral-geometric
threads provide context. The BGG question over projective space is not
merged with the earlier affine Serre--Quillen--Suslin problem. Related
vocabulary is never enough to establish identity of questions.

\GStalk{397}{1976-09-06}{I. M. Gelfand}{P000659}{A physics reading programme}
The opening question is \ru{Что надо знать из физики?}; an uncertain marginal attribution says ``following Sinai?''. The programme lists classical and Hamiltonian mechanics; quantum mechanics and scattering; relativity and field theory; and quantum field theory, including second quantization, Bose and Fermi statistics, functional integrals, perturbation series, Feynman diagrams, renormalization and Yang--Mills fields. Arnold, Landau--Lifshitz, Fermi, Weinberg and Hepp are reading suggestions. Consultation with physicists and a typed version are requested. No theorem, calculation or definite research problem is stated, so the programme is preserved without manufacturing mathematical occurrences. Gindikin begins on the same page; the catalogue's P000659--P000660 window overlaps that next record.


\GStalk{398}{1976-09-06}{S. G. Gindikin, with I. M. Gelfand}{P000659-P000662}{SL(2, R), integral geometry and analytic separation of series}
The group is represented as the determinant-one quadric, and the regular representation is studied through hyperplane and line integrals. Gelfand distinguishes separation into series from decomposition into individual irreducibles. The announced real-transform kernels isolate discrete and continuous parts; two complex domains then realize the positive and negative discrete parts as Hardy boundary values. The lecture ends with an explicit kernel and Arnold's conjecture identifying the domains by their fractional-linear action. The exact subsequent paper is \citet{gelfandgindikin1977complex}, received 6 June 1977. Its kernel has an additional nonconstant factor absent from the notebook. P000661--P000662 unmistakably continue this lecture, despite the catalogue's narrower page window.

\begin{GSitem}{1976-398.1}{E}{The determinant-one quadric of signature (2,2)}
\GSfield{Kind; source}{construction; \GSpage{P000659}}
\GSfield{The notes}{For \(g=\left(\begin{smallmatrix}\alpha&\beta\\\gamma&\delta\end{smallmatrix}\right)\), the notes identify \(SL(2,\mathbb R)\) with \(\alpha\delta-\beta\gamma=1\), then with \(x^2+y^2-z^2-t^2=1\).}
\GSfield{Statement}{The linear substitution \(x=(\alpha+\delta)/2\), \(z=(\alpha-\delta)/2\), \(y=(\beta-\gamma)/2\), \(t=(\beta+\gamma)/2\) is invertible and satisfies \(\det g=x^2+y^2-z^2-t^2\). Thus \(SL(2,\mathbb R)\) is precisely the nonsingular real level set of this quadratic form.}
\GSfield{Reconstruction}{\ours\ The explicit change of coordinates is supplied here; expanding the four squares verifies the identification.}
\GSfield{Fate}{This is the classical quadric model used again in \citet[pp.417--420]{gindikin2000integral}; it is background, not a 1976 novelty.}
\GSfield{Present status}{As of 10 November 2026: The identification is exact and elementary; no open question is attached to it.}
\GSfield{Assessment}{literal statement: correct; publication match: exact; verification: read the result; development: none; confidence: high}
\end{GSitem}
\begin{GSitem}{1976-398.2}{E}{Which linear symmetries come from left and right translations?}
\GSfield{Kind; source}{claim; \GSpage{P000659}}
\GSfield{The notes}{The formula is \(g\mapsto g_1^{-1}gg_2\). The next sentence says that \emph{any} linear transformation preserving the hyperboloid has this form.}
\GSfield{Statement}{For \(a,b\in SL(2,\mathbb R)\), \(X\mapsto aXb^{-1}\) gives a homomorphism onto \(SO_0(2,2)\), with kernel \(\{(I,I),(-I,-I)\}\). It does not exhaust \(O(2,2)\). For example \(X\mapsto X^{\mathsf T}\) preserves \(\det X\) but has determinant \(-1\) as a linear map on the four-dimensional matrix space.}
\GSfield{Reconstruction}{\ours\ The transpose counterexample checks the literal quantifier. The notebook's inverse placement defines a right action on points; pullback \(T_{g_1,g_2}f(X)=f(g_1^{-1}Xg_2)\) is a representation on functions. These conventions should not be conflated.}
\GSfield{Fate}{The left--right symmetry is standard in the real-quadric construction, as in \citet[pp.417--419]{gindikin2000integral}. The omitted identity-component restriction is a notebook error, not later progress.}
\GSfield{Present status}{As of 10 November 2026: The corrected identity-component statement holds. The unrestricted statement remains false.}
\GSfield{Assessment}{literal statement: false as written; publication match: partial; verification: read the result; development: none; confidence: high}
\end{GSitem}
\begin{GSitem}{1976-398.3}{E}{Stabilizers of points of the quadric}
\GSfield{Kind; source}{problem; \GSpage{P000659-P000660}}
\GSfield{The notes}{The stabilizer of the identity is the diagonal subgroup, followed by the exercise ``compute the stabilizer subgroups of other points''. }
\GSfield{Statement}{For \(x\in SL(2,\mathbb R)\), the pairs satisfying \(g_1^{-1}xg_2=x\) are exactly \(\{(xhx^{-1},h):h\in SL(2,\mathbb R)\}\). Hence every stabilizer is conjugate to the diagonal and \(G\simeq(G\times G)/\operatorname{diag}G\), with the action convention understood.}
\GSfield{Reconstruction}{\ours\ Multiply the fixed-point equation by \(g_1\) on the left and by \(x^{-1}\) on the right. This solves the exercise completely.}
\GSfield{Fate}{This is the elementary homogeneous-space structure underlying \citet[opening construction]{gindikin2000integral}; no independent research problem is created for the exercise.}
\GSfield{Present status}{As of 10 November 2026: Completely resolved by the displayed algebraic calculation.}
\GSfield{Assessment}{literal statement: correct; publication match: exact; verification: read the result; development: none; confidence: high}
\end{GSitem}
\begin{GSitem}{1976-398.4}{E}{Haar measure and the Hilbert space of the regular representation}
\GSfield{Kind; source}{construction; \GSpage{P000660}}
\GSfield{The notes}{On \(\beta\ne0\), coordinates are \((\alpha,\delta,\beta)\), and \(dg=d\alpha\,d\beta\,d\delta/|\beta|\). The space is \(\mathcal H=L^2(G,dg)\).}
\GSfield{Statement}{The measure \(\delta(\alpha\delta-\beta\gamma-1)\,d\alpha\,d\beta\,d\gamma\,d\delta\) is invariant under both determinant-one left and right multiplications. Integrating out \(\gamma=(\alpha\delta-1)/\beta\) gives the stated local Haar density, up to an overall positive constant. The left--right pullbacks are unitary on \(L^2(G)\).}
\GSfield{Reconstruction}{\ours\ The delta-function calculation is the coarea identity on the chart \(\beta\ne0\); it does not claim that this single chart covers all of \(G\).}
\GSfield{Fate}{The same invariant quadric measure is used in \citet[real geometry, pp.421--423]{gindikin2000integral}.}
\GSfield{Present status}{As of 10 November 2026: Standard Haar measure; the chart restriction and normalization are essential but routine.}
\GSfield{Assessment}{literal statement: correct; publication match: exact; verification: read the result; development: none; confidence: high}
\end{GSitem}
\begin{GSitem}{1976-398.5}{E}{Separating series before decomposing into irreducibles}
\GSfield{Kind; source}{problem; \GSpage{P000660}}
\GSfield{The notes}{The word ``irreducible'' is crossed out. Gelfand calls division into series algebraic and division of a series into representations transcendental; the example is explicitly said to be well known.}
\GSfield{Statement}{For the left--right regular representation of \(SL(2,\mathbb R)\), seek invariant orthogonal projectors onto the closed sum \(L^2_{\rm disc}(G)\) of square-integrable discrete-series matrix-coefficient spaces and its continuous-spectrum orthogonal complement, and then onto the holomorphic and antiholomorphic discrete parts. The first question is coarser than specifying the Plancherel measure and every irreducible constituent.}
\GSfield{Reconstruction}{The words ``algebraic'' and ``transcendental'' are a research preference, not a proved dichotomy. We formulate the definite projector problem using closed Hilbert subspaces.}
\GSfield{Fate}{The spectral decomposition predates the talk: \citet[entry R18(a)]{strichartz1973hyperboloids} records the 1973 hyperboloid analysis. The specifically complex-analytic implementation is the subsequent \citet{gelfandgindikin1977complex}; its occurrences below receive class A.}
\GSfield{Present status}{As of 10 November 2026: The separation problem for this group is solved. Different analytic realizations must be distinguished from the older existence of the Plancherel decomposition.}
\GSfield{Assessment}{literal statement: correct with missing hypotheses; publication match: exact; verification: indirect; development: none; confidence: high; problem \hyperref[prob:GP-1976-01]{GP-1976-01}}
\end{GSitem}
\begin{GSitem}{1976-398.6}{E}{Classification of real hyperplane sections}
\GSfield{Kind; source}{construction; \GSpage{P000660-P000661}}
\GSfield{The notes}{For \(\xi\alpha+\eta\delta-p\beta-q\gamma=\zeta\), the notes classify sections using \(\Delta=\xi\eta-pq\). The degenerate case mixes ``hyperbolic paraboloid'', ``hyperbolic cylinder'' and a two-parameter label; P000661 gives paraboloids three parameters.}
\GSfield{Statement}{Assume the normal covector is nonzero and identify proportional hyperplane equations. If \(\Delta<0\), the section is a one-sheeted hyperboloid. If \(\Delta>0\), it is one-sheeted, two-sheeted or a cone according as \(|\zeta|\) is greater than, less than or equal to \(2\sqrt\Delta\). If \(\Delta=0\), \(\zeta\ne0\) gives a hyperbolic paraboloid (three parameters), whereas \(\zeta=0\) gives a hyperbolic cylinder (two parameters). General nondegenerate sections have four parameters; tangent cones have three.}
\GSfield{Reconstruction}{\ours\ For the null normal, reduce to \(\alpha=\zeta\). When \(\zeta\ne0\), \(\delta=(1+\beta\gamma)/\zeta\); when \(\zeta=0\), \(\beta\gamma=-1\) and \(\delta\) is free. This independently resolves the conflated degenerate cases.}
\GSfield{Fate}{The section and horosphere geometry is treated in \citet[pp.418--423]{gindikin2000integral}; the underlying quadratic classification is elementary and predates the lecture.}
\GSfield{Present status}{As of 10 November 2026: Resolved with the null-level distinction. The two- and three-parameter families are different, not conflicting counts for one family.}
\GSfield{Assessment}{literal statement: false as written; publication match: partial; verification: read the result; development: none; confidence: high}
\end{GSitem}
\begin{GSitem}{1976-398.7}{E}{The hyperplane Radon transform on the quadric}
\GSfield{Kind; source}{construction; \GSpage{P000661}}
\GSfield{The notes}{The transform is \(\widehat\psi=\int_G\psi(g)\delta(\xi\alpha+\eta\delta-p\beta-q\gamma-\zeta)\,dg\), with \(\psi\in L^2(G)\) written without a domain qualification. The families \(A,B,D,E\) are then named.}
\GSfield{Statement}{For \(\psi\in C_c^\infty(G)\), the displayed expression defines a distribution in the hyperplane parameters by pushforward; at transverse sections it is the usual surface integral with the coarea density. It is homogeneous of degree \(-1\) under simultaneous nonzero real scaling of the five parameters, with absolute value in the scaling factor. For general \(L^2\) input, pointwise surface integration requires a specified extension or regularization.}
\GSfield{Reconstruction}{The test-function domain and distributional interpretation are supplied to avoid assigning a trace on a measure-zero surface to an arbitrary \(L^2\) equivalence class.}
\GSfield{Fate}{Compare the precise transform domains in \citet[complex and real geometry]{gindikin2000integral}. The definition is classical; its unrestricted \(L^2\) pointwise reading is not.}
\GSfield{Present status}{As of 10 November 2026: Well defined on test functions and in appropriate distributional extensions. The notebook does not specify one extension on all of \(L^2(G)\).}
\GSfield{Assessment}{literal statement: correct with missing hypotheses; publication match: partial; verification: read the result; development: none; confidence: high}
\end{GSitem}
\begin{GSitem}{1976-398.8}{F}{The unspecified range of the four- and three-parameter transforms}
\GSfield{Kind; source}{question; \GSpage{P000661}}
\GSfield{The notes}{The question is \ru{Как описать, какие ф-и 4-х и 3-х перем. появляются?} No support, decay, regularity or boundary conditions are attached.}
\GSfield{Statement}{A range problem would have to specify a source space \(X(G)\), one of the section families \(\mathcal F\), its density normalization, and a target space \(Y(\mathcal F)\), then characterize \(R_{\mathcal F}(X)\). These choices are missing, so there is no uniquely determined proposition here.}
\GSfield{Reconstruction}{We do not replace this question by a familiar Euclidean range theorem, or identify inversion with surjectivity onto a characterized range.}
\GSfield{Fate}{Targeted searches located \citet[main inversion construction]{gindikin2000integral}, but that does not identify the omitted function-space quantifiers. The 1977 Hardy theorem concerns a different, specified boundary-value problem.}
\GSfield{Present status}{As of 10 November 2026: Unidentified at the exact range-theorem level. No assertion that it remains open, and no new problem-index entry, is justified.}
\GSfield{Assessment}{literal statement: not determinable; publication match: thematic only; verification: unresolved lead; development: none; confidence: low}
\end{GSitem}
\begin{GSitem}{1976-398.9}{E}{Orthogonal discrete and continuous parts}
\GSfield{Kind; source}{theorem; \GSpage{P000661}}
\GSfield{The notes}{``Main theorem 1'' begins \(\mathcal H=\mathcal H_A\overset\circ+\mathcal H_B\); the spaces were introduced as kernels of the two kinds of hyperplane integration.}
\GSfield{Statement}{The unambiguous spectral statement is \(L^2(G)=L^2_{\rm disc}(G)\oplus L^2_{\rm cont}(G)\), an orthogonal decomposition invariant under both translations, where the first summand is the closed span of discrete-series matrix coefficients and the second its orthogonal complement. The identification of these summands with the notebook's integral kernels additionally requires precise domains for those transforms.}
\GSfield{Reconstruction}{The spectral statement is separated from the more delicate kernel realization in the next occurrence. The circle over the plus sign is read as an orthogonal direct sum, consistently with the later projector discussion.}
\GSfield{Fate}{The existence of the spectral decomposition precedes 1976; see the 1973 paper recorded in \citet[entry R18(a)]{strichartz1973hyperboloids}. The exact kernel formulations were not checked in that inaccessible article.}
\GSfield{Present status}{As of 10 November 2026: The spectral decomposition is established. Verification of this occurrence is indirect for the prior publication, not a claim to have read Strichartz's theorem.}
\GSfield{Assessment}{literal statement: correct with missing hypotheses; publication match: partial; verification: indirect; development: none; confidence: medium}
\end{GSitem}
\begin{GSitem}{1976-398.10}{E}{Common kernels of the section and ruling transforms}
\GSfield{Kind; source}{claim; \GSpage{P000661}}
\GSfield{The notes}{The notes assert \(\mathcal H_A=\mathcal H_D=\mathcal H_E=\mathcal H_C\), and label this the discrete series. Here \(C\) is the family of straight rulings.}
\GSfield{Statement}{A standard precise part is the cuspidal characterization: take smooth Harish--Chandra Schwartz functions on \(SL(2,\mathbb R)\) whose integrals \(\int_{\mathbb R}f(g_1^{-1}\left(\begin{smallmatrix}1&t\\0&1\end{smallmatrix}\right)g_2)\,dt\) vanish for all \(g_1,g_2\). Their closure in \(L^2(G)\) is \(L^2_{\rm disc}(G)\). The additional equality with every hyperplane-kernel realization is a distinct announced claim; the notebook omits the realizations needed to state it on all \(L^2\).}
\GSfield{Reconstruction}{The modern cuspidal interpretation supplies a well-defined part of the statement. We do not silently close the kernel on compactly supported functions, which need not recover the discrete spectrum.}
\GSfield{Fate}{The historical identification is consistent with the pre-talk harmonic analysis and \citet[real horospherical transform]{gindikin2000integral}. An exact source proving all four kernel equalities with the notebook's unspecified conventions was not verified.}
\GSfield{Present status}{As of 10 November 2026: The discrete/cuspidal identification is standard. The full notebook formulation receives only a partial publication match and medium confidence, not an invented unresolved-conjecture status.}
\GSfield{Assessment}{literal statement: correct with missing hypotheses; publication match: partial; verification: indirect; development: none; confidence: medium}
\end{GSitem}
\begin{GSitem}{1976-398.11}{E}{An ultrahyperbolic equation for the transform}
\GSfield{Kind; source}{claim; \GSpage{P000661}}
\GSfield{The notes}{The transformed function is said to satisfy a second-order ultrahyperbolic equation ``that can be written down'', but no operator is displayed.}
\GSfield{Statement}{For the test-function transform above, in the sense of distributions, \[ (\partial_\xi\partial_\eta-\partial_p\partial_q-\partial_\zeta^2)\widehat\psi=0. \] Indeed differentiation under the integral inserts \(\alpha\delta-\beta\gamma-1=0\). Additional range conditions do not follow merely from this equation.}
\GSfield{Reconstruction}{\ours\ The explicit PDE is reconstructed directly from the determinant equation; it is not a formula actually written in the notebook. The two minus signs in the \(p,q\) coefficients cancel in the mixed derivative.}
\GSfield{Fate}{This is the classical differential constraint of a Radon transform supported on a quadric; compare the transform setting in \citet[pp.418--423]{gindikin2000integral}. The derivation here checks only the asserted necessary PDE.}
\GSfield{Present status}{As of 10 November 2026: The reconstructed necessary condition is proved directly. It is not a sufficient range characterization.}
\GSfield{Assessment}{literal statement: correct; publication match: partial; verification: read the result; development: none; confidence: high}
\end{GSitem}
\begin{GSitem}{1976-398.12}{F}{The unexplained minimality of the families A, D and E}
\GSfield{Kind; source}{claim; \GSpage{P000661}}
\GSfield{The notes}{A separate sentence says \ru{Сем-ва A, D, E явл. минимальными}.}
\GSfield{Statement}{No order on families, admissible restriction, local completeness notion or minimal parameter count is supplied. Consequently ``minimal'' cannot be converted into a uniquely testable theorem from this sentence alone.}
\GSfield{Reconstruction}{Minimality under inclusion, minimal dimension for inversion, and minimality for detecting the continuous series are different assertions; none is selected editorially.}
\GSfield{Fate}{The exact-phrase and horospherical searches led to \citet{gindikin2000integral}, but no precise claim matching this unqualified sentence was identified.}
\GSfield{Present status}{As of 10 November 2026: Unidentified, rather than certified open or undeveloped.}
\GSfield{Assessment}{literal statement: not determinable; publication match: thematic only; verification: unresolved lead; development: none; confidence: low}
\end{GSitem}
\begin{GSitem}{1976-398.13}{E}{Straight rulings as translates of a unipotent subgroup}
\GSfield{Kind; source}{construction; \GSpage{P000661}}
\GSfield{The notes}{The examples are \(\left(\begin{smallmatrix}1&z\\0&1\end{smallmatrix}\right)\), \(\left(\begin{smallmatrix}h&z\\0&h^{-1}\end{smallmatrix}\right)\), and the general left--right translate; the ruling family has three parameters.}
\GSfield{Statement}{Every affine real line contained in \(\det X=1\) is \(A(I+tN)\), where \(A\in SL(2,\mathbb R)\), \(N\ne0\), \(\operatorname{tr}N=\det N=0\). Thus \(N^2=0\), and up to reparametrization the line is a left--right translate of the upper unipotent subgroup. The space of unoriented lines has dimension three.}
\GSfield{Reconstruction}{\ours\ Expand \(\det(A+tB)=1+t\operatorname{tr}(A^{-1}B)+t^2\det(A^{-1}B)\). At each \(A\), null directions form a one-dimensional projective cone; subtracting the one line parameter from \(3+1\) gives dimension 3.}
\GSfield{Fate}{These are the real horocycles in \citet[pp.418--422]{gindikin2000integral}; the classification follows by the elementary determinant calculation.}
\GSfield{Present status}{As of 10 November 2026: The construction and count are exact, with zero direction excluded.}
\GSfield{Assessment}{literal statement: correct; publication match: exact; verification: read the result; development: none; confidence: high}
\end{GSitem}
\begin{GSitem}{1976-398.14}{A}{The two complex domains and the Siegel coordinate chart}
\GSfield{Kind; source}{construction; \GSpage{P000661-P000662}}
\GSfield{The notes}{On the complex determinant-one quadric the inequality is \(\Im(\alpha\bar\beta)\Im(\delta\bar\beta)>(\Im\beta)^2\). The two signs are said to give domains with distinguished boundary \(G\); in three coordinates the boundary is called \(\mathbb R^3\).}
\GSfield{Statement}{On \(\beta\ne0\), put \(Z=\left(\begin{smallmatrix}\alpha/\beta&-1/\beta\\-1/\beta&\delta/\beta\end{smallmatrix}\right)\). The inequality is equivalent to \(\det\Im Z>0\), hence \(\Im Z\) is positive or negative definite. These are the two sign components in this chart. Their invariant continuation in \(SL(2,\mathbb C)\) has the real group as its distinguished boundary in the relevant Hardy realization. The chart omits \(Z_{12}=0\); neither this parametrization nor the phrase \(\mathbb R^3\) identifies the entire group with one Euclidean chart.}
\GSfield{Reconstruction}{\ours\ The coordinate calculation multiplies \(\det\Im Z\) by \(|\beta|^4\). ``Distinguished boundary'' requires a specified holomorphic function algebra or Hardy norm; it is not defined for arbitrary unbounded holomorphic functions solely by uniqueness.}
\GSfield{Fate}{\citet[\S 3, Proposition~5]{gelfandgindikin1977complex} publishes the same inequality and domain construction in 1977, after the 1976 talk.}
\GSfield{Present status}{As of 10 November 2026: The invariant domain realization is established; the global chart and function-class qualifications must accompany the notebook shorthand.}
\GSfield{Assessment}{literal statement: correct with missing hypotheses; publication match: partial; verification: read the result; development: none; confidence: high; problem \hyperref[prob:GP-1976-01]{GP-1976-01}}
\end{GSitem}
\begin{GSitem}{1976-398.15}{E}{A distinguished boundary is more than a uniqueness set}
\GSfield{Kind; source}{definition; \GSpage{P000662}}
\GSfield{The notes}{The parenthetical definition calls the skeleton the smallest boundary set from which a holomorphic function continuous on the closure is uniquely recovered.}
\GSfield{Statement}{For a unital uniform algebra \(\mathcal A\subset C(K)\) on a compact Hausdorff space, the Shilov boundary is the smallest closed set \(S\subset K\) such that \(\|f\|_K=\sup_S|f|\) for every \(f\in\mathcal A\). Mere uniqueness is weaker: for the disc algebra, every nontrivial closed boundary arc is a uniqueness set, whereas the Shilov boundary is the whole circle. Two disjoint such arcs show that a smallest uniqueness set need not exist.}
\GSfield{Reconstruction}{\ours\ Boundary vanishing on an arc gives analytic continuation by Schwarz reflection and then the identity theorem. Conversely, \((1+\bar\zeta z)/2\) peaks at any prescribed unit-circle point \(\zeta\), forcing every closed norming boundary to contain the whole circle. For the unbounded seminar domains one must specify the function class or compactification.}
\GSfield{Fate}{The classical minimal closed norming-boundary definition is recalled precisely in \citet[\S 1]{acosta2006boundaries}. The notebook's uniqueness-only definition is not equivalent to it.}
\GSfield{Present status}{As of 10 November 2026: The standard definition is established; the literal uniqueness-only definition is false in this generality.}
\GSfield{Assessment}{literal statement: false as written; publication match: partial; verification: read the result; development: annotated; confidence: high}
\end{GSitem}
\begin{GSitem}{1976-398.16}{A}{Hardy boundary values separate the two discrete series}
\GSfield{Kind; source}{theorem; \GSpage{P000662}}
\GSfield{The notes}{``Main theorem 2'' says \(\psi\in\mathcal H_A\) iff \(\psi=\psi_++\psi_-\), where the two functions continue analytically to \(D_+,D_-\). It also asserts zero intersection because everything is in \(L^2\).}
\GSfield{Statement}{Let \(H^2(D_\pm)\) be the invariant Hardy spaces with uniformly bounded \(L^2(G)\) norms on interior translates approaching the real distinguished boundary. Their boundary maps are isometries onto the holomorphic and antiholomorphic discrete-series summands \(H_\pm\). These summands are orthogonal, \(H_+\cap H_-=0\), and \(L^2_{\rm disc}(G)=H_+\oplus H_-\). Mere holomorphic continuation, without the Hardy bound and the specified boundary-value sense, is insufficient as a definition.}
\GSfield{Reconstruction}{The Hardy condition and boundary-limit interpretation are restored from the published theorem; they are not written in the notebook.}
\GSfield{Fate}{The precise statement appears in \citet[Theorem~3, Proposition~4 and \S 4]{gelfandgindikin1977complex}, received 6 June 1977.}
\GSfield{Present status}{As of 10 November 2026: Established for the specified Hardy spaces. The later continuous-series cohomological realization should not be substituted for this discrete-series theorem.}
\GSfield{Assessment}{literal statement: correct with missing hypotheses; publication match: exact; verification: read the result; development: none; confidence: high; problem \hyperref[prob:GP-1976-01]{GP-1976-01}}
\end{GSitem}
\begin{GSitem}{1976-398.17}{A}{The holomorphic projector kernel and its missing trace factor}
\GSfield{Kind; source}{claim; \GSpage{P000662}}
\GSfield{The notes}{The written kernel is \(c[2-\alpha_1\delta_2-\delta_1\alpha_2+\gamma_1\beta_2+\beta_1\gamma_2]^{-3/2}=c[2-\operatorname{tr}(g_1g_2^{-1})]^{-3/2}\). Both scan and transcriptions have just this factor.}
\GSfield{Statement}{For the invariant Haar normalization used in the published discrete-series projector, write \(t=\operatorname{tr}(g_1g_2^{-1})=\lambda+\lambda^{-1}\), choosing the holomorphic eigenvalue branch \(|\lambda|>1\) in the relevant semigroup. Up to a nonzero normalization constant and its consistent branch, the kernel is \[K=c\frac{\lambda^2}{(\lambda-1)^3(\lambda+1)}=c[(t-2)^3(t+2)]^{-1/2}.\] Its distributional boundary value gives the orthogonal projector. The extra factor \((t+2)^{-1/2}\) cannot be absorbed in \(c\).}
\GSfield{Reconstruction}{\oursComputed\ Summing the character series gives \[\frac{\sum_{n\ge1}n\lambda^{-n}}{\lambda-\lambda^{-1}}=\frac{\lambda^2}{(\lambda-1)^3(\lambda+1)}.\] Squaring and substituting \(t=\lambda+\lambda^{-1}\) verifies the trace expression independently. The full derivation and rational/numerical checks are preserved in the critical audit.}
\GSfield{Fate}{\citet[\S 4, Theorem~10, equations~18 and 22]{gelfandgindikin1977complex} gives the corrected kernel in 1977. This is a precise publication match to the construction, with a literal notebook error.}
\GSfield{Present status}{As of 10 November 2026: The corrected projector is established. The notebook's displayed formula is not the same invariant kernel even after constant renormalization.}
\GSfield{Assessment}{literal statement: false as written; publication match: partial; verification: read the result; development: annotated; confidence: high; problem \hyperref[prob:GP-1976-01]{GP-1976-01}}
\end{GSitem}
\begin{GSitem}{1976-398.18}{A}{The Fourier-integral technique behind the projector}
\GSfield{Kind; source}{method; \GSpage{P000662}}
\GSfield{The notes}{The concluding heading calls the ``submerged part'' Fourier-integral technique and describes \(K\) as an analogue of the Cauchy kernel.}
\GSfield{Statement}{In the verified construction, discrete-series characters are analytically continued into the semigroup and summed with their Plancherel weights to obtain a holomorphic reproducing kernel; its boundary value acts as a projection on \(L^2(G)\). This is a concrete implementation of the Cauchy-kernel analogy, not a separate general theorem about Fourier integral operators.}
\GSfield{Reconstruction}{The character-series procedure is supplied by the publication. The notebook itself records only the method label, with no phase, amplitude or symbol calculus.}
\GSfield{Fate}{\citet[\S 4, equations~15--22]{gelfandgindikin1977complex} supplies the precise construction in 1977.}
\GSfield{Present status}{As of 10 November 2026: The bounded method claim is identified. No broader Fourier-integral-operator theorem is attributed to this one sentence.}
\GSfield{Assessment}{literal statement: not determinable; publication match: partial; verification: read the result; development: none; confidence: high}
\end{GSitem}
\begin{GSitem}{1976-398.19}{A}{Arnold’s Möbius characterization of the two domains}
\GSfield{Kind; source}{conjecture; \GSpage{P000662}; V. I. Arnold}
\GSfield{The notes}{Arnold conjectures that \(D_+\) consists of transformations taking the disc into itself and \(D_-\) of transformations taking its exterior into itself. The last line says \(SL(2,\mathbb R)\) preserves the circle.}
\GSfield{Statement}{After Cayley conjugation from the real-projective-line/upper-half-plane model to the unit-circle model, the relevant open domains are the strict Möbius contraction semigroups: one sends the closed unit disc into the open disc, and the other sends the closed exterior in \(\mathbb{CP}^1\) into the open exterior. Their real edge is the circle-preserving form conjugate to \(SL(2,\mathbb R)\), modulo its central kernel. Without conjugation, the standard real matrices preserve \(\mathbb{RP}^1\), not the standard unit circle.}
\GSfield{Reconstruction}{The words ``strict'', the compactified exterior and the Cayley change of real form are required to distinguish the open domains from the semigroup boundary.}
\GSfield{Fate}{\citet[\S 3, Proposition~9]{gelfandgindikin1977complex} proves precisely this geometric identification in 1977.}
\GSfield{Present status}{As of 10 November 2026: Resolved by the 1977 proposition, with the model and strictness qualifications above.}
\GSfield{Assessment}{literal statement: correct with missing hypotheses; publication match: partial; verification: read the result; development: none; confidence: high; problem \hyperref[prob:GP-1976-02]{GP-1976-02}}
\end{GSitem}

\GStalk{399}{}{Academic-year heading}{P000659-P000664}{1976--77}
This catalogue entry labels the academic year and overlaps the opening lectures. It supplies no separately identifiable mathematical exposition. The actual occurrences in its page window belong to records 397--398 and 400--403 and are indexed there once.


\GStalk{400}{1976-09-13}{A. A. Kirillov}{P000663}{Superalgebras, supergroups and supermanifolds: announced heading}
Immediately below the date 13/IX-76, the manuscript gives Kirillov's name and this title, then explicitly switches to Fuchs and the Weil algebra. No intervening exposition is recorded. The title is retained as a programme heading; Fuchs's formulas and Kirillov's later intervention are assigned to records 401 and 402. The broad catalogue window P000662--P000664 must not be read as an independent three-page Kirillov lecture.


\GStalk{401}{1976-09-13}{D. B. Fuchs}{P000663-P000664}{The Weil algebra and its differential}
Fuchs constructs the graded algebra \(\Lambda\mathfrak g^*\otimes S\mathfrak g^*\), specifies its differential on generators, explains the signed Leibniz rule and gives the structure-constant formula on the symmetric generators. The concluding acyclicity statement omits the surviving degree-zero constants. Characteristic classes are announced for a future meeting, not developed here. This is established Weil--Cartan theory: \citet{cartan1950weil,cartan1950transgression}. The lecture begins on P000663 and ends before the explicit Kirillov intervention on P000664.

\begin{GSitem}{1976-401.1}{E}{The graded Weil algebra}
\GSfield{Kind; source}{definition; \GSpage{P000663}}
\GSfield{The notes}{\(W(\mathfrak g)=\Lambda^*(\mathfrak g')\otimes S^*(\mathfrak g')\), viewed as functions of \(2n\) variables, with \(\Lambda^p\otimes S^q\) in degree \(p+2q\).}
\GSfield{Statement}{For a finite-dimensional Lie algebra \(\mathfrak g\) over \(\mathbb R\) or \(\mathbb C\), \(W(\mathfrak g)\) is the free graded-commutative algebra on odd generators \(\theta^i\) of cohomological degree 1 and even generators \(u^i\) of degree 2, two copies of \(\mathfrak g^*\). Thus its underlying superdimension is \(n|n\); ``functions'' here means polynomials in the even variables and exterior polynomials in the odd ones.}
\GSfield{Reconstruction}{The field, finite-dimensionality and polynomial interpretation are supplied. No smooth or completed function algebra is implied.}
\GSfield{Fate}{The construction is already in \citet[Expos\'e~19, Weil-algebra construction]{cartan1950weil}, delivered 15 May 1950.}
\GSfield{Present status}{As of 10 November 2026: Standard definition. The earlier discussion of superfunctions in 1975-389.1 supplies context but is not the same construction.}
\GSfield{Assessment}{literal statement: correct; publication match: exact; verification: read the result; development: none; confidence: high}
\end{GSitem}
\begin{GSitem}{1976-401.2}{E}{The Chevalley differential sign convention}
\GSfield{Kind; source}{construction; \GSpage{P000663}}
\GSfield{The notes}{For \(h\in\Lambda^\ell\mathfrak g^*\), the displayed sum is \[\sum_{i<j}(-1)^{i+j-1}h([g_i,g_j],g_1,\ldots,\widehat g_i,\ldots,\widehat g_j,\ldots).\] In degree 1 it is \(h([g_1,g_2])\).}
\GSfield{Statement}{With one-based argument indices, define \(d_+h\) by exactly that alternating sum. It is the negative of the frequently used trivial-coefficient Chevalley--Eilenberg differential, so \(d_+^2=0\) by Jacobi and it obeys the graded Leibniz rule. Either overall sign convention computes the same cohomology.}
\GSfield{Reconstruction}{The remaining arguments retain their original order after the two omissions. The positive bracket sign is internally consistent and is not a transcription error.}
\GSfield{Fate}{Compare the exterior differential within the pre-existing Weil construction of \citet[Expos\'e~19]{cartan1950weil}; the overall sign is conventional.}
\GSfield{Present status}{As of 10 November 2026: Correct with the stated indexing and sign convention.}
\GSfield{Assessment}{literal statement: correct; publication match: exact; verification: read the result; development: none; confidence: high}
\end{GSitem}
\begin{GSitem}{1976-401.3}{E}{The unique Weil differential on both sets of generators}
\GSfield{Kind; source}{construction; \GSpage{P000663-P000664}}
\GSfield{The notes}{The differential sends \(h\otimes1\) to \((d_+h)\otimes1+1\otimes h\), is a graded derivation, and is said to extend uniquely with square zero. P000664 gives \(d(1\otimes h^k)=c_{ij}^{k}(h^i\otimes h^j)\).}
\GSfield{Statement}{If \([e_i,e_j]=c_{ij}^{k}e_k\), the notebook's sign convention gives \[d\theta^k=\tfrac12c_{ij}^{k}\theta^i\theta^j+u^k,\qquad du^k=c_{ij}^{k}\theta^iu^j.\] Together with \(d(ab)=d(a)b+(-1)^{\deg a}a\,d(b)\), these formulas define a unique degree-one derivation satisfying \(d^2=0\). The second formula is forced by \(d^2\theta^k=0\); Jacobi then gives \(d^2u^k=0\).}
\GSfield{Reconstruction}{\ours\ The factor \(1/2\) arises from summing over all ordered pairs \(i,j\); it disappears if one sums only \(i<j\). We have checked the sign of \(du\) against the notebook's positive exterior differential.}
\GSfield{Fate}{This is the Weil differential of \citet[Expos\'e~19]{cartan1950weil}, with an explicitly reconciled sign convention, not a new 1976 construction.}
\GSfield{Present status}{As of 10 November 2026: Correct. The two tensor factors in the structure-constant formula have degrees 1 and 2, respectively.}
\GSfield{Assessment}{literal statement: correct; publication match: exact; verification: read the result; development: none; confidence: high}
\end{GSitem}
\begin{GSitem}{1976-401.4}{E}{Acyclicity in positive degree, with constants surviving}
\GSfield{Kind; source}{theorem; \GSpage{P000664}}
\GSfield{The notes}{The ``main theorem'' says that the homologies of the complex are zero, without a qualification on degree.}
\GSfield{Statement}{For the unaugmented cochain algebra above, \(H^0(W(\mathfrak g),d)=\mathbb K\) and \(H^j(W(\mathfrak g),d)=0\) for \(j>0\). Equivalently, the augmentation \(W(\mathfrak g)\to\mathbb K\) is a quasi-isomorphism.}
\GSfield{Reconstruction}{\ours\ Put \(v^i=d\theta^i\). This is an invertible triangular replacement of the even generators; now \(d\theta^i=v^i\), \(dv^i=0\). The odd operator \(h=\sum_i\theta^i\partial/\partial v^i\) satisfies \(dh+hd=N\), where \(N\) counts the total number of generators. Division by positive \(N\) contracts every positive-weight component, but not the constant 1.}
\GSfield{Fate}{\citet[Expos\'e~20, \S 2, Theorem~1]{cartan1950transgression} states the precise acyclicity result in 1950.}
\GSfield{Present status}{As of 10 November 2026: The positive-degree theorem is established; literal vanishing in degree 0 is false. The notebook also says ``homologies'' although its differential raises degree.}
\GSfield{Assessment}{literal statement: false as written; publication match: partial; verification: read the result; development: annotated; confidence: high}
\end{GSitem}

\GStalk{402}{1976-09-13}{A. A. Kirillov}{P000664}{Transgression and invariant forms on compact Lie groups}
A short intervention explains the purpose of the Weil algebra in computing compact Lie-group cohomology. The intended correspondence is between primitive odd cohomology generators and indecomposable symmetric invariant polynomials. The notebook states a stronger, false dimension equality for all invariant forms and reverses the differential direction in its transgression shorthand. The surrounding Fuchs and Arnold lectures belong to separate records.

\begin{GSitem}{1976-402.1}{E}{Computing compact Lie-group cohomology from invariant forms}
\GSfield{Kind; source}{method; \GSpage{P000664}}
\GSfield{The notes}{Kirillov says that the algebra was devised for computing the homologies of simple compact Lie groups, and that invariant symmetric forms are easier to find than skew-symmetric ones.}
\GSfield{Statement}{For a compact connected Lie group \(K\) over \(\mathbb R\), averaging identifies de Rham cohomology with the cohomology of invariant forms; bi-invariant forms give the harmonic representatives for a bi-invariant metric. The cohomology algebra is exterior on primitive odd-degree generators. Invariant polynomials \((S\mathfrak k^*)^K\), together with transgression, determine their degrees. Simplicity is not needed for this real-coefficient statement.}
\GSfield{Reconstruction}{Connectedness, compactness and real coefficients are explicit. This does not assert an integral-cohomology computation, where torsion is invisible to differential forms.}
\GSfield{Fate}{\citet[Expos\'e~20, \S\S 3--4]{cartan1950transgression} supplies the transgression method before 1976.}
\GSfield{Present status}{As of 10 November 2026: Established real-coefficient theory. The method statement is distinct from the incorrect unrestricted counting rule in the next occurrence.}
\GSfield{Assessment}{literal statement: correct with missing hypotheses; publication match: exact; verification: read the result; development: none; confidence: high}
\end{GSitem}
\begin{GSitem}{1976-402.2}{E}{Primitive odd forms versus indecomposable symmetric invariants}
\GSfield{Kind; source}{claim; \GSpage{P000664}}
\GSfield{The notes}{The conclusion reads: as many invariant skew-symmetric forms of degree \(k\) as symmetric ones of degree \((k+1)/2\). A preceding sentence calls an odd form the differential of a symmetric expression.}
\GSfield{Statement}{Let \(I=(S\mathfrak k^*)^K\), \(I^+\) its positive-degree ideal, and \(Q(I)=I^+/(I^+)^2\). For compact connected \(K\), transgression identifies \(Q(I)_m\) with the primitive part of \(H^{2m-1}(K;\mathbb R)\), up to the usual choice of generators. A degree-\(m\) invariant polynomial has Weil degree \(2m\) and is \(d\) of a transgression element of degree \(2m-1\), whose exterior component represents the corresponding primitive class. The differential direction in the notebook must therefore be reversed.}
\GSfield{Reconstruction}{\ours\ For \(K=SU(2)\), the square of the Killing quadratic polynomial is a nonzero invariant symmetric polynomial of degree 4, whereas \(\Lambda^7\mathfrak{su}(2)^*=0\). This disproves the literal count at \(k=7\), while its indecomposable version is correct.}
\GSfield{Fate}{The corrected generator correspondence is \citet[Expos\'e~20, \S 4, Theorem~2]{cartan1950transgression}.}
\GSfield{Present status}{As of 10 November 2026: Established for primitive and indecomposable generators. The unrestricted equality and the displayed differential direction are false as written.}
\GSfield{Assessment}{literal statement: false as written; publication match: partial; verification: read the result; development: annotated; confidence: high}
\end{GSitem}

\GStalk{403}{1976-09-13}{V. I. Arnold}{P000664-P000666}{Point equivalence of second-order equations and projective geometry}
Arnold explicitly revisits nineteenth-century work of Lie, Tresse and Liouville and Cartan's 1924 projective connections. The lecture separates point linearization, pairs of line fields in a contact three-manifold, projective parameters along solutions, and inequivalent projective structures on a circle. It then normalizes the quadratic transverse terms, obtains a relative 5/2-density, and constructs projective scalar invariants from its reciprocal power. A damaged global question and Gelfand's proposed cohomological connection are recorded without inventing missing quantifiers. The local normalization calculations can be checked, but the stated flatness implication needs additional hypotheses.

\begin{GSitem}{1976-403.1}{E}{When can a scalar second-order ODE be straightened by point transformations?}
\GSfield{Kind; source}{question; \GSpage{P000664}}
\GSfield{The notes}{For \(y''=F(x,y,y')\), Arnold asks whether solution curves can be straightened. The page explicitly cites Tresse 1896 and Acta Mathematica 1894, and contrasts general equations with \(y''=0\).}
\GSfield{Statement}{For a smooth scalar equation on an open set in \((x,y,p)\), local equivalence under a nonsingular point change \((x,y)\mapsto(X,Y)\), on a chart where \(dX/dx\ne0\), to \(Y_{XX}=0\) holds exactly when the two relative invariants vanish identically there: \[I=F_{pppp},\quad H=D^2F_{pp}-4DF_{yp}-F_pDF_{pp}+6F_{yy}-3F_yF_{pp}+4F_pF_{yp},\qquad D=\partial_x+p\partial_y+F\partial_p.\] Vanishing at one point is insufficient; these are local identities on an open set.}
\GSfield{Reconstruction}{The displayed criterion and smooth regular setting are modern precision supplied from the classical equivalence theory, not formulas written on P000664.}
\GSfield{Fate}{\citet[\S 2, \S 4.1]{kruglikov2009tresse} reconstructs the Lie--Liouville--Tresse criterion and its nineteenth-century provenance. Cartan's 1924 treatment reaches the same flatness test geometrically \citep[\S VIII, no.24]{cartan1924projective}.}
\GSfield{Present status}{As of 10 November 2026: The local point-linearization problem is solved. The 2008 preprint/2009 exposition is not the first circulation of the classical result.}
\GSfield{Assessment}{literal statement: correct; publication match: exact; verification: read the result; development: none; confidence: high; problem \hyperref[prob:GP-1976-03]{GP-1976-03}}
\end{GSitem}
\begin{GSitem}{1976-403.2}{E}{Every linear second-order equation is locally point-flat}
\GSfield{Kind; source}{theorem; \GSpage{P000664}}
\GSfield{The notes}{The notes state that \(y''+a(x)y'+b(x)y=0\) becomes \(y''=0\) after a change of independent variable and a change linear in \(y\); configuration theorems such as Desargues are then transferred to the solution family.}
\GSfield{Statement}{Let \(a,b\) be smooth on an interval and let \(u,v\) be independent solutions, with \(v\ne0\) on a smaller interval. Then \(X=u/v\), \(Y=y/v\) is a local point diffeomorphism because \(X'=W/v^2\ne0\), and every solution becomes \(Y=c_1X+c_2\). Any incidence theorem formulated for straight-line configurations transfers to configurations lying wholly in this coordinate neighbourhood; no global straightening across zeros or monodromy is asserted.}
\GSfield{Reconstruction}{\ours\ The ratio-of-solutions proof also states the actual domain restriction. The doubtful Russian phrase translated ``linear menu'' is not used to infer extra content.}
\GSfield{Fate}{This classical fact is recalled in \citet[introduction]{kruglikov2009tresse}; Arnold presents it explicitly as old material.}
\GSfield{Present status}{As of 10 November 2026: Completely established locally. The incidence-theorem remark is a consequence of transporting points and curves by this diffeomorphism.}
\GSfield{Assessment}{literal statement: correct; publication match: exact; verification: read the result; development: none; confidence: high}
\end{GSitem}
\begin{GSitem}{1976-403.3}{E}{Cartan geometry of a general second-order equation}
\GSfield{Kind; source}{construction; \GSpage{P000665}}
\GSfield{The notes}{Cartan 1924--25 is invoked: solutions are geodesics of a projective connection. The next lines pass to the space of line elements \((x,y,y')\).}
\GSfield{Statement}{A regular second-order equation defines a path geometry on the three-dimensional contact-element space and has a canonical normal Cartan connection of projective path-geometry type. It descends to an ordinary projective equivalence class of torsion-free affine connections on the two-dimensional point space precisely when \(F\) is cubic or lower in \(p=y'\). Thus a projective connection on contact elements and one on points are different notions.}
\GSfield{Reconstruction}{We preserve the passage to contact elements rather than automatically declaring the general-\(F\) sentence false. The surface-connection reading requires the cubic condition recorded later on P000666.}
\GSfield{Fate}{\citet[introduction and \S\S VII--VIII, especially nos.22--24]{cartan1924projective} explicitly distinguishes the two geometries and proves the general contact-element construction.}
\GSfield{Present status}{As of 10 November 2026: Established; the distinction prevents importing an affine-geodesic claim for arbitrary \(F\).}
\GSfield{Assessment}{literal statement: correct with missing hypotheses; publication match: exact; verification: read the result; development: none; confidence: high}
\end{GSitem}
\begin{GSitem}{1976-403.4}{E}{Two transverse line fields spanning a contact distribution}
\GSfield{Kind; source}{construction; \GSpage{P000665}}
\GSfield{The notes}{Arnold says that reducing a pair of direction fields in \(\mathbb R^3\) under diffeomorphisms is equivalent to the ODE problem, then remarks that their span is a contact structure.}
\GSfield{Statement}{Let \(L,V\) be an ordered pair of smooth rank-one distributions on a three-manifold, transverse at the point, with \(L\oplus V\) contact. Locally choose contact coordinates \((x,y,p)\) in which \(V=\langle\partial_p\rangle\) and \(L=\langle\partial_x+p\partial_y+F\partial_p\rangle\). Local diffeomorphisms preserving the ordered pair correspond to point transformations on the local leaf space of \(V\), with their prolongations. Arbitrary pairs whose span is integrable are excluded.}
\GSfield{Reconstruction}{The rank, ordering, contact and local leaf-space assumptions are restored explicitly. Exchanging the two fields gives the dual equation rather than the same point-equivalence problem.}
\GSfield{Fate}{The contact-element formulation and duality are in \citet[\S VII and \S VIII, no.23]{cartan1924projective}.}
\GSfield{Present status}{As of 10 November 2026: The regular contact version is exact. Without the contact hypothesis the announced equivalence is false.}
\GSfield{Assessment}{literal statement: correct with missing hypotheses; publication match: exact; verification: read the result; development: none; confidence: high}
\end{GSitem}
\begin{GSitem}{1976-403.5}{E}{A projective parameter along each solution}
\GSfield{Kind; source}{construction; \GSpage{P000665-P000666}}
\GSfield{The notes}{After choosing one solution as \(y=0\), the linear part is straightened by \(\widetilde x=f(x)\), \(\widetilde y=g(x)y\). The function \(f\) satisfies a Schwarz-type equation, and a projective structure is said to arise on every solution.}
\GSfield{Statement}{Along a smooth solution \(y_0\), infinitesimal transverse variations satisfy \(\eta''=F_p(x,y_0,y_0')\eta'+F_y(x,y_0,y_0')\eta\). If \(u,v\) are independent Jacobi solutions, \([u(x):v(x)]\in\mathbb{RP}^1\) gives a local projective coordinate; changing the Jacobi basis changes it by a Möbius map. For \(\eta''+a\eta'+b\eta=0\), its affine ratio satisfies \(S(f)=2b-a'-a^2/2\), where \(S(f)=f^{(3)}/f'-\tfrac32(f''/f')^2\).}
\GSfield{Reconstruction}{\ours\ Removing the first derivative by \(\eta=e^{-\frac12\int a}z\) gives \(z''+(b-a'/2-a^2/4)z=0\); the ratio identity proves the Schwarzian formula. The construction is local along a nonsingular solution.}
\GSfield{Fate}{This is the projective parametrization inherent in the classical connection theory \citep[\S\S VII--VIII]{cartan1924projective}; the notebook explains it through linearization rather than a new existence theorem.}
\GSfield{Present status}{As of 10 November 2026: Established locally. Global holonomy and the topology of closed solutions are additional data, not determined by this local construction.}
\GSfield{Assessment}{literal statement: correct; publication match: partial; verification: read the result; development: none; confidence: high}
\end{GSitem}
\begin{GSitem}{1976-403.6}{E}{Inequivalent projective structures on the circle}
\GSfield{Kind; source}{construction; \GSpage{P000665}}
\GSfield{The notes}{The standard \(\mathbb{RP}^1\) and its pullback under a degree-two covering \(S^1\to\mathbb{RP}^1\) are called inequivalent. Arnold proposes examining infinitesimal automorphisms \(\partial_x,x\partial_x,x^2\partial_x\).}
\GSfield{Statement}{For each positive integer \(n\), pull back the projective atlas on \(\mathbb{RP}^1\) by the degree-\(n\) covering. The resulting circle structures are pairwise inequivalent under projective-structure-preserving diffeomorphisms. In an angular coordinate their infinitesimal projective automorphism fields are spanned by \(\partial_\theta,\sin(n\theta)\partial_\theta,\cos(n\theta)\partial_\theta\). A nonzero such field can have at most \(2n\) distinct zeros, and the maximum is achieved.}
\GSfield{Reconstruction}{\ours\ The zero-count maximum is preserved by diffeomorphisms, distinguishing \(n=1\) from \(n=2\) as the notebook requests. The abstract Lie algebra is \(\mathfrak{sl}_2\) in both cases; its realization on the circle, not its abstract isomorphism type, distinguishes them.}
\GSfield{Fate}{This is an elementary consequence of the classical projective action; \citet[projective frame construction]{cartan1924projective} supplies the historical framework. No complete classification of all circle projective structures is claimed here.}
\GSfield{Present status}{As of 10 November 2026: The two stated examples, and their covering generalization, are rigorously distinguished by the displayed argument.}
\GSfield{Assessment}{literal statement: correct; publication match: partial; verification: read the result; development: annotated; confidence: high}
\end{GSitem}
\begin{GSitem}{1976-403.7}{F}{The damaged global projective-structure question}
\GSfield{Kind; source}{question; \GSpage{P000665}}
\GSfield{The notes}{The line marked ``Problem'' contains two cancellations and survives only as a question involving ``the equation'', ``globally'' and ``locally projective space''. The scan does not recover the missing relation.}
\GSfield{Statement}{No unique global problem is recoverable. It might concern developing maps, holonomy, realization of a projective circle as a solution, or global equivalence of path geometries; these have different objects and answers.}
\GSfield{Reconstruction}{The list of possible readings explains the uncertainty and is not a proposed reconstruction of Arnold's question.}
\GSfield{Fate}{Searches of the exact local-projective vocabulary and Cartan/Arnold sources did not resolve the deleted words. \citet{cartan1924projective} supplies context only.}
\GSfield{Present status}{As of 10 November 2026: Unidentified from the source; no open-problem or nondevelopment claim is made.}
\GSfield{Assessment}{literal statement: not determinable; publication match: thematic only; verification: unresolved lead; development: none; confidence: low}
\end{GSitem}
\begin{GSitem}{1976-403.8}{E}{Removing the mixed quadratic transverse terms}
\GSfield{Kind; source}{method; \GSpage{P000666}}
\GSfield{The notes}{After linear straightening, \(y''=a(x)y^2+b(x)yy'+c(x)y'^2+\cdots\). Coordinate changes \(\widetilde x=x+yf(x,y)\), \(\widetilde y=y(1+yg(x,y))\) are said to kill \(b,c\).}
\GSfield{Statement}{Work locally near \(y=p=0\) with zero constant and linear transverse terms. Up to terms of total degree at least 3 in \((y,p)\), use \(X=x+f(x)y\), \(Y=y+g(x)y^2\). The new quadratic coefficients are \(a+g''\), \(b+4g'-f''\), \(c+2g-2f'\). Choosing \(g=f'-c/2\) and \(f''=(2c'-b)/3\) removes the last two and leaves \[A=a+\tfrac16c''-\tfrac13b'.\] The coordinate Jacobian equals 1 along the reference solution.}
\GSfield{Reconstruction}{\oursComputed\ The explicit coefficient formulas, local integrations for \(f,g\), and the stated jet order are editorial calculations. No convergence of an infinite normal-form process is inferred.}
\GSfield{Fate}{The residual coefficient is \(H/12\) on the normalized solution, where \(H\) is the classical Tresse invariant reviewed in \citet[\S 2.1]{kruglikov2009tresse}. This identifies the old invariant behind the notebook calculation.}
\GSfield{Present status}{As of 10 November 2026: The finite-order elimination is verified algebraically. The normalization does not itself settle all higher-order terms.}
\GSfield{Assessment}{literal statement: correct; publication match: partial; verification: read the result; development: annotated; confidence: high}
\end{GSitem}
\begin{GSitem}{1976-403.9}{E}{The relative 5/2-density along a solution}
\GSfield{Kind; source}{construction; \GSpage{P000666}}
\GSfield{The notes}{The remaining coefficient is written \(a(x)(dx)^{5/2}\), defined up to a multiplicative constant depending on the solution.}
\GSfield{Statement}{Let \(A\) be the residual quadratic coefficient after the preceding normalization. On an oriented projective coordinate interval, the residual transformations preserving the flat variational equation have \(X=\phi(x)\), \(Y=k\sqrt{\phi'(x)}\,y+O(y^2)\), with \(\phi\) Möbius, \(\phi'>0\), \(k\ne0\). After renormalizing mixed quadratic terms, \(\widetilde A=A/(k(\phi')^{5/2})\), so \(A(dx)^{5/2}\) is determined up to the constant factor \(k^{-1}\). For unoriented real coordinates use the corresponding density bundle.}
\GSfield{Reconstruction}{\ours\ The weight is checked by the transformation of the normalized equation. The notation denotes a tensor density, not an ordinary differential form of noninteger exterior degree; its scale is not canonically fixed.}
\GSfield{Fate}{It is the solutionwise form of the classical relative invariant \(H\) in \citet[\S 2.1, weights]{kruglikov2009tresse}. An exact later Arnold publication of this particular normalization was not verified; the invariant identification, rather than that unverified lead, supports class E.}
\GSfield{Present status}{As of 10 November 2026: The local relative-density statement is well defined with orientation and scale specified.}
\GSfield{Assessment}{literal statement: correct with missing hypotheses; publication match: partial; verification: read the result; development: annotated; confidence: medium}
\end{GSitem}
\begin{GSitem}{1976-403.10}{E}{Why vanishing on one solution does not imply formal straightening}
\GSfield{Kind; source}{claim; \GSpage{P000666}}
\GSfield{The notes}{The recorded theorem says that if \(a(x)(dx)^{5/2}\equiv0\), the equation can be straightened in formal series. The page has just fixed a single solution; no ``all solutions'' or cubic hypothesis is stated in the theorem line.}
\GSfield{Statement}{Taken along one chosen solution, the implication is false: \(y''=y^3\) has the solution \(y=0\), zero linear and quadratic transverse parts, and hence zero displayed density there, but its Tresse invariant is \(H=36y\), a nonzero formal series. It is therefore not formally point-equivalent to \(Y''=0\) near that solution. A sufficient corrected condition is the simultaneous formal vanishing of \(I,H\); in the cubic-in-\(p\) class, vanishing of \(H\) on an open set suffices because \(I=0\).}
\GSfield{Reconstruction}{\ours\ This counterexample audits the strongest literal interpretation supported by the neighbouring sentences. We do not claim it disproves an unstated theorem about a family of all solutions with additional assumptions.}
\GSfield{Fate}{The obstruction is already part of classical point equivalence \citep[\S 2 and \S 4.1]{kruglikov2009tresse}; the notebook's abbreviated implication is not a new unresolved conjecture.}
\GSfield{Present status}{As of 10 November 2026: False under the one-solution reading. The source does not determine which extra quantifiers Arnold intended.}
\GSfield{Assessment}{literal statement: false as written; publication match: partial; verification: read the result; development: annotated; confidence: high; problem \hyperref[prob:GP-1976-03]{GP-1976-03}}
\end{GSitem}
\begin{GSitem}{1976-403.11}{E}{The normal form with a squared first derivative factor}
\GSfield{Kind; source}{method; \GSpage{P000666}}
\GSfield{The notes}{Further diffeomorphisms are said to lead to \(y''=y'^2A(x,y,y')\).}
\GSfield{Statement}{Every regular smooth scalar second-order equation admits this form locally near a chosen initial direction after a point transformation. First choose a local foliation by solutions and use its leaf label as \(y\), obtaining \(F(x,y,0)=0\). Then put \(X=\phi(x,y)\), \(Y=y\), with \(\phi_x\ne0\) and \(\phi_{xx}=F_p(x,y,0)\phi_x\). The transformed right side has both value and first \(P\)-derivative zero at \(P=0\), hence equals \(P^2A(X,Y,P)\) by smooth division.}
\GSfield{Reconstruction}{\ours\ The foliation exists by the local ODE existence theorem and a transverse initial curve. The second coordinate change is obtained by integrating a linear equation for \(\phi_x\), choosing it nonzero. This proof supplies the precise local scope omitted in the notes.}
\GSfield{Fate}{This is a local coordinate normalization within the classical path-geometry equivalence problem \citep[\S\S VII--VIII]{cartan1924projective}, not a criterion for flatness.}
\GSfield{Present status}{As of 10 November 2026: The local finite normalization is verified. It does not make the arbitrary coefficient \(A\) vanish and imposes no global reduction.}
\GSfield{Assessment}{literal statement: correct; publication match: partial; verification: read the result; development: annotated; confidence: high}
\end{GSitem}
\begin{GSitem}{1976-403.12}{E}{Cubic slope equations form an invariant projective class}
\GSfield{Kind; source}{theorem; \GSpage{P000666}}
\GSfield{The notes}{The equation \(y''=P_3(y')\), with coefficients depending on \(x,y\), is said to be invariant under plane diffeomorphisms and is associated with R. Liouville.}
\GSfield{Statement}{The class \(F=A_0(x,y)+A_1(x,y)p+A_2(x,y)p^2+A_3(x,y)p^3\) is preserved by local point transformations wherever the transformed independent variable is valid. Exactly these equations describe the unparametrized geodesics of torsion-free affine connections on a surface, modulo projective equivalence \(\widehat\nabla_XY=\nabla_XY+\alpha(X)Y+\alpha(Y)X\).}
\GSfield{Reconstruction}{The coefficients may vary on the surface; ``third degree'' means degree at most 3. The graph-coordinate restriction excludes directions sent to vertical ones.}
\GSfield{Fate}{\citet[introduction, \S VIII no.24]{cartan1924projective} explicitly gives the cubic criterion; \citet[\S 4.1]{kruglikov2009tresse} supplies a modern account of Lie and Liouville's earlier work.}
\GSfield{Present status}{As of 10 November 2026: Established projective differential geometry, already classical in 1976.}
\GSfield{Assessment}{literal statement: correct; publication match: exact; verification: read the result; development: none; confidence: high}
\end{GSitem}
\begin{GSitem}{1976-403.13}{E}{A vector field from a nonvanishing relative density}
\GSfield{Kind; source}{construction; \GSpage{P000666}}
\GSfield{The notes}{The scalar-invariant discussion begins with \(v(x)=a(x)^{-2/5}\partial_x\).}
\GSfield{Statement}{On an oriented interval where the real density coefficient \(A\) is nowhere zero, choose a consistent fifth-root convention and set \(v=A^{-2/5}\partial_x\). Under coordinate changes this is a vector field. Rescaling the original density by a nonzero constant rescales \(v\) by a constant; at zeros of \(A\) the formula is singular and provides no smooth invariant field.}
\GSfield{Reconstruction}{The nonvanishing and branch qualifications, and the residual constant ambiguity, are necessary. They prevent treating this as a globally defined absolute invariant of every equation.}
\GSfield{Fate}{This is algebra on the classical relative invariant identified in \citet[\S 2.1]{kruglikov2009tresse}, not an independent classification theorem.}
\GSfield{Present status}{As of 10 November 2026: Well defined only on the stated nonvanishing locus, up to the inherited constant.}
\GSfield{Assessment}{literal statement: correct with missing hypotheses; publication match: exact; verification: read the result; development: none; confidence: high}
\end{GSitem}
\begin{GSitem}{1976-403.14}{E}{The projective scalar I2 and the adjoint quadratic form}
\GSfield{Kind; source}{construction; \GSpage{P000666}}
\GSfield{The notes}{The notebook says that the projective group acts on two-jets by the adjoint representation and writes \(I_2=2v''v-v'^2\).}
\GSfield{Statement}{For a specified smooth vector field \(v(x)\partial_x\) on a projective line, \(I_2(v)=2vv''-(v')^2\) is a scalar invariant under Möbius changes. More generally, if \(X=\phi(x)\), \(V(X)=\phi'(x)v(x)\), then \[2VV_{XX}-V_X^2=2vv''-v'^2+2v^2S(\phi).\] The quadratic polynomial vector fields form \(\mathfrak{sl}_2\) under the usual bracket; evaluating their two-jet at a point gives an isomorphism, and \(I_2\) is the corresponding discriminant/Killing quadratic form up to normalization.}
\GSfield{Reconstruction}{\ours\ The transformation identity proves invariance since \(S(\phi)=0\) for Möbius maps. A pointed two-jet includes a base point; the notebook's adjoint statement refers to the induced identification of its three coefficient coordinates with a quadratic projective vector field.}
\GSfield{Fate}{This is classical projective invariant theory; the Cartan framework predates the talk \citep{cartan1924projective}. The displayed identity is independently checked in the audit rather than attributed to an unread formula in a later text.}
\GSfield{Present status}{As of 10 November 2026: Correct for a specified vector field and projective coordinate changes. Its value scales by \(c^2\) under \(v\mapsto cv\), so it is not absolute for the unnormalized ODE density.}
\GSfield{Assessment}{literal statement: correct with missing hypotheses; publication match: partial; verification: read the result; development: annotated; confidence: high}
\end{GSitem}
\begin{GSitem}{1976-403.15}{E}{Invariant differentiation and the higher In}
\GSfield{Kind; source}{construction; \GSpage{P000666}}
\GSfield{The notes}{The recurrence is \(I_3=vI_2'\), then \(I_n=vI_{n-1}'\).}
\GSfield{Statement}{The operator \(\nabla_v=v(x)d/dx\) is coordinate invariant. Starting from the Möbius scalar \(I_2\), define \(I_n=\nabla_v^{n-2}I_2\) for \(n\ge3\); each is a scalar invariant of the specified vector field on the projective line. If \(v\) is rescaled by a constant \(c\), then \(I_n\) is rescaled by \(c^n\). No completeness or functional independence of this list is claimed.}
\GSfield{Reconstruction}{\ours\ Coordinate invariance follows from \(V\partial_X=v\partial_x\). Induction proves the scaling law, recording exactly the ambiguity inherited from the density.}
\GSfield{Fate}{Invariant differentiation is part of the classical differential-invariant calculus reviewed in \citet[\S 1--2]{kruglikov2009tresse}. The notebook's narrower recurrence is verified directly here.}
\GSfield{Present status}{As of 10 November 2026: The recurrence is correct with its scale qualification; it is not a complete classification of second-order equations.}
\GSfield{Assessment}{literal statement: correct; publication match: partial; verification: read the result; development: annotated; confidence: high}
\end{GSitem}
\begin{GSitem}{1976-403.16}{F}{Gelfand’s proposed connection with Gelfand–Fuchs cohomology}
\GSfield{Kind; source}{conjecture; \GSpage{P000666}; I. M. Gelfand}
\GSfield{The notes}{The final line says \ru{Гипотеза Гельфанда: связано с когомологиями Гельфанда--Фукса}. Pommaret is then mentioned as having worked on differential invariants.}
\GSfield{Statement}{The source specifies neither a Lie algebra/module, a cohomological degree, a cocycle nor a map relating cohomology to the displayed invariants. Consequently no definite conjecture with verifiable quantifiers can be extracted.}
\GSfield{Reconstruction}{The cohomological terminology is retained as a proposed connection. It is not automatically identified with the Virasoro cocycle, the Schwarzian cocycle or an arbitrary theorem about formal vector fields.}
\GSfield{Fate}{Targeted searches of the 5/2-density, projective-vector-field invariants and Gelfand--Fuchs terminology did not identify an exact statement. \citet{kruglikov2009tresse} matches the differential-invariant side only.}
\GSfield{Present status}{As of 10 November 2026: Unidentified as a precise conjecture. No C/D classification or development claim is warranted.}
\GSfield{Assessment}{literal statement: not determinable; publication match: thematic only; verification: unresolved lead; development: none; confidence: low}
\end{GSitem}

\GStalk{404}{1976-09-20}{S. P. Novikov}{P000667-P000669}{Lie superalgebras and the Wess--Zumino multiplet}
Novikov explicitly recommends Ogievetsky's UFN review and a Zumino review, with Volkov and Akulov noted in the margin. He defines Lie superalgebras, uses real Majorana spinors to construct supertranslations, and records Gelfand's request to adjoin Lorentz generators. The Wess--Zumino chiral multiplet, its auxiliary fields, Lagrangian and infinitesimal transformations occupy the remainder. The material predates the talk, as the cited 1975 review confirms. Several faithful notebook errors require emendation: Gol'fand's name, the scalar label on the spinor, an uncontracted kinetic index, and incompatible coefficients between the two action formulas.

\begin{GSitem}{1976-404.1}{E}{Lie superalgebra signs and parity}
\GSfield{Kind; source}{definition; \GSpage{P000667-P000668}}
\GSfield{The notes}{\(L=L_0\oplus L_1\), \([a,b]=(-1)^{da\,db+1}[b,a]\), and Jacobi ``with signs''. The first definition writes \(da=\dim a\), while the continuation correctly writes \(\operatorname{deg}a\).}
\GSfield{Statement}{Over a field of characteristic zero, a Lie superalgebra is a \(\mathbb Z/2\)-graded vector space with a bilinear degree-zero bracket, satisfying for homogeneous elements \([a,b]=-(-1)^{|a||b|}[b,a]\) and \((-1)^{|a||c|}[a,[b,c]]+(-1)^{|b||a|}[b,[c,a]]+(-1)^{|c||b|}[c,[a,b]]=0\). The number \(|a|\) is parity, not vector-space dimension.}
\GSfield{Reconstruction}{The complete Jacobi formula and homogeneous-element qualification are supplied. The subsequent manuscript line decides the intended meaning of \(da\).}
\GSfield{Fate}{\citet[\S 2]{ogievetskymezincescu1975superfields} gives the pre-talk algebraic framework. The seminar had already treated Lie superalgebras in 1975-389.5.}
\GSfield{Present status}{As of 10 November 2026: Standard definition; the dimensional notation is a faithful local error resolved by the next page.}
\GSfield{Assessment}{literal statement: false as written; publication match: partial; verification: read the result; development: none; confidence: high}
\end{GSitem}
\begin{GSitem}{1976-404.2}{E}{The associative supercommutator}
\GSfield{Kind; source}{construction; \GSpage{P000667-P000668}}
\GSfield{The notes}{Even-even and even-odd brackets are commutators; odd-odd brackets are anticommutators in an associative matrix algebra.}
\GSfield{Statement}{For an associative \(\mathbb Z/2\)-graded algebra \(A\) with \(A_iA_j\subset A_{i+j}\), \([a,b]_s=ab-(-1)^{|a||b|}ba\) defines a Lie superalgebra. Any graded linear subspace closed under this bracket is a Lie subsuperalgebra. A matrix example requires a chosen grading, such as the block grading of \(\operatorname{End}(V_0\oplus V_1)\). Closure under associative multiplication is not required for the subsuperalgebra.}
\GSfield{Reconstruction}{The grading of the ambient associative algebra is made explicit; arbitrary ungraded matrix subspaces with an assigned parity need not work.}
\GSfield{Fate}{This standard realization is used in \citet[\S 2]{ogievetskymezincescu1975superfields}; see also the earlier occurrence 1975-389.5.}
\GSfield{Present status}{As of 10 November 2026: Established directly by expanding the graded Jacobi identity.}
\GSfield{Assessment}{literal statement: correct; publication match: exact; verification: read the result; development: none; confidence: high}
\end{GSitem}
\begin{GSitem}{1976-404.3}{E}{Real Majorana spinors in signature (-,+,+,+)}
\GSfield{Kind; source}{construction; \GSpage{P000668}}
\GSfield{The notes}{The metric is \(\operatorname{diag}(-1,1,1,1)\), the spinor space is explicitly corrected to real dimension 4, and four real matrices satisfy \(\gamma_\mu\gamma_\nu+\gamma_\nu\gamma_\mu=2\eta_{\mu\nu}\). A barred spinor is written by multiplication by \(\gamma_0\).}
\GSfield{Statement}{For this signature, the real Clifford algebra \(\mathrm{Cl}(3,1)\) is isomorphic to \(M_4(\mathbb R)\) and acts on a four-dimensional real Majorana module \(S\). One may choose \(\gamma_0\) antisymmetric and \(\gamma_i\) symmetric. A charge-conjugation bilinear pairing must be specified when converting a column spinor to a row; for a compatible real basis one can write \(\bar s=s^{\mathsf T}C\). The covariance formulas depend on this convention.}
\GSfield{Reconstruction}{The transpose and charge-conjugation matrix distinguish the algebraic pairing from an unexplained row/column product. ``Approximately quaternions'' is only an analogy, not the Clifford-algebra identification.}
\GSfield{Fate}{The Majorana conventions used for the component theory are discussed in \citet[\S 3 and Appendix]{ogievetskymezincescu1975superfields}.}
\GSfield{Present status}{As of 10 November 2026: The real four-component construction is correct for the stated signature; barred-index formulas require a pairing convention.}
\GSfield{Assessment}{literal statement: correct with missing hypotheses; publication match: exact; verification: read the result; development: none; confidence: high}
\end{GSitem}
\begin{GSitem}{1976-404.4}{E}{The four-dimensional supertranslation algebra}
\GSfield{Kind; source}{construction; \GSpage{P000668}}
\GSfield{The notes}{After crossing out ``Poincar\'e'', the even part is translations; the odd part is spinors, with \([L_0,L_1]=0\) and \([\bar Q_i,\bar Q_j]=-2(\gamma^\mu)_{ij}P_\mu\).}
\GSfield{Statement}{Let \(V=\mathbb R^{1,3}\) and \(S\) a real Majorana spin module. A nonzero \(\operatorname{Spin}(1,3)\)-equivariant symmetric map \(\beta:S\otimes S\to V\) defines a Lie superalgebra \(V\oplus S\) by \([s,t]=\beta(s,t)\), with \(V\) central. In a Majorana matrix formula the symmetric matrices are \(C\gamma^\mu\), with indices placed consistently. Equivalently, the complex Weyl notation has \(\{Q_\alpha,\bar Q_{\dot\beta}\}=2\sigma^\mu_{\alpha\dot\beta}P_\mu\), \(\{Q,Q\}=\{\bar Q,\bar Q\}=0\).}
\GSfield{Reconstruction}{The invariant pairing and normalizations replace the notebook's unqualified two-lowered-index gamma matrix. Because odd brackets land in the central even space, the three-odd Jacobi identity is immediate.}
\GSfield{Fate}{The pre-talk construction appears in \citet[\S 2]{ogievetskymezincescu1975superfields}; the 1971 algebra is reproduced in \citet[Appendix, equation~1]{likhtman2001supersymmetry}.}
\GSfield{Present status}{As of 10 November 2026: Established supertranslation algebra. It has four even translations and four real odd generators, before Lorentz generators are added.}
\GSfield{Assessment}{literal statement: correct with missing hypotheses; publication match: partial; verification: read the result; development: none; confidence: high}
\end{GSitem}
\begin{GSitem}{1976-404.5}{E}{Extending supertranslations to the Poincaré superalgebra}
\GSfield{Kind; source}{problem; \GSpage{P000668}; I. M. Gelfand}
\GSfield{The notes}{Gelfand asks to extend the Poincar\'e algebra to a superalgebra. The preceding historical sentence names ``Gelfand and Likhtman''. }
\GSfield{Statement}{Adjoin \(\mathfrak{so}(1,3)\) to the supertranslation algebra by its vector action on \(V\) and spin action \(\rho\) on \(S\): \(\mathfrak s=\mathfrak{so}(1,3)\ltimes(V\oplus S)\), \([M,s]=\rho(M)s\). The \(\mathfrak{so}(1,3)\)-equivariance of \(\beta\) is exactly the mixed Jacobi condition. Thus \(\mathfrak s_0=\mathfrak{iso}(1,3)\), \(\mathfrak s_1=S\), giving the four-dimensional \(N=1\) super-Poincar\'e algebra without central charges.}
\GSfield{Reconstruction}{The request is interpreted as the existence/construction problem actually written, not a classification of all possible extensions. The historical name is Yu. A. Gol'fand, not the seminar's I. M. Gelfand.}
\GSfield{Fate}{\citet[Appendix, equation~1 and release note]{likhtman2001supersymmetry} reproduces a preprint passed for printing 12 April 1971; its references identify the Gol'fand--Likhtman 1971 publication. The requested extension was therefore already known before this talk.}
\GSfield{Present status}{As of 10 November 2026: Solved by the displayed semidirect construction; not a 1976 open problem. The 2001 reprint date is separated from the 1971 circulation and publication.}
\GSfield{Assessment}{literal statement: correct; publication match: exact; verification: read the result; development: none; confidence: high; problem \hyperref[prob:GP-1976-04]{GP-1976-04}}
\end{GSitem}
\begin{GSitem}{1976-404.6}{E}{The off-shell Wess--Zumino chiral multiplet}
\GSfield{Kind; source}{definition; \GSpage{P000668}}
\GSfield{The notes}{The list is \(A\) scalar, \(B\) pseudoscalar, \(\psi\) again labelled scalar, and auxiliary \(F\) scalar, \(G\) pseudoscalar. A bracket groups the five fields as a supermultiplet.}
\GSfield{Statement}{The four-dimensional \(N=1\) chiral multiplet comprises two real bosonic fields \(A,B\), one four-real-component Majorana fermion \(\psi\), and two real auxiliary bosons \(F,G\). Off shell it has four real bosonic and four real fermionic component functions; the auxiliary fields have no kinetic terms. The scalar label on \(\psi\) is incompatible with its gamma-matrix kinetic term and supersymmetry variations.}
\GSfield{Reconstruction}{The spinor correction is compelled both by the following notebook formulas and the named model. The off-shell component count does not count propagating degrees of freedom.}
\GSfield{Fate}{\citet[\S 3, chiral superfield; \S 4(b)]{ogievetskymezincescu1975superfields} describes the pre-existing Wess--Zumino model, introduced in 1974.}
\GSfield{Present status}{As of 10 November 2026: Established multiplet; \(\psi\) must be a spinor. No experimental discovery of a supersymmetric multiplet is implied.}
\GSfield{Assessment}{literal statement: false as written; publication match: partial; verification: read the result; development: none; confidence: high}
\end{GSitem}
\begin{GSitem}{1976-404.7}{E}{A consistent supersymmetric action and the notebook’s coefficient defects}
\GSfield{Kind; source}{construction; \GSpage{P000668-P000669}}
\GSfield{The notes}{Two formulas are named the Wess--Zumino Lagrangian, one after and one before eliminating \(F,G\). The latter has \(\bar\Psi\gamma^\beta\partial_\mu\Psi\) and a fermion mass coefficient differing by a factor 2 from the former.}
\GSfield{Statement}{One consistent modern convention uses metric \((+---)\), a complex scalar \(\phi\), Weyl fermion \(\chi\), complex auxiliary field \(\mathcal F\), and holomorphic \(W(\phi)=m\phi^2/2+g\phi^3/3\): \[\mathcal L=\partial_\mu\bar\phi\,\partial^\mu\phi+i\bar\chi\bar\sigma^\mu\partial_\mu\chi+\bar{\mathcal F}\mathcal F+\bigl(\mathcal F W'(\phi)-\tfrac12W''(\phi)\chi\chi+\text{h.c.}\bigr).\] Here spinor contractions use the standard antisymmetric two-spinor pairing. This action has the intended chiral supersymmetry; it is not a term-for-term endorsement of the notebook's inconsistent real-component display.}
\GSfield{Reconstruction}{The complex normalization, opposite metric convention and \(W\) are expressly modern reconstruction. In the notebook, the kinetic indices must contract; moreover eliminating bosonic auxiliary fields cannot change an untouched fermion mass coefficient from \(-im\) to \(-im/2\).}
\GSfield{Fate}{The intended component construction is verified in \citet[\S 4(b), equations~4.12--4.23]{ogievetskymezincescu1975superfields}, already in the December 1975 review.}
\GSfield{Present status}{As of 10 November 2026: The standard model is established. The literal pair of notebook Lagrangians is inconsistent and should not be used as a checked supersymmetric action.}
\GSfield{Assessment}{literal statement: false as written; publication match: partial; verification: read the result; development: none; confidence: high}
\end{GSitem}
\begin{GSitem}{1976-404.8}{E}{Eliminating the auxiliary fields and auditing the potential}
\GSfield{Kind; source}{method; \GSpage{P000668-P000669}}
\GSfield{The notes}{Because no derivatives of \(F,G\) occur, their Euler equations are said to recover the preceding action. Its cubic scalar term is positive and its quartic coefficient is written ambiguously as \(g_2^2\).}
\GSfield{Statement}{Using exactly the real auxiliary terms printed on P000669, \[\mathcal L_{\rm aux}=\tfrac12(F^2+G^2)+m(FA+GB)+g\{F(A^2-B^2)+2GAB\},\] their equations uniquely give \(F=-mA-g(A^2-B^2)\), \(G=-mB-2gAB\). Substitution yields \[-\tfrac12m^2(A^2+B^2)-mgA(A^2+B^2)-\tfrac12g^2(A^2+B^2)^2.\] Thus the cubic sign in P000668 is inconsistent with P000669; \(g^2/2\) is the forced quartic coefficient in these conventions.}
\GSfield{Reconstruction}{\oursComputed\ Completing both squares independently verifies the substitution. The scan has no clear fraction bar at \(g_2^2\), so the quartic change is an evidence-based emendation, not a confident new paleographic reading.}
\GSfield{Fate}{Auxiliary elimination is part of the pre-existing construction \citep[\S 4(b), equations~4.16 and 4.23]{ogievetskymezincescu1975superfields}. The comparison here checks the notebook internally.}
\GSfield{Present status}{As of 10 November 2026: The algebraic elimination is exact. It diagnoses a faithful cubic-sign error and an ambiguous quartic notation; it does not repair all fermionic conventions automatically.}
\GSfield{Assessment}{literal statement: false as written; publication match: partial; verification: read the result; development: annotated; confidence: high}
\end{GSitem}
\begin{GSitem}{1976-404.9}{E}{Infinitesimal supersymmetry and invariance of the action}
\GSfield{Kind; source}{construction; \GSpage{P000669}}
\GSfield{The notes}{Five variations of \(A,B,\Psi,F,G\) are displayed, with \(\delta A=i\bar\alpha\Psi\), derivative terms in \(\delta\Psi\), and derivatives of \(\Psi\) in \(\delta F,\delta G\). The text says that the Lagrangian is invariant.}
\GSfield{Statement}{In the consistent complex convention of the previous occurrence, take a constant Grassmann-odd Weyl parameter \(\epsilon\). The transformations \(\delta\phi=\sqrt2\epsilon\chi\), \(\delta\chi=i\sqrt2\sigma^\mu\bar\epsilon\,\partial_\mu\phi+\sqrt2\epsilon\mathcal F\), \(\delta\mathcal F=i\sqrt2\bar\epsilon\bar\sigma^\mu\partial_\mu\chi\), and their conjugates close off shell on translations. The variation of \(\mathcal L\) is a total derivative, so the action is invariant when boundary terms vanish. It is not necessary for the density to be pointwise unchanged.}
\GSfield{Reconstruction}{Odd parameters, the use of an even variation operator and the total-derivative qualification are supplied. A term-by-term sign identification with the notebook's barred-Majorana conventions is not asserted.}
\GSfield{Fate}{\citet[\S 1, equation~1.1; \S 4(b)]{ogievetskymezincescu1975superfields} verifies this pre-talk off-shell symmetry and component interpretation.}
\GSfield{Present status}{As of 10 November 2026: Established classical supersymmetry. The named action and multiplet are identified exactly, while the notebook's unexplained sign conventions receive only a partial formula match.}
\GSfield{Assessment}{literal statement: correct with missing hypotheses; publication match: partial; verification: read the result; development: none; confidence: high}
\end{GSitem}

\GStalk{405}{1976-09-20}{A. A. Migdal}{P000669}{Why quarks suggest local colour gauge fields}
Migdal's short discussion begins after Novikov's variations on P000669. It moves from three-quark nucleons and quark--antiquark mesons to short-distance freedom, confinement, relativistic locality and the colour/flavour matrix. No QCD field equation is actually written. Its force-law and uncertainty-principle slogans are not mathematical confinement proofs. The matrix uses three light flavours and three colours; it should not be interpreted as nine independent flavours or as an exhaustive 1976 particle inventory. Bernstein's dated 18/X-76 lecture starts only on P000670 and is outside this record.

\begin{GSitem}{1976-405.1}{E}{Valence quarks in nucleons and mesons}
\GSfield{Kind; source}{claim; \GSpage{P000669}}
\GSfield{The notes}{A proton is drawn with three dots labelled ``3 quarks''; the last sentence says mesons are made from a quark and an antiquark.}
\GSfield{Statement}{In the minimal valence-quark model, a proton has flavour content \(uud\), a neutron \(udd\), and ordinary mesons have \(q\bar q\) content. Colour singlets occur through \(\Lambda^3\mathbb C^3\cong\mathbb C\) for three quarks and the contraction \(\mathbb C^3\otimes(\mathbb C^3)^*\to\mathbb C\) for a quark--antiquark pair. These are valence classifications, not assertions that interacting hadrons contain exactly this number of quantum excitations.}
\GSfield{Reconstruction}{The representation-theoretic singlet maps make the claim precise; gluons, sea pairs and nonminimal hadrons are not excluded.}
\GSfield{Fate}{The pre-talk constituent picture is recalled by \citet[\S I]{wilson1974confinement}; the modern qualification is explicit in \citet[\S 15.1]{amslercredegrand2024quark}. It continues 1975-377.1.}
\GSfield{Present status}{As of 10 November 2026: The valence terminology remains standard with the stated limits; it is not an exact Fock-space decomposition theorem.}
\GSfield{Assessment}{literal statement: correct with missing hypotheses; publication match: exact; verification: read the result; development: none; confidence: high}
\end{GSitem}
\begin{GSitem}{1976-405.2}{E}{Short-distance freedom seen through electromagnetic probes}
\GSfield{Kind; source}{claim; \GSpage{P000669}}
\GSfield{The notes}{Quarks are said to be ``free inside'' and visible when the proton is probed by gamma quanta.}
\GSfield{Statement}{The controlled theoretical statement is asymptotic freedom: for four-dimensional \(SU(3)\) gauge theory with \(N_f\) Dirac fermions in the fundamental representation and no additional scalar matter, \[\mu\frac{dg}{d\mu}=-\frac{11-2N_f/3}{16\pi^2}g^3+O(g^5).\] For \(N_f\le16\), the weak-coupling ultraviolet flow is toward \(g=0\). Deep-inelastic electromagnetic scattering probes short-distance charged constituents, with calculable scaling violations; it does not show exactly free quarks at all distances inside a hadron.}
\GSfield{Reconstruction}{The beta-function formula and regime are supplied as the precise interpretation of the physical shorthand, not as a computation performed at this talk.}
\GSfield{Fate}{\citet[equations~8--9]{grosswilczek1973ultraviolet} published the result 25 June 1973, received 27 April 1973, before the lecture.}
\GSfield{Present status}{As of 10 November 2026: Established perturbative ultraviolet behaviour. It does not by itself prove long-distance confinement.}
\GSfield{Assessment}{literal statement: correct with missing hypotheses; publication match: partial; verification: read the result; development: none; confidence: high}
\end{GSitem}
\begin{GSitem}{1976-405.3}{E}{Confinement as a pre-existing physical programme}
\GSfield{Kind; source}{claim; \GSpage{P000669}}
\GSfield{The notes}{The juxtaposed assertion is that quarks cannot be taken out of a nucleon and that the binding is associated with colour.}
\GSfield{Statement}{The intended physical hypothesis is absence of isolated colour-charged asymptotic particles, while observable hadron states are colour singlets. A precise theorem would additionally require a defined quantum gauge theory, its physical Hilbert space and a specified confinement criterion. Those data are not supplied by this page. Cutoff-lattice strong-coupling confinement and continuum QCD confinement are different statements.}
\GSfield{Reconstruction}{This records the prior physical programme without replacing the notes by a new request to solve the continuum Yang--Mills mass-gap problem.}
\GSfield{Fate}{\citet[abstract and \S\S I--II]{wilson1974confinement} had already formulated a gauge-theoretic confinement mechanism and its lattice-cutoff limitation. \citet[\S 15.1]{amslercredegrand2024quark} still phrases general QCD confinement as a belief.}
\GSfield{Present status}{As of 10 November 2026: Class E marks pre-talk exposition, not a completed rigorous continuum proof. The notebook provides no unique additional unresolved mathematical scope to index as a new C/D problem.}
\GSfield{Assessment}{literal statement: not determinable; publication match: partial; verification: read the result; development: none; confidence: medium}
\end{GSitem}
\begin{GSitem}{1976-405.4}{E}{A growing potential is not necessarily a growing force}
\GSfield{Kind; source}{claim; \GSpage{P000669}}
\GSfield{The notes}{The emphatic sentence is:\par \ru{В теории силы должны расти с расстояниями!}}
\GSfield{Statement}{In a static-source model, a confining potential \(V(R)\sim\sigma R\), \(\sigma>0\), has force magnitude \(|V'(R)|\sim\sigma\), not a force diverging with separation. Formally an area law \(\langle W(R,T)\rangle\sim e^{-\sigma RT}\) yields this linear potential via \(V(R)=-\lim_{T\to\infty}T^{-1}\log\langle W(R,T)\rangle\), if that limit exists. With dynamical light quarks the asymptotic static-source problem also involves string breaking.}
\GSfield{Reconstruction}{\ours\ Differentiating a linear potential suffices to refute the necessity of increasing force. The Wilson-loop limit is an explanatory reconstruction with its existence assumption explicit.}
\GSfield{Fate}{The area-law mechanism is already in \citet[\S II and \S IV]{wilson1974confinement}; modern model and string-breaking qualifications appear in \citet[\S\S 15.7--15.8]{amslercredegrand2024quark}.}
\GSfield{Present status}{As of 10 November 2026: The literal force requirement is false. The appropriate statement concerns energy cost and the specified static-source regime.}
\GSfield{Assessment}{literal statement: false as written; publication match: partial; verification: read the result; development: annotated; confidence: high}
\end{GSitem}
\begin{GSitem}{1976-405.5}{F}{The uncertainty-principle slogan does not specify a stability theorem}
\GSfield{Kind; source}{remark; \GSpage{P000669}}
\GSfield{The notes}{A short sentence says that by the uncertainty principle the quarks cannot merge.}
\GSfield{Statement}{An uncertainty relation such as \(\Delta x\,\Delta p\ge\hbar/2\) gives a kinetic localization cost only after a state and Hamiltonian are specified. Whether that cost prevents collapse depends on the attractive interaction, particle statistics and relativistic regime. None of these quantitative data occurs in the sentence.}
\GSfield{Reconstruction}{No particular many-body stability inequality or confining Hamiltonian is inferred.}
\GSfield{Fate}{The exact stability assertion could not be identified in the searched confinement and quark-model sources; \citet[\S 15.7]{amslercredegrand2024quark} discusses models but cannot supply the missing Hamiltonian.}
\GSfield{Present status}{As of 10 November 2026: Unidentified as a mathematical result. The sentence is physical intuition, not an independent proof of confinement or noncollapse.}
\GSfield{Assessment}{literal statement: not determinable; publication match: thematic only; verification: unresolved lead; development: none; confidence: low}
\end{GSitem}
\begin{GSitem}{1976-405.6}{E}{Local Lorentz-covariant gauge dynamics instead of a prescribed force law}
\GSfield{Kind; source}{method; \GSpage{P000669}}
\GSfield{The notes}{Relativistic invariance is said to obstruct a direct potential law; pointlike probing is offered as evidence for locality, and the proposed response is to construct Lagrangians and gauge fields.}
\GSfield{Statement}{A local classical model on Minkowski space uses an \(SU(3)\) connection \(A_\mu\), its curvature \(F_{\mu\nu}\), and Dirac fields \(q_f\) in the colour fundamental representation, with \(\mathcal L=-\tfrac14F^a_{\mu\nu}F^{a\mu\nu}+\sum_f\bar q_f(i\gamma^\mu D_\mu-m_f)q_f\). This is local, Lorentz covariant and invariant under smooth local colour changes. Small observed constituent size is evidence for a useful local-field model; it is not a proof of the exact quantum microcausality axiom.}
\GSfield{Reconstruction}{The explicit Lagrangian is a labelled modern reconstruction of the proposed method; no action formula appears in Migdal's short notes.}
\GSfield{Fate}{The covariant-derivative and curvature construction predates QCD: \citet[equations~2--11]{yangmills1954conservation}. Its colour application is already proposed in \citet[p.1345]{grosswilczek1973ultraviolet}.}
\GSfield{Present status}{As of 10 November 2026: Established classical gauge-theory framework. Defining its full nonperturbative continuum quantum dynamics is a separate task not formulated here.}
\GSfield{Assessment}{literal statement: correct with missing hypotheses; publication match: partial; verification: read the result; development: none; confidence: high}
\end{GSitem}
\begin{GSitem}{1976-405.7}{E}{Colour and flavour are two different indices}
\GSfield{Kind; source}{construction; \GSpage{P000669}}
\GSfield{The notes}{The \(3\times3\) array has columns \(p,n,\lambda\) and rows 1,2,3. The vertical group is \(SU(3)\) colour, the horizontal flavour group is annotated ``unknown'', and the text says there are nine quarks.}
\GSfield{Statement}{For the three light flavours, write \(q^a_f\), \(a=1,2,3\), \(f=u,d,s\). This gives nine colour--flavour components before spinor indices; it does not give nine flavours. A local colour matrix acts on \(a\), while a global flavour matrix acts on \(f\). With equal light-quark masses and electromagnetism omitted, vector \(SU(3)\) flavour is a symmetry; unequal masses break it. The full modern quark inventory has six flavours, each with three colours.}
\GSfield{Reconstruction}{The traditional \(p,n,\lambda\) labels are interpreted as the light-flavour triplet, not physical proton/neutron/Lambda fields in the same gauge multiplet. The ``unknown'' annotation does not specify a unique unresolved group-classification problem.}
\GSfield{Fate}{\citet[p.1345]{grosswilczek1973ultraviolet} already separates colour from global flavour symmetries. \citet[\S\S 15.1--15.2,15.5.1]{amslercredegrand2024quark} gives modern terminology. The seminar had already discussed four flavours and three colours in 1975-367.5.}
\GSfield{Present status}{As of 10 November 2026: The nine-component light-flavour model is identifiable. Treating it as a complete inventory in 1976 would contradict the earlier seminar's four-flavour discussion; the shorthand is therefore not silently expanded into such a claim.}
\GSfield{Assessment}{literal statement: correct with missing hypotheses; publication match: partial; verification: read the result; development: none; confidence: high}
\end{GSitem}
\begin{GSitem}{1976-405.8}{E}{Local colour gauge symmetry}
\GSfield{Kind; source}{definition; \GSpage{P000669}}
\GSfield{The notes}{The last conceptual distinction says that the colour group corresponds to a much stronger symmetry, namely gauge symmetry.}
\GSfield{Statement}{The gauge group consists locally of smooth maps \(U:M\to SU(3)\), acting by \(q_f\mapsto Uq_f\) and, for \(D=d+A\), \(A\mapsto UAU^{-1}-(dU)U^{-1}\). Then \(F\mapsto UFU^{-1}\). Constant colour rotations form a subgroup of this local transformation group. ``Stronger'' means freedom to choose the colour frame independently at each spacetime point, with the connection compensating derivatives.}
\GSfield{Reconstruction}{This specifies the mathematical content of the comparison; gauge transformations describe redundancy of field variables and are not simply an extra independent global flavour action.}
\GSfield{Fate}{\citet[isotopic gauge transformation, equations~1--5]{yangmills1954conservation} gives the local construction in 1954; the colour specialization predates the talk.}
\GSfield{Present status}{As of 10 November 2026: Established local gauge symmetry. No new gauge-group extension problem is posed by the sentence.}
\GSfield{Assessment}{literal statement: correct; publication match: exact; verification: read the result; development: none; confidence: high}
\end{GSitem}

\GStalk{406}{1976-10-18}{I. N. Bernstein}{P000670-P000674}{Superdomains, Berezin integration, and Stokes boundaries}
The dated continuation begins on P000670, after a separate question about factors. Bernstein explains superfunctions, parity-preserving morphisms, nilpotent Taylor expansion, and why a boundary requires an embedded subsupermanifold in addition to its ordinary body. The circle and half-line examples lead to the Stokes problem, soon printed by Bernstein and Leites. P000673--P000674 explicitly form a retrospective summary of earlier super talks copied from I.~Dorfman's notebook; their definitions and linear algebra are indexed here with that provenance, rather than assigned a new seminar date. The catalogue's P000668--P000675 span also includes adjoining talks.

\begin{GSitem}{1976-406.1}{F}{Factors and matrices of continuous order}
\GSfield{Kind; source}{problem; \GSpage{P000670}}
\GSfield{The notes}{Before Bernstein's heading the page asks: \ru{записать \(\mathrm{II}_1\)-факторы через матрицы непр. порядка}, followed by a doubtful \ru{не сделанная}. No attribution, definition, or proposed equivalence is supplied.}
\GSfield{Statement}{The surviving wording requests some representation of type \(\mathrm{II}_1\) factors by ``matrices of continuous order.'' It does not specify the matrix object, a topology, an approximation property, or whether every factor is intended. No determinate modern theorem or conjecture can therefore be extracted.}
\GSfield{Reconstruction}{``Continuous order'' is retained as an unidentified expression. Interpreting it as finite-dimensional microstate approximation, an embedding into an ultrapower, or continuous geometry would add a substantive condition absent from the page.}
\GSfield{Fate}{Searches for the Russian phrase, continuous-order matrices, factors, and Gelfand did not identify an exact publication. The connection with known matrix approximation problems is only a possible theme and is not used as a solution claim.}
\GSfield{Present status}{As of November 2026, the exact question remains unidentified; this is an evidence limitation, not an assertion that a specified operator-algebra problem is open.}
\GSfield{Assessment}{literal statement: not determinable; publication match: none found; verification: unresolved lead; development: none; confidence: low}
\end{GSitem}
\begin{GSitem}{1976-406.2}{E}{Supercommutative functions and their parity}
\GSfield{Kind; source}{definition; \GSpage{P000670-P000671; P000673-P000674}}
\GSfield{The notes}{The notes write \(C^\infty(\mathbb R^p)\otimes\Lambda(\xi_1,\ldots,\xi_q)\), expand in monomials with exponents zero or one, and set parity equal to total odd degree modulo two. The retrospective page gives \(A=A_0\oplus A_1\) and \(ab=(-1)^{p(a)p(b)}ba\).}
\GSfield{Statement}{For an open \(U\subset\mathbb R^p\), set \(\mathcal O(U)=C^\infty(U;\mathbb R)\otimes\Lambda(\mathbb R^q)\). Every element has a unique expansion \(f=\sum_{I\subseteq\{1,\ldots,q\}}f_I(u)\xi^I\), with increasing indices in \(\xi^I\). For homogeneous \(a,b\), \(ab=(-1)^{|a||b|}ba\); the grading is by \(|I|\bmod2\), not by polynomial degree in the even variables.}
\GSfield{Reconstruction}{Homogeneity in the sign rule and the increasing order of exterior generators are explicit here. The final exponent \(\alpha_k\) on P000670 is a faithful index slip for \(\alpha_q\). Earlier-index continuity: 1975-389.1. These links identify a recurring statement or construction; only the explicitly indexed inverse problem reuses a problem ID.}
\GSfield{Fate}{Already printed in \citet[\S2, p.~505]{berezinleites1975supermanifolds}; the same definitions appear in \citet[\S\S1--2]{bernsteinleites1977integral}. The 1975 paper was received on 24 March 1975.}
\GSfield{Present status}{As of November 2026, this is the standard local algebra of a smooth supermanifold.}
\GSfield{Assessment}{literal statement: correct with missing hypotheses; publication match: exact; verification: read the result; development: none; confidence: high}
\end{GSitem}
\begin{GSitem}{1976-406.3}{E}{Superdomains as locally ringed spaces}
\GSfield{Kind; source}{definition; \GSpage{P000670; P000674}}
\GSfield{The notes}{A superdomain is an open set \(U\subset\mathbb R^p\) with the sheaf \(V\mapsto C^\infty(V)\otimes\Lambda(\xi_1,\ldots,\xi_q)\). The sketches illustrate an ordinary interval carrying this sheaf and an underlying one-point space.}
\GSfield{Statement}{A smooth superdomain of dimension \(p|q\) is \((U,\mathcal O_U)\), where \(U\subset\mathbb R^p\) is open and \(\mathcal O_U=C^\infty_U\otimes\Lambda(\mathbb R^q)\). A smooth supermanifold is a locally ringed space locally isomorphic to one of these models; its body is obtained by quotienting the nilpotent ideal generated by the odd functions.}
\GSfield{Reconstruction}{The local-ring condition and the body quotient are supplied to make precise the ringed-space shorthand. The notes only use the local split model; they do not assert a canonical global splitting. Earlier-index continuity: 1975-389.1. These links identify a recurring statement or construction; only the explicitly indexed inverse problem reuses a problem ID.}
\GSfield{Fate}{This is the construction of \citet[Definition~1, pp.~505--506]{berezinleites1975supermanifolds}, before this talk; compare the compact formulation in \citet[\S\S2, 8]{bernsteinleites1977integral}.}
\GSfield{Present status}{As of November 2026, the local ringed-space formulation is established and carries no unresolved classification claim here.}
\GSfield{Assessment}{literal statement: correct with missing hypotheses; publication match: exact; verification: read the result; development: none; confidence: high}
\end{GSitem}
\begin{GSitem}{1976-406.4}{E}{Morphisms determined by coordinate pullbacks}
\GSfield{Kind; source}{theorem; \GSpage{P000670; P000674}}
\GSfield{The notes}{A map covers a map of the bodies and induces a unital algebra homomorphism in the reverse direction preserving parity. The retrospective page says that it suffices to specify \(v_i^*,\eta_j^*\), and gives \(v=u+\xi_1\xi_2,\eta_1=u\xi_1,\eta_2=\xi_2\).}
\GSfield{Statement}{A morphism \(\Phi:U^{p|q}\to V^{m|n}\) is a body map together with a parity-preserving local unital homomorphism of sheaves \(\Phi^\sharp:\mathcal O_V\to\Phi_*\mathcal O_U\). It is uniquely specified by even sections \(a_1,\ldots,a_m\) and odd sections \(\beta_1,\ldots,\beta_n\), provided the body tuple \((a_{1,0},\ldots,a_{m,0})\) takes values in \(V\). Smooth functions of the \(a_i\) are evaluated by their finite nilpotent Taylor expansion.}
\GSfield{Reconstruction}{The body-image condition and the sheaf/locality requirement are essential additions. The example with \(\eta_1=u\xi_1\) is a morphism but is not an invertible coordinate change where \(u=0\). Earlier-index continuity: 1975-389.3. These links identify a recurring statement or construction; only the explicitly indexed inverse problem reuses a problem ID.}
\GSfield{Fate}{Morphisms were defined in \citet[Definition~3 and coordinate changes, p.~506]{berezinleites1975supermanifolds}; the exact coordinate determination is stated in \citet[\S3]{bernsteinleites1977integral}.}
\GSfield{Present status}{As of November 2026, the coordinate theorem is standard; invertibility requires an invertible body differential and invertible odd linear block.}
\GSfield{Assessment}{literal statement: correct with missing hypotheses; publication match: exact; verification: read the result; development: none; confidence: high}
\end{GSitem}
\begin{GSitem}{1976-406.5}{E}{Endomorphisms of a point with two odd directions}
\GSfield{Kind; source}{construction; \GSpage{P000670}}
\GSfield{The notes}{For \(p=0,q=2\) the underlying space is a point. The allowed images of \(\xi_1,\xi_2\) are arbitrary linear combinations of the odd generators; a term proportional to \(\xi_1\xi_2\) is crossed out because parity must be preserved.}
\GSfield{Statement}{Every parity-preserving unital \(\mathbb R\)-algebra endomorphism of \(\Lambda(\xi_1,\xi_2)\) is determined by \(\xi_1\mapsto a\xi_1+b\xi_2\), \(\xi_2\mapsto c\xi_1+d\xi_2\). Any real \(a,b,c,d\) are allowed; the endomorphism is an automorphism exactly when \(ad-bc\ne0\). A quadratic term cannot occur in the image of an odd generator.}
\GSfield{Reconstruction}{The automorphism criterion is an elementary consequence, not a claim explicitly written on the page. It clarifies the distinction between an endomorphism and a coordinate change.}
\GSfield{Fate}{This is the \(0|2\) instance of the pre-existing coordinate-morphism formalism in \citet[Definition~3]{berezinleites1975supermanifolds}; compare \citet[\S3]{bernsteinleites1977integral}.}
\GSfield{Present status}{As of November 2026, this finite exterior-algebra example is completely determined by its odd degree-one part.}
\GSfield{Assessment}{literal statement: correct; publication match: exact; verification: read the result; development: none; confidence: high}
\end{GSitem}
\begin{GSitem}{1976-406.6}{E}{The nilpotent Taylor formula for a smooth pullback}
\GSfield{Kind; source}{question; \GSpage{P000671; P000674}}
\GSfield{The notes}{The classroom question is \ru{куда переходит \(\sin v\)?} for \(v=u+\xi_1\xi_2\). The next display supplies \(g_{00}(u+\xi_1\xi_2)=g_{00}(u)+\xi_1\xi_2g'_{00}(u)\).}
\GSfield{Statement}{If \(\varepsilon=\xi_1\xi_2\), then \(\varepsilon^2=0\) and \(g(u+\varepsilon)=g(u)+\varepsilon g'(u)\) for every smooth \(g\). Thus \(\Phi^*(\sin v)=\sin u+\xi_1\xi_2\cos u\). More generally, if even nilpotents \(n_i\) generate a nilpotent ideal, \(g(u+n)=\sum_\alpha (\partial^\alpha g)(u)n^\alpha/\alpha!\), a finite sum.}
\GSfield{Reconstruction}{The sine answer and the multi-variable version spell out the Taylor rule rather than assume analyticity of \(g\). This repeats the earlier seminar item 1975-389.7, not a new unsolved problem. Earlier-index continuity: 1975-389.7. These links identify a recurring statement or construction; only the explicitly indexed inverse problem reuses a problem ID.}
\GSfield{Fate}{The construction is implicit in the 1975 coordinate formalism and explicit in \citet[\S3]{bernsteinleites1977integral}; the present notes themselves answer the posed question.}
\GSfield{Present status}{As of November 2026, the question has the displayed exact algebraic answer for smooth functions.}
\GSfield{Assessment}{literal statement: correct; publication match: exact; verification: read the result; development: none; confidence: high}
\end{GSitem}
\begin{GSitem}{1976-406.7}{E}{Even and left odd partial derivatives}
\GSfield{Kind; source}{definition; \GSpage{P000671; P000673-P000674}}
\GSfield{The notes}{Even derivatives differentiate coefficients. The odd rules include \(\partial_{\xi_1}(\xi_1\xi_2)=\xi_2\), \(\partial_{\xi_2}(\xi_1\xi_2)=-\xi_1\), and the sign from the number of preceding generators.}
\GSfield{Statement}{On the ordered monomial \(f(u)\xi_1^{\alpha_1}\cdots\xi_q^{\alpha_q}\), \(\partial_{u_i}\) acts on \(f\), whereas the left derivative is \[\partial_{\xi_j}(f\xi^\alpha)=\alpha_j(-1)^{\alpha_1+\cdots+\alpha_{j-1}}f\xi_1^{\alpha_1}\cdots\widehat{\xi_j^{\alpha_j}}\cdots\xi_q^{\alpha_q}.\] It satisfies \(\partial_{\xi_j}(ab)=(\partial_{\xi_j}a)b+(-1)^{|a|}a\partial_{\xi_j}b\) for homogeneous \(a\).}
\GSfield{Reconstruction}{The left-derivative convention is made explicit; switching to right derivatives changes matrix-coordinate conventions and must be done consistently.}
\GSfield{Fate}{These pre-existing Grassmann differentiation rules are written explicitly in \citet[\S4]{bernsteinleites1977integral}; the earlier density formula in \citet[\S3]{berezinleites1975supermanifolds} instead labels a right-derivative convention.}
\GSfield{Present status}{As of November 2026, the formulas are standard once the left/right convention is fixed.}
\GSfield{Assessment}{literal statement: correct; publication match: exact; verification: read the result; development: none; confidence: high}
\end{GSitem}
\begin{GSitem}{1976-406.8}{E}{The Berezinian and the reciprocal formula in the recap}
\GSfield{Kind; source}{theorem; \GSpage{P000673}}
\GSfield{The notes}{The first formula gives \(\det(A_{11})^{-1}\det(A_{22}-A_{21}A_{11}^{-1}A_{12})\). Lower down the same page gives \(\operatorname{Ber}J=\det(D)^{-1}\det(A-BD^{-1}C)\) and multiplicativity.}
\GSfield{Statement}{For an invertible even block matrix \(J=\left(\begin{smallmatrix}A&B\\C&D\end{smallmatrix}\right)\) over a supercommutative real or complex algebra, with invertible even diagonal blocks, \[\operatorname{Ber}J=\frac{\det(A-BD^{-1}C)}{\det D}=\frac{\det A}{\det(D-CA^{-1}B)}.\] It is an even unit and \(\operatorname{Ber}(J_1J_2)=\operatorname{Ber}J_1\operatorname{Ber}J_2\).}
\GSfield{Reconstruction}{The first notebook expression is the inverse of the usual Berezinian. This is a faithful mathematical error, not an OCR correction: on a diagonal \(1|1\) matrix \(\operatorname{diag}(a,d)\), it gives \(d/a\), whereas the lower formula gives \(a/d\). The independent check is recorded in the critical annotations.}
\GSfield{Fate}{The density multiplier already appears in \citet[formula~(3)]{berezinleites1975supermanifolds}; \citet[\S1]{bernsteinleites1977integral} states the Berezinian formula and multiplicativity and credits Berezin's 1971 formula. Thus this is exposition with an internal reciprocal inconsistency.}
\GSfield{Present status}{As of November 2026, the corrected Berezinian formula is established; the two literal notebook definitions cannot be used simultaneously.}
\GSfield{Assessment}{literal statement: false as written; publication match: exact; verification: read the result; development: annotated; confidence: high}
\end{GSitem}
\begin{GSitem}{1976-406.9}{E}{Supertrace, logarithm, and supertranspose}
\GSfield{Kind; source}{claim; \GSpage{P000673}}
\GSfield{The notes}{The recap writes \(\operatorname{str}J=\operatorname{tr}A-\operatorname{tr}D\), \(\det J=\exp\operatorname{str}\log J\), ordinary cyclicity of supertrace, and \(\left(\begin{smallmatrix}A&B\\C&D\end{smallmatrix}\right)^{*}=\left(\begin{smallmatrix}A^*&C^*\\-B^*&D^*\end{smallmatrix}\right)\).}
\GSfield{Statement}{For an even matrix \(J\), the supertrace is \(\operatorname{tr}A-\operatorname{tr}D\). Whenever a compatible formal or analytic logarithm exists, \(\operatorname{Ber}J=\exp(\operatorname{str}\log J)\); equivalently \(\operatorname{Ber}(e^X)=e^{\operatorname{str}X}\) for even \(X\). For homogeneous endomorphisms \(X,Y\), \(\operatorname{str}(XY)=(-1)^{|X||Y|}\operatorname{str}(YX)\). With the stated even-matrix convention the supertranspose is \(J^{st}=\left(\begin{smallmatrix}A^t&C^t\\-B^t&D^t\end{smallmatrix}\right)\).}
\GSfield{Reconstruction}{The ungraded cyclic identity holds for the even matrices under discussion; it would be false without that restriction. The logarithm formula uses the lower, usual Berezinian, not the first reciprocal determinant. The star is read as transpose, not complex adjoint; the page gives no conjugation convention.}
\GSfield{Fate}{These are algebraic identities belonging to the same older Berezinian formalism. The opened \citet[\S1]{bernsteinleites1977integral} verifies the determinant normalization, but does not itself print all the supertrace and transpose assertions: the publication match for this combined recital is partial, and the identities have been checked algebraically here.}
\GSfield{Present status}{As of November 2026, the identities hold with the displayed parity and logarithm restrictions.}
\GSfield{Assessment}{literal statement: correct with missing hypotheses; publication match: partial; verification: read the result; development: annotated; confidence: high}
\end{GSitem}
\begin{GSitem}{1976-406.10}{E}{Berezinian volume elements rather than a top ordinary form}
\GSfield{Kind; source}{definition; \GSpage{P000671; P000674}}
\GSfield{The notes}{The notes say that one integrates densities transforming by the Berezinian, and define \(\rho=fD(x)\), \(D(y)/D(x)=\operatorname{Ber}J\). They contrast these with ordinary differential forms.}
\GSfield{Statement}{For an oriented smooth \(p|q\) supermanifold, the Berezinian line has local generator \(D(u,\xi)\) and transition law \(D(y)=\operatorname{Ber}(\partial y/\partial x)D(x)\), with coefficients pulled back to a common chart. Thus a section satisfies \(f_x=(f_y\circ\Phi)\operatorname{Ber}D\Phi\). For \(q>0\), the usual super differential-form algebra has no maximal form degree, since the differentials of odd coordinates commute in exterior degree. A Berezinian section is therefore not a highest ordinary differential form.}
\GSfield{Reconstruction}{Orientation and coefficient pullback are explicit here. The statement ``forms cannot be integrated'' is read as the lack of this canonical top-form construction, not a prohibition of every integration operation on superforms.}
\GSfield{Fate}{The density law is already \citet[Definition~5 and formula~(3)]{berezinleites1975supermanifolds}; \citet[\S\S5--7]{bernsteinleites1977integral} distinguishes volume, differential, and integral forms.}
\GSfield{Present status}{As of November 2026, the distinction is standard and is essential to the Stokes formulation below.}
\GSfield{Assessment}{literal statement: correct with missing hypotheses; publication match: exact; verification: read the result; development: none; confidence: high}
\end{GSitem}
\begin{GSitem}{1976-406.11}{E}{Compactly supported Berezin integration}
\GSfield{Kind; source}{theorem; \GSpage{P000672; P000674}}
\GSfield{The notes}{The integral is the ordinary integral of the coefficient of \(\xi_1\cdots\xi_q\). The recap attributes coordinate independence to Berezin and explicitly requires compact support.}
\GSfield{Statement}{Fix an orientation of \(U\subset\mathbb R^p\) and the odd convention \(\int\xi_1\cdots\xi_q\,D\xi=1\). For a compactly supported Berezinian section \(\rho=fD(u,\xi)\), define \(\int_{U^{p|q}}\rho=\int_U f_{\{1,\ldots,q\}}(u)\,du\). This is invariant under orientation-compatible super-coordinate changes and extends by partitions of unity to oriented smooth supermanifolds.}
\GSfield{Reconstruction}{The fixed odd ordering removes the apparent discrepancy between different conventions for \(d\xi_1d\xi_2\). Compact support is inside the open body; allowing support to reach a missing boundary invalidates the unqualified change-of-variables statement. Earlier-index continuity: 1975-389.2. These links identify a recurring statement or construction; only the explicitly indexed inverse problem reuses a problem ID.}
\GSfield{Fate}{This is \citet[Theorem~1]{berezinleites1975supermanifolds}, published in 1975. \citet[\S5]{bernsteinleites1977integral} restates it and gives Rudakov's boundary counterexample when the support hypothesis is removed.}
\GSfield{Present status}{As of November 2026, the theorem is established; boundary integrals require the extra geometric data in the following items.}
\GSfield{Assessment}{literal statement: correct; publication match: exact; verification: read the result; development: none; confidence: high}
\end{GSitem}
\begin{GSitem}{1976-406.12}{E}{Regular ideals and embedded even boundaries}
\GSfield{Kind; source}{definition; \GSpage{P000671-P000672}}
\GSfield{The notes}{A subsupermanifold is described by an ideal of functions; a boundary is cut out by an even superfunction \(f=0\). The notes emphasize that it carries the structure of a subsupermanifold.}
\GSfield{Statement}{An embedded codimension \(r|s\) subsupermanifold is locally given by \(r\) even and \(s\) odd functions whose body Jacobian has the corresponding full even and odd ranks, so that they extend to a super-coordinate system. In particular an even \(f\) with \(df_0\ne0\) on \(f_0^{-1}(0)\) defines a codimension \(1|0\) boundary. An arbitrary ideal defines a ringed subspace, not necessarily a subsupermanifold.}
\GSfield{Reconstruction}{The necessary rank condition is supplied; it cannot be omitted from the ideal description.}
\GSfield{Fate}{The regular-equation result predates the talk: \citet[Theorem~2]{berezinleites1975supermanifolds}. The boundary-specific formulation is \citet[\S\S8, 11]{bernsteinleites1977integral}.}
\GSfield{Present status}{As of November 2026, regularity separates smooth embedded boundaries from singular quotient spaces.}
\GSfield{Assessment}{literal statement: correct with missing hypotheses; publication match: exact; verification: read the result; development: none; confidence: high}
\end{GSitem}
\begin{GSitem}{1976-406.13}{A}{A closed superdomain includes a choice of superboundary}
\GSfield{Kind; source}{construction; \GSpage{P000671-P000672}}
\GSfield{The notes}{The same body \([0,\infty)\) has boundaries \(u=0\) and \(u+c\xi_1\xi_2=0\). The disc similarly has the two boundaries \(u_1^2+u_2^2-1=0\) and \(u_1^2+u_2^2-1-\xi_1\xi_2=0\).}
\GSfield{Statement}{A closed superdomain consists of a closed region \(F\) with smooth body boundary and an embedded \((p-1)|q\) subsupermanifold with body \(\partial F\). Distinct ideals \((u)\) and \((u+c\xi_1\xi_2)\) on the same ambient \(1|2\) superdomain give distinct such boundaries when \(c\ne0\). Berezin integration is defined in boundary-adapted charts; the body alone does not determine it.}
\GSfield{Reconstruction}{This is extra geometric boundary data, not a choice between two unequal ordinary point sets. With a compactly supported \(g\), the naive change \(v=u+c\xi_1\xi_2\) changes a top-coefficient half-line integral by \(c\int_0^\infty g'_0(u)du=-cg_0(0)\) if both coordinate half-lines are held fixed.}
\GSfield{Fate}{The exact closed-domain definition is \citet[\S11]{bernsteinleites1977integral}, submitted 6 October 1976 and printed in 1977. The submission date precedes this talk; a separate public-circulation date has not been established. Class A records publication chronology rather than priority of this seminar.}
\GSfield{Present status}{As of November 2026, the extra boundary embedding and the resulting boundary term are established.}
\GSfield{Assessment}{literal statement: correct with missing hypotheses; publication match: exact; verification: read the result; development: none; confidence: high}
\end{GSitem}
\begin{GSitem}{1976-406.14}{A}{Straightening the nilpotently displaced circle}
\GSfield{Kind; source}{claim; \GSpage{P000672}}
\GSfield{The notes}{To normalize \(u_1^2+u_2^2-1-\xi_1\xi_2=0\), the notebook writes \((u_1-\tfrac12\xi_1\xi_2)^2+u_2^2-1=0\). The scan confirms the missing denominator is not supplied in the handwriting.}
\GSfield{Statement}{Let \(\varepsilon=\xi_1\xi_2\). On the chart \(u_1\ne0\), the substitution \(v_1=u_1-\varepsilon/(2u_1)\), \(v_2=u_2\), with unchanged odd coordinates, gives exactly \(v_1^2+v_2^2-1=u_1^2+u_2^2-1-\varepsilon\). An analogous chart uses \(u_2\ne0\). Such charts cover the body circle.}
\GSfield{Reconstruction}{Direct expansion of the notebook substitution gives \(u_1^2+u_2^2-1-u_1\varepsilon\), so the written normalization is false. The local division by \(u_1\) is an editorial correction with its domain stated, not a claimed global coordinate formula.}
\GSfield{Fate}{The boundary-straightening framework is exactly the one used in \citet[\S11]{bernsteinleites1977integral}, published in 1977 after submission on 6 October 1976. That short paper does not print this particular circle substitution, so the publication match is partial.}
\GSfield{Present status}{As of November 2026, the local corrected substitution is exact because \(\varepsilon^2=0\); the literal notebook formula fails.}
\GSfield{Assessment}{literal statement: false as written; publication match: partial; verification: read the result; development: annotated; confidence: high}
\end{GSitem}
\begin{GSitem}{1976-406.15}{A}{A boundary-preserving half-line change of variables}
\GSfield{Kind; source}{claim; \GSpage{P000672}}
\GSfield{The notes}{For \(v=u+u\xi_1\xi_2\), the notes observe that \(u=0\) becomes \(v=0\), then begin a check using \(g_{11}\) and mark \(Dv/Du\). The displayed expansion does not yet include the transformed full superfunction and Berezinian.}
\GSfield{Statement}{Put \(\varepsilon=\xi_1\xi_2\), \(v=u(1+\varepsilon)\), and \(\eta_i=\xi_i\). For \(G(v,\eta)=g_0(v)+g_1(v)\eta_1+g_2(v)\eta_2+g_{12}(v)\eta_1\eta_2\), one has \(\operatorname{Ber}D\Phi=1+\varepsilon\). The top coefficient of \(\Phi^*(G D(v,\eta))\) is \(g_{12}(u)+(ug_0(u))'\). If the coefficients are smooth at zero and compactly supported toward infinity, its half-line integral is \(\int_0^\infty g_{12}(u)du\).}
\GSfield{Reconstruction}{The omitted lower coefficient and Jacobian are restored. The extra term integrates to \([ug_0(u)]_0^\infty=0\), which directly verifies the intended boundary-preserving invariance; expanding only \(g_{11}\) is not that proof.}
\GSfield{Fate}{This is the local model of \citet[\S11]{bernsteinleites1977integral}, received 6 October 1976 and printed in 1977; the explicit coefficient check here is a critical annotation, not a new theorem.}
\GSfield{Present status}{As of November 2026, the corrected calculation establishes the claimed invariance under its stated support hypotheses.}
\GSfield{Assessment}{literal statement: correct with missing hypotheses; publication match: partial; verification: read the result; development: annotated; confidence: high}
\end{GSitem}
\begin{GSitem}{1976-406.16}{A}{Stokes formula for a nilpotently displaced disc}
\GSfield{Kind; source}{problem; \GSpage{P000671-P000672}}
\GSfield{The notes}{The text asks \ru{Теорема Стокса --- ?}, \ru{Что интегр. по границе?}, and finally \ru{теор. Стокса для круга \(u_1^2+u_2^2+c\xi_1\xi_2-1=0\)}. It names Bernstein and Leites in connection with boundary integration.}
\GSfield{Statement}{Let \(M\) be an oriented smooth supermanifold of dimension \(p|q\), and let \(\mathcal D\) be a closed superdomain with embedded even boundary \(\partial\mathcal D\), oriented by the outward-normal convention. For a compactly supported integral form \(\omega\) of degree \(p-1\), \[\int_{\mathcal D}d\omega=\int_{\partial\mathcal D}\iota^*\omega.\] In particular this applies in \(\mathbb R^{2|2}\) to the disc with boundary ideal \((u_1^2+u_2^2+c\xi_1\xi_2-1)\), since its body differential is nonzero on the unit circle.}
\GSfield{Reconstruction}{Integral forms, their differential, and the boundary restriction map supply the missing objects. They are defined in the later publication; the notebook itself poses the problem and does not record a complete general Stokes theorem.}
\GSfield{Fate}{The exact answer is \citet[\S12, Stokes' formula, with \S\S7, 9, 11]{bernsteinleites1977integral}. The manuscript was submitted on 6 October 1976, before the 18 October seminar, and printed in 1977; no earlier public-circulation date has been verified.}
\GSfield{Present status}{As of November 2026, this boundary problem is solved in the stated integral-form category. It does not furnish a remaining class C or D research problem.}
\GSfield{Assessment}{literal statement: correct with missing hypotheses; publication match: exact; verification: read the result; development: none; confidence: high; problem \hyperref[prob:GP-1976-10]{GP-1976-10}}
\end{GSitem}
\begin{GSitem}{1976-406.17}{E}{Berezin and Marinov on classical Grassmann mechanics}
\GSfield{Kind; source}{remark; \GSpage{P000672}}
\GSfield{The notes}{A bibliographic aside, with a question mark after Marinov's name, refers to a work on classical mechanics in Fermi and Bose variables. No Hamiltonian, bracket, or motion equation is written.}
\GSfield{Statement}{The identifiable subject is classical dynamics on an algebra containing both commuting and anticommuting variables. The notes do not specify a model or a theorem, and no particular equations of spinning-particle mechanics are attributed to the speaker here.}
\GSfield{Reconstruction}{The 1975 announcement and 1976 report/1977 full paper are plausible bibliographic referents. Identification of the topic is stronger than identification of which version the speaker meant.}
\GSfield{Fate}{The author's institutional bibliography lists \citet{berezinmarinov1975classical}, before the talk. INSPIRE records \citet{berezinmarinov1977particle} as report ITEP-43-1976 and publication in 1977. Only bibliographic metadata was checked for this aside; no unrecorded mechanics result is inferred.}
\GSfield{Present status}{As of November 2026, the bibliographic topic is identified, while the exact version referred to and any intended mathematical proposition remain unspecified.}
\GSfield{Assessment}{literal statement: not determinable; publication match: thematic only; verification: indirect; development: none; confidence: medium}
\end{GSitem}

\GStalk{407}{1976-10-25}{V. N. Popov [catalogue; attribution conflict]}{P000675-P000677}{Connections, free Yang--Mills fields, and Ricci-flat metrics}
The notes begin on P000675 with Gelfand's request for an elementary explanation of a connection and a reference to Yang and Mills. They develop gauge covariance, the principal-bundle equation \(H\delta_h\Phi=0\), its independence of the auxiliary vertical metric, and the affine-frame realization of the vacuum Einstein equations. The last two results match D.~A.~Popov's 1975 paper and D.~A.~Popov--L.~I.~Daikhin's 1975 note exactly. The supplied catalogue calls the speaker V.~N.~Popov, while the handwritten heading gives only Popov; that conflict is retained explicitly. Mathematical identification of the papers does not prove a correction of the speaker's initials. The identified notes occupy P000675--P000677, narrower than the catalogue's P000673--P000679 packet.

\begin{GSitem}{1976-407.1}{E}{What is a connection?}
\GSfield{Kind; source}{question; \GSpage{P000675-P000676}; I. M. Gelfand}
\GSfield{The notes}{Gelfand asks \ru{Что такое связность? (для первокурсников и Полякова)} and lists Nomizu, Sternberg, Bishop--Crittenden, and Landau--Lifshitz. Later the talk represents a connection by a Lie-algebra-valued one-form \(\Omega\) on a principal bundle.}
\GSfield{Statement}{For a smooth right principal \(G\)-bundle \(\pi:P\to M\), a connection is a smooth right-equivariant horizontal subbundle \(H\subset TP\) with \(TP=H\oplus\ker d\pi\). Equivalently it is \(\Omega\in\Omega^1(P;\mathfrak g)\) such that \(\Omega(a^\#)=a\) for each fundamental vertical field and \(R_g^*\Omega=\operatorname{Ad}(g^{-1})\Omega\). Local potentials are \(s^*\Omega\) for local sections \(s\).}
\GSfield{Reconstruction}{The two connection axioms and their equivalent horizontal-space interpretation are supplied. Gelfand's request is pedagogical, not a new research question; no specific edition of the reading-list books is inferred.}
\GSfield{Fate}{The exact principal-connection interpretation is recalled in \citet[\S1, Definition~1]{popovdaikhin1975einstein} and in \citet[\S1]{popov1975yangmills}, both published before the seminar.}
\GSfield{Present status}{As of November 2026, the question is answered by the standard principal-connection definition.}
\GSfield{Assessment}{literal statement: correct with missing hypotheses; publication match: exact; verification: read the result; development: none; confidence: high}
\end{GSitem}
\begin{GSitem}{1976-407.2}{E}{The massless Dirac equation used to introduce gauge coupling}
\GSfield{Kind; source}{definition; \GSpage{P000675}}
\GSfield{The notes}{The opening equation is the massless Dirac equation for an ordinary four-component spinor. No interaction term is included in that displayed equation.}
\GSfield{Statement}{On four-dimensional Minkowski space with a fixed spin structure, choose matrices \(\gamma^\mu\) satisfying \(\gamma^\mu\gamma^\nu+\gamma^\nu\gamma^\mu=2g^{\mu\nu}I\). A free massless spinor satisfies \(i\gamma^\mu\partial_\mu\Psi=0\), equivalently the equation without the factor \(i\). A gauge coupling in a representation \(R\) replaces \(\partial_\mu\) by \(D_\mu=\partial_\mu+R_*(A_\mu)\).}
\GSfield{Reconstruction}{The Clifford relation and a definite signature convention are needed to interpret the written symbols. The covariant derivative is explanatory reconstruction; the notebook does not yet state a coupled Dirac--Yang--Mills system.}
\GSfield{Fate}{The free Dirac equation was classical long before 1976. The local-symmetry motivation for replacing the derivative is the construction in \citet[``Isotopic gauge transformation,'' pp.~192--193]{yangmills1954conservation}; the publication match is thematic for the notebook's isolated massless equation.}
\GSfield{Present status}{As of November 2026, this is established background, and the displayed free spinor equation does not by itself couple matter to the later gauge field.}
\GSfield{Assessment}{literal statement: correct with missing hypotheses; publication match: thematic only; verification: read the result; development: none; confidence: high}
\end{GSitem}
\begin{GSitem}{1976-407.3}{E}{Curvature and the source-free Yang--Mills equations}
\GSfield{Kind; source}{theorem; \GSpage{P000675}}
\GSfield{The notes}{The notes give \(\partial_\nu F_a^{\mu\nu}+f^b_{ac}F_b^{\mu\nu}A_\nu^c=0\), the nonlinear curvature formula, and the Killing form \(K_{ab}=f_{ac}^{d}f_{bd}^{c}\) for lowering Lie-algebra indices.}
\GSfield{Statement}{Let \((M,g)\) be pseudo-Riemannian, \(G\) semisimple, and \(A=A_\mu^aT_a dx^\mu\) a local principal connection, with \([T_b,T_c]=f^a_{bc}T_a\). Then \[F=dA+\tfrac12[A,A],\qquad F^a_{\mu\nu}=\partial_\mu A^a_\nu-\partial_\nu A^a_\mu+f^a_{bc}A^b_\mu A^c_\nu.\] The Euler--Lagrange equation of the quadratic curvature action, using a nondegenerate invariant bilinear form, is \(D_\mu F^{\mu\nu}=0\). On Minkowski space the covariant base derivative reduces to \(\partial_\mu\); Lie-algebra dual indices can be identified by the Killing form.}
\GSfield{Reconstruction}{The invariant tensor notation fixes the ambiguous handwritten structure-constant indices. The factor \(1/2\) in the antisymmetrized notebook component expression gives exactly the displayed commutator term. On a curved base one must include the Levi-Civita connection or the equivalent density divergence.}
\GSfield{Fate}{The gauge-field equations originate in \citet[pp.~193--194]{yangmills1954conservation}, received 28 June and published 1 October 1954. The geometric semisimple formulation is \citet[Proposition~3 and equation~(37)]{popov1975yangmills}; it was received in July 1974 and printed in 1975.}
\GSfield{Present status}{As of November 2026, these are the established classical source-free equations; no assertion about quantum existence, a mass gap, or confinement is made here.}
\GSfield{Assessment}{literal statement: correct with missing hypotheses; publication match: exact; verification: read the result; development: none; confidence: high}
\end{GSitem}
\begin{GSitem}{1976-407.4}{E}{Gelfand's reference on quarks}
\GSfield{Kind; source}{remark; \GSpage{P000675}; I. M. Gelfand}
\GSfield{The notes}{Gelfand refers readers to Zel'dovich's UFN article \ru{Кварки для пешеходов}.}
\GSfield{Statement}{This is a reading recommendation about the classification of elementary particles and the quark model. It contains no new mathematical claim in the seminar record.}
\GSfield{Reconstruction}{The abbreviated spoken title is identified with Zel'dovich's 1965 review, ``Classification of elementary particles and quarks `for pedestrians'.'' No result from that review is imported into the notebook.}
\GSfield{Fate}{Publisher and Math-Net metadata identify \citet{zeldovich1965quarks}, published in 1965, as the exact bibliographic referent.}
\GSfield{Present status}{As of November 2026, the reading reference is identified; it should be linked as background rather than indexed as an open problem.}
\GSfield{Assessment}{literal statement: correct; publication match: exact; verification: publisher metadata checked; development: none; confidence: high}
\end{GSitem}
\begin{GSitem}{1976-407.5}{E}{The distinction between a gauge field and quark matter}
\GSfield{Kind; source}{claim; \GSpage{P000675}}
\GSfield{The notes}{The notebook says \ru{Ур-ия Янга--Миллса описывают кварки} immediately after the source-free gauge equations and the earlier free Dirac equation.}
\GSfield{Statement}{The displayed unknown \(A\) is a Lie-algebra-valued gauge potential. Quark fields are spinor-valued matter fields carrying a color representation; a classical model with such matter requires a covariant Dirac equation and a current term \(D_\mu F^{\mu\nu}=j^\nu\). Consequently the source-free equation written here describes the gauge sector, not the full coupled quark dynamics.}
\GSfield{Reconstruction}{The phrase may be broad physical motivation, but as a literal identification of the displayed equations with quark fields it is false. Neither confinement nor a statement about individual observed quarks follows.}
\GSfield{Fate}{The original \citet[pp.~192--194]{yangmills1954conservation} distinguishes the gauge potential from the matter wavefunction and its current; the note's free-field terminology is also explicit in \citet[Definitions~1--3]{popovdaikhin1975einstein}. The correction concerns that pre-existing distinction.}
\GSfield{Present status}{As of November 2026, the literal field identification requires correction; no unresolved research assertion is extracted from the motivational sentence.}
\GSfield{Assessment}{literal statement: false as written; publication match: partial; verification: read the result; development: annotated; confidence: high}
\end{GSitem}
\begin{GSitem}{1976-407.6}{E}{The abelian reduction and gradient gauge transformations}
\GSfield{Kind; source}{theorem; \GSpage{P000675}}
\GSfield{The notes}{For \(G=U(1)\), the notes obtain Maxwell's equations and write \(A_\mu\mapsto A_\mu+\partial_\mu\alpha\). The antisymmetric matrix separates electric and magnetic components.}
\GSfield{Statement}{For an abelian connection, \(F=dA\) locally, so \(dF=0\) identically and the source-free Yang--Mills equation is \(d{*}F=0\). These are the vacuum Maxwell equations. A local gauge change \(A\mapsto A+d\alpha\) leaves \(F\) unchanged; globally \(A\) need not be a single one-form when the \(U(1)\)-bundle is nontrivial.}
\GSfield{Reconstruction}{The Bianchi equation supplies the homogeneous half of Maxwell's system. Metric raising of both indices is required: the formula with derivatives of \(A_\nu\) directly defines \(F_{\mu\nu}\), not \(F^{\mu\nu}\) in Minkowski coordinates. Electric/magnetic signs depend on the metric and time-coordinate convention.}
\GSfield{Fate}{The comparison with electromagnetism is explicit in \citet[pp.~192--193]{yangmills1954conservation} and \citet[\S1, after Definition~3]{popovdaikhin1975einstein}, both before this talk.}
\GSfield{Present status}{As of November 2026, the abelian reduction is exact, subject to the local-potential and index conventions stated here.}
\GSfield{Assessment}{literal statement: correct with missing hypotheses; publication match: exact; verification: read the result; development: none; confidence: high}
\end{GSitem}
\begin{GSitem}{1976-407.7}{E}{Nonabelian gauge covariance}
\GSfield{Kind; source}{theorem; \GSpage{P000676}}
\GSfield{The notes}{With \(B_\mu=T_aA_\mu^a\), the notes write \(B'_\mu=S^{-1}B^\mu S+S^{-1}\partial_\mu S\), \(F'_{\mu\nu}=S^{-1}F_{\mu\nu}S\), and assert invariance of the equations.}
\GSfield{Statement}{For a smooth \(S:U\to G\), fix \(A'=S^{-1}AS+S^{-1}dS\). Then \(F'=S^{-1}FS\) and \(D'({*}F')=S^{-1}D({*}F)S\). Therefore \(D{*}F=0\) is invariant under changes of local trivialization. In the component formula the same covariant index \(\mu\) must occur on both \(A'_\mu\) and \(A_\mu\).}
\GSfield{Reconstruction}{The isolated raised \(\mu\) in the conjugated potential is a faithful index error. The exterior equation fixes all local-sign conventions and explains precisely what is gauge invariant: the equation is covariant, while the curvature itself is conjugated.}
\GSfield{Fate}{These are \citet[equations~(3)--(5)]{yangmills1954conservation}, with coupling and Hermitian conventions converted to a Lie-algebra-valued connection, and \citet[equations~(5), (20)]{popov1975yangmills}.}
\GSfield{Present status}{As of November 2026, gauge covariance is established; the notebook requires only the indicated index correction.}
\GSfield{Assessment}{literal statement: false as written; publication match: exact; verification: read the result; development: annotated; confidence: high}
\end{GSitem}
\begin{GSitem}{1976-407.8}{E}{A metric on the total space of a principal bundle}
\GSfield{Kind; source}{construction; \GSpage{P000676}}
\GSfield{The notes}{Given \(g\) on \(M\), a connection \(\Omega\), and an arbitrary nondegenerate quadratic form \(K\) on \(\mathfrak g\), the notes set \(h=\pi^*g+K(\Omega,\Omega)\), observing that it is generally not invariant.}
\GSfield{Statement}{Let \(g\) be nondegenerate and symmetric on \(TM\), and \(K\) a fixed nondegenerate symmetric form on \(\mathfrak g\). The formula \[h_p(X,Y)=g_{\pi(p)}(d\pi X,d\pi Y)+K(\Omega X,\Omega Y)\] defines a smooth nondegenerate pseudo-Riemannian metric on \(P\): horizontal and vertical directions are orthogonal, with metrics \(g\) and \(K\). It is right \(G\)-invariant exactly when \(K\) is \(\operatorname{Ad}(G)\)-invariant.}
\GSfield{Reconstruction}{Symmetry is made explicit because the source calls \(K\) a quadratic form. Semisimplicity is not needed for this auxiliary metric construction; it is needed only if one insists on taking the Killing form.}
\GSfield{Fate}{The semisimple metric is in \citet[equation~(22)]{popov1975yangmills}; the arbitrary-group extension using an arbitrary nondegenerate form appears explicitly in \citet[\S1, before Lemma~1]{popovdaikhin1975einstein}, printed in 1975.}
\GSfield{Present status}{As of November 2026, the construction is valid without an invariant vertical metric, as the notes correctly emphasize.}
\GSfield{Assessment}{literal statement: correct; publication match: exact; verification: read the result; development: none; confidence: high}
\end{GSitem}
\begin{GSitem}{1976-407.9}{E}{Horizontal curvature and the structural identity}
\GSfield{Kind; source}{definition; \GSpage{P000676}}
\GSfield{The notes}{Curvature is defined by \(\Phi=H d\Omega=d\Omega+\tfrac12[\Omega,\Omega]\), where \(H\) takes the horizontal component. A horizontal form vanishes when one argument is vertical.}
\GSfield{Statement}{For a principal connection, \(\Phi=d\Omega+\tfrac12[\Omega\wedge\Omega]\) is horizontal and \(\operatorname{Ad}\)-equivariant. Hence it is equivalently a two-form on \(M\) valued in the adjoint bundle. Its covariant exterior derivative vanishes: \(D\Phi=0\), the Bianchi identity.}
\GSfield{Reconstruction}{The bracket includes the wedge product and the Lie bracket. The Bianchi identity is stated here to distinguish curvature geometry from the additional field equation; it is implicit in the later comparison with Maxwell's equations.}
\GSfield{Fate}{These standard connection identities are used in \citet[\S1]{popov1975yangmills} and restated in \citet[\S1]{popovdaikhin1975einstein}, both preceding the seminar.}
\GSfield{Present status}{As of November 2026, the curvature and Bianchi identities are established independently of any field equation.}
\GSfield{Assessment}{literal statement: correct; publication match: exact; verification: read the result; development: none; confidence: high}
\end{GSitem}
\begin{GSitem}{1976-407.10}{E}{The projected codifferential is independent of the vertical metric}
\GSfield{Kind; source}{theorem; \GSpage{P000676-P000677}}
\GSfield{The notes}{The notes propose \(H\delta_h\Phi=0\) and emphasize that it does not depend on \(K\). A parenthetical sentence makes the broader assertion for equivariant horizontal vector-valued forms.}
\GSfield{Statement}{For the metric \(h\) above and a principal-connection curvature \(\Phi\), \(H\delta_h\Phi\) is horizontal, \(\operatorname{Ad}\)-equivariant, and independent of the chosen nondegenerate \(K\). Thus it defines an intrinsic equation \(\mathcal D\Phi=0\) for any Lie group. For semisimple \(G\), this coincides with the Euler--Lagrange Yang--Mills equation formed using the Killing metric.}
\GSfield{Reconstruction}{The exact verified theorem is for the curvature form. The wider notebook parenthesis about arbitrary equivariant horizontal forms is not needed for the field equation and is not asserted here without specifying its representation and codifferential conventions. No invariant \(K\) is imposed on the curvature theorem: that would improperly weaken the result.}
\GSfield{Fate}{The semisimple equivalence is \citet[Proposition~3]{popov1975yangmills}. The arbitrary-group, arbitrary-\(K\) assertion is exactly \citet[Lemma~1 and Definition~3, p.~791]{popovdaikhin1975einstein}, received 26 June 1975 and published in 1975. It is therefore E, not a post-seminar development.}
\GSfield{Present status}{As of November 2026, the verified curvature equation is established for arbitrary Lie groups; the scope of the broader parenthetical claim is kept separate.}
\GSfield{Assessment}{literal statement: correct with missing hypotheses; publication match: partial; verification: read the result; development: none; confidence: high}
\end{GSitem}
\begin{GSitem}{1976-407.11}{E}{Projection of the bundle equation to a local gauge}
\GSfield{Kind; source}{theorem; \GSpage{P000677}}
\GSfield{The notes}{For a section \(\zeta\), the notes set \(\theta=\zeta^*\Omega\) and \(\zeta^*\Phi=T_aF^a_{\mu\nu}dx^\mu\wedge dx^\nu\), then write \(\delta_M\zeta^*\Phi+*^{-1}[\theta,*\zeta^*\Phi]=0\).}
\GSfield{Statement}{In a local section, \(A=s^*\Omega\), \(F=s^*\Phi=dA+\tfrac12[A,A]\), and the projected equation is \(D_A{*}F=0\), or \(D_{A,\mu}F^{\mu\nu}=0\). Equivalently it is a covariant codifferential equation \(\delta_A F=0\), with the signature-dependent Hodge sign included in the chosen definition of \(\delta_A\). With a sum over all ordered indices, \(F=\tfrac12T_aF^a_{\mu\nu}dx^\mu\wedge dx^\nu\).}
\GSfield{Reconstruction}{The notebook formula suppresses either a factor \(1/2\) or a restriction \(\mu<\nu\). The invariant exterior formulation avoids declaring its Hodge sign wrong without a convention. Existence of a bi-invariant metric is needed for the usual invariant-action interpretation, not for the arbitrary-group projected equation itself.}
\GSfield{Fate}{The local formula is \citet[equations~(36)--(37)]{popov1975yangmills} and \citet[equations~(2)--(3)]{popovdaikhin1975einstein}; both exact formulations were available in 1975.}
\GSfield{Present status}{As of November 2026, the local and bundle equations agree after fixing normalization and codifferential conventions.}
\GSfield{Assessment}{literal statement: correct with missing hypotheses; publication match: exact; verification: read the result; development: none; confidence: high}
\end{GSitem}
\begin{GSitem}{1976-407.12}{E}{The affine Cartan connection induced by a metric}
\GSfield{Kind; source}{construction; \GSpage{P000677}}
\GSfield{The notes}{For the affine frame bundle and \(G=\operatorname{Aff}(\mathbb R^4)\), the notes characterize \(\bar\Omega\) by \(i^*\bar\Omega=(\Omega^g,\chi)\), where \(\chi\) is the displacement or solder form.}
\GSfield{Statement}{For a pseudo-Riemannian \(n\)-manifold, let \(L(M)\) and \(A(M)\) be its linear and affine frame bundles. The Levi-Civita connection \(\Omega^g\) induces the affine connection uniquely characterized by \(i^*\bar\Omega=(\Omega^g,\chi)\) in \(\mathfrak{gl}_n\ltimes\mathbb R^n\). Its curvature restricts to \((R^g,T^g)=(R^g,0)\), because the Levi-Civita connection is torsion free.}
\GSfield{Reconstruction}{The dimension-general formulation and torsion-free condition clarify the four-dimensional example. ``Cartan connection'' follows the affine-frame convention of the source; it is not an arbitrary affine gauge potential with independent translational part.}
\GSfield{Fate}{This construction is written explicitly in \citet[\S2, p.~792, before equation~(5)]{popovdaikhin1975einstein}, published in 1975 after receipt on 26 June.}
\GSfield{Present status}{As of November 2026, the induced affine connection and its curvature decomposition are established.}
\GSfield{Assessment}{literal statement: correct with missing hypotheses; publication match: exact; verification: read the result; development: none; confidence: high}
\end{GSitem}
\begin{GSitem}{1976-407.13}{E}{Vacuum Einstein equations as the affine curvature equation}
\GSfield{Kind; source}{theorem; \GSpage{P000677}}
\GSfield{The notes}{On the metric-induced Cartan connections, \(H\delta\Phi=0\) is identified with \(\operatorname{Ric}=0\). The notes underline that these are first order in the connection.}
\GSfield{Statement}{For the affine connection \(\bar\Omega\) induced by a pseudo-Riemannian metric as above, \[H\delta_h\bar\Phi=0\quad\Longleftrightarrow\quad\operatorname{Ric}(g)=0.\] In the splitting \(\mathfrak{aff}_n=\mathfrak{gl}_n\ltimes\mathbb R^n\), the translational component of the projected equation is, up to the fixed Hodge convention, contraction of the Ricci tensor; the linear component is the covariant-divergence equation for curvature.}
\GSfield{Reconstruction}{The equivalence holds on the specified metric-induced, torsion-free affine connections. It is false if one replaces them merely by linear-frame Levi-Civita connections: their Yang--Mills equation alone gives harmonic curvature, not necessarily Ricci flatness. The assertion is about vacuum equations with zero cosmological constant.}
\GSfield{Fate}{This is exactly \citet[Theorem~2 and equation~(5), p.~792]{popovdaikhin1975einstein}, printed in 1975. The source explicitly contrasts its affine result with Yang's linear-frame equation.}
\GSfield{Present status}{As of November 2026, the equivalence is settled with those constraints; it is not an equivalence between all unrestricted Yang--Mills solutions and all Einstein metrics.}
\GSfield{Assessment}{literal statement: correct with missing hypotheses; publication match: exact; verification: read the result; development: none; confidence: high}
\end{GSitem}
\begin{GSitem}{1976-407.14}{E}{Why the second-order connection equation follows from Ricci flatness}
\GSfield{Kind; source}{claim; \GSpage{P000677}}
\GSfield{The notes}{A final diagram links \(R_{ij}=0\) with the equation \(R_{ij,k}=R_{ik,j}\); the second-order equations are called consequences of the first-order equations.}
\GSfield{Statement}{For a Levi-Civita connection, the contracted differential Bianchi identity expresses the divergence of Riemann curvature as the antisymmetrized covariant derivative of Ricci, with index signs depending on the curvature convention. Hence \(\operatorname{Ric}=0\) implies \(\nabla_j\operatorname{Ric}_{ik}=\nabla_k\operatorname{Ric}_{ij}\) and \(D^*R=0\). Ricci uses one derivative of the connection coefficients, whereas its covariant derivative uses two. In metric variables their orders are two and three respectively.}
\GSfield{Reconstruction}{The derivative-order assertion is relative to connection coefficients and to the metric-induced constraint. The implication is one-way for the linear-frame equation; a parallel nonzero Ricci tensor also satisfies the displayed derivative identity.}
\GSfield{Fate}{The consequence and the fact that the translational component is fixed are stated immediately after \citet[Theorem~2]{popovdaikhin1975einstein}. This is the same pre-1976 observation, not an additional field equation discovered at the talk.}
\GSfield{Present status}{As of November 2026, the implication follows from the differential Bianchi identity; its converse requires the affine translational equation.}
\GSfield{Assessment}{literal statement: correct with missing hypotheses; publication match: exact; verification: read the result; development: none; confidence: high}
\end{GSitem}

\GStalk{408}{1976-11-22}{I. N. Bernstein}{P000678-P000680}{Euler products, compact-group heat kernels, and affine Lie algebras}
Bernstein explains Euler and Jacobi identities, Macdonald's root-system generalization, and Kostant's expression through characters and the Casimir. A shifted spectral zeta function is introduced, but its displayed reflection equation lacks the necessary Mellin completion. Olshanetsky intervenes on radial Laplacians and quantum integrability. Bernstein then gives the Cartan-matrix construction and the affine \(A_2\) loop-algebra example. The date 22 November is supplied by the catalogue; the corresponding handwritten start on P000678 is undated and P000680 ends with a New Year greeting. This evidence does not establish additional sessions within the packet. The mathematical notes are P000678--P000680, rather than the catalogue's wider P000677--P000681 span.

\begin{GSitem}{1976-408.1}{E}{Euler's pentagonal-number identity}
\GSfield{Kind; source}{theorem; \GSpage{P000678}}
\GSfield{The notes}{The page sets \(\varphi(x)=\prod_{n\ge1}(1-x^n)\) and writes \(\varphi(x)=\sum_{n\in\mathbb Z}(-1)^n x^{(3n^2+n)/2}\).}
\GSfield{Statement}{In \(\mathbb Z[[q]]\), and analytically for \(|q|<1\), \[\prod_{m=1}^\infty(1-q^m)=\sum_{n\in\mathbb Z}(-1)^nq^{n(3n+1)/2}.\] Replacing \(n\) by \(-n\) gives the familiar \(n(3n-1)/2\) convention.}
\GSfield{Reconstruction}{The formal-series or unit-disc domain is added; the plus sign in the notebook exponent is correct, not a transcription error. Earlier-index continuity: 1975-366.1. These links identify a recurring statement or construction; only the explicitly indexed inverse problem reuses a problem ID.}
\GSfield{Fate}{This is the classical Euler identity discussed in \citet[\S7, formula~(23)]{demazure1977identites}. That exposition is dated June 1976, although its collected volume was printed in 1977; it therefore documents pre-talk circulation independently of the much older identity.}
\GSfield{Present status}{As of November 2026, the identity is established in both the formal and analytic senses stated.}
\GSfield{Assessment}{literal statement: correct; publication match: exact; verification: read the result; development: none; confidence: high}
\end{GSitem}
\begin{GSitem}{1976-408.2}{E}{Jacobi's cubic Euler-product identity}
\GSfield{Kind; source}{theorem; \GSpage{P000678}}
\GSfield{The notes}{Jacobi's formula is \(\varphi(x)^3=\sum_{n\ge0}(-1)^n(2n+1)x^{n(n+1)/2}\), described as arising from the theta function.}
\GSfield{Statement}{For \(|q|<1\), or as a formal series, \[\prod_{m\ge1}(1-q^m)^3=\sum_{n\ge0}(-1)^n(2n+1)q^{n(n+1)/2}.\] It is the rank-one specialization of the affine denominator identity, equivalently a specialization or derivative of the Jacobi triple product.}
\GSfield{Reconstruction}{The source's lower limit zero and the coefficient \(2n+1\) are retained. This item identifies the identity itself; it does not assert that the notebook contains the derivation. Earlier-index continuity: 1975-366.2. These links identify a recurring statement or construction; only the explicitly indexed inverse problem reuses a problem ID.}
\GSfield{Fate}{\citet[Introduction, equations~(0.6)--(0.7)]{macdonald1972affine} explicitly recovers Jacobi's old identity as the \(A_1\) case. The paper was printed in 1972, well before the seminar.}
\GSfield{Present status}{As of November 2026, the identity and its rank-one root-system interpretation are established.}
\GSfield{Assessment}{literal statement: correct; publication match: exact; verification: read the result; development: none; confidence: high}
\end{GSitem}
\begin{GSitem}{1976-408.3}{E}{Partitions and the reciprocal Euler product}
\GSfield{Kind; source}{claim; \GSpage{P000678}}
\GSfield{The notes}{The notes write \(\varphi(x)^{-1}=\sum p(n)x^n\) and describe \(p(n)\) as partitions into nonnegative summands.}
\GSfield{Statement}{Let \(p(n)\) count unordered partitions of \(n\) into positive integer parts, with \(p(0)=1\). Then \[\prod_{m\ge1}(1-q^m)^{-1}=\sum_{n\ge0}p(n)q^n.\] For each \(m\), the factor \(\sum_{r\ge0}q^{mr}\) records the multiplicity of the positive part \(m\).}
\GSfield{Reconstruction}{Allowing arbitrary zero parts as distinct summands would make the count infinite. The word ``nonnegative'' is harmless only if zero padding is ignored; the precise convention uses positive parts. Earlier-index continuity: 1975-366.1. These links identify a recurring statement or construction; only the explicitly indexed inverse problem reuses a problem ID.}
\GSfield{Fate}{The partition function and its use in eta identities are already discussed in \citet[Introduction, pp.~5--6 of the Russian translation]{macdonald1972affine}. The coefficient interpretation follows directly from the geometric-series factors and is old expository material.}
\GSfield{Present status}{As of November 2026, the generating function is established with positive parts or an explicitly equivalent zero-padding convention.}
\GSfield{Assessment}{literal statement: false as written; publication match: partial; verification: read the result; development: annotated; confidence: high}
\end{GSitem}
\begin{GSitem}{1976-408.4}{E}{Ramanujan's tau function and the eta normalization}
\GSfield{Kind; source}{definition; \GSpage{P000678}}
\GSfield{The notes}{The definition is \(x\varphi(x)^{24}=\sum\tau(n)x^n\), followed by a remark about a combinatorial definition and a reference to Dedekind's ``theta'' function.}
\GSfield{Statement}{For \(q=e^{2\pi i\tau}\), \(\operatorname{Im}\tau>0\), Dedekind's function is \(\eta(\tau)=q^{1/24}\prod_{m\ge1}(1-q^m)\). The discriminant cusp form is \(\Delta_{\rm mod}(\tau)=\eta(\tau)^{24}=\sum_{n\ge1}\tau(n)q^n\). A direct finite signed combinatorial interpretation of \(\tau(n)\) is the sum of \((-1)^{\sum_j|S_j|}\) over 24 finite subsets \(S_j\subset\mathbb Z_{>0}\) satisfying \(\sum_j\sum_{m\in S_j}m=n-1\).}
\GSfield{Reconstruction}{The notation \(\Delta_{\rm mod}\) avoids collision with the Laplacian later on the page. The signed subset formula is the coefficient expansion of the displayed product, supplied here; it is not claimed to be the specific combinatorial formula intended in the shorthand remark. The Dedekind function is \(\eta\), whereas the earlier Jacobi reference is indeed to theta. Earlier-index continuity: 1975-366.3, 1975-370.3. These links identify a recurring statement or construction; only the explicitly indexed inverse problem reuses a problem ID.}
\GSfield{Fate}{\citet[Introduction, \(A_4\) example and Dyson's tau formula]{macdonald1972affine} discusses precisely the 24th-power application before the talk. \citet[\S1(b)]{demazure1977identites}, a June 1976 exposition, gives the \(SU(5)\) example.}
\GSfield{Present status}{As of November 2026, the defining series is established; no Ramanujan conjecture is being posed by this definition.}
\GSfield{Assessment}{literal statement: correct with missing hypotheses; publication match: exact; verification: read the result; development: none; confidence: high}
\end{GSitem}
\begin{GSitem}{1976-408.5}{E}{The invariant Casimir and the compact-group Laplacian}
\GSfield{Kind; source}{construction; \GSpage{P000678}}
\GSfield{The notes}{For compact semisimple \(G\), a Killing-orthonormal basis \(X_i\) is used to write \(\Delta=\sum X_i^2\), asserted to be two-sided invariant and scalar on irreducible representations.}
\GSfield{Statement}{Let \(G\) be a connected compact semisimple Lie group and choose \(X_1,\ldots,X_d\) orthonormal for minus the Killing form. The operator \(\Delta=\sum_i X_i^2\), acting as left-invariant vector fields, is a nonpositive bi-invariant Laplacian. Its negative \(C=-\Delta\) is the positive Casimir and acts on a finite-dimensional irreducible representation of highest weight \(\lambda\) by \(c_\lambda=(\lambda,\lambda+2\rho)\), using the dual Killing normalization.}
\GSfield{Reconstruction}{Minus the Killing form, the Laplace sign, and the normalization on each simple factor are explicit. Without these choices the exponents in the next eta identity and the shift \(d/24\) are not determined.}
\GSfield{Fate}{The Killing-normalized Casimir is the one in \citet[\S3]{demazure1977identites} and in \citet[Theorem~0.1 and its preceding notation]{kostant2003euler}, which restates Kostant's 1976 result. The June 1976 Bourbaki lecture verifies that this formulation circulated before the seminar.}
\GSfield{Present status}{As of November 2026, the operator construction is standard; the notebook needs a consistent sign to use it as an exponent.}
\GSfield{Assessment}{literal statement: correct with missing hypotheses; publication match: exact; verification: read the result; development: none; confidence: high}
\end{GSitem}
\begin{GSitem}{1976-408.6}{E}{Kostant's character expansion of an Euler-product power}
\GSfield{Kind; source}{theorem; \GSpage{P000678-P000679}}
\GSfield{The notes}{For \(k=\dim G\), the notes write \(\varphi(x)^k=\sum_V c_V\dim(V)x^{\Delta|_V}\), with \(c_V=\chi_V(a)\). The example table gives dimensions 3, 24, and 10 for \(SU(2)\), \(SU(5)\), and the symplectic type \(C_2\).}
\GSfield{Statement}{Let \(G\) be compact, connected, semisimple, and simply connected, \(d=\dim G\), and identify the real weight space \(i\operatorname{Lie}(T)\) with its dual by the Killing form. For \(a_\rho=\exp(2\pi i\,2\rho)\), where \(\rho\) is half the sum of positive roots under that identification, \[\prod_{m\ge1}(1-q^m)^d=\sum_{\lambda\in P_+}\chi_\lambda(a_\rho)\dim(V_\lambda)q^{c_\lambda}.\] Here \(c_\lambda\) is the positive Casimir eigenvalue; \(\chi_\lambda(a_\rho)\in\{-1,0,1\}\). The equality holds for \(0<q<1\); only integral Casimir exponents have nonzero character coefficient.}
\GSfield{Reconstruction}{The notes do not specify \(a\), the sign of \(\Delta\), or the covering group. These are supplied from the exact theorem. \(Sp(4)\) in the table is interpreted as the complex matrix notation for type \(C_2\), whose compact form is often denoted \(Sp(2)\); its dimension is 10. Earlier-index continuity: 1975-366.4, 1975-370.18. These links identify a recurring statement or construction; only the explicitly indexed inverse problem reuses a problem ID.}
\GSfield{Fate}{The exact expansion is restated by its author in \citet[Theorem~0.1, citing his 1976 Theorem~3.1]{kostant2003euler}. Its pre-talk formulation was checked directly in \citet[\S3, formulas~(15)--(16)]{demazure1977identites}, dated June 1976. The original 1976 publisher full text was inaccessible; no claim of having read that original is made.}
\GSfield{Present status}{As of November 2026, the character identity is established; the extra hypotheses and normalization prevent a false literal Laplacian-sign reading.}
\GSfield{Assessment}{literal statement: correct with missing hypotheses; publication match: exact; verification: read the result; development: none; confidence: high}
\end{GSitem}
\begin{GSitem}{1976-408.7}{E}{The distinguished regular conjugacy class in SU(n)}
\GSfield{Kind; source}{construction; \GSpage{P000679}}
\GSfield{The notes}{For \(SU(n)\), the notes describe a unique regular conjugacy class whose \(n\)th power is scalar, represented by \(c\operatorname{diag}(\xi,\xi^2,\ldots,\xi^n)\), with \(c\) chosen for determinant one.}
\GSfield{Statement}{Take a primitive \(n\)th root \(\zeta\) and choose \(c\in U(1)\) so that \(c^n\zeta^{n(n+1)/2}=1\). Then \(a=c\operatorname{diag}(\zeta,\ldots,\zeta^n)\in SU(n)\), has distinct eigenvalues, and \(a^n\) is scalar. Conversely, if \(a\in SU(n)\) has distinct eigenvalues and scalar \(n\)th power, its eigenvalues form the full set of roots of one scalar; the determinant condition makes it conjugate to this element. Different permissible choices of \(c\) merely permute that set.}
\GSfield{Reconstruction}{Primitivity of \(\zeta\) and the distinction between one element and one conjugacy class are explicit. The short converse is an elementary reconstruction from the written characterization.}
\GSfield{Fate}{This is the type \(A\) principal-element case of the already circulating Kostant description, recorded in \citet[\S3]{demazure1977identites}; the canonical type-\(\rho\) element is also defined before \citet[Theorem~0.1]{kostant2003euler}.}
\GSfield{Present status}{As of November 2026, the conjugacy-class characterization is exact under the distinct-eigenvalue condition.}
\GSfield{Assessment}{literal statement: correct; publication match: exact; verification: read the result; development: none; confidence: high}
\end{GSitem}
\begin{GSitem}{1976-408.8}{E}{The spectral zeta reflection requires completion}
\GSfield{Kind; source}{claim; \GSpage{P000679}}
\GSfield{The notes}{The notes define \(\zeta_a(s)\) as the central convolution kernel of \((-\Delta+d/24)^{-s}\) evaluated at \(a\), then assert \(\zeta_a(s)=\zeta_a(d/2-s)\). The scan shows no Gamma or power-of-\(2\pi\) factor.}
\GSfield{Statement}{With the Killing-normalized positive Casimir \(C=-\Delta\), normalized Haar measure, and the distinguished element \(a=a_\rho\), initially for \(\operatorname{Re}s>d/2\), \[Z_a(s)=\sum_\lambda\frac{\chi_\lambda(a)\dim V_\lambda}{(c_\lambda+d/24)^s}.\] Its completion \(\Lambda_a(s)=(2\pi)^{-s}\Gamma(s)Z_a(s)\) satisfies \[\Lambda_a(s)=\int_0^\infty t^{s-1}\eta(it)^d\,dt,\qquad \Lambda_a(s)=\Lambda_a(d/2-s).\] The integral is entire in \(s\), owing to exponential decay at both ends after eta inversion.}
\GSfield{Reconstruction}{Kostant's character identity gives the heat-kernel integrand. The classical transformation \(\eta(i/t)=\sqrt t\,\eta(it)\) and the substitution \(t\mapsto1/t\) prove the completed reflection. The raw symmetry is false: \(Z_a(0)=0\), while \(Z_a(d/2)=(2\pi)^{d/2}\Lambda_a(d/2)/\Gamma(d/2)>0\) for \(d>0\). This is a normalization correction, not a new zeta theorem. Earlier-index continuity: 1975-366.5, 1975-370.15. These links identify a recurring statement or construction; only the explicitly indexed inverse problem reuses a problem ID.}
\GSfield{Fate}{The heat-kernel identification is tied to Kostant's 1976 result, restated in \citet[Theorem~0.1]{kostant2003euler}. The eta inversion used to check the missing completion is the classical formula in \citet[23.18.5--23.18.6]{nist2026dlmf}. The exact raw reflection in the notebook has no valid publication match.}
\GSfield{Present status}{As of November 2026, the completed reflection is established; the literal uncompleted equation needs the displayed correction.}
\GSfield{Assessment}{literal statement: false as written; publication match: partial; verification: read the result; development: annotated; confidence: high}
\end{GSitem}
\begin{GSitem}{1976-408.9}{F}{The unidentified 39th- and 45th-power formulas}
\GSfield{Kind; source}{remark; \GSpage{P000679}}
\GSfield{The notes}{A short line says that there are also formulas for \(\varphi^{39}\) and \(\varphi^{45}\), followed by the Latin-script name ``Van Esh.'' No sum, lattice, coefficient rule, or normalization is given.}
\GSfield{Statement}{The only definite content is the claimed existence of unspecified identities for two Euler-product powers. Since every such power has an elementary coefficient expansion, the mathematical content depends on the missing special form of the asserted formulas.}
\GSfield{Reconstruction}{The plausible author identification is A.~G. (Bram) van Asch, not a new theorem attributed to a literal ``Van Esh.'' His 1976 paper and the 1973 announcement cited by Demazure are leads; the exact 39/45 assertions were not recovered from an opened full text.}
\GSfield{Fate}{\citet[\S1 and bibliography item~7]{demazure1977identites} identifies van Asch's direct proof and a 1973 announcement. Publisher metadata verifies \citet{vanasch1976modular}, received 28 September 1975 and published June 1976. Searches for the name variants and the exponents did not verify the two particular formulas; the accessible Demazure text does not state them.}
\GSfield{Present status}{As of November 2026, the author and likely literature are partly identified, but the exact notebook item is still not determinate enough to assert a solved or open problem.}
\GSfield{Assessment}{literal statement: not determinable; publication match: thematic only; verification: unresolved lead; development: none; confidence: medium}
\end{GSitem}
\begin{GSitem}{1976-408.10}{E}{The affine denominator and homological explanation of Macdonald identities}
\GSfield{Kind; source}{method; \GSpage{P000679}}
\GSfield{The notes}{Bernstein credits Kac with a proof using infinite-dimensional Lie algebras and gives the precise Garland--Lepowsky citation, Inventiones 34 (1976), 37--76.}
\GSfield{Statement}{For a symmetrizable Kac--Moody algebra with positive roots \(\Delta_+\), root multiplicities \(\operatorname{mult}\alpha\), Weyl group \(W\), and a weight \(\rho\) satisfying \(\rho(h_i)=1\), the denominator identity is \[\prod_{\alpha\in\Delta_+}(1-e^{-\alpha})^{\operatorname{mult}\alpha}=\sum_{w\in W}(-1)^{\ell(w)}e^{w\rho-\rho}.\] It is an identity in the completion by the positive root cone. In the affine case, the imaginary-root factors account for the Euler-product factor; specialization yields Macdonald's identities.}
\GSfield{Reconstruction}{The formula, the symmetrizability hypothesis, the root multiplicities, and the formal completion are reconstructed from the cited proof tradition, not displayed in the notebook. They explain the role of infinite-dimensional homology rather than merely naming it.}
\GSfield{Fate}{The exact bibliography is \citet{garlandlepowsky1976homology}, received 7 January and published February 1976. Its publisher metadata was checked; the denominator statement and explanation were read in \citet[\S\S4--6, especially (19), (21)]{demazure1977identites}, delivered June 1976. Hence the method was already public before this seminar.}
\GSfield{Present status}{As of November 2026, the affine denominator interpretation is established, including the multiplicities of imaginary roots.}
\GSfield{Assessment}{literal statement: correct with missing hypotheses; publication match: exact; verification: read the result; development: none; confidence: high}
\end{GSitem}
\begin{GSitem}{1976-408.11}{A}{Radial Laplacians become inverse-sine-square Hamiltonians}
\GSfield{Kind; source}{construction; \GSpage{P000679}; M. A. Olshanetsky}
\GSfield{The notes}{Olshanetsky writes that the radial Laplace operator on a symmetric space becomes \(\tfrac12\sum\partial_{t_k}^2+U\), with \(U=g\sum_{k>\ell}\sin^{-2}(t_k-t_\ell)\).}
\GSfield{Statement}{On a regular radial chamber of a Riemannian symmetric space, let \(J(t)\) be the radial volume density and \(\Delta_{\rm rad}=J^{-1}\operatorname{div}(J\nabla)\). Conjugation gives \[-\tfrac12J^{1/2}\Delta_{\rm rad}J^{-1/2}=-\tfrac12\Delta_{\mathfrak a}+\tfrac14\Delta_{\mathfrak a}\log J+\tfrac18|\nabla\log J|^2.\] In compact type \(A\), \(J=\prod_{i<j}|\sin(t_i-t_j)|^m\); the resulting potential, up to an additive constant and metric scale, is \(m(m-2)\sum_{i<j}\sin^{-2}(t_i-t_j)/4\).}
\GSfield{Reconstruction}{The conjugating factor, chamber, kinetic sign, and geometric coupling are supplied. A general symmetric space has its own restricted root system and root multiplicities; the single pairwise \(A\)-type expression is not its universal formula. In this geometric derivation the coupling is fixed by multiplicities, not arbitrary.}
\GSfield{Fate}{The exact radial-conjugation construction appears in \citet[\S3, Lemma~1 and formula~(3.6)]{olshanetskyperelomov1978radial}, received 16 June 1977 and published in 1978. Class A follows the intervention's author and the subsequent joint paper; older many-particle examples alone would not establish this geometric match.}
\GSfield{Present status}{As of November 2026, the radial construction is established with group multiplicities and a specified radial domain.}
\GSfield{Assessment}{literal statement: correct with missing hypotheses; publication match: exact; verification: read the result; development: none; confidence: high}
\end{GSitem}
\begin{GSitem}{1976-408.12}{A}{A complete commuting family of quantum integrals from radial operators}
\GSfield{Kind; source}{claim; \GSpage{P000679}; M. A. Olshanetsky}
\GSfield{The notes}{Olshanetsky adds that there is a complete system of observables ``in the quantum sense.'' No domains or specific higher operators are written.}
\GSfield{Statement}{For the Hamiltonians obtained from symmetric-space radial parts with the group values of the coupling constants, the radial invariant differential operators form a commutative algebra. Conjugating this algebra by the same \(J^{1/2}\) gives commuting differential operators including the Hamiltonian and having \(r=\dim\mathfrak a\) algebraically independent principal symbols. This is formal quantum complete integrability on the regular chamber.}
\GSfield{Reconstruction}{The precise assertion concerns commuting differential expressions and independence of symbols. Strong commutation of chosen self-adjoint realizations at collision walls would require domain conditions absent from the notes and is not inferred.}
\GSfield{Fate}{This is the group-coupling assertion of \citet[\S4, Theorem~2 and the preceding commutativity corollary]{olshanetskyperelomov1978radial}, received in 1977 and printed in 1978. The paper gives explicit generators for several root systems in \S5.}
\GSfield{Present status}{As of November 2026, the formal commuting algebra is established under those hypotheses; operator-domain issues are outside this recorded claim.}
\GSfield{Assessment}{literal statement: correct with missing hypotheses; publication match: exact; verification: read the result; development: none; confidence: high}
\end{GSitem}
\begin{GSitem}{1976-408.13}{E}{Generalized Cartan matrices}
\GSfield{Kind; source}{definition; \GSpage{P000679}}
\GSfield{The notes}{The notebook requires integer entries, \(a_{ii}=2\), symmetric occurrence of zeros, but writes \(a_{ij}<0\) for every \(i\ne j\).}
\GSfield{Statement}{A generalized Cartan matrix is an integer matrix \(A=(a_{ij})\) with \(a_{ii}=2\), \(a_{ij}\le0\) for \(i\ne j\), and \(a_{ij}=0\) if and only if \(a_{ji}=0\). It is symmetrizable if \(DA\) is symmetric for a positive diagonal matrix \(D\).}
\GSfield{Reconstruction}{The strict inequality is a faithful error: it contradicts the adjacent comment about zeros and would exclude disconnected diagrams and most ordinary Dynkin diagrams. Symmetrizability is additional structure, not part of the three basic conditions.}
\GSfield{Fate}{The conditions are classical before the talk: see \citet[Chapter~II, \S1, conditions (M2)--(M5)]{kac1968graded} and \citet[\S5]{demazure1977identites}.}
\GSfield{Present status}{As of November 2026, the corrected nonpositive-off-diagonal definition is standard.}
\GSfield{Assessment}{literal statement: false as written; publication match: exact; verification: read the result; development: annotated; confidence: high}
\end{GSitem}
\begin{GSitem}{1976-408.14}{E}{Chevalley--Serre generators and the missing conventions}
\GSfield{Kind; source}{construction; \GSpage{P000680}}
\GSfield{The notes}{There are three families \(h_i,e_i,f_i\), called ``\(n\) generators.'' The general brackets use the signs opposite to the later \(\mathfrak{sl}_2\) example, and the \([e,f]\) line repeats the subscript \(i\). Only the \(e\)-Serre relation is explicitly displayed; a maximal-ideal quotient is also mentioned.}
\GSfield{Statement}{Over \(\mathbb C\), the derived Kac--Moody algebra for a generalized Cartan matrix is generated by \(e_i,f_i,h_i\) with \[[h_i,h_j]=0,\quad[h_i,e_j]=a_{ij}e_j,\quad[h_i,f_j]=-a_{ij}f_j,\quad[e_i,f_j]=\delta_{ij}h_i,\] together with \((\operatorname{ad}e_i)^{1-a_{ij}}e_j=(\operatorname{ad}f_i)^{1-a_{ij}}f_j=0\) for \(i\ne j\). The full Kac--Moody algebra uses a realization with Cartan dimension \(2n-\operatorname{rank}A\); for singular \(A\), it is larger than the algebra generated by these \(3n\) elements.}
\GSfield{Reconstruction}{The statement fixes a standard sign convention consistent with the notebook's explicit \(\mathfrak{sl}_2\) brackets, supplies the missing \(f\)-relations and \(i\ne j\), and distinguishes the derived from the full algebra. In a contragredient maximal-ideal construction, the quotient ideal must meet the chosen Cartan subalgebra trivially; an unrestricted ``maximal normal subalgebra'' is not the definition.}
\GSfield{Fate}{The finite/affine defining relations are in \citet[Chapter~II, Proposition~13, equations~(39)]{kac1968graded}; the generalized construction in the form being explained is \citet[\S5]{demazure1977identites}, dated June 1976.}
\GSfield{Present status}{As of November 2026, the corrected presentation is established; the literal mixed signs and indices are not a consistent presentation of the examples that follow.}
\GSfield{Assessment}{literal statement: false as written; publication match: exact; verification: read the result; development: annotated; confidence: high}
\end{GSitem}
\begin{GSitem}{1976-408.15}{E}{Finite type and the positive-definite criterion}
\GSfield{Kind; source}{theorem; \GSpage{P000680}}
\GSfield{The notes}{The matrix \((2)\) is associated with \(SU(2)\) and \(\left(\begin{smallmatrix}2&-1\\-1&2\end{smallmatrix}\right)\) with \(SU(3)\). For symmetric \(A\), finite dimensionality is asserted equivalent to positive definiteness.}
\GSfield{Statement}{For a symmetric generalized Cartan matrix \(A\), the associated complex Kac--Moody algebra is finite dimensional exactly when \(A\) is positive definite. In that case it is semisimple, with simple factors given by the connected finite Dynkin diagrams. The matrices \((2)\) and \(A_2=\left(\begin{smallmatrix}2&-1\\-1&2\end{smallmatrix}\right)\) give \(\mathfrak{sl}_2(\mathbb C)\) and \(\mathfrak{sl}_3(\mathbb C)\), whose compact real forms integrate to \(SU(2)\) and \(SU(3)\).}
\GSfield{Reconstruction}{The notebook calls the constructed Lie algebra a group. Distinguishing the complex algebra, compact real form, and group avoids tensoring a Lie group with a polynomial ring in the next item. The positive-definite criterion concerns the symmetric matrix here; for symmetrizable matrices it concerns a positive diagonal symmetrization.}
\GSfield{Fate}{The finite diagrams and their contrast with affine diagrams are classified in \citet[Chapter~II, \S\S4--7, Tables~0--3]{kac1968graded}; the finite presentation is also recalled in \citet[\S5]{demazure1977identites}. Both precede the seminar.}
\GSfield{Present status}{As of November 2026, the finite-type criterion and the two examples are established.}
\GSfield{Assessment}{literal statement: correct with missing hypotheses; publication match: exact; verification: read the result; development: none; confidence: high}
\end{GSitem}
\begin{GSitem}{1976-408.16}{E}{Affine A2 as a central extension of the loop algebra}
\GSfield{Kind; source}{theorem; \GSpage{P000680}}
\GSfield{The notes}{For \(A=\left(\begin{smallmatrix}2&-1&-1\\-1&2&-1\\-1&-1&2\end{smallmatrix}\right)\), the quotient by a one-dimensional center is written \(SU(3)\otimes\mathbb C[x,x^{-1}]\), with \([x^kg,x^\ell g']=x^{k+\ell}[g,g']\).}
\GSfield{Statement}{The derived affine algebra of type \(A_2^{(1)}\) is \(\widehat{\mathfrak{sl}}_3'=(\mathfrak{sl}_3(\mathbb C)\otimes\mathbb C[t,t^{-1}])\oplus\mathbb Cc\), with \[[X\otimes t^m,Y\otimes t^n]=[X,Y]\otimes t^{m+n}+m\delta_{m+n,0}\kappa(X,Y)c,\] for a fixed nonzero invariant form \(\kappa\). Quotienting by \(\mathbb Cc\) gives the loop algebra and exactly the bracket written in the notebook. The full affine algebra additionally has the degree derivation \(d\), with \([d,X\otimes t^m]=mX\otimes t^m\).}
\GSfield{Reconstruction}{The tensor factor is the complex Lie algebra, not \(SU(3)\). The central cocycle and the derived/full distinction are reconstruction explaining why the notebook's center quotient is correct in its intended convention.}
\GSfield{Fate}{\citet[Chapter~II, \S7, loop construction, Lemma~22 and Proposition~13]{kac1968graded} gives the affine realization modulo center. The notebook's conclusion is an exposition of that pre-existing result.}
\GSfield{Present status}{As of November 2026, the central extension, quotient, and derivation distinction are established.}
\GSfield{Assessment}{literal statement: correct with missing hypotheses; publication match: exact; verification: read the result; development: none; confidence: high}
\end{GSitem}
\begin{GSitem}{1976-408.17}{E}{The scope of the positive-semidefinite loop classification}
\GSfield{Kind; source}{claim; \GSpage{P000680}}
\GSfield{The notes}{After the affine example, the notes say \ru{все нестрог. опред. получаются так}.}
\GSfield{Statement}{If a symmetric generalized Cartan matrix is indecomposable, positive semidefinite, and singular, it has affine type and corank one; its derived algebra modulo center is the corresponding untwisted loop algebra of finite type \(A,D,E\), including the rank-one affine case. If indecomposability is dropped, positive-semidefinite matrices decompose into finite and affine blocks. In the broader symmetrizable setting, twisted loop algebras also occur.}
\GSfield{Reconstruction}{``Non-strictly definite'' is interpreted as positive semidefinite in the immediately preceding symmetric setting. It cannot mean every matrix that fails positive definiteness: indefinite Kac--Moody algebras generally have neither this corank nor this loop realization.}
\GSfield{Fate}{The distinction between finite, affine, and other growth is the content of \citet[Chapter~II, Theorem~1, \S7 and Tables~0--3]{kac1968graded}. This prior classification supplies the missing scope.}
\GSfield{Present status}{As of November 2026, the restricted affine statement is settled; the unrestricted reading would be false.}
\GSfield{Assessment}{literal statement: correct with missing hypotheses; publication match: exact; verification: read the result; development: none; confidence: high}
\end{GSitem}

\GStalk{409}{}{-}{P000680-P000681}{YEAR 1977}
This catalogue record is the archival year heading between the late-1976 material and the dated Manin notes beginning on P000681. It supplies no independent mathematical statement. The wider catalogue packet P000679--P000681 contains the neighboring talks; no content is invented for the heading.


\GStalk{410}{1977-02-14}{Yu. I. Manin}{P000681-P000684}{Spectral curves, rank-two solitons, and filtered bimodules}
Manin begins with a reference to Calogero's boomeron, then describes a fourth-order Lax system and an explicitly singular, periodically oscillating rank-two pulse. The algebraic part passes from commuting operators to a spectral curve and eigenfunction bundles, contrasts rank-one Jacobian motion with higher-rank bundle moduli, and gives the filtered-bimodule construction and its inverse. Two further connections yield evolution and zero-curvature equations. The material has an unusually close later match in Manin's 1978 survey, including a sign error in the flatness formula that survives into the printed English translation. The last sentence announces Duistermaat for the next meeting; it does not contain that lecture.

\begin{GSitem}{1976-410.1}{E}{Calogero and the boomeron reference}
\GSfield{Kind; source}{remark; \GSpage{P000681}}
\GSfield{The notes}{The opening reads \ru{Калоджеро: бумерон}, beside a sketch of a wave turning back. No equation or parameter identification is attached to this heading.}
\GSfield{Statement}{The identifiable bibliographical referent is the returning-soliton construction named ``boomeron'' by Calogero and Degasperis. The heading alone does not specify which member of their coupled evolution systems is intended, or identify it with the singular pulse subsequently displayed by Manin.}
\GSfield{Reconstruction}{The sketch is read only as a change of propagation direction. The returning boomeron and Manin's periodically singular geyseron are not identified solely from their juxtaposition.}
\GSfield{Fate}{\citet{calogerodegasperis1976boomeron} appeared in the July 1976 issue and was received on 6 May 1976, before this lecture. The publisher's displayed 2008 date is online publication, not the original issue date. Only the publication metadata and title were checked, not an exact match of the subsequent formulas.}
\GSfield{Present status}{As of November 2026, the named reference is identified, while the notebook supplies no separate mathematical assertion about its exact solution.}
\GSfield{Assessment}{literal statement: not determinable; publication match: thematic only; verification: publisher metadata checked; development: none; confidence: high}
\end{GSitem}
\begin{GSitem}{1976-410.2}{A}{The fourth-order Lax commutator and two incorrect coefficients}
\GSfield{Kind; source}{claim; \GSpage{P000681}}
\GSfield{The notes}{The notes set \(L=D^4+vD^2+wD+z\), \(P=D^2+AD+u\), \(D=\partial_x\), and \(L_t=[L,P]\). They print \(v_t=v_{xx}-Av_x-2w_x\), \(w_t=2v_{xxx}+2vv_x-w_{xx}-Aw_x-2z_x\), and the displayed third equation below.}
\GSfield{Statement}{For scalar smooth or meromorphic coefficients on a connected coordinate domain, write \([L,P]=LP-PL\). Absence of the \(D^4\) and \(D^3\) terms forces \(A_x=0\) and \(u=v/2+c_0(t)\). The equivalent system is
 \[\begin{aligned}v_t&=2v_{xx}-Av_x-2w_x,\\ w_t&=2v_{xxx}+vv_x-w_{xx}-Aw_x-2z_x,\\ z_t&=\tfrac12v_{xxxx}+\tfrac12vv_{xx}+\tfrac12wv_x-z_{xx}-Az_x.\end{aligned}\]
 The scalar summand \(c_0(t)\) commutes with \(L\) and has no effect.}
\GSfield{Reconstruction}{The notebook's coefficient of \(v_{xx}\) is 1 instead of 2, and its coefficient of \(vv_x\) is 2 instead of 1. Both are legible in the scan and faithfully transcribed. They cannot both be explained by reversing the commutator or rescaling time. Independent Leibniz-rule expansion and application to a test exponential give the displayed correction; this is a critical annotation, not a new integrable-system result.}
\GSfield{Fate}{The same fourth/second-order system is printed with the correct relative coefficients in \citet[Chapter~III, \S1.5, pp.~79--80; \S6.3, equation~(17), p.~96]{manin1978algebraic}. That source uses \(L_t=[P,L]\) and explicit time/speed scale factors, so a convention change is needed before comparing the right-hand sides. The Russian survey appeared in 1978, the cited English translation in 1979.}
\GSfield{Present status}{As of November 2026, the commutator expansion is settled; the two handwritten coefficients are false under the stated convention.}
\GSfield{Assessment}{literal statement: false as written; publication match: exact; verification: read the result; development: annotated; confidence: high}
\end{GSitem}
\begin{GSitem}{1976-410.3}{A}{The recorded rank-two pulse and the missing solution data}
\GSfield{Kind; source}{construction; \GSpage{P000681}}
\GSfield{The notes}{With \(\xi=x+ct\), the notebook gives
 \[\mu_1=2\frac{\cosh2\xi-\cos2\xi+2\sin\sigma t}{\cosh2\xi+\cos2\xi-2\cos\sigma t},\qquad \mu_2=2\frac{\sinh2\xi-\sin2\xi}{\cosh2\xi+\cos2\xi-2\cos\sigma t},\]
 then \(v=-4\mu_{2x}\), \(\sigma w=4\mu_2\mu_{2x}+6\mu_{2xx}+4\mu_{1x}\). It omits \(z\), calling it inessential, and leaves the relation of \(\sigma\) to \(A\) unspecified.}
\GSfield{Statement}{The displayed expressions define meromorphic functions wherever their common denominator is nonzero. They do not, by themselves, specify a complete solution \((v,w,z)\) of a particular Lax system: \(z\), the parameter relation, and a consistent time/sign normalization must also be given. A complete rank-two one-soliton family with this hyperbolic/trigonometric structure is supplied in Manin's later construction.}
\GSfield{Reconstruction}{No missing \(z\) or parameter relation is invented. In the printed family, writing \(\theta=4a^2\omega t\), the denominator is \(\cosh(2a\xi)+\cos(2a\xi)+2\cos\theta\), the prefactor of \(\mu_2\) is \(-2a\), and the formula is \(w=-6\mu_{2xx}-4\mu_{1x}+4\mu_2\mu_{2x}\). These differences prevent a literal formula-for-formula identification without an explicit change of variables, which the notes do not record.}
\GSfield{Fate}{\citet[Chapter~III, \S1.5, p.~80; \S6.4, equations~(18)--(22) and following display, p.~96]{manin1978algebraic} gives the construction and complete coefficients. The precise structural match supports identification of the family, while the recorded normalization and omitted component justify only a partial publication match.}
\GSfield{Present status}{As of November 2026, the published family is established; the truncated notebook expression is not certified as a complete solution in its literal normalization.}
\GSfield{Assessment}{literal statement: not determinable; publication match: partial; verification: read the result; development: none; confidence: medium}
\end{GSitem}
\begin{GSitem}{1976-410.4}{A}{Periodic singularities of the displayed denominator}
\GSfield{Kind; source}{claim; \GSpage{P000681}}
\GSfield{The notes}{Beside the pulse sketch the notes say that every \(2\pi/\sigma\) it ``spits,'' with the names geyseron and shooting soliton, and describe explosive instability.}
\GSfield{Statement}{For real \(\xi,t\) and \(\sigma\ne0\), the notebook denominator \(\Delta=\cosh2\xi+\cos2\xi-2\cos\sigma t\) is nonnegative and vanishes exactly when \(\xi=0\) and \(\sigma t\in2\pi\mathbb Z\). At those times \(\mu_2=4/\xi+O(\xi^3)\), hence \(v=-4\partial_x\mu_2=16/\xi^2+O(\xi^2)\). The positive temporal period is \(2\pi/|\sigma|\).}
\GSfield{Reconstruction}{The check follows from \(\cosh2\xi+\cos2\xi=2\sum_{j\ge0}(2\xi)^{4j}/(4j)!\). This verifies an actual pole of the recorded expression and its recurrence, but not nonlinear instability in a specified norm, which would require a perturbation theorem. For \(\sigma=0\), the formula instead has a persistent singular center and the stated period is undefined.}
\GSfield{Fate}{\citet[Chapter~III, \S1.5, p.~80]{manin1978algebraic} describes exactly the same qualitative geyser behavior for the published normalization, whose singular times are phase-shifted. The denominator and Laurent-series check here is independent of the omitted \(z\).}
\GSfield{Present status}{As of November 2026, periodic pole formation in the displayed functions is verified; no stronger stability theorem is implied by the notebook wording.}
\GSfield{Assessment}{literal statement: correct with missing hypotheses; publication match: partial; verification: read the result; development: annotated; confidence: high}
\end{GSitem}
\begin{GSitem}{1976-410.5}{E}{Commuting operators, spectral curves, and eigenfunction sheaves}
\GSfield{Kind; source}{method; \GSpage{P000681-P000682}}
\GSfield{The notes}{Manin first constructs \([L_1,P_1]=0\), then solves \(L_1\psi=\lambda\psi\), \(P_1\psi=\mu\psi\) over an affine algebraic curve \(X\). The sketch puts the common-eigenfunction space over \((\lambda,\mu)\), and the next page extends the bundle to a completion.}
\GSfield{Statement}{For commuting scalar ordinary differential operators of positive orders over a differential field of characteristic zero with algebraically closed constant field, the Burchnall--Chaundy relation gives a nonzero polynomial \(Q\) over the constants with \(Q(L_1,P_1)=0\). On a component of the resulting spectral curve where the joint-eigenspace dimension is constant, these eigenspaces form the spectral bundle. Compactification and normalization, with a specified point at infinity and suitable local data, give the algebro-geometric spectral data. At singularities one must allow a torsion-free sheaf rather than automatically a vector bundle.}
\GSfield{Reconstruction}{The scalar, constant-field, generic-rank and completion qualifications are supplied. The sketch is a fibre diagram, not evidence that every singular joint spectrum has an everywhere locally free eigenbundle. The matrix-coefficient case later in the talk is handled through the bimodule axioms, not by an unrestricted extension of the scalar theorem. Earlier-index continuity: 1975-391.1. These links identify a recurring statement or construction; only the explicitly indexed inverse problem reuses a problem ID.}
\GSfield{Fate}{\citet[\S2, proposition on p.~293]{krichever1976zakharov} already states the polynomial relation and rank-one inverse construction; it was received 8 October and presented 24 October 1975, then published in 1976. \citet[Theorem~1.1, pp.~176--177]{krichever1978commutative} explicitly attributes the relation to Burchnall and Chaundy. This continues earlier item 1975-391.1.}
\GSfield{Present status}{As of November 2026, the qualified spectral-curve method is established; the singular-fibre caveat remains necessary.}
\GSfield{Assessment}{literal statement: correct with missing hypotheses; publication match: exact; verification: read the result; development: none; confidence: high}
\end{GSitem}
\begin{GSitem}{1976-410.6}{E}{Rank-one evolution as translation on the Jacobian}
\GSfield{Kind; source}{theorem; \GSpage{P000682}}
\GSfield{The notes}{For a one-dimensional common-eigenfunction space the page writes \(t\longrightarrow\) a path on the Jacobian. The previous page identifies the rank-one restriction in Krichever's construction.}
\GSfield{Statement}{Let \(C\) be a smooth projective complex curve, \(p_\infty\) a marked point with local coordinate, and choose nonspecial rank-one Baker--Akhiezer divisor data. The commuting flows associated with prescribed principal parts at \(p_\infty\) evolve the degree-zero class of the spectral line bundle by \(a(t)=a(0)+tU\) in \(\operatorname{Jac}(C)\), where \(U\) is determined by the corresponding normalized second-kind differential. Several commuting times give a sum of such translations. The local reconstruction is used away from its exceptional theta divisor.}
\GSfield{Reconstruction}{The notebook says only ``path''; linear translation and the hypotheses are the precise established rank-one statement. Identifying \(\operatorname{Pic}^d(C)\) with the Jacobian requires a reference line bundle. This is not asserted for every path of arbitrary higher-rank bundles. Earlier-index continuity: 1975-384.3, 1975-385.8. These links identify a recurring statement or construction; only the explicitly indexed inverse problem reuses a problem ID.}
\GSfield{Fate}{The rank-one construction is already \citet[\S1, Theorem~1 and Corollaries~1--2]{krichever1976zakharov}; \citet[Chapter~III, \S4, especially the concluding Jacobian linearization]{manin1978algebraic} presents the corresponding bimodule proof. The topic continues the earlier rank-one inverse problem GP-1975-14.}
\GSfield{Present status}{As of November 2026, the rank-one linearization is established under the stated spectral-data hypotheses.}
\GSfield{Assessment}{literal statement: correct with missing hypotheses; publication match: exact; verification: read the result; development: none; confidence: high}
\end{GSitem}
\begin{GSitem}{1976-410.7}{B}{The recurring higher-rank inverse construction}
\GSfield{Kind; source}{question; \GSpage{P000681-P000682}}
\GSfield{The notes}{The notes contrast \(\dim\{\psi\}=1\) with \(r>1\), refer to Drinfeld's nonexplicit description and difficulties, and replace the Jacobian by a moduli space of bundles. They do not give a higher-rank reconstruction formula here.}
\GSfield{Statement}{The precise recurring inverse problem is to recover commuting scalar ordinary differential operators of generic spectral rank \(r>1\) from appropriate pointed-curve and rank-\(r\) spectral-bundle data. In the nonsingular-curve, general-position setting, one must include a local parameter, a normalized matrix divisor (Tyurin data), and the additional functional parameters; the curve and an unframed bundle alone do not determine the operators.}
\GSfield{Reconstruction}{The inverse-problem formulation is reconstructed from the contrast with the preceding construction and from earlier occurrences 1975-391.3--4. It is narrower than a claim that every holomorphic path in every bundle-moduli space yields a Lax solution. The data and genericity here follow the later effective reconstruction, not an unrecorded theorem attributed to Manin. Earlier-index continuity: 1975-391.3, 1975-391.4. These links identify a recurring statement or construction; only the explicitly indexed inverse problem reuses a problem ID.}
\GSfield{Fate}{\citet[Proposition~2]{drinfeld1977commutative}, received 14 April 1976 and printed in the first 1977 issue, gives the abstract geometric correspondence; the precise public issue date relative to 14 February is not established. \citet[\S2, Lemma~2.2 and Theorem~2.3, pp.~181--182]{krichever1978commutative} gives an effective reconstruction for smooth spectral curves in general position. Its introduction expressly corrects serious errors in the 1977 implementation; it was submitted 15 February 1977, the day after this talk, and printed in 1978. The class B assignment refers to that effective reconstruction by another author, not to discovery of the older abstract correspondence.}
\GSfield{Present status}{As of November 2026, the stated generic inverse problem is solved. Singular spectral curves, special strata, and global real regularity are separate questions, not residual open claims assigned to this notebook.}
\GSfield{Assessment}{literal statement: not determinable; publication match: partial; verification: read the result; development: none; confidence: medium; problem \hyperref[prob:GP-1975-26]{GP-1975-26}}
\end{GSitem}
\begin{GSitem}{1976-410.8}{F}{Hyperelliptic curves and the two-quadrics shorthand}
\GSfield{Kind; source}{claim; \GSpage{P000682}}
\GSfield{The notes}{Immediately after the higher-rank moduli remark, the page says \ru{Гиперэлл. кривая} \(\longrightarrow\) \ru{пересеч. 2-х квадрик в} \(\mathbb C^5\). It does not specify a genus, rank, determinant, or which moduli space the arrow denotes.}
\GSfield{Statement}{A precise plausible referent is the classical theorem that for a smooth projective complex curve \(C\) of genus two and a fixed line bundle \(\Lambda\) of odd degree, the moduli space of stable rank-two bundles with determinant \(\Lambda\) is a smooth intersection of two quadrics in \(\mathbb P^5\). The notebook does not supply enough data to identify that theorem uniquely, and it writes affine \(\mathbb C^5\), not projective \(\mathbb P^5\).}
\GSfield{Reconstruction}{Genus two, rank two, fixed odd determinant, stability, and projective ambient space are all candidate reconstruction, not recovered words. A general hyperelliptic curve is not itself a three-dimensional intersection of two quadrics. An affine chart of the classical moduli threefold could explain \(\mathbb C^5\), but no chart is specified.}
\GSfield{Fate}{\citet[\S10, p.~546]{ramanan1980vector} states the hyperelliptic moduli/quadric construction and explicitly attributes the genus-two case to Newstead (1968) and Narasimhan--Ramanan (1969); the publisher metadata for \citet{narasimhanramanan1969moduli} was also checked. This is a strong historical candidate, not an exact identification of the under-specified arrow.}
\GSfield{Present status}{As of November 2026, the candidate classical theorem is settled, but the precise object intended by this notebook shorthand remains uncertain.}
\GSfield{Assessment}{literal statement: not determinable; publication match: thematic only; verification: read the result; development: none; confidence: medium}
\end{GSitem}
\begin{GSitem}{1976-410.9}{F}{The unqualified question about analytic paths}
\GSfield{Kind; source}{question; \GSpage{P000682}}
\GSfield{The notes}{The separate two-word question \ru{Аналит. пути?} follows the moduli and two-quadrics remarks.}
\GSfield{Statement}{No definite analytic-path problem is specified: the notes give neither a vector field, an initial condition, a regularity category, a domain of continuation, nor a chosen moduli stratum. The question cannot be assigned a precise existence, integrability, or completeness assertion from this wording alone.}
\GSfield{Reconstruction}{It may concern explicit higher-rank flows on a bundle-moduli space, but that interpretation is not enough to select one theorem. In particular, generic local reconstruction and global nonsingular continuation are different mathematical questions.}
\GSfield{Fate}{The neighboring sources \citet[Chapter~III, \S\S1.4, 3--6]{manin1978algebraic} and \citet[\S\S2--3]{krichever1978commutative} discuss higher-rank construction and bundle deformations, but neither provides an exact identifiable target for the isolated wording. No open-problem status is inferred from failure of identification.}
\GSfield{Present status}{As of November 2026, the intended analytic question is unidentified; this is not evidence for an unresolved theorem.}
\GSfield{Assessment}{literal statement: not determinable; publication match: thematic only; verification: unresolved lead; development: none; confidence: low}
\end{GSitem}
\begin{GSitem}{1976-410.10}{A}{Nodal spectral curves and soliton solutions}
\GSfield{Kind; source}{method; \GSpage{P000682}}
\GSfield{The notes}{The notes describe a soliton as a nonstationary solution obtained by using a stationary spectral curve with one double point and distinct tangents.}
\GSfield{Statement}{In the algebro-geometric soliton construction, a rational spectral curve with an ordinary node is obtained by identifying two distinct smooth points of its normalization. A spectral line bundle, or higher-rank bundle with specified gluing matrix, together with marked-point asymptotics and compatible commuting connections, can produce one-soliton-type solutions. More nodes give corresponding gluing parameters and multisoliton families. The node alone does not determine a solution or ensure real nonsingularity.}
\GSfield{Reconstruction}{The normalization, gluing, spectral data and compatibility assumptions are supplied to replace the overbroad descriptive definition. The actual double point in the note has two branches; a cusp would require a different local construction.}
\GSfield{Fate}{\citet[Chapter~III, \S1.4; \S5, especially \S\S5.1--5.5 and 5.12; \S6.1]{manin1978algebraic} gives the precise nodal gluing construction and identifies the number of pairs of glued points with the soliton count. The rank-one degeneration was older; this class A entry concerns the rank-two/matrix construction being explained by the speaker and published in his 1978 survey.}
\GSfield{Present status}{As of November 2026, the stated spectral-data construction is established; the literal node-only criterion needs the additional data and admissibility conditions.}
\GSfield{Assessment}{literal statement: correct with missing hypotheses; publication match: partial; verification: read the result; development: none; confidence: high}
\end{GSitem}
\begin{GSitem}{1976-410.11}{A}{Differential-operator algebra and compatible evolution conventions}
\GSfield{Kind; source}{definition; \GSpage{P000682}}
\GSfield{The notes}{The algebraic setting is a characteristic-zero field \(k\), a \(k\)-algebra \(\mathcal B\), and commuting derivations \(\partial,\partial_1,\partial_2\). The analytic example is \(M_\ell(\mathcal B_0)\), with meromorphic coefficients. The notebook lists stationary Lax, nonstationary Lax and Zakharov--Shabat equations.}
\GSfield{Statement}{Let \(\mathcal B\) be a unital associative central \(k\)-algebra with commuting \(k\)-derivations \(d,d_1,d_2\). The Ore algebra \(\mathcal B[D;d]\) consists of finite sums \(\sum b_jD^j\), with multiplication fixed by \(Db=bD+d(b)\). The derivations \(d_i\) extend coefficientwise and satisfy \(d_iD=0\). Stationarity is \([L,P]=0\). With the later connection convention \(\nabla_1\mathbf1=P\mathbf1\), \(\nabla_2\mathbf1=L\mathbf1\), the compatible equations are \(d_1L=[P,L]\) and \(d_1L-d_2P=[P,L]\).}
\GSfield{Reconstruction}{The Ore multiplication rule is made explicit, especially because matrix coefficients need not commute. The notebook's initial \(d_1L=[L,P]\) is another admissible convention after replacing the evolution operator by its negative, but it cannot be mixed unchanged with the later positive connection definitions. The derivative-plus sign in its third equation is corrected in the final item.}
\GSfield{Fate}{The algebraic setup is exactly \citet[Chapter~III, \S2.1, p.~81]{manin1978algebraic}; the connections are treated in Theorems~2.5--2.6. This packages older Lax and zero-curvature equations into the speaker's subsequent bimodule exposition.}
\GSfield{Present status}{As of November 2026, the algebra and the convention-dependent equations are established; a consistent choice of signs is necessary.}
\GSfield{Assessment}{literal statement: correct with missing hypotheses; publication match: exact; verification: read the result; development: annotated; confidence: high}
\end{GSitem}
\begin{GSitem}{1976-410.12}{A}{Krichever--Drinfeld filtered bimodules}
\GSfield{Kind; source}{definition; \GSpage{P000682-P000683}}
\GSfield{The notes}{The structures are a left \(\mathcal B\)-module and right \(\mathcal O\)-module \(\mathcal M\), a distinguished \(\mathbf1\), an exhaustive filtration, and a connection \(\nabla\). The actions are faithful and commute; \(\mathcal M_{-1}=0\), \(\mathcal M_0=\mathcal B\mathbf1\), and every graded piece has rank one. The final line gives \(i\ge-1\) for both the filtration inclusion and the induced isomorphism.}
\GSfield{Statement}{Take \(\mathcal O\) to be the coordinate ring of a geometrically integral affine curve over \(k\). Require commuting faithful left \(\mathcal B\) and right \(\mathcal O\) actions, an exhaustive filtration by left submodules \(0=\mathcal M_{-1}\subset\mathcal M_0\subset\cdots\), with \(\mathbf1\) a free generator of \(\mathcal M_0\), and each \(\mathcal M_i/\mathcal M_{i-1}\) free of rank one for \(i\ge0\). Require \(\nabla(bm)=d(b)m+b\nabla m\), \(\nabla(m\varphi)=(\nabla m)\varphi\), \(\nabla\mathcal M_i\subset\mathcal M_{i+1}\), and induced isomorphisms \(\mathcal M_i/\mathcal M_{i-1}\to\mathcal M_{i+1}/\mathcal M_i\) for \(i\ge0\).}
\GSfield{Reconstruction}{The induced maps on graded pieces are \(\mathcal B\)-linear because the extra Leibniz term drops to the lower filtration level. The isomorphism cannot start at \(i=-1\): its source would be zero (or involve an undefined \(\mathcal M_{-2}\)) while its target is \(\mathcal M_0\ne0\). The degree-raising inclusion alone is harmless at \(-1\). Earlier-index continuity: 1975-385.1. These links identify a recurring statement or construction; only the explicitly indexed inverse problem reuses a problem ID.}
\GSfield{Fate}{These are the structures and axioms of \citet[Chapter~III, \S2.2, p.~81]{manin1978algebraic}, a close publication match to the notebook. The induced-isomorphism range is interpreted as \(i\ge0\); no nonexistent initial graded isomorphism is retained.}
\GSfield{Present status}{As of November 2026, the corrected axioms give the filtered differential-module model used in the reconstruction theorem.}
\GSfield{Assessment}{literal statement: false as written; publication match: exact; verification: read the result; development: annotated; confidence: high}
\end{GSitem}
\begin{GSitem}{1976-410.13}{A}{The regular bimodule attached to commuting operators}
\GSfield{Kind; source}{construction; \GSpage{P000683}}
\GSfield{The notes}{For commuting \(L,P\), with \(\bar k[L,P]\) having no zero divisors, the notes put \(\mathcal M=\mathcal B[D]\), \(\mathcal O=k[L,P]\), use left and right multiplication, filter by order, and set \(\nabla Q=DQ\).}
\GSfield{Statement}{Given an injective \(k\)-algebra homomorphism \(i:\mathcal O\hookrightarrow\mathcal B[D;d]\), where \(\mathcal O\) is the coordinate ring of a geometrically integral affine curve, take \(\mathcal M=\mathcal B[D;d]\), left multiplication by \(\mathcal B\), right multiplication by \(i(\mathcal O)\), \(\mathbf1=1\), \(\mathcal M_i=\{\sum_{j=0}^i b_jD^j\}\), and \(\nabla Q=DQ\). These satisfy all the preceding axioms.}
\GSfield{Reconstruction}{For the special choice \(\mathcal O=k[L,P]\), one must also ensure that it is an affine curve ring. Absence of zero divisors alone is not enough: taking scalar constant \(L,P\) gives dimension zero. Assuming a geometrically integral curve algebra at the outset avoids claiming an unjustified transcendence-degree assertion for arbitrary matrix operators.}
\GSfield{Fate}{This is exactly \citet[Chapter~III, \S2.2, p.~81]{manin1978algebraic}, which starts from an embedding of an already specified affine curve ring. The printed formulation supplies the missing dimension hypothesis in the notebook shorthand.}
\GSfield{Present status}{As of November 2026, the construction is established with the curve-algebra hypothesis; the domain condition by itself is insufficient.}
\GSfield{Assessment}{literal statement: correct with missing hypotheses; publication match: exact; verification: read the result; development: annotated; confidence: high}
\end{GSitem}
\begin{GSitem}{1976-410.14}{A}{Recovering the commuting algebra from a cyclic filtered module}
\GSfield{Kind; source}{theorem; \GSpage{P000683}}
\GSfield{The notes}{The first theorem recovers an embedding \(\mathcal O\to\mathcal B[D]\) from the bimodule. The diagrams first identify differential operators generated by \(\mathcal B\) and \(\nabla\), then define \(L_\varphi\mathbf1=\mathbf1\varphi\).}
\GSfield{Statement}{For a bimodule satisfying the corrected axioms above, the map \(\mathcal B[D;d]\to\mathcal M\), \(Q\mapsto Q(\nabla)\mathbf1\), is an isomorphism of filtered left modules. For every \(\varphi\in\mathcal O\) there is a unique \(L_\varphi\) with \(L_\varphi\mathbf1=\mathbf1\varphi\); the map \(\varphi\mapsto L_\varphi\) is an injective \(k\)-algebra homomorphism, and the regular-bimodule construction recovers \(\mathcal M\).}
\GSfield{Reconstruction}{The filtration shows inductively that \(\mathbf1,\nabla\mathbf1,\ldots,\nabla^i\mathbf1\) form a left \(\mathcal B\)-basis of \(\mathcal M_i\); this proves existence and uniqueness rather than assuming cyclicity. Right \(\mathcal O\)-linearity of \(\nabla\) makes the right action commute with every differential operator, and faithfulness proves injectivity. The diagrams express these maps; no extra unrecorded geometric correspondence is needed. Earlier-index continuity: 1975-383.10. These links identify a recurring statement or construction; only the explicitly indexed inverse problem reuses a problem ID.}
\GSfield{Fate}{The notebook's Theorem~1 becomes \citet[Chapter~III, Theorem~2.4, pp.~82--83]{manin1978algebraic}, including the observation that the cyclic-module part does not use \(\mathcal O\). This is an exact later publication of the speaker's presentation.}
\GSfield{Present status}{As of November 2026, the algebraic reconstruction is established for the stated axioms.}
\GSfield{Assessment}{literal statement: correct; publication match: exact; verification: read the result; development: none; confidence: high}
\end{GSitem}
\begin{GSitem}{1976-410.15}{A}{A second connection produces Lax evolution}
\GSfield{Kind; source}{theorem; \GSpage{P000683}}
\GSfield{The notes}{The second theorem introduces \(\nabla_1\) commuting with \(\nabla\) and the right \(\mathcal O\)-action, puts \(\nabla_1\mathbf1=P\mathbf1\), and indicates a Lax equation with an uncertain commutator sign. Its written Leibniz line is \(\nabla_1(bm)=\partial_1b+\nabla_1m\), omitting both multiplicative factors.}
\GSfield{Statement}{Let \(d_1\) commute with \(d\), and let \(\nabla_1\) satisfy \(\nabla_1(bm)=d_1(b)m+b\nabla_1m\), \([\nabla_1,\nabla]=0\), and right \(\mathcal O\)-linearity. Define \(P\) by \(\nabla_1\mathbf1=P\mathbf1\). Then every recovered \(L_\varphi\) satisfies \(d_1L_\varphi=[P,L_\varphi]\).}
\GSfield{Reconstruction}{Applying \(\nabla_1\) to \(L_\varphi\mathbf1=\mathbf1\varphi\) gives \((d_1L_\varphi+L_\varphi P)\mathbf1=PL_\varphi\mathbf1\); cyclic freeness removes \(\mathbf1\). This fixes the sign for the stated convention. The earlier notebook convention \(d_1L=[L,P]\) requires replacing its evolution operator by \(-P\). The missing Leibniz factors are a faithful error, not a translation issue.}
\GSfield{Fate}{This is \citet[Chapter~III, Theorem~2.5, p.~83]{manin1978algebraic}, with the exact sign \([P,L_\varphi]\) verified in the printed formula and proof.}
\GSfield{Present status}{As of November 2026, the theorem is established after restoring the Leibniz rule and fixing the evolution-operator convention.}
\GSfield{Assessment}{literal statement: false as written; publication match: exact; verification: read the result; development: annotated; confidence: high}
\end{GSitem}
\begin{GSitem}{1976-410.16}{A}{Three commuting connections and the shared zero-curvature sign error}
\GSfield{Kind; source}{theorem; \GSpage{P000684}}
\GSfield{The notes}{The third theorem drops all mention of \(\mathcal O\), assumes pairwise commuting \(\nabla,\nabla_1,\nabla_2\), puts \(\nabla_1\mathbf1=P\mathbf1\), \(\nabla_2\mathbf1=L\mathbf1\), and writes \(\partial_2P+\partial_1L=\pm[L,P]\). Both the plus sign and the explicit sign uncertainty are visible in the scan.}
\GSfield{Statement}{Let the filtered left \(\mathcal B\)-module, cyclic element and \(d\)-connection satisfy the preceding axioms with \(\mathcal O\) omitted. Let \(d,d_1,d_2\) commute and let their connections \(\nabla,\nabla_1,\nabla_2\) commute pairwise. Define \(P,L\) by the notebook's two equations. Then
 \[d_1L-d_2P=[P,L].\]
 Equivalently, \([d_1-P,d_2-L]=0\) as differential operators acting on a common auxiliary wavefunction.}
\GSfield{Reconstruction}{On the cyclic vector, flatness gives \(0=(d_1L+LP-d_2P-PL)\mathbf1\). Thus the derivative terms have opposite signs, and the bracket sign is fixed. As a concrete check, take \(\mathcal B=k[t,z]\), \(d=0\), \(\mathcal M=\mathcal B[D]\), \(\nabla=D\), \(\nabla_1=\partial_t+R_z\), \(\nabla_2=\partial_z+R_t\), where \(R\) is right multiplication. These connections commute and give \(P=z,L=t\); the correct difference is zero, whereas the notebook's sum is 2 and either commutator sign is zero.}
\GSfield{Fate}{\citet[Chapter~III, Theorem~2.6, p.~83]{manin1978algebraic} is an exact structural match, but its printed English statement also has a plus sign: with its swapped labels \(\nabla_1\mathbf1=L\mathbf1\), \(\nabla_2\mathbf1=P\mathbf1\), it prints \(d_1P+d_2L=[P,L]\). Its own flatness proof requires \(d_1P-d_2L=[L,P]\). The scan of that printed page was inspected. Thus publication confirms the theorem's provenance but does not validate the literal formula; this is a faithful error shared by the notebook and the cited translation.}
\GSfield{Present status}{As of November 2026, the corrected flatness identity is established; the literal plus-sign formula is false even for scalar commuting multiplication operators.}
\GSfield{Assessment}{literal statement: false as written; publication match: exact; verification: read the result; development: annotated; confidence: high}
\end{GSitem}

\GStalk{411}{1977-02-21}{J. J. Duistermaat}{P000685-P000689}{Morse indices, Lagrangian intersections, and the spherical trace formula}
The first part explains the finite Morse index of a regular variational problem with general endpoint constraints through a path of Lagrangian subspaces. The general-boundary theory was publicly presented in May 1974. The second part passes from spherical convolution to a distributional Selberg formula and joint spectral asymptotics; its identifiable continuation is the 1979 paper with Kolk and Varadarajan. The notebook suppresses transversality and normalization conventions, indexes orbital terms by elements rather than conjugacy classes, and ends with a remainder exponent not supported by the matched theorem. These are assessed separately.

\begin{GSitem}{1976-411.1}{E}{Variational problem with a submanifold of endpoint pairs}
\GSfield{Kind; source}{definition; \GSpage{P000685}}
\GSfield{The notes}{The action is \(I(\gamma)=\int_0^T f(\gamma,\dot\gamma)\,dt\), with \((\gamma(0),\gamma(T))\in R\subset\mathbb R^n\times\mathbb R^n\). Fixed endpoints and the diagonal \(R\) are examples.}
\GSfield{Statement}{Let \(R\) be smooth and \(f\) be \(C^2\) near a smooth critical path. Allowed variations satisfy \((v(0),v(T))\in T_{(\gamma(0),\gamma(T))}R\). Stationarity consists of the Euler--Lagrange equation and \((-p(0),p(T))\in N^*R\), where \(p=f_{\dot\gamma}\).}
\GSfield{Reconstruction}{The momentum condition and regularity assumptions are additions. An arbitrary smooth Lagrangian need not define a \(C^2\) action on all of \(H^1\); the Hessian along the smooth path is used here.}
\GSfield{Fate}{The setup appears in \citet[\S1]{duistermaat1974morse}.}
\GSfield{Present status}{As of November 2026: this is standard regular variational calculus.}
\GSfield{Assessment}{literal statement: correct with missing hypotheses; publication match: exact; verification: read the result; development: none; confidence: high}
\end{GSitem}
\begin{GSitem}{1976-411.2}{E}{Positive Legendre condition gives a finite Morse index}
\GSfield{Kind; source}{theorem; \GSpage{P000685}}
\GSfield{The notes}{Positivity of \(f_{\dot\gamma\dot\gamma}\) is said to make the number of negative squares finite.}
\GSfield{Statement}{On a compact time interval, suppose \(A(t)=f_{\dot\gamma\dot\gamma}(\gamma(t),\dot\gamma(t))\) is continuous and positive definite. The second variation on the closed \(H^1\) space of admissible endpoint variations, including the boundary term from \(R\), has finite negative index.}
\GSfield{Reconstruction}{Compactness makes the Legendre bound uniform. Lower-order terms and endpoint traces are compact relative to the positive derivative form after adding an \(L^2\) term.}
\GSfield{Fate}{See \citet[\S1]{duistermaat1974morse}.}
\GSfield{Present status}{As of November 2026: the assertion is valid under these hypotheses.}
\GSfield{Assessment}{literal statement: correct with missing hypotheses; publication match: exact; verification: read the result; development: none; confidence: high}
\end{GSitem}
\begin{GSitem}{1976-411.3}{E}{Legendre transform and the symplectic variational flow}
\GSfield{Kind; source}{construction; \GSpage{P000686}}
\GSfield{The notes}{Euler--Lagrange equations become Hamilton equations; \(\Phi(t)\) is the derivative of the flow along the extremal.}
\GSfield{Statement}{An invertible velocity Hessian gives the local Legendre transform \(p=f_v(q,v)\), \(H(q,p)=p\cdot v-f(q,v)\). In a symplectic trivialization along the extremal, the linearized flow satisfies \(\Phi(0)=1\) and \(\Phi(t)^*\sigma=\sigma\).}
\GSfield{Reconstruction}{Only local invertibility is needed.}
\GSfield{Fate}{See \citet[\S1]{duistermaat1974morse}.}
\GSfield{Present status}{As of November 2026: the Hamiltonian reduction is established.}
\GSfield{Assessment}{literal statement: correct with missing hypotheses; publication match: exact; verification: read the result; development: none; confidence: high}
\end{GSitem}
\begin{GSitem}{1976-411.4}{E}{Graphs of symplectic maps as Lagrangian planes}
\GSfield{Kind; source}{construction; \GSpage{P000686}}
\GSfield{The notes}{The page writes \(\varphi(t)=\operatorname{graph}\varphi(t)\) and says that \(\sigma\oplus(-\sigma)\) vanishes on the graph.}
\GSfield{Statement}{For symplectic \(V\) of dimension \(2n\), \(\Lambda(t)=\{(v,\Phi(t)v):v\in V\}\) is Lagrangian in \((V\oplus V,\sigma\oplus(-\sigma))\) exactly when \(\Phi(t)\) is symplectic.}
\GSfield{Reconstruction}{Upper-case \(\Phi\) is the map and \(\Lambda\) the subspace.}
\GSfield{Fate}{See \citet[\S\S1--3]{duistermaat1974morse}.}
\GSfield{Present status}{As of November 2026: the equivalence is elementary and exact.}
\GSfield{Assessment}{literal statement: correct; publication match: exact; verification: read the result; development: none; confidence: high}
\end{GSitem}
\begin{GSitem}{1976-411.5}{E}{Symmetric-form chart on the Lagrangian Grassmannian}
\GSfield{Kind; source}{construction; \GSpage{P000686}}
\GSfield{The notes}{For \(\alpha\cap\beta=0\) and \(\gamma\cap\beta=0\), write \(\gamma=\operatorname{graph}\Gamma\), \(\Gamma:\alpha\to\beta\), and \(Q(u,v)=\sigma(\Gamma u,v)\).}
\GSfield{Statement}{Relative to \(V=\alpha\oplus\beta\), every plane transverse to \(\beta\) is such a graph. It is Lagrangian exactly when \(Q\) is symmetric, and \(\gamma\cap\alpha=\ker Q\) under the graph identification.}
\GSfield{Reconstruction}{The phrase that \(\sigma\) vanishes on \(\Gamma\) means on its graph.}
\GSfield{Fate}{See \citet[\S2]{duistermaat1974morse}.}
\GSfield{Present status}{As of November 2026: the coordinate description is standard.}
\GSfield{Assessment}{literal statement: correct with missing hypotheses; publication match: exact; verification: read the result; development: none; confidence: high}
\end{GSitem}
\begin{GSitem}{1976-411.6}{E}{Maslov crossings and independence of the auxiliary complement}
\GSfield{Kind; source}{construction; \GSpage{P000686-P000687}}
\GSfield{The notes}{At a simple crossing \(\gamma(t)\cap\alpha\ne0\), take the jump of the quadratic-form index; it does not depend on the transverse \(\beta\).}
\GSfield{Statement}{For a Lagrangian path with endpoints transverse to \(\alpha\), signed regular crossings of \(\Sigma_\alpha=\{\Lambda:\Lambda\cap\alpha\ne0\}\) define an integer. Symmetric-form charts give compatible local counts; homotopies with fixed transverse endpoints preserve the integer.}
\GSfield{Reconstruction}{Perturbation defines nonregular crossings. A symplectic sign convention fixes the sign.}
\GSfield{Fate}{See \citet[\S2]{duistermaat1974morse}.}
\GSfield{Present status}{As of November 2026: the Maslov index is established; endpoint conventions remain essential.}
\GSfield{Assessment}{literal statement: correct with missing hypotheses; publication match: exact; verification: read the result; development: none; confidence: high}
\end{GSitem}
\begin{GSitem}{1976-411.7}{E}{Topology of the Lagrangian Grassmannian and the singular Maslov cycle}
\GSfield{Kind; source}{claim; \GSpage{P000687}}
\GSfield{The notes}{The notes call \(\Sigma_\alpha\) an embedded two-sided hypersurface and infer \(\pi_1(\operatorname{Lag})=\mathbb Z\); its complement is identified with quadratic forms.}
\GSfield{Statement}{For \(N\geq1\), \(\operatorname{Lag}(\mathbb R^{2N})=U(N)/O(N)\) has fundamental group \(\mathbb Z\). Transverse planes form a contractible symmetric-matrix chart. The one-dimensional-intersection stratum is a cooriented smooth hypersurface; higher-intersection strata are singular when \(N\geq2\).}
\GSfield{Reconstruction}{The whole cycle is not smooth. The complement description alone is not a proof of the fundamental group.}
\GSfield{Fate}{Compare \citet[\S2]{duistermaat1974morse}.}
\GSfield{Present status}{As of November 2026: the topological conclusion is correct, but the literal smoothness assertion is not.}
\GSfield{Assessment}{literal statement: false as written; publication match: partial; verification: read the result; development: annotated; confidence: high}
\end{GSitem}
\begin{GSitem}{1976-411.8}{E}{Endpoint correction for an open Lagrangian path}
\GSfield{Kind; source}{definition; \GSpage{P000687}}
\GSfield{The notes}{The proposed index is \([\gamma:\alpha]+\operatorname{ind}Q\{\gamma(T),\alpha,\gamma(0)\}\); in dimension one it is described by completing an arc on a circle.}
\GSfield{Statement}{An open-path Maslov index needs an endpoint normalization. For triples admitting the graph chart, a triple quadratic form measures the change of reference plane. The Hörmander identity supplies a reference-independent correction. Degenerate endpoints require a specified extension.}
\GSfield{Reconstruction}{The notes define neither all nontransverse triple forms nor an unambiguous orientation of the circle completion.}
\GSfield{Fate}{Compare \citet[\S2, formula (2.4)]{duistermaat1974morse}.}
\GSfield{Present status}{As of November 2026: the construction is known, but some notebook conventions are undetermined.}
\GSfield{Assessment}{literal statement: not determinable; publication match: partial; verification: read the result; development: annotated; confidence: high}
\end{GSitem}
\begin{GSitem}{1976-411.9}{E}{The conormal Lagrangian for a boundary condition}
\GSfield{Kind; source}{construction; \GSpage{P000687}}
\GSfield{The notes}{The notes define \(R^\perp\) with reversal of the second momentum, write \(\rho=T(R^\perp)\), and call it a fibre of the normal bundle.}
\GSfield{Statement}{Under \((x,y;\xi,\eta)\mapsto(x,\xi;y,-\eta)\), \(N^*R\) is Lagrangian in the doubled phase space. Its tangent at the endpoint covector is the required Lagrangian \(2n\)-plane \(\rho\).}
\GSfield{Reconstruction}{This is a tangent to the total conormal bundle; a fibre has dimension \(2n-\dim R\).}
\GSfield{Fate}{See \citet[\S1]{duistermaat1974morse}.}
\GSfield{Present status}{As of November 2026: the geometry is correct after the tangent/fibre distinction.}
\GSfield{Assessment}{literal statement: false as written; publication match: partial; verification: read the result; development: annotated; confidence: high}
\end{GSitem}
\begin{GSitem}{1976-411.10}{E}{General-boundary Morse--Maslov index formula}
\GSfield{Kind; source}{theorem; \GSpage{P000686-P000687}}
\GSfield{The notes}{The equality is \(\operatorname{ind}_R\gamma=\operatorname{ind}(\varphi)+i(I,\rho,\varphi(T))+\dim\pi(\rho\cap I)-2n\), with \(i=(\operatorname{ind}+\operatorname{null})Q\).}
\GSfield{Statement}{Under regularity and the positive Legendre condition, the Morse index is computable from flow graphs, the boundary conormal Lagrangian, and endpoint forms. The fixed-endpoint specialization is \(\operatorname{ind}\gamma=\sum_{0<t<T}\dim\{J:J\text{ is Jacobi},\,J(0)=J(t)=0\}\), with endpoint nullity separate.}
\GSfield{Reconstruction}{The general equality is retained as an announcement. Its absolute constant and triple-index convention cannot be recovered from the transverse-chart definition: \(\varphi(0)=I\), and \(\rho\cap I\) may be nonzero. The precise fixed-endpoint specialization is an editorial reconstruction.}
\GSfield{Fate}{See \citet[\S3, formulas (3.1)--(3.4)]{duistermaat1974morse}.}
\GSfield{Present status}{As of November 2026: the general theorem is established; the exact notebook normalization is partially matched.}
\GSfield{Assessment}{literal statement: not determinable; publication match: partial; verification: read the result; development: annotated; confidence: high}
\end{GSitem}
\begin{GSitem}{1976-411.11}{E}{Spherical convolution and kernels on a compact quotient}
\GSfield{Kind; source}{construction; \GSpage{P000687-P000688}}
\GSfield{The notes}{For \(X=K\backslash G/\Gamma\), the kernel is \(K_\varphi(x,y)=\sum_{\gamma\in\Gamma}\varphi(x\gamma y^{-1})\), with bi-\(K\)-invariant \(\varphi\).}
\GSfield{Statement}{Let \(G\) be connected semisimple with finite centre, \(K\) maximal compact, and \(\Gamma\) torsion-free cocompact. For \(\varphi\in C_c^\infty(K\backslash G/K)\), the locally finite kernel is smooth on \(X\times X\), hence trace class. Invariant differential operators arise from suitable distributions supported at the identity coset.}
\GSfield{Reconstruction}{Torsion-freeness makes \(X\) a manifold. In the bi-\(K\)-invariant realization on \(G\), identity-coset support means support on \(K\).}
\GSfield{Fate}{The setup is reviewed in \citet[\S\S2--3]{duistermaatkolkvaradarajan1979spectra}.}
\GSfield{Present status}{As of November 2026: the compact-quotient construction is established.}
\GSfield{Assessment}{literal statement: correct with missing hypotheses; publication match: exact; verification: read the result; development: none; confidence: high}
\end{GSitem}
\begin{GSitem}{1976-411.12}{E}{Commuting spherical operators and a common eigenbasis}
\GSfield{Kind; source}{theorem; \GSpage{P000688}}
\GSfield{The notes}{The spherical convolution algebra is commutative with involution; its compact operators have common orthonormal eigenvectors \(e_j\).}
\GSfield{Statement}{For the Riemannian symmetric pair \((G,K)\), bi-\(K\)-invariant convolution is commutative. On compact \(X\), spherical smoothing operators and an invariant elliptic Laplacian have a complete orthonormal simultaneous eigenbasis. Eigenvalue functionals are multiplicative characters.}
\GSfield{Reconstruction}{An elliptic operator or approximate identity supplies completeness; compactness alone would leave a possible common kernel.}
\GSfield{Fate}{See \citet[\S3]{duistermaatkolkvaradarajan1979spectra}.}
\GSfield{Present status}{As of November 2026: the joint decomposition is standard.}
\GSfield{Assessment}{literal statement: correct with missing hypotheses; publication match: exact; verification: read the result; development: none; confidence: high}
\end{GSitem}
\begin{GSitem}{1976-411.13}{E}{Trace formula and the grouping into conjugacy classes}
\GSfield{Kind; source}{theorem; \GSpage{P000688-P000689}}
\GSfield{The notes}{The notes equate \(\sum_j\lambda_j(\varphi)\), the diagonal integral, and a sum of individual \(\int_X\varphi(x\gamma x^{-1})\,dx\).}
\GSfield{Statement}{The trace is the diagonal integral. Its geometric expansion is a sum over \(\Gamma\)-conjugacy classes with centralizer-volume factors and orbital integrals. An individual \(\varphi(x\gamma x^{-1})\) generally does not descend to \(X\); an ungrouped sum can be integrated over a specified fundamental domain.}
\GSfield{Reconstruction}{Class regrouping repairs the quotient-invariance gap. One can instead integrate over a specified fundamental domain, but the notes do not say so.}
\GSfield{Fate}{The normalized side is in \citet[Theorem 5.1]{duistermaatkolkvaradarajan1979spectra}.}
\GSfield{Present status}{As of November 2026: the trace identity is established with this interpretation.}
\GSfield{Assessment}{literal statement: false as written; publication match: partial; verification: read the result; development: annotated; confidence: high}
\end{GSitem}
\begin{GSitem}{1976-411.14}{E}{Abel transform and the Euclidean spherical transform}
\GSfield{Kind; source}{construction; \GSpage{P000688}}
\GSfield{The notes}{The map is \(f(a)=e^{\langle\rho,a\rangle}\int_N\varphi(an)\,dn\); the margin calls \(A\) a vector space.}
\GSfield{Statement}{Use \(G=KAN\), write \(a=\exp H\), and put \(\rho=\tfrac12\sum_{\alpha>0}m_\alpha\alpha\). With compatible Haar measures, \(\mathcal A\varphi(H)=e^{\rho(H)}\int_N\varphi(\exp H\,n)\,dn\) identifies spherical convolution with convolution on \(C_c^\infty(\mathfrak a)^W\).}
\GSfield{Reconstruction}{Restricted-root multiplicities and exponential coordinates are additions.}
\GSfield{Fate}{See \citet[\S\S2--3]{duistermaatkolkvaradarajan1979spectra}.}
\GSfield{Present status}{As of November 2026: this is classical, with measure-dependent normalization.}
\GSfield{Assessment}{literal statement: correct with missing hypotheses; publication match: exact; verification: read the result; development: none; confidence: high}
\end{GSitem}
\begin{GSitem}{1976-411.15}{E}{Joint spectral parameters and multiplicities}
\GSfield{Kind; source}{definition; \GSpage{P000688-P000689}}
\GSfield{The notes}{Multiplicativity gives \(\lambda_j(\varphi)=\widetilde f(\mu_j)\), with \(\mu_j\in A_\mathbb C^*\); these form the spectrum.}
\GSfield{Statement}{A spherical eigencharacter is evaluation at a parameter in \(\mathfrak a_\mathbb C^*/W\). Joint eigenspaces are counted with dimension. Representatives before the Weyl quotient are noncanonical, and nonprincipal parameters need not lie on the real spectral axis.}
\GSfield{Reconstruction}{Multiplicities and the Weyl quotient are needed for the counting formula.}
\GSfield{Fate}{See \citet[\S3]{duistermaatkolkvaradarajan1979spectra}.}
\GSfield{Present status}{As of November 2026: the parametrization is established.}
\GSfield{Assessment}{literal statement: correct with missing hypotheses; publication match: exact; verification: read the result; development: none; confidence: high}
\end{GSitem}
\begin{GSitem}{1976-411.16}{A}{Poisson form of the spherical trace identity}
\GSfield{Kind; source}{theorem; \GSpage{P000689}}
\GSfield{The notes}{The trace becomes \(\sum_j\widetilde f(\mu_j)=\sum_\gamma\langle T_\gamma,f\rangle\), for a locally finite family.}
\GSfield{Statement}{For compact \(X\), the joint spectral distribution is a sum of \(W\)-invariant distributions indexed by conjugacy classes, with centralizer-volume weights. Nonidentity terms avoid a common neighbourhood of zero. The identity term is quotient volume times the Fourier transform of spherical Plancherel measure, with fixed Haar normalization.}
\GSfield{Reconstruction}{Weights and normalization are supplied from the paper.}
\GSfield{Fate}{See \citet[Theorem 5.1]{duistermaatkolkvaradarajan1979spectra}, received July 1978 and published 1979.}
\GSfield{Present status}{As of November 2026: the distributional formula is proved.}
\GSfield{Assessment}{literal statement: correct with missing hypotheses; publication match: exact; verification: read the result; development: none; confidence: high}
\end{GSitem}
\begin{GSitem}{1976-411.17}{A}{Closed-geodesic length and support of the geometric distribution}
\GSfield{Kind; source}{claim; \GSpage{P000689}}
\GSfield{The notes}{The boxed assertion reads \(\ell(\gamma)=\operatorname{dist}(\operatorname{supp}T_\gamma,0)\).}
\GSfield{Statement}{For nontrivial \(c\), the matched theorem places \(\operatorname{supp}T_c\) in a Weyl union of affine subspaces at distance equal to the minimal translation length \(\ell(c)\). For real regular Iwasawa classes, including rank one, \(T_c\) is a nonzero sum of Dirac masses on that Weyl orbit, giving equality of distances.}
\GSfield{Reconstruction}{A general support inclusion gives a lower bound, not itself equality. Equality is verified here in the explicit regular case.}
\GSfield{Fate}{See \citet[Theorems 5.1--5.2 and Proposition 5.15]{duistermaatkolkvaradarajan1979spectra}.}
\GSfield{Present status}{As of November 2026: the regular case is established; the universal equality is partially verified.}
\GSfield{Assessment}{literal statement: not determinable; publication match: partial; verification: read the result; development: annotated; confidence: high}
\end{GSitem}
\begin{GSitem}{1976-411.18}{A}{Joint Weyl law and the unsupported remainder exponent}
\GSfield{Kind; source}{claim; \GSpage{P000689}}
\GSfield{The notes}{The last line gives a ball count with main term \(\operatorname{vol}(X)\int_{\operatorname{ball}(R)}\beta\) and \(O(R^{n-r-1})\); the middle exponent letter is uncertain.}
\GSfield{Statement}{For bounded open principal spectral \(\Omega\) with smooth boundary, the matched result is \(\sum_{\lambda\in R\Omega}m(\lambda)=c_G\operatorname{vol}(X)\int_{R\Omega}\beta+O(R^{n-1})\), \(n=\dim X\). Complementary parameters are separately controlled; rank one gives \(O(R^{n-1}/\log R)\).}
\GSfield{Reconstruction}{These are verified conclusions, not a claim that the extra exponent symbol is absent from the scan.}
\GSfield{Fate}{See \citet[Theorems 8.8--8.9 and 9.1]{duistermaatkolkvaradarajan1979spectra}.}
\GSfield{Present status}{As of November 2026: the classical bounds are established; the stronger exponent was not verified.}
\GSfield{Assessment}{literal statement: not determinable; publication match: partial; verification: read the result; development: annotated; confidence: high}
\end{GSitem}

\GStalk{412}{1977-02-28}{Yu. I. Manin}{P000690-P000693}{Filtered bimodules, nodal spectral curves, and soliton equations}
This continuation recapitulates the bimodule reconstruction of commuting operators, Lax evolutions, and zero-curvature equations, then constructs the modules explicitly on an affine line with pairs of points identified. The block interpolation matrix, filtration, connection matrices, and operator orders closely match Chapter III, sections 2 and 5, of Manin's 1978 Russian exposition (English translation 1979). Several signs and indices in the notebook are demonstrably inconsistent with that construction. The final question about the behaviour of an \(N\)-soliton solution specifies neither an equation nor a limiting regime; it is retained without manufacturing a new open problem.

\begin{GSitem}{1976-412.1}{E}{Matrix differential-operator algebra and three compatibility equations}
\GSfield{Kind; source}{definition; \GSpage{P000690}}
\GSfield{The notes}{The coefficients are \(\mathcal B=M_\ell(\mathcal B_0)\), where \(\mathcal B_0\) consists of meromorphic germs in \(x\), \(x,t\), or \(x,t,y\). The targets are \([L,P]=0\), \(L_t=[P,L]\), and a displayed Zakharov--Shabat equation.}
\GSfield{Statement}{Form the Ore algebra \(\mathcal B[\partial_x]\) with \(\partial_x b=b\partial_x+\partial_x(b)\). The stationary equation is commutation; the Lax equation is compatibility of \(L\psi=\lambda\psi\) and \(\psi_t=P\psi\). Two evolution equations \(\psi_y=L\psi\), \(\psi_t=P\psi\) have compatibility \(L_t-P_y=[P,L]\).}
\GSfield{Reconstruction}{Commutators mean \([A,B]=AB-BA\). The sign correction is justified again in the dedicated zero-curvature item.}
\GSfield{Fate}{The framework is summarized in \citet[Ch.~III, \S2]{manin1978algebraic}.}
\GSfield{Present status}{As of November 2026: these operator equations are standard.}
\GSfield{Assessment}{literal statement: correct with missing hypotheses; publication match: exact; verification: read the result; development: none; confidence: high}
\end{GSitem}
\begin{GSitem}{1976-412.2}{A}{Filtered bimodule data and the omitted generation axioms}
\GSfield{Kind; source}{construction; \GSpage{P000690}}
\GSfield{The notes}{The notes specify a \((\mathcal B,\mathcal O)\)-bimodule \(M\), a filtration, a distinguished \(\mathbf1\), and \(\nabla(bm)=b_xm+b\nabla m\), referring backwards for further axioms.}
\GSfield{Statement}{A sufficient precise version is an exhaustive filtration \(M_0=0\subset M_1\subset\cdots\), \(M_1=\mathcal B\mathbf1\), with each quotient free of rank one, and \(\nabla M_i\subset M_{i+1}\) inducing isomorphisms on consecutive quotients. The right action of the commutative algebra \(\mathcal O\) commutes with \(\nabla\) and the left \(\mathcal B\)-action, and is faithful.}
\GSfield{Reconstruction}{These are the generation and commutation hypotheses needed below, supplied from the referenced earlier axioms and the matched exposition. The notebook alone does not repeat them.}
\GSfield{Fate}{See \citet[Ch.~III, \S2.2]{manin1978algebraic}, Russian original 1978, English 1979.}
\GSfield{Present status}{As of November 2026: the reconstruction data are established.}
\GSfield{Assessment}{literal statement: correct with missing hypotheses; publication match: exact; verification: read the result; development: none; confidence: high}
\end{GSitem}
\begin{GSitem}{1976-412.3}{A}{A free cyclic differential module produces commuting operators}
\GSfield{Kind; source}{theorem; \GSpage{P000690}}
\GSfield{The notes}{The lemma makes \(M\) a free rank-one \(\mathcal B[\partial_x]\)-module, with \(\partial_x\) acting as \(\nabla\); \(\mathbf1\varphi=L_\varphi\mathbf1\).}
\GSfield{Statement}{Under the preceding filtration axioms, \(\mathbf1,\nabla\mathbf1,\ldots,\nabla^{i-1}\mathbf1\) are a \(\mathcal B\)-basis of \(M_i\). Thus evaluation at \(\mathbf1\) identifies \(\mathcal B[\partial_x]\) with \(M\). Each right multiplication by \(\varphi\in\mathcal O\) determines \(L_\varphi\), and \(\varphi\mapsto L_\varphi\) is an injective algebra map; its image is commutative.}
\GSfield{Reconstruction}{Faithfulness gives injectivity. This does not say that arbitrary operators commuting with one chosen \(L\) automatically form this algebra. The scalar specialization \(\ell=1\) is the filtered-module reconstruction mechanism in GP-1975-26; the occurrence does not replace its pointed-curve and local-data hypotheses by a bare curve and bundle. Matrix coefficients and the later explicit nodal examples are additional scope, not a changed definition of that earlier problem.}
\GSfield{Fate}{The precise reconstruction is \citet[Ch.~III, Theorem 2.4]{manin1978algebraic}.}
\GSfield{Present status}{As of November 2026: the cyclic-module theorem is proved.}
\GSfield{Assessment}{literal statement: correct with missing hypotheses; publication match: exact; verification: read the result; development: none; confidence: high; problem \hyperref[prob:GP-1975-26]{GP-1975-26}}
\end{GSitem}
\begin{GSitem}{1976-412.4}{A}{A commuting time connection yields a Lax equation}
\GSfield{Kind; source}{theorem; \GSpage{P000690}}
\GSfield{The notes}{The notes add \(\nabla_t\), require \([\nabla_x,\nabla_t]=0\), and obtain \(L_t=[P,L]\).}
\GSfield{Statement}{Let \(\nabla_t\) extend \(\partial_t\), commute with \(\nabla_x\), and commute with the right \(\mathcal O\)-action. Define \(P\) by \(\nabla_t\mathbf1=P\mathbf1\). Then \(\partial_tL_\varphi=[P,L_\varphi]\).}
\GSfield{Reconstruction}{Compatibility with the spectral-ring action is needed in addition to the displayed commutator.}
\GSfield{Fate}{See \citet[Ch.~III, Theorem 2.5]{manin1978algebraic}.}
\GSfield{Present status}{As of November 2026: the stated construction is established.}
\GSfield{Assessment}{literal statement: correct with missing hypotheses; publication match: exact; verification: read the result; development: none; confidence: high}
\end{GSitem}
\begin{GSitem}{1976-412.5}{A}{Zero curvature and the sign in the Zakharov--Shabat equation}
\GSfield{Kind; source}{claim; \GSpage{P000690-P000691}}
\GSfield{The notes}{The page has \(L_t+P_y=[P,L]\), while \(\nabla_t\mathbf1=P\mathbf1\), \(\nabla_y\mathbf1=L\mathbf1\), and the three connections commute.}
\GSfield{Statement}{With the displayed connection convention and \([A,B]=AB-BA\), the resulting identity is \(L_t-P_y=[P,L]\). Indeed \([\nabla_t-P,\nabla_y-L]\mathbf1=(-L_t+P_y+[P,L])\mathbf1=0\); cyclic freeness makes the operator vanish.}
\GSfield{Reconstruction}{This is a mathematical emendation, not a claim that the scan has a minus. Independent visual inspection of the published English Theorem~2.6, p.~83, confirms a printed plus rather than an OCR artifact. Its connection conventions likewise require a difference; the notebook and publication both need a mathematical sign correction.}
\GSfield{Fate}{Compare \citet[Ch.~III, Theorem 2.6 and its proof]{manin1978algebraic}.}
\GSfield{Present status}{As of November 2026: zero curvature is established with the corrected sign.}
\GSfield{Assessment}{literal statement: false as written; publication match: partial; verification: read the result; development: annotated; confidence: high}
\end{GSitem}
\begin{GSitem}{1976-412.6}{A}{An affine nodal spectral curve by pairwise identification}
\GSfield{Kind; source}{construction; \GSpage{P000691}}
\GSfield{The notes}{The ring is \(\mathcal O=\{f\in k[\lambda]:f(\alpha_i)=f(\beta_i),\,1\leq i\leq N\}\).}
\GSfield{Statement}{For \(2N\) pairwise distinct complex numbers \(\alpha_i,\beta_i\), this subring is the coordinate ring of the affine line with the \(N\) pairs identified as ordinary nodes. Its normalization is \(\mathbb A^1\). For a real model the node set and identifications must be invariant under conjugation.}
\GSfield{Reconstruction}{Pairwise distinctness of all \(2N\) points is stronger than the notebook's explicitly written \(\alpha_i\ne\alpha_j\), and excludes unintended multiple identifications.}
\GSfield{Fate}{The exact construction is \citet[Ch.~III, \S5.1]{manin1978algebraic}.}
\GSfield{Present status}{As of November 2026: this spectral-curve construction is established.}
\GSfield{Assessment}{literal statement: correct with missing hypotheses; publication match: exact; verification: read the result; development: none; confidence: high}
\end{GSitem}
\begin{GSitem}{1976-412.7}{A}{Gluing a trivial module over the normalization}
\GSfield{Kind; source}{construction; \GSpage{P000691}}
\GSfield{The notes}{Inside \(N=\mathcal B[\lambda]^r\), the module is defined by \(m(\alpha_i)=g_i m(\beta_i)\); the left \(\mathcal B\)-action is componentwise right multiplication by \(b^t\).}
\GSfield{Statement}{Write the ambient module as \(\mathcal N=\mathcal B[\lambda]^r\) to distinguish it from the number of nodes. The left action \(b\cdot(m_j)=(m_jb^t)\) commutes with left block matrices \(g_i\in GL_r(\mathcal B_\mathbb C)\). The gluing equations define a left \(\mathcal B\), right \(\mathcal O\) sub-bimodule, the algebraic analogue of sections of a bundle with prescribed nodal gluing.}
\GSfield{Reconstruction}{The transpose is essential for associativity of this left action; it is not a transcription error.}
\GSfield{Fate}{See \citet[Ch.~III, \S\S5.1--5.2]{manin1978algebraic}.}
\GSfield{Present status}{As of November 2026: the module model is established.}
\GSfield{Assessment}{literal statement: correct; publication match: exact; verification: read the result; development: none; confidence: high}
\end{GSitem}
\begin{GSitem}{1976-412.8}{A}{Reality conditions at conjugate nodes}
\GSfield{Kind; source}{claim; \GSpage{P000691}}
\GSfield{The notes}{For real \(k\), the notes put \(\alpha_i=\overline{\beta_i}\) and \(\overline{g_i}=g_i^{-1}\).}
\GSfield{Statement}{Work first over the complexification of the real coefficient algebra. At conjugate paired nodes, conjugating \(m(\alpha_i)=g_i m(\beta_i)\) gives the same condition precisely when \(\overline{g_i}g_i=1\). Taking the fixed real subspace then descends the module.}
\GSfield{Reconstruction}{The points cannot all be required to lie in \(\mathbb R\) as well as to be distinct conjugate pairs. The scalar field in the initial point list must be complexified.}
\GSfield{Fate}{These conditions appear in \citet[Ch.~III, \S5.1]{manin1978algebraic}.}
\GSfield{Present status}{As of November 2026: the real descent condition is valid with complexified node parameters.}
\GSfield{Assessment}{literal statement: correct with missing hypotheses; publication match: exact; verification: read the result; development: none; confidence: high}
\end{GSitem}
\begin{GSitem}{1976-412.9}{A}{The staircase filtration on polynomial columns}
\GSfield{Kind; source}{definition; \GSpage{P000691}}
\GSfield{The notes}{In \(N_{ar+k}\), the first \(k\) components have degree at most \(a\), the remaining components degree at most \(a-1\), \(0\leq k<r\).}
\GSfield{Statement}{For \(a\geq0\), define \(\mathcal N_{ar+k}\) by those bounds, interpreting degree at most \(-1\) as zero. Then \(\mathcal N_0=0\), the filtration is exhaustive, and each successive quotient is a free left \(\mathcal B\)-module of rank one.}
\GSfield{Reconstruction}{The boundary convention \(a=0\) is explicit; it matters for the later shift.}
\GSfield{Fate}{See \citet[Ch.~III, \S5.3]{manin1978algebraic}.}
\GSfield{Present status}{As of November 2026: the filtration is established.}
\GSfield{Assessment}{literal statement: correct; publication match: exact; verification: read the result; development: none; confidence: high}
\end{GSitem}
\begin{GSitem}{1976-412.10}{A}{The block interpolation matrix and nondegeneracy}
\GSfield{Kind; source}{theorem; \GSpage{P000691-P000692}}
\GSfield{The notes}{The block matrix has entries \(\alpha_i^j I_r-\beta_i^jg_i\), \(i=1,\ldots,N\), \(j=0,\ldots,N-1\); the notebook requires nonzero determinant.}
\GSfield{Statement}{Assume \(G=(\alpha_i^jI_r-\beta_i^jg_i)_{i,j}\) is invertible in \(M_{rN}(\mathcal B)\). If \(M_{(i)}=M\cap\mathcal N_i\), then \(M_{(rN)}=0\), and \(M_{(i)}/M_{(i-1)}\to\mathcal N_i/\mathcal N_{i-1}\) is an isomorphism for \(i>rN\).}
\GSfield{Reconstruction}{Over a matrix coefficient algebra the clean condition is invertibility. A scalar determinant is taken only after expanding into matrices over the commutative coefficient ring and restricting to a locus where it is a unit.}
\GSfield{Fate}{See \citet[Ch.~III, \S5.4 and Lemma 5.5]{manin1978algebraic}.}
\GSfield{Present status}{As of November 2026: the interpolation lemma is proved.}
\GSfield{Assessment}{literal statement: correct with missing hypotheses; publication match: exact; verification: read the result; development: none; confidence: high}
\end{GSitem}
\begin{GSitem}{1976-412.11}{A}{The sign and indexing of the filtration shift}
\GSfield{Kind; source}{claim; \GSpage{P000692}}
\GSfield{The notes}{After \(M_{(rN)}=0\), the definition reads \(M_i=M_{(i-rN)}\), but the generator is then placed in \(M_1\).}
\GSfield{Statement}{With the notebook's one-based convention \(M_0=0\), \(M_1\) generated by \(\mathbf1\), the correct shift is \(M_i=M_{(i+rN)}\). With the paper's generator in \(M_0\), it is \(M_i=M_{(i+rN+1)}\).}
\GSfield{Reconstruction}{The negative shift leaves the supposed first generator space zero and is incompatible with the preceding vanishing lemma. The one-unit difference between conventions must not be mistaken for another error.}
\GSfield{Fate}{Compare \citet[Ch.~III, Lemma 5.5]{manin1978algebraic}.}
\GSfield{Present status}{As of November 2026: the corrected positive shift is definite.}
\GSfield{Assessment}{literal statement: false as written; publication match: partial; verification: read the result; development: annotated; confidence: high}
\end{GSitem}
\begin{GSitem}{1976-412.12}{A}{Normalized cyclic vector from finite linear equations}
\GSfield{Kind; source}{construction; \GSpage{P000692}}
\GSfield{The notes}{The normalized vector is \(\mathbf1=\sum_{j=0}^{N-1}m_j\lambda^j+\lambda^Ne_1\), with a displayed system using \(m_i\) inside the sum over \(j\).}
\GSfield{Statement}{Its coefficients uniquely solve \(\sum_{j=0}^{N-1}(\alpha_i^jI_r-\beta_i^jg_i)m_j=-(\alpha_i^NI_r-\beta_i^Ng_i)e_1\), \(1\leq i\leq N\), under invertibility of \(G\).}
\GSfield{Reconstruction}{The summation coefficient must be \(m_j\), not the row-indexed \(m_i\).}
\GSfield{Fate}{The normalized vector is \citet[Ch.~III, \S5.6, equation (5)]{manin1978algebraic}.}
\GSfield{Present status}{As of November 2026: the finite linear construction is established.}
\GSfield{Assessment}{literal statement: false as written; publication match: partial; verification: read the result; development: annotated; confidence: high}
\end{GSitem}
\begin{GSitem}{1976-412.13}{A}{Connection preserving the nodal gluing equations}
\GSfield{Kind; source}{theorem; \GSpage{P000692}}
\GSfield{The notes}{For \(\nabla_x=\partial_x+d_x(\lambda)\), the notes give \(\partial_xg_i=g_id_x(\beta_i)-d_x(\alpha_i)g_i\).}
\GSfield{Statement}{Let \(d_x\in M_r(\mathcal B[\lambda])\), acting on columns on the left. The displayed differential equations imply \(\nabla_xM\subset M\); differentiating the gluing relation proves this directly. Such matrices are \(\mathcal B[\lambda]\)-linear endomorphisms of the ambient module under the transpose action.}
\GSfield{Reconstruction}{The notebook's equality with all \(\mathcal B\)-linear endomorphisms is too broad: commuting with multiplication by \(\lambda\) is additionally required.}
\GSfield{Fate}{This is \citet[Ch.~III, Lemma 5.8]{manin1978algebraic}.}
\GSfield{Present status}{As of November 2026: the preservation criterion is established.}
\GSfield{Assessment}{literal statement: correct with missing hypotheses; publication match: exact; verification: read the result; development: none; confidence: high}
\end{GSitem}
\begin{GSitem}{1976-412.14}{A}{Explicit evolution of the gluing matrices}
\GSfield{Kind; source}{construction; \GSpage{P000692}}
\GSfield{The notes}{For constant \(d_x(\lambda)\), \(g_i(x)=e^{-xd_x(\alpha_i)}g_i(0)e^{xd_x(\beta_i)}\).}
\GSfield{Statement}{If \(d_x\) is independent of \(x\), the displayed formula solves the gluing-preservation equation for every invertible initial \(g_i(0)\). Matrix exponentials are required even when \(r=1\) if the coefficient algebra has matrix size \(\ell>1\).}
\GSfield{Reconstruction}{The marginal assertion that \(r=1\) gives ordinary scalar exponentials additionally requires \(\ell=1\), or scalar commuting coefficients.}
\GSfield{Fate}{See \citet[Ch.~III, \S5.9]{manin1978algebraic}.}
\GSfield{Present status}{As of November 2026: the explicit evolution is established.}
\GSfield{Assessment}{literal statement: correct with missing hypotheses; publication match: exact; verification: read the result; development: none; confidence: high}
\end{GSitem}
\begin{GSitem}{1976-412.15}{A}{Flatness of the parameter connections and two-time exponential signs}
\GSfield{Kind; source}{claim; \GSpage{P000692-P000693}}
\GSfield{The notes}{The notes write \(\partial_xd_t-\partial_td_x=[d_x,d_t]\), and later put \(+td_t(\alpha_i)\) in the left exponential as well as \(+td_t(\beta_i)\) on the right.}
\GSfield{Statement}{For \(\nabla_x=\partial_x+d_x\), \(\nabla_t=\partial_t+d_t\), flatness is \(\partial_xd_t-\partial_td_x+[d_x,d_t]=0\). For constant commuting matrices the compatible solution is \[g_i(x,t)=e^{-xd_x(\alpha_i)-td_t(\alpha_i)}g_i(0,0)e^{xd_x(\beta_i)+td_t(\beta_i)}.\]}
\GSfield{Reconstruction}{Both signs follow by differentiation. Choosing \(d_t,d_y\) as scalar polynomials in \(d_x\) ensures commutation; arbitrary constants need not commute.}
\GSfield{Fate}{Compare \citet[Ch.~III, \S5.12, equations (9)--(10)]{manin1978algebraic}.}
\GSfield{Present status}{As of November 2026: the equations hold with these corrected signs.}
\GSfield{Assessment}{literal statement: false as written; publication match: partial; verification: read the result; development: annotated; confidence: high}
\end{GSitem}
\begin{GSitem}{1976-412.16}{A}{Graded cyclicity and the order of a spectral operator}
\GSfield{Kind; source}{theorem; \GSpage{P000693}}
\GSfield{The notes}{The lemma \(\operatorname{ord}L_\varphi=r\deg\varphi\) ends with an unfinished condition. A matrix pattern and the example \(\begin{pmatrix}0&\lambda\\b&0\end{pmatrix}\) follow.}
\GSfield{Statement}{Assume \(\nabla_x\) increases filtration by one and is invertible on successive graded pieces. Multiplication by a polynomial \(\varphi\in\mathcal O\) of degree \(d\) raises the staircase filtration by exactly \(rd\), giving \(\operatorname{ord}L_\varphi=rd\). For \(r=2\), the displayed matrix has the required graded property when \(b\) is invertible.}
\GSfield{Reconstruction}{The unfinished hypothesis is supplied from the adjacent graded-isomorphism condition, not guessed solely from the desired order. For general \(r\), a cyclic lower subdiagonal with invertible entries and an invertible \(\lambda\)-coefficient in the upper-right corner suffices.}
\GSfield{Fate}{See \citet[Ch.~III, Lemma 5.10 and \S\S5.11--5.13]{manin1978algebraic}.}
\GSfield{Present status}{As of November 2026: the conditional order formula is proved.}
\GSfield{Assessment}{literal statement: correct with missing hypotheses; publication match: exact; verification: read the result; development: none; confidence: high}
\end{GSitem}
\begin{GSitem}{1976-412.17}{A}{A quadratic function on the nodal curve and an operator of order twice the rank}
\GSfield{Kind; source}{construction; \GSpage{P000693}}
\GSfield{The notes}{If \(\alpha_i+\beta_i\) is independent of \(i\), there is \(\psi\in\mathcal O\) of degree two, hence an operator of order \(2r\).}
\GSfield{Statement}{If every pair is distinct and \(\alpha_i+\beta_i=s\), then \(\psi(\lambda)=\lambda^2-s\lambda\) has equal values at paired points. Under the graded-cyclicity hypotheses, \(L_\psi\) has order \(2r\), including order four for \(r=2\). Conversely, a genuine quadratic polynomial constant on all pairs forces their sums to agree.}
\GSfield{Reconstruction}{The explicit polynomial and converse follow by factoring \(\psi(\alpha)-\psi(\beta)\).}
\GSfield{Fate}{The degree/order application appears in \citet[Ch.~III, \S5.13 and \S6]{manin1978algebraic}.}
\GSfield{Present status}{As of November 2026: the construction is established.}
\GSfield{Assessment}{literal statement: correct; publication match: exact; verification: read the result; development: none; confidence: high}
\end{GSitem}
\begin{GSitem}{1976-412.18}{E}{The rank-one nodal construction includes KdV multisolitons}
\GSfield{Kind; source}{construction; \GSpage{P000691-P000693}}
\GSfield{The notes}{The notes identify the rank-one case with \(N\)-soliton solutions of KdV and its hierarchy.}
\GSfield{Statement}{Take scalar coefficients \(\ell=1\), spectral rank \(r=1\), paired spectral points with a common sum, and the nodal gluing data of the preceding construction. On any parameter domain where the interpolation determinant is nonzero, the quadratic spectral function gives a second-order scalar differential operator. Choosing the compatible cubic time flow yields the KdV Lax equation after the usual normalization; its nodal specialization includes the classical multisoliton family.}
\GSfield{Reconstruction}{This does not assert that every arbitrary gluing matrix yields a globally smooth real soliton. The real multisoliton subclass needs real spectral choices and a nonvanishing real denominator; the notebook states neither global regularity nor those restrictions.}
\GSfield{Fate}{The identification is explicit in \citet[Ch.~III, \S5.13]{manin1978algebraic}. The rank-one multisoliton family is treated there as known, whereas the uniform higher-rank module construction is the newer part.}
\GSfield{Present status}{As of November 2026: the stated local construction includes the established rank-one soliton family.}
\GSfield{Assessment}{literal statement: correct with missing hypotheses; publication match: exact; verification: read the result; development: none; confidence: high}
\end{GSitem}
\begin{GSitem}{1976-412.19}{F}{The unqualified question about N-soliton behaviour}
\GSfield{Kind; source}{question; \GSpage{P000693}}
\GSfield{The notes}{The talk ends with \ru{Вопрос: поведение \(N\)-солит. решения?}.}
\GSfield{Statement}{No equation, spectral rank, reality or nonsingularity conditions, time direction, or requested observable is specified in the final question.}
\GSfield{Reconstruction}{It cannot be equated with the earlier two-soliton phase-shift problem GP-1973-54 or promoted to a new higher-rank scattering conjecture merely because of the preceding examples.}
\GSfield{Fate}{The surrounding soliton construction is matched by \citet[Ch.~III, \S5.13 and \S6]{manin1978algebraic}, but no definite final question was identified.}
\GSfield{Present status}{As of November 2026: the final question is too underspecified to support an open-problem claim.}
\GSfield{Assessment}{literal statement: not determinable; publication match: thematic only; verification: read the result; development: annotated; confidence: medium}
\end{GSitem}

\GStalk{413}{1977-03-14}{A. M. Gabrielov}{P000694-P000696}{Pontryagin cycles, changes of sections, and hypersimplicial chains}
Gelfand motivates a combinatorial theory by local Euler-characteristic formulas and the difficulty of noncommuting rotations. Gabrielov then describes the first Pontryagin class by a corank-two degeneracy locus, reduces the chain calculations to pairs of sections of a real plane bundle, and organizes their cobordisms by hypersimplices. The notes identify MacPherson's February 1977 Bourbaki exposition, which was published in the 1978 proceedings. The core constructions were already in the two Gabrielov--Gelfand--Losik papers of 1975. The final diagrams sketch low-dimensional chain assignments rather than a new, fully specified open problem.

\begin{GSitem}{1976-413.1}{F}{The programme of an elementary geometry of polyhedra}
\GSfield{Kind; source}{remark; \GSpage{P000694}; I. M. Gelfand}
\GSfield{The notes}{Gelfand says that an elementary geometry of higher-dimensional polyhedra remains to be made; he contrasts commutative planar rotations with noncommutative \(SO(3)\).}
\GSfield{Statement}{The contrast concerns the search for local combinatorial characteristic-class formulas. Noncommutativity of \(SO(3)\) is a definite fact, but the assertion that three-dimensional polyhedra have no substantive theory specifies no mathematical obstruction or question.}
\GSfield{Reconstruction}{No universal nonexistence statement or orphan problem is extracted from this programme. The claim that named mathematicians learned from the GGL work is recorded as Gelfand's recollection, not independently established intellectual genealogy.}
\GSfield{Fate}{The contemporary programme is described in \citet[introduction, p.~105]{macpherson1978combinatorial}; it is not an exact theorem matching the rhetoric.}
\GSfield{Present status}{As of November 2026: the programme is identifiable, but no definite open proposition is stated.}
\GSfield{Assessment}{literal statement: not determinable; publication match: thematic only; verification: read the result; development: none; confidence: medium}
\end{GSitem}
\begin{GSitem}{1976-413.2}{E}{The two-faces condition does not characterize a manifold}
\GSfield{Kind; source}{claim; \GSpage{P000694}; I. M. Gelfand}
\GSfield{The notes}{Every \((n-1)\)-simplex is said to bound exactly two \(n\)-simplices.}
\GSfield{Statement}{For a finite pure simplicial complex, this is the codimension-one incidence condition of a pseudomanifold without boundary. To obtain a PL manifold one must also impose the appropriate sphere links at every simplex; incidence alone is insufficient.}
\GSfield{Reconstruction}{The notebook's use of this condition as a description of the intended polyhedral manifold needs the link condition. MacPherson additionally formulates his flattenability hypotheses.}
\GSfield{Fate}{Compare \citet[introduction, pp.~105--106]{macpherson1978combinatorial}.}
\GSfield{Present status}{As of November 2026: the distinction between pseudomanifolds and PL manifolds is standard.}
\GSfield{Assessment}{literal statement: correct with missing hypotheses; publication match: exact; verification: read the result; development: annotated; confidence: high}
\end{GSitem}
\begin{GSitem}{1976-413.3}{E}{Assigning local vertex weights for the Euler characteristic}
\GSfield{Kind; source}{question; \GSpage{P000694}; I. M. Gelfand}
\GSfield{The notes}{Gelfand asks how numbers at vertices can sum to the Euler characteristic and invokes angular defects.}
\GSfield{Statement}{For a finite simplicial complex \(K\), the vertex weights
 \[w(v)=\sum_{\sigma\ni v}\frac{(-1)^{\dim\sigma}}{\dim\sigma+1}\]
 satisfy \(\sum_vw(v)=\chi(K)\). For a closed piecewise Euclidean triangulated surface, the angle defects
 \[\kappa(v)=2\pi-\sum_{\sigma\ni v}\theta_{\sigma,v}\]
 satisfy \(\sum_v\kappa(v)=2\pi\chi(K)\).}
\GSfield{Reconstruction}{These are precise reconstructions of the local-counting question. The notebook's external-angle sentence does not identify which angles it sums and should not be read as the displayed \(2\pi\) formula without explanation.}
\GSfield{Fate}{The classical Euler formula is recalled in \citet[introduction, p.~105]{macpherson1978combinatorial}.}
\GSfield{Present status}{As of November 2026: the definite local Euler-characteristic question has elementary established answers.}
\GSfield{Assessment}{literal statement: not determinable; publication match: partial; verification: read the result; development: annotated; confidence: high}
\end{GSitem}
\begin{GSitem}{1976-413.4}{E}{Two degree-four characteristic numbers}
\GSfield{Kind; source}{claim; \GSpage{P000694}; I. M. Gelfand}
\GSfield{The notes}{Gelfand says that four-dimensional polyhedra have two invariants, the Euler characteristic and one more.}
\GSfield{Statement}{For a closed oriented smooth four-manifold, the characteristic numbers in view are \(\langle e(TX),[X]\rangle=\chi(X)\) and \(\langle p_1(TX),[X]\rangle=3\,\operatorname{sign}(X)\). These do not classify four-manifolds and are not their only topological invariants.}
\GSfield{Reconstruction}{The statement is interpreted within the local characteristic-class programme, not as a classification claim.}
\GSfield{Fate}{The Euler/Pontryagin context is explicit in \citet[introduction, p.~105]{macpherson1978combinatorial}.}
\GSfield{Present status}{As of November 2026: the characteristic-number interpretation is established.}
\GSfield{Assessment}{literal statement: correct with missing hypotheses; publication match: partial; verification: read the result; development: annotated; confidence: high}
\end{GSitem}
\begin{GSitem}{1976-413.5}{E}{The first Pontryagin cycle as a corank-two degeneracy locus}
\GSfield{Kind; source}{construction; \GSpage{P000694-P000695}}
\GSfield{The notes}{Four generic vector fields on a four-manifold give \(Z=\{x:\operatorname{rk}(e_1,e_2,e_3,e_4)\leq2\}\), a codimension-four signed set.}
\GSfield{Statement}{Let \(E\) be a real rank-four bundle on a smooth closed oriented four-manifold. A generic bundle map \(\mathbb R^4\to E\) meets the rank-two stratum transversely and avoids lower ranks. The resulting zero-cycle, with the standard determinantal coorientation, represents the Poincaré dual of \(p_1(E)\). Apply this to \(E=TX\).}
\GSfield{Reconstruction}{Genericity means transversality to rank strata; the rank-at-most-one locus has codimension nine and is avoided. The sign convention is normalized by \(p_1\) of an oriented plane bundle being its Euler class squared.}
\GSfield{Fate}{The degeneracy-cycle framework underlies \citet[\S1]{gabrielovgelfandlosik1975computation}.}
\GSfield{Present status}{As of November 2026: the construction is classical and was expository here.}
\GSfield{Assessment}{literal statement: correct with missing hypotheses; publication match: exact; verification: read the result; development: none; confidence: high}
\end{GSitem}
\begin{GSitem}{1976-413.6}{E}{Canonical coorientation and independence of the signed count}
\GSfield{Kind; source}{construction; \GSpage{P000695}}
\GSfield{The notes}{The rank degeneracy set in \(\operatorname{Hom}(\mathbb R^m,\mathbb R^n)\) is said to carry a canonical coorientation; the signed count does not depend on the chosen fields.}
\GSfield{Statement}{For a square map of corank two, the normal space to the rank stratum is \(\operatorname{Hom}(\ker A,\operatorname{coker}A)\), with both factors two-dimensional. Its determinant transforms by even powers, so a fixed universal orientation convention is independent of bases. A generic homotopy of sections gives a compact oriented cobordism of the zero-cycles.}
\GSfield{Reconstruction}{The assertion is restricted to the corank-two square situation actually used; arbitrary determinantal strata need not have a canonical coorientation.}
\GSfield{Fate}{See the section-change framework of \citet[\S1]{gabrielovgelfandlosik1975computation}.}
\GSfield{Present status}{As of November 2026: the signed homology class is independent of the generic fields.}
\GSfield{Assessment}{literal statement: correct with missing hypotheses; publication match: exact; verification: read the result; development: none; confidence: high}
\end{GSitem}
\begin{GSitem}{1976-413.7}{E}{Two sections of a plane bundle represent the square of the Euler class}
\GSfield{Kind; source}{construction; \GSpage{P000695}}
\GSfield{The notes}{For a two-dimensional real bundle \(E\), \(Z=\{e_1=e_2=0\}\) is identified with its first Pontryagin cycle.}
\GSfield{Statement}{For an oriented real rank-two bundle over a closed oriented four-manifold, a transverse section \((e_1,e_2)\) of \(E\oplus E\) has zero-cycle dual to \(e(E\oplus E)=e(E)^2=p_1(E)\).}
\GSfield{Reconstruction}{Orientation is supplied so that the ordinary integral Euler class notation is unambiguous. The canonical orientation of \(E\oplus E\) also admits a formulation without orienting \(E\), but that extension is not needed here.}
\GSfield{Fate}{This is the plane-bundle model used in \citet[\S1]{gabrielovgelfandlosik1975computation}.}
\GSfield{Present status}{As of November 2026: the model is established.}
\GSfield{Assessment}{literal statement: correct with missing hypotheses; publication match: exact; verification: read the result; development: none; confidence: high}
\end{GSitem}
\begin{GSitem}{1976-413.8}{E}{A one-chain relating cycles after replacing one section}
\GSfield{Kind; source}{construction; \GSpage{P000695}}
\GSfield{The notes}{The first identity is \(Z(e_2,e_3)-Z(e_1,e_3)=\partial Z^1(e_1,e_2;e_3)\), with \(e_3=0\), \(e_1=\lambda e_2\), \(\lambda>0\).}
\GSfield{Statement}{For generic sections of a real plane bundle, use the compactified parameter \(\lambda\in[0,\infty]\). The incidence locus \(e_3=0,\ e_1=\lambda e_2\), projected as an oriented chain to the base, has boundary the difference of the two zero-cycles, with sign fixed by the parameter orientation.}
\GSfield{Reconstruction}{The chain is the closure or pushforward of the parameterized locus; the open set with \(\lambda>0\) alone has no included boundary points. Transversality and compactification are explicit additions.}
\GSfield{Fate}{The section-change calculus is developed in \citet{gabrielovgelfandlosik1975calculus}.}
\GSfield{Present status}{As of November 2026: the cobordism identity is established after orientation conventions.}
\GSfield{Assessment}{literal statement: correct with missing hypotheses; publication match: exact; verification: read the result; development: none; confidence: high}
\end{GSitem}
\begin{GSitem}{1976-413.9}{E}{A two-chain filling the cyclic sum of section changes}
\GSfield{Kind; source}{construction; \GSpage{P000695}}
\GSfield{The notes}{Adding three displayed one-chain boundaries gives a closed curve; it is said to bound the locus \(e_1=\lambda e_2=\mu e_3\), \(\lambda,\mu>0\).}
\GSfield{Statement}{For generic triples, compactify the positive proportionality parameters and take the corresponding oriented incidence chain. Its codimension-two projection has boundary equal to the cyclic sum of the three codimension-three section-change chains, with compatible signs.}
\GSfield{Reconstruction}{The parameter compactification and signs are necessary for a chain identity; the notebook gives the geometric idea rather than a complete orientation convention.}
\GSfield{Fate}{This is part of the published hypersimplicial chain calculus \citep{gabrielovgelfandlosik1975calculus}.}
\GSfield{Present status}{As of November 2026: the geometric chain relation is established.}
\GSfield{Assessment}{literal statement: correct with missing hypotheses; publication match: partial; verification: read the result; development: none; confidence: high}
\end{GSitem}
\begin{GSitem}{1976-413.10}{E}{Hypersimplices and their two kinds of triangular faces}
\GSfield{Kind; source}{definition; \GSpage{P000695}}
\GSfield{The notes}{A hypersimplex is the convex hull of centres of faces of a simplex. Vertices carry index sets; the octahedron has triangular faces of two types.}
\GSfield{Statement}{The standard hypersimplex is \(\Delta(k,m)=\{x\in[0,1]^m:\sum_i x_i=k\}\), the convex hull of indicator vectors of \(k\)-subsets. Scaling by \(1/k\) identifies these with barycentres of \((k-1)\)-faces. Its faces are obtained by fixing coordinates to zero or one; \(\Delta(2,4)\) is an octahedron whose triangular faces arise in the two ways.}
\GSfield{Reconstruction}{The notebook's appended fixed indices correspond to coordinates fixed to one.}
\GSfield{Fate}{See \citet[\S1, pp.~107--109]{macpherson1978combinatorial}.}
\GSfield{Present status}{As of November 2026: the combinatorial definition is standard.}
\GSfield{Assessment}{literal statement: correct; publication match: exact; verification: read the result; development: none; confidence: high}
\end{GSitem}
\begin{GSitem}{1976-413.11}{E}{Boundary-compatible chains assigned to hypersimplices}
\GSfield{Kind; source}{problem; \GSpage{P000695-P000696}}
\GSfield{The notes}{The task is to construct a chain for each hypersimplex with the required properties; vertices represent double-zero loci and edges represent the one-section changes.}
\GSfield{Statement}{Construct rational chains \(C(H)\) in spaces of vector configurations, extending the specified vertex and edge assignments, so that \(\partial C(H)=C(\partial H)\), with the appropriate codimension and compatibility under restriction to faces. In the first-Pontryagin construction these chains must descend through the configuration and hypersimplicial complexes.}
\GSfield{Reconstruction}{The boundary equation is the strongest definite compatibility indicated by the preceding examples. The full list of admissibility and locality properties is supplied by the published construction, not repeated by the notebook.}
\GSfield{Fate}{The chain systems already exist in \citet[\S18, Proposition 14; \S\S27--28]{gabrielovgelfandlosik1975calculus}; compare \citet[\S3]{macpherson1978combinatorial}.}
\GSfield{Present status}{As of November 2026: the first-Pontryagin chain problem was solved before this talk; no new unsolved variant is specified.}
\GSfield{Assessment}{literal statement: correct with missing hypotheses; publication match: partial; verification: read the result; development: none; confidence: high; problem \hyperref[prob:GP-1976-20]{GP-1976-20}}
\end{GSitem}
\begin{GSitem}{1976-413.12}{F}{The two triangular chain assignments and the handwritten one-half}
\GSfield{Kind; source}{remark; \GSpage{P000696}}
\GSfield{The notes}{One triangle with labels \(14,24,34\) points to \(\tfrac12(e_4=0)\); the other, labelled \(12,13,23\), points to positive proportionality of \(e_1,e_2,e_3\).}
\GSfield{Statement}{These drawings distinguish the two triangular face types. The second has the incidence interpretation in the preceding two-chain item. The first shows a rational coefficient \(1/2\), but the notes do not specify the oriented region, multiplicity function, or chain category needed to interpret \(\tfrac12\{e_4=0\}\) as a complete boundary-compatible chain.}
\GSfield{Reconstruction}{The PNG confirms the coefficient. It is not corrected to an integer and is not promoted to a fully defined chain formula.}
\GSfield{Fate}{The surrounding rational chain construction is identifiable in \citet[\S3]{macpherson1978combinatorial}, but the isolated drawn assignment is only partially specified.}
\GSfield{Present status}{As of November 2026: the diagram belongs to a known construction; its literal chain convention remains undetermined.}
\GSfield{Assessment}{literal statement: not determinable; publication match: partial; verification: read the result; development: annotated; confidence: medium}
\end{GSitem}

\GStalk{414}{1977-03-21}{M. I. Graev}{P000697-P000698}{Range equations, differential-form inversion, and complexes of lines}
Graev develops the complex line transform in three-space, its range equations, and inversion by integrating a closed \((1,1)\)-form over a cycle of directions. This matches the explicitly cited 1967 Gelfand--Graev--Shapiro paper. The classification of admissible complexes matches the 1968 Gelfand--Graev paper, including its algebraic and genericity conditions. The final discussion of linear complexes is a separate projective-geometric construction: nondegenerate and special complexes must be distinguished, and it is not automatically an example of the preceding admissibility theorem.

\begin{GSitem}{1976-414.1}{E}{The affine complex line transform}
\GSfield{Kind; source}{definition; \GSpage{P000697}}
\GSfield{The notes}{A line is \(x_1=\alpha_1x_3+\beta_1,\ x_2=\alpha_2x_3+\beta_2\), and \(\varphi(\alpha,\beta)\) is the integral of \(f\) against \(dx_3\,d\bar x_3\).}
\GSfield{Statement}{For \(f\in\mathcal S(\mathbb C^3)\), define \(\varphi(\alpha,\beta)=\int_{\mathbb C}f(\alpha_1t+\beta_1,\alpha_2t+\beta_2,t)\,dA(t)\). This chart excludes directions parallel to \(x_3=0\); changing the line parameter transforms the density by the corresponding squared Jacobian.}
\GSfield{Reconstruction}{The area-measure convention and chart limitation are explicit additions. A global scalar realization uses the compensating density factor of the cited paper.}
\GSfield{Fate}{See \citet[\S1, equations (1)--(4)]{gelfand1967planes}.}
\GSfield{Present status}{As of November 2026: the line transform is classical.}
\GSfield{Assessment}{literal statement: correct with missing hypotheses; publication match: exact; verification: read the result; development: none; confidence: high}
\end{GSitem}
\begin{GSitem}{1976-414.2}{E}{Reconstruction from all line integrals and reduction to a plane}
\GSfield{Kind; source}{question; \GSpage{P000697}}
\GSfield{The notes}{The notes ask how to recover \(f\) from \(\varphi\), call the data redundant, and restrict to lines in a plane.}
\GSfield{Statement}{For a fixed complex two-plane through \(x\), the supplied line integrals over its lines are the complex Radon transform of \(f\) restricted to that plane. Its inversion recovers \(f(x)\); taking all such planes gives uniqueness of the full line transform on Schwartz functions.}
\GSfield{Reconstruction}{The restriction is to lines contained in the chosen plane, not a projection of \(f\) onto that plane.}
\GSfield{Fate}{The precise cycle inversion is \citet[\S8, Theorem 2]{gelfand1967planes}.}
\GSfield{Present status}{As of November 2026: the reconstruction problem is solved.}
\GSfield{Assessment}{literal statement: correct with missing hypotheses; publication match: exact; verification: read the result; development: none; confidence: high; problem \hyperref[prob:GP-1976-21]{GP-1976-21}}
\end{GSitem}
\begin{GSitem}{1976-414.3}{E}{John equation and the missing global range hypotheses}
\GSfield{Kind; source}{claim; \GSpage{P000697}}
\GSfield{The notes}{In the real case \(\varphi_{\alpha_1\beta_2}=\varphi_{\alpha_2\beta_1}\), and the equation is said to be sufficient.}
\GSfield{Statement}{For the real line transform in \(\mathbb R^3\), differentiation under the integral gives this equation. It characterizes the range within the global smooth line-data class whose density-normalized derivatives decay faster than every inverse power of the line's distance from zero, uniformly in compact direction charts, with the correct coordinate covariance.}
\GSfield{Reconstruction}{A local smooth solution in one affine chart, with no decay or global extension, is not covered by sufficiency. The condition \(k<n-1\) in the paper holds for \(k=1,n=3\).}
\GSfield{Fate}{See \citet[\S7, Theorem 1 and its real-case footnote]{gelfand1967planes}.}
\GSfield{Present status}{As of November 2026: the range theorem is established in its stated global class.}
\GSfield{Assessment}{literal statement: correct with missing hypotheses; publication match: exact; verification: read the result; development: none; confidence: high}
\end{GSitem}
\begin{GSitem}{1976-414.4}{F}{A wavefront-set heuristic from singular plane integrals}
\GSfield{Kind; source}{method; \GSpage{P000697}}
\GSfield{The notes}{The page suggests that \((x_0,\xi_0)\in WF(u)\) when integration over parallel planes in direction \(\xi_0\) has a singularity.}
\GSfield{Statement}{For a localized distribution, Radon data can detect conormal singularities through the incidence canonical relation. An unlocalized integral over a whole plane does not identify which point caused a singularity, and cancellations can erase contributions. Thus a pointwise criterion requires localization near \(x_0\), a convention making \(\xi_0\) a normal covector, and a precise microlocal implication.}
\GSfield{Reconstruction}{These are necessary qualifications, not a theorem present in the notebook. No unique equivalence is reconstructed from the single sentence.}
\GSfield{Fate}{The related characteristic geometry is explained in \citet[\S\S2--3]{guillemin1985line}; the notebook's particular wavefront criterion is not matched.}
\GSfield{Present status}{As of November 2026: the heuristic is recognizable, but its literal pointwise assertion is insufficiently specified.}
\GSfield{Assessment}{literal statement: not determinable; publication match: thematic only; verification: read the result; development: annotated; confidence: medium}
\end{GSitem}
\begin{GSitem}{1976-414.5}{E}{Differential forms on the space of lines through a point}
\GSfield{Kind; source}{construction; \GSpage{P000697}}
\GSfield{The notes}{In the complex case \(\mathcal H\varphi=\sum_{i,j}\varphi_{\beta_i\bar\beta_j}\,d\alpha_i\wedge d\bar\alpha_j\); in the real case \(\mathcal H\varphi=\sum_i\varphi_{\beta_i}\,d\alpha_i\). Substitute \(\beta_i=x_i-\alpha_i x_3\).}
\GSfield{Statement}{The substitution pulls the relevant differentiated data back to the pencil of lines through \(x\). In the indicated chart it gives a complex \((1,1)\)-form, or a real one-form. With the line-density transformation law these expressions patch to forms on the projective space of directions.}
\GSfield{Reconstruction}{Global patching and measure normalization are implicit in the handwritten chart formula.}
\GSfield{Fate}{See \citet[\S\S4--6]{gelfand1967planes}; differential-form generalization in \citet[\S1]{gelfand1969forms}.}
\GSfield{Present status}{As of November 2026: the differential-form construction is established.}
\GSfield{Assessment}{literal statement: correct with missing hypotheses; publication match: exact; verification: read the result; development: none; confidence: high}
\end{GSitem}
\begin{GSitem}{1976-414.6}{E}{The closed-form characterization of the complex line-transform range}
\GSfield{Kind; source}{theorem; \GSpage{P000697}}
\GSfield{The notes}{The theorem is \(\varphi=Rf\Longleftrightarrow d(\mathcal H_x\varphi)=0\) for every \(x\).}
\GSfield{Statement}{Let \(\varphi\) be globally smooth normalized line data in the rapidly decreasing class: all derivatives in local vector-bundle trivializations of the affine Grassmannian decay rapidly in the transverse displacement, uniformly over directions. Then \(\varphi=Rf\) for some \(f\in\mathcal S(\mathbb C^3)\) exactly when the associated \((1,1)\)-form is closed on the pencil through every \(x\).}
\GSfield{Reconstruction}{The global class is essential and is not contained in the notebook's displayed equivalence.}
\GSfield{Fate}{This is \citet[\S7, Theorem 1]{gelfand1967planes}, published ten years before the seminar.}
\GSfield{Present status}{As of November 2026: the characterization is proved.}
\GSfield{Assessment}{literal statement: correct with missing hypotheses; publication match: exact; verification: read the result; development: none; confidence: high; problem \hyperref[prob:GP-1976-21]{GP-1976-21}}
\end{GSitem}
\begin{GSitem}{1976-414.7}{E}{Cycle inversion and the real dimension of the integrating cycle}
\GSfield{Kind; source}{theorem; \GSpage{P000697}}
\GSfield{The notes}{The inversion formula is \(\int_\gamma\mathcal H_x\varphi=c_\gamma f(x)\); the cycle is called one-dimensional.}
\GSfield{Statement}{For data in the range and a closed real two-cycle \(\gamma\) in the complex projective direction space \(\mathbb{CP}^2\), the formula holds with a constant depending only on its homology class and density normalization. A projective line is a cycle with nonzero coefficient and gives reconstruction.}
\GSfield{Reconstruction}{A \((1,1)\)-form has real degree two. The notebook's dimension is interpretable as complex dimension one, not real dimension one. A nontrivial cycle also cannot be contained entirely in the contractible affine direction chart.}
\GSfield{Fate}{See \citet[\S8, Theorem 2, pp.~29--30 of the Russian original]{gelfand1967planes}.}
\GSfield{Present status}{As of November 2026: cycle inversion is established with the corrected dimension convention.}
\GSfield{Assessment}{literal statement: correct with missing hypotheses; publication match: exact; verification: read the result; development: annotated; confidence: high}
\end{GSitem}
\begin{GSitem}{1976-414.8}{E}{Why the analogous real local period gives zero}
\GSfield{Kind; source}{claim; \GSpage{P000697}}
\GSfield{The notes}{For the real case the notes say \(C_\gamma=0\), and inversion is nonlocal.}
\GSfield{Statement}{For real Schwartz line data in the range, the closed real one-form on the pencil has zero periods: \(H^1_{\mathrm{dR}}(\mathbb{RP}^2)=0\). This particular local differential-form period cannot reconstruct \(f(x)\); a real inversion uses a different, nonlocal operator.}
\GSfield{Reconstruction}{The conclusion is about this period formula, not a claim that the real line transform lacks an inverse.}
\GSfield{Fate}{The distinction follows from the real range form and complex inversion theorem in \citet[\S\S7--8]{gelfand1967planes}.}
\GSfield{Present status}{As of November 2026: the obstruction to this particular local period is established.}
\GSfield{Assessment}{literal statement: correct with missing hypotheses; publication match: partial; verification: read the result; development: annotated; confidence: high}
\end{GSitem}
\begin{GSitem}{1976-414.9}{E}{Restricted data and the definition of an admissible complex}
\GSfield{Kind; source}{question; \GSpage{P000697}}
\GSfield{The notes}{For a three-dimensional family \(K\) of lines, the notes ask when knowing \(\mathcal H\varphi\) on \(K\) permits recovery of \(\varphi\).}
\GSfield{Statement}{In the matched 1968 formulation, \(K\) is a closed irreducible algebraic three-dimensional family of complex lines covering a dense open subset of \(\mathbb C^3\). Admissibility means the inversion form on \(K\cap G_x\) can be computed by a differential operator from restricted scalar data \(\varphi|_K\), independently of an extension off \(K\). Reconstruction additionally requires a cycle with nonzero inversion coefficient.}
\GSfield{Reconstruction}{The notebook's question supplies differentiated data, whereas the published definition asks whether those derivatives descend from scalar restricted data. These distinct inverse problems must not be identified without this correction.}
\GSfield{Fate}{See \citet[\S2, Definitions 1--2 and formula (7)]{gelfandgraev1968complexes}.}
\GSfield{Present status}{As of November 2026: the classical admissibility problem is formulated and solved in that setting.}
\GSfield{Assessment}{literal statement: correct with missing hypotheses; publication match: partial; verification: read the result; development: annotated; confidence: high; problem \hyperref[prob:GP-1976-22]{GP-1976-22}}
\end{GSitem}
\begin{GSitem}{1976-414.10}{E}{Admissible line complexes come from a curve or a surface}
\GSfield{Kind; source}{theorem; \GSpage{P000697}}
\GSfield{The notes}{The asserted alternatives are all lines meeting a curve or all lines tangent to a surface.}
\GSfield{Statement}{For the algebraic admissibility notion just stated, a complex of lines in \(\mathbb C^3\), near a generic regular line and its regular critical point, is either the family meeting a local curve or the family tangent to a local surface. The general-dimensional theorem imposes distinct critical points and general-position tangent hyperplanes; the paper states the three-dimensional dichotomy as its corollary.}
\GSfield{Reconstruction}{This is not a classification of arbitrary smooth three-parameter families with any conceivable injectivity property. Singular global branches and the data model require the paper's hypotheses.}
\GSfield{Fate}{See \citet[\S5, Theorem 6 and corollary, pp.~49--51]{gelfandgraev1968complexes}; compare \citet[\S3]{guillemin1985line}.}
\GSfield{Present status}{As of November 2026: the identified classification predates the talk.}
\GSfield{Assessment}{literal statement: correct with missing hypotheses; publication match: exact; verification: read the result; development: none; confidence: high; problem \hyperref[prob:GP-1976-22]{GP-1976-22}}
\end{GSitem}
\begin{GSitem}{1976-414.11}{E}{Plücker coordinates and the equation of a linear complex}
\GSfield{Kind; source}{definition; \GSpage{P000698}}
\GSfield{The notes}{For a plane spanned by \(x,y\in k^4\), \(p_{ij}=x_i y_j-x_j y_i\); a linear complex satisfies \(\sum c_{ij}p_{ij}=0\).}
\GSfield{Statement}{A projective line in \(\mathbb P^3_k\), \(k=\mathbb R\) or \(\mathbb C\), is a two-plane in \(k^4\). Its decomposable bivector \(x\wedge y\) gives a point of the Klein quadric in \(\mathbb P(\Lambda^2k^4)\). A linear complex is its hyperplane section, equivalently the two-planes isotropic for an alternating form \(\omega\).}
\GSfield{Reconstruction}{The plane in the notebook is the vector two-plane representing a projective line.}
\GSfield{Fate}{The Grassmannian framework is recalled in \citet[\S\S1--2]{guillemin1985line}; this is classical projective geometry.}
\GSfield{Present status}{As of November 2026: the construction is established.}
\GSfield{Assessment}{literal statement: correct; publication match: exact; verification: read the result; development: none; confidence: high}
\end{GSitem}
\begin{GSitem}{1976-414.12}{E}{Symplectic polarity attached to a nondegenerate linear complex}
\GSfield{Kind; source}{construction; \GSpage{P000697-P000698}}
\GSfield{The notes}{Each point is associated with a plane containing the complex's lines through it; each line is associated with another line meeting those lines.}
\GSfield{Statement}{If \(\omega\) is nondegenerate, the plane of complex lines through \(p=[v]\) is \(\mathbb P(v^\perp)\). A line \(\mathbb P(L)\) has polar line \(\mathbb P(L^\perp)\). Every isotropic line meeting \(\mathbb P(L)\) also meets \(\mathbb P(L^\perp)\).}
\GSfield{Reconstruction}{Nondegeneracy is needed for a polarity. If \(v\in\ker\omega\), all lines through \(p\) satisfy the complex equation, and no single projective plane contains them.}
\GSfield{Fate}{This is the alternating-form realization of the classical line geometry underlying \citet[\S2]{guillemin1985line}.}
\GSfield{Present status}{As of November 2026: the stated polarity is valid in the nondegenerate case.}
\GSfield{Assessment}{literal statement: correct with missing hypotheses; publication match: partial; verification: read the result; development: annotated; confidence: high}
\end{GSitem}
\begin{GSitem}{1976-414.13}{E}{The standard linear complex and its isomorphism class}
\GSfield{Kind; source}{claim; \GSpage{P000698}}
\GSfield{The notes}{The example is \(x_1=\alpha x_3+\beta,\ x_2=\gamma x_3+\alpha\); the notes say any other complex is isomorphic to it.}
\GSfield{Statement}{In homogeneous coordinates this is the affine chart of \(p_{14}+p_{23}=0\), a rank-four alternating form. All nondegenerate alternating forms on \(k^4\) are linearly equivalent. Rank-two forms give a different projective orbit: their complexes consist of lines meeting the projectivized kernel.}
\GSfield{Reconstruction}{The universal assertion must be restricted to nondegenerate linear complexes. It does not classify all nonlinear complexes and does not prove admissibility for restricted scalar line data.}
\GSfield{Fate}{The linear-algebra normal form supplies the correction; compare the admissibility distinction in \citet[\S2]{gelfandgraev1968complexes}.}
\GSfield{Present status}{As of November 2026: the nondegenerate normal form is established; the unqualified assertion is false.}
\GSfield{Assessment}{literal statement: false as written; publication match: partial; verification: read the result; development: annotated; confidence: high}
\end{GSitem}
\begin{GSitem}{1976-414.14}{F}{An unspecified Legendre transformation}
\GSfield{Kind; source}{remark; \GSpage{P000698}}
\GSfield{The notes}{The single line says \ru{Преобр. Лежандра есть}.}
\GSfield{Statement}{No source space, target space, generating function, or transformation formula is given. Contact or projective duality associated with the line complex is a possible interpretation, but is not uniquely determined.}
\GSfield{Reconstruction}{No specific Legendre transform is inserted into the notebook.}
\GSfield{Fate}{The phrase was compared with the cited integral-geometry papers and the line-complex literature; no exact transformation was identified.}
\GSfield{Present status}{As of November 2026: the intended transformation remains unidentifiable from this record.}
\GSfield{Assessment}{literal statement: not determinable; publication match: none found; verification: unresolved lead; development: none; confidence: low}
\end{GSitem}

\GStalk{416}{1977-03-28}{I. M. Gelfand}{P000699}{Poncelet closure and Figiel’s short proof of Dvoretzky’s theorem}
The short intervention above Vershik's heading states the Poncelet porism for two conics, asks for the closing condition, and points to Griffiths's 1976 paper. It then identifies Figiel's 1976 Compositio proof of Dvoretzky's theorem. This material precedes record 415 on the manuscript page despite the catalogue record-number order.

\begin{GSitem}{1976-416.1}{E}{When does the Poncelet polygon close?}
\GSfield{Kind; source}{problem; \GSpage{P000699}}
\GSfield{The notes}{A polygon is inscribed in one conic and tangent to another; Gelfand asks when such a polygon closes, citing Griffiths.}
\GSfield{Statement}{For two nonsingular complex plane conics whose incidence curve of a point on the outer conic and a tangent to the inner conic is smooth, that curve is genus one. The successive-tangent operation is translation by a fixed \(\delta\) on it. A polygon closes after \(m\) steps exactly when \(m\delta=0\); exact period \(m\) requires \(\delta\) to have exact order \(m\). Nested real ovals select the real branch and direction of iteration.}
\GSfield{Reconstruction}{Smoothness of the incidence curve excludes degenerate conic pairs. The elliptic-curve translation is the precise reconstruction behind the notes' mention of Jacobi's elliptic integrals.}
\GSfield{Fate}{See \citet[\S I(d), pp.~345--348]{griffiths1976abel}, published before the talk.}
\GSfield{Present status}{As of November 2026: the closure criterion is classical and solved.}
\GSfield{Assessment}{literal statement: correct with missing hypotheses; publication match: exact; verification: read the result; development: none; confidence: high; problem \hyperref[prob:GP-1976-23]{GP-1976-23}}
\end{GSitem}
\begin{GSitem}{1976-416.2}{E}{Poncelet’s porism}
\GSfield{Kind; source}{theorem; \GSpage{P000699}}
\GSfield{The notes}{If the construction closes from one point, it closes from every other point in the same number of steps.}
\GSfield{Statement}{For the nondegenerate conic pair and consistent tangent iteration above, the existence of a closed \(m\)-step orbit means the elliptic translation is \(m\)-torsion. Therefore its \(m\)th iterate is the identity for every admissible starting point. Exceptional points of the algebraic construction are understood by extension.}
\GSfield{Reconstruction}{The same branch and a nondegenerate conic pair are required; changing tangent directions arbitrarily is not this map.}
\GSfield{Fate}{The porism is stated in \citet[\S I(d), formula (1.37) and discussion]{griffiths1976abel}.}
\GSfield{Present status}{As of November 2026: the theorem is classical.}
\GSfield{Assessment}{literal statement: correct with missing hypotheses; publication match: exact; verification: read the result; development: none; confidence: high; problem \hyperref[prob:GP-1976-23]{GP-1976-23}}
\end{GSitem}
\begin{GSitem}{1976-416.3}{E}{Invariant measure and circle-rotation coordinates}
\GSfield{Kind; source}{claim; \GSpage{P000699}}
\GSfield{The notes}{The notebook writes that the Poncelet theorem implies an invariant measure and an isomorphism with a circle rotation.}
\GSfield{Statement}{For nested smooth real conics with a regular real Poncelet map, the elliptic-curve parameter gives an invariant positive measure on the relevant circle component. Integrating and normalizing this measure gives a coordinate conjugating the map to a rigid rotation. Rational rotation number is the closure criterion.}
\GSfield{Reconstruction}{The construction uses the elliptic geometry and a full-support nonatomic invariant density. The mere logical statement of a porism does not supply such a density or the conjugacy.}
\GSfield{Fate}{The elliptic-parameter construction is in \citet[\S I(d)]{griffiths1976abel}.}
\GSfield{Present status}{As of November 2026: the geometric rotation description is established, with the stated regularity.}
\GSfield{Assessment}{literal statement: correct with missing hypotheses; publication match: partial; verification: read the result; development: annotated; confidence: high}
\end{GSitem}
\begin{GSitem}{1976-416.4}{E}{Figiel’s short proof of Dvoretzky’s theorem}
\GSfield{Kind; source}{remark; \GSpage{P000699}}
\GSfield{The notes}{Gelfand points to the latest Compositio issue and attributes the simplest proof of Dvoretzky's theorem to Figiel.}
\GSfield{Statement}{For every integer \(k\geq2\) and \(\varepsilon>0\), there is \(N(k,\varepsilon)\) such that every real or complex normed space of dimension at least \(N(k,\varepsilon)\) has a \(k\)-dimensional subspace with Banach--Mazur distance at most \(1+\varepsilon\) from the corresponding Euclidean space.}
\GSfield{Reconstruction}{This is a bibliographically identifiable pointer; the notebook contains no details of the proof and makes no quantitative claim about \(N\).}
\GSfield{Fate}{The exact source is \citet[theorem, p.~297]{figiel1976short}, Compositio 33(3), 1976.}
\GSfield{Present status}{As of November 2026: the theorem and the cited proof predate the talk.}
\GSfield{Assessment}{literal statement: correct; publication match: exact; verification: read the result; development: none; confidence: high}
\end{GSitem}

\GStalk{415}{1977-03-28}{A. M. Vershik}{P000699-P000701}{Probability, the Young graph, and Plancherel asymptotics}
The manuscript starts Vershik's talk below Gelfand's intervention on the page headed 28 March 1977, without another date. The catalogue's 4 April date is therefore corrected here, with the catalogue value preserved in the provenance. The exposition moves from Young's finite representation theory and branching paths to Plancherel growth, the limit shape, Ulam's longest increasing subsequence problem, and a proposed fluctuation theorem. The exact limit-shape and first-row results are in the 1977 Vershik--Kerov announcement, received in December 1976. The fluctuation line is not an exact statement of the later Kerov theorem: both the displayed fourth-root scaling and the white-noise description require correction.

\begin{GSitem}{1976-415.1}{E}{Symmetric groups, partitions, and Young diagrams}
\GSfield{Kind; source}{definition; \GSpage{P000699}}
\GSfield{The notes}{The notes define \(S_n\), finitary \(S_\infty\), and a descending array of left-aligned rows called a Young tableau.}
\GSfield{Statement}{\(S_\infty=\bigcup_nS_n\) under the inclusions fixing the last letter. A partition \(\lambda_1\geq\lambda_2\geq\cdots\geq0\), \(\sum\lambda_i=n\), specifies a Young diagram. A standard tableau additionally labels its cells \(1,\ldots,n\), increasingly in rows and columns.}
\GSfield{Reconstruction}{The distinction between an unfilled diagram and a labelled tableau is editorial terminology; the notebook uses the Russian word for tableau for both.}
\GSfield{Fate}{This is the classical framework recalled in \citet[\S\S1,5]{vershikokounkov2005approach}; the asymptotic setting is explicit in \citet[p.~1024]{vershikkerov1977asymptotics}.}
\GSfield{Present status}{As of November 2026: these definitions are standard.}
\GSfield{Assessment}{literal statement: correct; publication match: exact; verification: read the result; development: none; confidence: high}
\end{GSitem}
\begin{GSitem}{1976-415.2}{E}{The common constituent of row and column induced representations}
\GSfield{Kind; source}{construction; \GSpage{P000699-P000700}}
\GSfield{The notes}{For the row subgroup \(G_1\) and column subgroup \(G_2\), the notes induce the trivial and alternating representations and select their one common irreducible constituent.}
\GSfield{Statement}{Over \(\mathbb C\), let \(R_\lambda\) and \(C_\lambda\) permute the labelled cells within each row and column. The representations \(\operatorname{Ind}_{R_\lambda}^{S_n}{\bf1}\) and \(\operatorname{Ind}_{C_\lambda}^{S_n}\operatorname{sgn}\) have just one common irreducible constituent, \(V^\lambda\), with multiplicity one in both. The resulting \(V^\lambda\) exhaust the irreducible representations.}
\GSfield{Reconstruction}{The field and multiplicity are made explicit. A different labelling gives conjugate subgroups and the same representation class.}
\GSfield{Fate}{The traditional Young construction is recalled in \citet[Introduction]{vershikokounkov2005approach}; the paper's new inductive proof is later, but the theorem itself is classical.}
\GSfield{Present status}{As of November 2026: this is established characteristic-zero representation theory.}
\GSfield{Assessment}{literal statement: correct with missing hypotheses; publication match: exact; verification: read the result; development: none; confidence: high}
\end{GSitem}
\begin{GSitem}{1976-415.3}{E}{Hook-length formula for the degree of an irreducible}
\GSfield{Kind; source}{theorem; \GSpage{P000700}}
\GSfield{The notes}{The page gives \(n!/\prod h\) and dates the hook formula to 1955; its hook description mentions the cells to the right and below.}
\GSfield{Statement}{For \(\lambda\vdash n\), \(d_\lambda=\dim V^\lambda=n!/\prod_{(i,j)\in\lambda}h_{ij}\), where \(h_{ij}=\lambda_i-j+\lambda'_j-i+1\). The corner cell itself is included.}
\GSfield{Reconstruction}{The \(+1\) is essential. Interpreting the notebook's verbal description as excluding the corner would give zero hook lengths at removable cells.}
\GSfield{Fate}{The original publication is \citet{framerobinsonthrall1954hook}, dated 1954 in the opened publisher record; the formula also appears in \citet[p.~1024]{vershikkerov1977asymptotics}. The original full article was not accessible in this run.}
\GSfield{Present status}{As of November 2026: the formula is classical; the notebook date is one year late.}
\GSfield{Assessment}{literal statement: correct with missing hypotheses; publication match: exact; verification: read the result; development: annotated; confidence: high}
\end{GSitem}
\begin{GSitem}{1976-415.4}{E}{Multiplicity-free branching by removing a cell}
\GSfield{Kind; source}{theorem; \GSpage{P000700}}
\GSfield{The notes}{One sentence says that every irreducible of \(S_{n-1}\) occurs once, but the following line restricts to diagrams preceding \(\lambda\) in the Young graph.}
\GSfield{Statement}{\(\operatorname{Res}^{S_n}_{S_{n-1}}V^\lambda\cong\bigoplus_{\mu\nearrow\lambda}V^\mu\), where \(\mu\nearrow\lambda\) means that \(\lambda/\mu\) consists of one removable cell. No other irreducible occurs.}
\GSfield{Reconstruction}{The unrestricted wording is corrected using the immediately following notebook qualification.}
\GSfield{Fate}{The classical branching theorem is proved in \citet[Theorems~2.9 and 5.8]{vershikokounkov2005approach}.}
\GSfield{Present status}{As of November 2026: the predecessor restriction is indispensable.}
\GSfield{Assessment}{literal statement: false as written; publication match: exact; verification: read the result; development: annotated; confidence: high}
\end{GSitem}
\begin{GSitem}{1976-415.5}{E}{Paths, standard tableaux, and a Gelfand--Tsetlin basis}
\GSfield{Kind; source}{construction; \GSpage{P000700}}
\GSfield{The notes}{Paths from the empty diagram to \(\lambda\) index basis vectors and count the degree of the representation.}
\GSfield{Statement}{Successive multiplicity-free restrictions along \(S_1\subset\cdots\subset S_n\) decompose \(V^\lambda\) into one-dimensional joint eigenspaces for the centers of the subgroup algebras. Paths to \(\lambda\) index these lines, hence \(d_\lambda=f^\lambda\), the number of standard tableaux. Choosing a nonzero vector in each line gives a basis, unique only up to independent rescalings.}
\GSfield{Reconstruction}{A canonical basis of lines does not choose phases or scales of vectors.}
\GSfield{Fate}{See \citet[\S1 and Theorem~5.8]{vershikokounkov2005approach}.}
\GSfield{Present status}{As of November 2026: the classical branching-path interpretation is exact.}
\GSfield{Assessment}{literal statement: correct with missing hypotheses; publication match: exact; verification: read the result; development: none; confidence: high}
\end{GSitem}
\begin{GSitem}{1976-415.6}{E}{Counting paths between two fixed Young diagrams}
\GSfield{Kind; source}{question; \GSpage{P000700}}
\GSfield{The notes}{After calling combinatorial counting unsolved, the notes ask for the number of paths from a diagram of size \(n\) to a specified diagram of size \(m<n\).}
\GSfield{Statement}{For \(\mu\vdash m\), \(\lambda\vdash n\), \(\mu\subseteq\lambda\), the multiplicity of \(V^\mu\) in \(\operatorname{Res}^{S_n}_{S_m}V^\lambda\) is \(f^{\lambda/\mu}\). Equivalently, it counts standard tableaux of the skew shape. For \(r\geq\ell(\lambda)\), padding both partitions with zeros,
 \[
 f^{\lambda/\mu}=(n-m)!\det\left[\frac1{(\lambda_i-\mu_j-i+j)!}\right]_{i,j=1}^r,
 \]
 with \(1/k!=0\) for \(k<0\).}
\GSfield{Reconstruction}{The determinant and the precise skew-tableau formulation are supplied here. The notebook does not specify a stronger demand such as a positive product formula or uniform asymptotics; such demands are not manufactured as open problems.}
\GSfield{Fate}{Iterated branching gives the multiplicity interpretation. The classical Aitken determinant follows from Jacobi--Trudi, as explicitly derived in \citet[\S1, equations~(1.1)--(1.2)]{stanley2011skew}; the latter is a checked later exposition, not a claim that the formula originated in 2011.}
\GSfield{Present status}{As of November 2026: the literal finite counting problem has an exact determinant; no precise unsolved strengthening is recorded.}
\GSfield{Assessment}{literal statement: correct; publication match: exact; verification: read the result; development: none; confidence: high; problem \hyperref[prob:GP-1976-24]{GP-1976-24}}
\end{GSitem}
\begin{GSitem}{1976-415.7}{E}{The assertion that total path counting is unsolved}
\GSfield{Kind; source}{claim; \GSpage{P000700}}
\GSfield{The notes}{The notebook says \ru{Комбинаторная задача подсчета не решена} and mentions only Frobenius's formula, despite having just written the hook formula.}
\GSfield{Statement}{For paths from the empty diagram, the hook formula already counts them. For two fixed endpoints, the skew determinant in the preceding item gives an exact count. Without an additional restriction on the permitted formula or asymptotic regime, the stated unsolved assertion is false.}
\GSfield{Reconstruction}{The two possible literal readings are distinguished; neither supports an identified open problem.}
\GSfield{Fate}{Compare \citet{framerobinsonthrall1954hook} and \citet[equation~(1.2)]{stanley2011skew}.}
\GSfield{Present status}{As of November 2026: the claim cannot be retained as evidence of an open counting problem.}
\GSfield{Assessment}{literal statement: false as written; publication match: partial; verification: read the result; development: annotated; confidence: high}
\end{GSitem}
\begin{GSitem}{1976-415.8}{E}{First-row length detected by invariant vectors}
\GSfield{Kind; source}{theorem; \GSpage{P000700}}
\GSfield{The notes}{The first row is characterized by restriction to symmetric subgroups and the appearance of an invariant vector.}
\GSfield{Statement}{For the standard inclusion \(S_m\subset S_n\), \((V^\lambda)^{S_m}\ne0\) if and only if \(m\leq\lambda_1\). Thus \(\lambda_1=\max\{m:(V^\lambda)^{S_m}\ne0\}\).}
\GSfield{Reconstruction}{This follows by iterating the branching theorem: the one-row partition \((m)\) must be a subdiagram of \(\lambda\).}
\GSfield{Fate}{This is an elementary consequence of the classical branching rule, verified in \citet[Theorem~5.8]{vershikokounkov2005approach}.}
\GSfield{Present status}{As of November 2026: the representation-theoretic characterization is exact.}
\GSfield{Assessment}{literal statement: correct; publication match: exact; verification: read the result; development: none; confidence: high}
\end{GSitem}
\begin{GSitem}{1976-415.9}{E}{Finite Plancherel measure}
\GSfield{Kind; source}{definition; \GSpage{P000700}}
\GSfield{The notes}{The regular representation is described by functions on a symbol resembling \(S^n\); weights are \((\dim\lambda)^2/n!\).}
\GSfield{Statement}{\(\mathbb C[S_n]\cong\bigoplus_{\lambda\vdash n}V^\lambda\otimes(V^\lambda)^*\); therefore \(\sum_\lambda d_\lambda^2=n!\), and \(\mathrm{Pl}_n(\lambda)=d_\lambda^2/n!\) is a probability measure on partitions of \(n\).}
\GSfield{Reconstruction}{The group is \(S_n\), not a sphere or a Cartesian power.}
\GSfield{Fate}{See \citet[p.~1024]{vershikkerov1977asymptotics}.}
\GSfield{Present status}{As of November 2026: this is the classical finite-group Plancherel measure.}
\GSfield{Assessment}{literal statement: correct with missing hypotheses; publication match: exact; verification: read the result; development: none; confidence: high}
\end{GSitem}
\begin{GSitem}{1976-415.10}{A}{The coherent Plancherel measure on infinite Young paths}
\GSfield{Kind; source}{construction; \GSpage{P000700}}
\GSfield{The notes}{The notes pass to infinite paths and conditional probabilities for adding a cell.}
\GSfield{Statement}{Give a specified length-\(n\) path ending at \(\lambda\) mass \(d_\lambda/n!\). The resulting coherent probability measure has level marginal \(d_\lambda^2/n!\). Its growth transition is
 \[
 p(\lambda,\nu)=\frac{d_\nu}{(n+1)d_\lambda}\quad(\lambda\nearrow\nu,\ |\lambda|=n).
 \]
 The dimension identity for induction makes the transition probabilities sum to one.}
\GSfield{Reconstruction}{These precise cylinder and transition probabilities expand the notebook shorthand.}
\GSfield{Fate}{The infinite path measure is described in \citet[p.~1025, before Theorem~2]{vershikkerov1977asymptotics}. The paper was received 14 December 1976 and presented 17 December 1976, then printed in 1977; those editorial dates do not establish an earlier public preprint.}
\GSfield{Present status}{As of November 2026: the coherent growth process is established.}
\GSfield{Assessment}{literal statement: correct with missing hypotheses; publication match: exact; verification: read the result; development: none; confidence: high}
\end{GSitem}
\begin{GSitem}{1976-415.11}{F}{Gelfand asks for an infinite-dimensional Plancherel formulation}
\GSfield{Kind; source}{question; \GSpage{P000700}; I. M. Gelfand}
\GSfield{The notes}{Gelfand asks for a continuous, infinite-dimensional analogue of Plancherel's formula.}
\GSfield{Statement}{For finitary \(S_\infty\), a precise interpretation in the subsequent discussion is the spectral measure of the Young-path commutative subalgebra in the regular representation. The notebook does not specify a unique general Plancherel theorem for arbitrary infinite-dimensional groups.}
\GSfield{Reconstruction}{A type-I direct integral over irreducible representations is not silently substituted for this question: the regular von Neumann algebra of \(S_\infty\) is a type \(\mathrm{II}_1\) factor.}
\GSfield{Fate}{The specific path-space realization and the type \(\mathrm{II}_1\) factor are explicit in \citet[p.~1025]{vershikkerov1977asymptotics}; this supplies only a partial match to the unqualified request.}
\GSfield{Present status}{As of November 2026: the specific construction is known; the broader requested object is not uniquely determined.}
\GSfield{Assessment}{literal statement: not determinable; publication match: partial; verification: read the result; development: none; confidence: medium}
\end{GSitem}
\begin{GSitem}{1976-415.12}{A}{Plancherel limit shape in half-rotated coordinates}
\GSfield{Kind; source}{theorem; \GSpage{P000701}}
\GSfield{The notes}{The page displays uniform convergence in probability to \(\frac2\pi(X\arcsin X+\sqrt{1-X^2})\), but interchanges the sum and difference in the coordinate definition and omits the exterior formula.}
\GSfield{Statement}{For nonnegative row-column coordinates \(x,y\), set \(X=(x-y)/2\), \(Y=(x+y)/2\), and let \(\widetilde\lambda(X)\) be the upper boundary of the rotated diagram. Under \(\mathrm{Pl}_n\),
 \[
 \sup_X\left|n^{-1/2}\widetilde\lambda(\sqrt n X)-\Omega(X)\right|
 \longrightarrow0
 \quad\hbox{in probability},
 \]
 where \(\Omega(X)=\frac2\pi(X\arcsin X+\sqrt{1-X^2})\) for \(|X|\leq1\), and \(\Omega(X)=|X|\) otherwise.}
\GSfield{Reconstruction}{The corrected coordinates and the exterior branch are checked against the 1977 paper. Their normalization has half the coordinates used in many later accounts.}
\GSfield{Fate}{This is \citet[Theorem~1, pp.~1024--1025]{vershikkerov1977asymptotics}, received December 1976 and published in 1977.}
\GSfield{Present status}{As of November 2026: the limit-shape theorem is established with these conventions.}
\GSfield{Assessment}{literal statement: false as written; publication match: exact; verification: read the result; development: annotated; confidence: high}
\end{GSitem}
\begin{GSitem}{1976-415.13}{A}{Almost-sure convergence along coherent paths}
\GSfield{Kind; source}{theorem; \GSpage{P000701}}
\GSfield{The notes}{The notes separately assert convergence almost everywhere for infinite diagrams, interpreted as paths.}
\GSfield{Statement}{For almost every infinite path under the coherent Plancherel measure, the uniformly rescaled boundary at level \(n\) converges uniformly to \(\Omega\). The assertion uses that coupling of the marginals, not an unspecified coupling of independent samples.}
\GSfield{Reconstruction}{The ambient probability space is made explicit.}
\GSfield{Fate}{This is \citet[Theorem~2, p.~1025]{vershikkerov1977asymptotics}.}
\GSfield{Present status}{As of November 2026: almost-sure pathwise convergence is established.}
\GSfield{Assessment}{literal statement: correct with missing hypotheses; publication match: exact; verification: read the result; development: none; confidence: high}
\end{GSitem}
\begin{GSitem}{1976-415.14}{A}{Asymptotic first-row length}
\GSfield{Kind; source}{theorem; \GSpage{P000701}}
\GSfield{The notes}{The displayed corollary has \(\lambda_1/(2\sqrt n)\) near one and is called a solution of Ulam's problem.}
\GSfield{Statement}{For every \(\varepsilon>0\),
 \(\mathrm{Pl}_n\{|\lambda_1/(2\sqrt n)-1|<\varepsilon\}\to1\). The same first-order statement holds for the first column by transposition.}
\GSfield{Reconstruction}{The missing terminal probability limit is supplied. Uniform boundary convergence alone does not control arbitrarily thin outer rows; the original proof also invokes a bound on expected first-row length.}
\GSfield{Fate}{See \citet[Corollary~1 and its proof, p.~1026]{vershikkerov1977asymptotics}.}
\GSfield{Present status}{As of November 2026: the first-order law is proved; it is not itself a fluctuation theorem.}
\GSfield{Assessment}{literal statement: correct with missing hypotheses; publication match: exact; verification: read the result; development: none; confidence: high; problem \hyperref[prob:GP-1976-25]{GP-1976-25}}
\end{GSitem}
\begin{GSitem}{1976-415.15}{B}{The proposed white-noise central limit theorem}
\GSfield{Kind; source}{claim; \GSpage{P000701}}
\GSfield{The notes}{The notes say \ru{Уклонения дают белый шум} and display fourth roots in both the argument and prefactor of an uncentered boundary functional; the test-function class and limiting distribution are unspecified.}
\GSfield{Statement}{A verified later formulation uses full rotated coordinates \(x=2X\), profile \(\bar\lambda_n(x)=2n^{-1/2}\widetilde\lambda(\sqrt n x/2)\), and \(\Omega_{\rm full}(x)=2\Omega(x/2)\). Put \(\Delta_n(x)=\frac{\sqrt n}{2}(\bar\lambda_n(x)-\Omega_{\rm full}(x))\). For any finite collection of polynomials \(p\), their integrals against \(\Delta_n\) converge jointly to those against the generalized Gaussian field
 \[
 \Delta(2\cos\theta)=\frac1\pi\sum_{k\geq2}\frac{\xi_k}{\sqrt k}\sin(k\theta),\qquad 0<\theta<\pi,
 \]
 with independent standard Gaussian \(\xi_k\). This field is not spatial white noise.}
\GSfield{Reconstruction}{This is a later mathematically precise substitute, not a transcription repair or a claim that the 1977 speaker proved Kerov's theorem. The mode weights \(k^{-1/2}\), the vanishing area fluctuation, and polynomial-test topology are part of the correction.}
\GSfield{Fate}{Kerov announced the central limit theorem in 1993. \citet[Introduction, Theorem~7.1, and \S9]{ivanovolshanski2002kerov} give a full account, published in the 2002 volume and posted to arXiv in 2003; their introduction distinguishes the 1993 announcement and 1999 unpublished notes.}
\GSfield{Present status}{As of November 2026: the later generalized-Gaussian theorem is proved; the notebook fourth-root formula is not a correct statement of it.}
\GSfield{Assessment}{literal statement: false as written; publication match: partial; verification: read the result; development: annotated; confidence: high}
\end{GSitem}
\begin{GSitem}{1976-415.16}{F}{An unidentified stationarity assertion}
\GSfield{Kind; source}{remark; \GSpage{P000701}}
\GSfield{The notes}{The words \ru{марковость и стационарность проверяются} immediately follow the fluctuation formula, without naming the process or action.}
\GSfield{Statement}{Plancherel growth is Markov with the transition probabilities above, but the unrescaled shape process is not stationary: the number of cells is deterministically \(n\). A stationary process of rescaled fluctuations would require an additional time parametrization and an identified law, neither of which is given.}
\GSfield{Reconstruction}{The known Markov construction does not identify what stationarity meant in this position on the page.}
\GSfield{Fate}{The path construction is checked in \citet[p.~1025]{vershikkerov1977asymptotics}; the fluctuation source \citet[\S7]{ivanovolshanski2002kerov} does not identify this notebook sentence.}
\GSfield{Present status}{As of November 2026: the Markov part has a precise interpretation, while the stationarity claim remains unidentified.}
\GSfield{Assessment}{literal statement: not determinable; publication match: partial; verification: read the result; development: none; confidence: medium}
\end{GSitem}
\begin{GSitem}{1976-415.17}{E}{Row insertion connects permutations and Plancherel diagrams}
\GSfield{Kind; source}{method; \GSpage{P000701}}
\GSfield{The notes}{A bumping map into diagrams is credited to Knuth (1970) and a 1961 author written as Shenstead.}
\GSfield{Statement}{Schensted insertion associates to a permutation of \(n\) a pair of standard tableaux of the same shape \(\lambda\). This correspondence is bijective, so a uniform permutation has insertion shape distributed by \(d_\lambda^2/n!\). Its longest increasing subsequence has length \(\lambda_1\).}
\GSfield{Reconstruction}{The 1961 author is C. Schensted. Knuth's 1970 paper extends the correspondence; the distinction between a pair of tableaux and just a diagram is restored.}
\GSfield{Fate}{See \citet[Theorems~1--3]{schensted1961longest} and \citet[\S\S2--3]{knuth1970permutations}.}
\GSfield{Present status}{As of November 2026: the classical insertion mechanism exactly transfers the Plancherel law to increasing subsequences.}
\GSfield{Assessment}{literal statement: correct with missing hypotheses; publication match: exact; verification: read the result; development: none; confidence: high}
\end{GSitem}
\begin{GSitem}{1976-415.18}{A}{Ulam's problem for a random real sequence}
\GSfield{Kind; source}{question; \GSpage{P000701}}
\GSfield{The notes}{For uniform \(x_i\in[0,1]\), the notes ask for the greatest increasing-subsequence length and state \(2\sqrt n\). The displayed maximum is over \(r\), but the subsequence indices end at \(n\).}
\GSfield{Statement}{Let \(X_1,\ldots,X_n\) be independent uniform random variables on \([0,1]\), and let \(L_n=\max\{r:\exists i_1<\cdots<i_r,\ X_{i_1}<\cdots<X_{i_r}\}\). Then \(L_n/(2\sqrt n)\to1\) in probability. The same holds for any independent identically distributed continuous law, since its ranks form a uniform permutation.}
\GSfield{Reconstruction}{Independence, continuity, and the terminal index \(r\) are supplied. With continuous laws, strict and weak inequalities agree almost surely.}
\GSfield{Fate}{The uniform version is \citet[Corollary~2, p.~1027]{vershikkerov1977asymptotics}; the rank reduction uses \citet[Theorems~1 and 3]{schensted1961longest}.}
\GSfield{Present status}{As of November 2026: the first-order Ulam problem is solved.}
\GSfield{Assessment}{literal statement: correct with missing hypotheses; publication match: exact; verification: read the result; development: none; confidence: high; problem \hyperref[prob:GP-1976-25]{GP-1976-25}}
\end{GSitem}

\GStalk{417}{1977-04-04}{A. Tyurin}{P000702-P000704}{Packets of flat representations and exponential and Schwarzian periods}
The scan identifies only A. Tyurin; the catalogue gives A. S., whereas the exactly matching 1977 paper and 1978 survey are by A. N. Tyurin. This identification is made from the publications, not from an expanded manuscript initial. The talk compares flat representations having isomorphic underlying holomorphic bundles, explicitly solves the rank-one packet problem by logarithmic differentials, and turns to Schwarzian derivatives and projective monodromy. The GL(n) packet dimension in the notes is the fixed-determinant dimension, and the universal assertion about any two irreducible representations is false, already for unitary characters. Both require mathematical annotations rather than silent transcription repairs.

\begin{GSitem}{1976-417.1}{E}{The marked surface group}
\GSfield{Kind; source}{definition; \GSpage{P000702}}
\GSfield{The notes}{The group \(FG_g\) has generators \(a_1,\ldots,a_g,b_1,\ldots,b_g\) and relation \(\prod_i[a_i,b_i]=1\).}
\GSfield{Statement}{For a compact connected oriented surface \(\Sigma_g\), this is a presentation of \(\pi_1(\Sigma_g)\). A marked Riemann surface specifies an identification of its fundamental group with this fixed group; the orientation determines the integral intersection form on its abelianization.}
\GSfield{Reconstruction}{Marking matters because the same abstract representations are compared while the complex structure varies.}
\GSfield{Fate}{See \citet[Ch.~II, \S2, equation~(2.2.6) and Definition~2.2.6]{tyurin1978periods}. The survey was received 19 April 1977 and published in the 1978 volume; its introduction dates the underlying lectures to 1973.}
\GSfield{Present status}{As of November 2026: the group and marking are standard.}
\GSfield{Assessment}{literal statement: correct; publication match: exact; verification: read the result; development: none; confidence: high}
\end{GSitem}
\begin{GSitem}{1976-417.2}{E}{Dimension of the general linear character variety}
\GSfield{Kind; source}{claim; \GSpage{P000702}}
\GSfield{The notes}{The notes subtract the relation and conjugacy dimensions to obtain \(2(g-1)n^2+2\), calling the quotient an algebraic variety with singularities.}
\GSfield{Statement}{For \(g\geq2\), the irreducible locus of
 \(\operatorname{Hom}(\pi_1(\Sigma_g),GL_n(\mathbb C))/\!/GL_n(\mathbb C)\)
 is smooth of complex dimension \(2(g-1)n^2+2\). A point of the affine quotient represents a semisimple conjugacy class; singular and reducible strata are not covered by the displayed tangent-dimension count.}
\GSfield{Reconstruction}{At an irreducible representation, \(H^0(\Sigma_g,\operatorname{ad}\rho)\) and \(H^2\) each have dimension one. Euler characteristic gives \(\dim H^1=2+(2g-2)n^2\). This cohomological calculation checks the notebook's generic parameter count.}
\GSfield{Fate}{The flat-bundle and representation interpretation is \citet[Ch.~II, \S2]{tyurin1978periods}; the explicit tangent-space argument is an editorial reconstruction of the familiar generic count.}
\GSfield{Present status}{As of November 2026: the generic dimension is correct; the unrestricted singular quotient needs the stated qualification.}
\GSfield{Assessment}{literal statement: correct with missing hypotheses; publication match: partial; verification: read the result; development: annotated; confidence: high}
\end{GSitem}
\begin{GSitem}{1976-417.3}{E}{The flat bundle associated to a representation}
\GSfield{Kind; source}{construction; \GSpage{P000702}}
\GSfield{The notes}{The notes form \(E_\rho=U(X)\times V^n/\pi_1(X)\) and declare two representations analytically equivalent when their underlying bundles are isomorphic.}
\GSfield{Statement}{For a representation \(\rho:\pi_1(X)\to GL_n(\mathbb C)\), quotient the trivial bundle on the universal cover by the compatible deck and linear action. Constant transition functions give a flat connection and a holomorphic bundle \(E_\rho\). Forgetting the flat connection sends the conjugacy class of \(\rho\) to the holomorphic isomorphism class of \(E_\rho\).}
\GSfield{Reconstruction}{The deck-action convention must be used consistently; changing it changes \(\rho\) to the corresponding inverse convention, not the underlying correspondence.}
\GSfield{Fate}{This is \citet[Ch.~II, \S2, equations~(2.2.3)--(2.2.5) and Definition~2.2.7]{tyurin1978periods}.}
\GSfield{Present status}{As of November 2026: the construction and equivalence relation are standard.}
\GSfield{Assessment}{literal statement: correct; publication match: exact; verification: read the result; development: none; confidence: high}
\end{GSitem}
\begin{GSitem}{1976-417.4}{E}{A packet is a fiber of forgetting the flat structure}
\GSfield{Kind; source}{definition; \GSpage{P000702}}
\GSfield{The notes}{For fixed \(X,E\), a packet \(H(X,E)\) is a subset of the representation space.}
\GSfield{Statement}{\(H(X,E)\) consists of conjugacy classes of representations whose associated holomorphic bundle is isomorphic to \(E\). Equivalently, it is the set of holomorphic connections on \(E\), modulo the action of \(\operatorname{Aut}_{\rm hol}(E)\). On a complex curve a holomorphic connection is flat because its holomorphic curvature is a two-form.}
\GSfield{Reconstruction}{The quotient by bundle automorphisms is necessary before calling this a space of representation classes.}
\GSfield{Fate}{See \citet[Ch.~II, Definition~2.2.8 and Exercises~1--5 to \S2, pp.~199, 204]{tyurin1978periods}. The survey explicitly refers the connection theory back to Atiyah's earlier work.}
\GSfield{Present status}{As of November 2026: the packet interpretation is established.}
\GSfield{Assessment}{literal statement: correct with missing hypotheses; publication match: exact; verification: read the result; development: none; confidence: high}
\end{GSitem}
\begin{GSitem}{1976-417.5}{E}{The packet dimension mixes GL(n) and fixed determinant}
\GSfield{Kind; source}{claim; \GSpage{P000702}}
\GSfield{The notes}{An arrow labels \(H(X,E)\) as an affine space of dimension \((n^2-1)(g-1)\), although the ambient representations are in \(GL_n\).}
\GSfield{Statement}{If \(E\) is a stable degree-zero rank-\(n\) bundle on a compact curve of genus \(g\geq2\), its holomorphic connections form an affine torsor for \(H^0(X,\operatorname{End}E\otimes K_X)\), of dimension \(n^2(g-1)+1\). Scalar automorphisms act trivially. If one also fixes the connection on \(\det E\), the modeling space is \(H^0(X,\operatorname{End}_0E\otimes K_X)\), of dimension \((n^2-1)(g-1)\).}
\GSfield{Reconstruction}{The dimensions follow from Riemann--Roch and \(h^0(\operatorname{End}E)=1\). In rank one the GL(1) packet has dimension \(g\), whereas the notebook formula gives zero; this directly exposes the missing determinant constraint.}
\GSfield{Fate}{The affine difference of connections and the dimension exercise are \citet[Ch.~II, \S2, Exercises~1 and 8, p.~204]{tyurin1978periods}. The corrected Riemann--Roch calculation is supplied here.}
\GSfield{Present status}{As of November 2026: the displayed formula is valid only for the fixed-determinant connection problem.}
\GSfield{Assessment}{literal statement: false as written; publication match: partial; verification: read the result; development: annotated; confidence: high}
\end{GSitem}
\begin{GSitem}{1976-417.6}{F}{An ambiguous question about related bundles and representations}
\GSfield{Kind; source}{question; \GSpage{P000702}}
\GSfield{The notes}{A short line asks \ru{Расслоения связ. \(\Rightarrow\) что можно сказать о предст.?}; the word abbreviated as \ru{связ.} is uncertain.}
\GSfield{Statement}{No specific relation between the two bundles is given. If holomorphic isomorphism is intended, the question asks about membership in a common packet; if deformation or another relation is intended, the notebook does not supply a precise condition or requested conclusion.}
\GSfield{Reconstruction}{The surrounding definition suggests packet equivalence, but that suggestion is not substituted for the uncertain word.}
\GSfield{Fate}{The packet framework of \citet[Definitions~2.2.7--2.2.8]{tyurin1978periods} is relevant, but does not identify this isolated question's exact quantifiers.}
\GSfield{Present status}{As of November 2026: the local wording remains insufficient for a separate problem.}
\GSfield{Assessment}{literal statement: not determinable; publication match: thematic only; verification: read the result; development: none; confidence: low}
\end{GSitem}
\begin{GSitem}{1976-417.7}{E}{Any two irreducibles need not become analytically equivalent}
\GSfield{Kind; source}{claim; \GSpage{P000702}}
\GSfield{The notes}{The manuscript states \ru{Любые два неприв. предст. эквивал. относит. нек-рой рим. пов.} and explains that they enter one packet.}
\GSfield{Statement}{The literal assertion is false. Fix \(g\geq2\), \(0<\theta<2\pi\), and rank-one unitary characters \(\chi_0=1\), \(\chi_1(a_1)=e^{i\theta}\), with \(\chi_1=1\) on the remaining generators. Both are irreducible and satisfy the surface relation, but for every complex structure on \(\Sigma_g\) their underlying holomorphic line bundles are nonisomorphic.}
\GSfield{Reconstruction}{Proof: an isomorphism would give a nowhere-zero holomorphic section of the unitary flat line bundle with character \(\chi_1\chi_0^{-1}\). In unitary parallel local frames, its log norm is harmonic; compactness makes it constant. Each local holomorphic coefficient then has constant modulus, hence is constant. The section is parallel, forcing the character to be trivial, a contradiction. An independent universal-cover derivation reaches the same conclusion. This is a correction of a false literal assertion, not a new open problem.}
\GSfield{Fate}{\citet[Ch.~II, \S2, Exercise~10, p.~204]{tyurin1978periods} states that a packet cannot contain two distinct unitary representation classes. The 1973-lecture provenance establishes that the surrounding theory was not new in the 1978 printing; no claim is made that this exact notebook sentence is reproduced in the publication.}
\GSfield{Present status}{As of November 2026: the universal assertion is disproved by the checked rank-one example.}
\GSfield{Assessment}{literal statement: false as written; publication match: partial; verification: read the result; development: annotated; confidence: high}
\end{GSitem}
\begin{GSitem}{1976-417.8}{E}{The affine fibration exists only over an appropriate bundle locus}
\GSfield{Kind; source}{claim; \GSpage{P000702}}
\GSfield{The notes}{The notes draw \(R_C\to S_n(X)\), label it an affine bundle, and call \(S_n(X)\) the moduli of bundles.}
\GSfield{Statement}{Let \(X\) be a smooth projective complex curve of genus \(g\geq2\). Over the moduli of stable degree-zero bundles, forgetting a holomorphic connection gives a torsor modeled on the cotangent bundle, with fiber dimension \(n^2(g-1)+1\); fixing determinant and its connection gives the trace-free version. The whole character variety does not map to stable-bundle moduli in this way: an irreducible flat bundle can have unstable underlying holomorphic bundle.}
\GSfield{Reconstruction}{This is a restricted de Rham moduli statement. In a presentation using the Betti character variety, Riemann--Hilbert is analytic, so the notebook does not establish an algebraic affine bundle on the full Betti quotient.}
\GSfield{Fate}{The difference-of-connections model is \citet[Ch.~II, \S2, Exercises~1--5]{tyurin1978periods}; the moduli-locus and torsor formulation make the suppressed hypotheses explicit.}
\GSfield{Present status}{As of November 2026: the restricted torsor picture is standard; the global unrestricted diagram is too strong.}
\GSfield{Assessment}{literal statement: false as written; publication match: partial; verification: read the result; development: annotated; confidence: high}
\end{GSitem}
\begin{GSitem}{1976-417.9}{E}{Describe the packet of a fixed line bundle}
\GSfield{Kind; source}{problem; \GSpage{P000703}}
\GSfield{The notes}{The typical problem is to describe \(H(X,L)\subset\mathbb C^{*\,2g}\) for a fixed one-dimensional holomorphic bundle.}
\GSfield{Statement}{For a compact connected Riemann surface \(X\), the packet is empty unless \(\deg L=0\). If \(\deg L=0\), choose one flat connection \(\nabla_0\); all others are \(\nabla_0+\eta\), \(\eta\in H^0(X,K_X)\). With a fixed sign convention for periods, the character set is a coset
 \[
 H(X,L)=\chi_0\left\{\gamma\longmapsto
       \exp\!\int_\gamma\eta:\eta\in H^0(X,K_X)\right\}.
 \]}
\GSfield{Reconstruction}{Degree zero, the additive connection parameter, and the period convention are supplied. The next notebook calculations give an explicit divisor version of this description.}
\GSfield{Fate}{The precise packet question is answered by \citet[Ch.~II, \S2, Exercises~1--5 and 12; \S3]{tyurin1978periods}; the logarithmic construction is also the introduction to \citet{tyurin1977periods}.}
\GSfield{Present status}{As of November 2026: the rank-one packet is explicitly described; there is no residual open inverse problem here.}
\GSfield{Assessment}{literal statement: correct with missing hypotheses; publication match: exact; verification: read the result; development: none; confidence: high; problem \hyperref[prob:GP-1976-26]{GP-1976-26}}
\end{GSitem}
\begin{GSitem}{1976-417.10}{E}{Constant transition functions and degree zero}
\GSfield{Kind; source}{theorem; \GSpage{P000703}}
\GSfield{The notes}{The line bundle is represented by a cocycle in \(\mathcal O_X^*\), and comes from a character if a representative has constant transition functions.}
\GSfield{Statement}{The forgetful map \(H^1(X,\mathbb C^*)\to H^1(X,\mathcal O_X^*)\) has image \(\operatorname{Pic}^0(X)\). Thus a holomorphic line bundle on a compact Riemann surface admits a flat structure exactly when its degree is zero; many flat structures can have the same holomorphic image.}
\GSfield{Reconstruction}{This is a statement about existence of a constant cocycle representative, not that an arbitrarily chosen holomorphic transition cocycle is already constant.}
\GSfield{Fate}{See \citet[Ch.~II, \S2, Definitions~2.2.1--2.2.2 and Exercise~5]{tyurin1978periods}.}
\GSfield{Present status}{As of November 2026: the line-bundle criterion is established.}
\GSfield{Assessment}{literal statement: correct with missing hypotheses; publication match: exact; verification: read the result; development: none; confidence: high}
\end{GSitem}
\begin{GSitem}{1976-417.11}{E}{A meromorphic section gives an integral-residue differential}
\GSfield{Kind; source}{construction; \GSpage{P000703}}
\GSfield{The notes}{When transition functions are constant, the local \(df_i/f_i\) glue to one meromorphic differential. Its poles are simple and its residues are integers.}
\GSfield{Statement}{For a nonzero meromorphic section of a flat line bundle, \(\omega=df_i/f_i\) is global. At a zero or pole of local order \(m\in\mathbb Z\), it has residue \(m\). Hence \(\sum_x\operatorname{Res}_x\omega=0\), and its residue divisor is the divisor of that section, of degree zero.}
\GSfield{Reconstruction}{The residue sum and degree condition use compactness. The notation \([\omega]\) here records the residue divisor, not the full differential.}
\GSfield{Fate}{This construction is explicitly given in \citet[Introduction, pp.~1359--1360]{tyurin1977periods} and \citet[Ch.~II, \S3]{tyurin1978periods}.}
\GSfield{Present status}{As of November 2026: the logarithmic differential construction is classical and exact.}
\GSfield{Assessment}{literal statement: correct; publication match: exact; verification: read the result; development: none; confidence: high}
\end{GSitem}
\begin{GSitem}{1976-417.12}{E}{Changing a section changes the residue divisor by a principal divisor}
\GSfield{Kind; source}{construction; \GSpage{P000703}}
\GSfield{The notes}{The notes say that two representatives give equivalent divisors.}
\GSfield{Statement}{Two nonzero meromorphic sections of the same holomorphic line bundle differ by a nonzero meromorphic function \(h\). Their divisors differ by \(\operatorname{div}(h)\); for the same flat connection the associated differentials differ by \(dh/h\). Exponentiated periods of \(dh/h\) equal one.}
\GSfield{Reconstruction}{Holding the flat connection fixed is essential for the displayed differential comparison.}
\GSfield{Fate}{This is the divisor and exponential-integral correspondence in \citet[Introduction]{tyurin1977periods}.}
\GSfield{Present status}{As of November 2026: the gauge and divisor equivalence is established.}
\GSfield{Assessment}{literal statement: correct with missing hypotheses; publication match: exact; verification: read the result; development: none; confidence: high}
\end{GSitem}
\begin{GSitem}{1976-417.13}{E}{Constructing every character in a packet from a divisor}
\GSfield{Kind; source}{construction; \GSpage{P000703}}
\GSfield{The notes}{Fix a divisor, construct a differential with its principal parts, and solve \(df/f=\omega+\omega_0\) on the universal cover, with \(\omega_0\) holomorphic.}
\GSfield{Statement}{Let \(D=\sum m_pp\) be a degree-zero divisor representing \(L\). There exists a meromorphic one-form with only simple poles of residues \(m_p\); all such forms differ by \(H^0(X,K_X)\). Their exponential integrals are meromorphic on \(U(X)\), because integer residues make the local exponent single-valued at each pole. Their deck multipliers give exactly \(H(X,L)\), in a compatible divisor/character convention.}
\GSfield{Reconstruction}{The residue theorem's degree-zero condition is the missing existence hypothesis. Local meromorphic integrability of this first-order equation really does follow from the integer residues.}
\GSfield{Fate}{See \citet[Introduction, pp.~1359--1360]{tyurin1977periods} and \citet[Ch.~II, \S3]{tyurin1978periods}.}
\GSfield{Present status}{As of November 2026: this solves the stated rank-one packet problem.}
\GSfield{Assessment}{literal statement: correct with missing hypotheses; publication match: exact; verification: read the result; development: none; confidence: high; problem \hyperref[prob:GP-1976-26]{GP-1976-26}}
\end{GSitem}
\begin{GSitem}{1976-417.14}{E}{The trivial packet is the exponential of holomorphic periods}
\GSfield{Kind; source}{construction; \GSpage{P000703}}
\GSfield{The notes}{The zero fiber is described as exponentials of primitives of holomorphic differentials.}
\GSfield{Statement}{The homomorphism \(H^0(X,K_X)\to\operatorname{Hom}(\pi_1(X),\mathbb C^*)\), \(\eta\mapsto(\gamma\mapsto e^{\int_\gamma\eta})\), has trivial kernel and image the packet of \(\mathcal O_X\). Indeed, a kernel element gives a single-valued nowhere-zero holomorphic function \(e^{\int\eta}\) on compact \(X\), hence a constant, so \(\eta=0\). Every other nonempty line packet is a coset.}
\GSfield{Reconstruction}{The injectivity proof identifies the affine dimension \(g\) and provides a second check on the erroneous GL(1) dimension in P000702.}
\GSfield{Fate}{Compare \citet[Ch.~II, \S2, Exercise~12(a), pp.~204--205]{tyurin1978periods}.}
\GSfield{Present status}{As of November 2026: the trivial packet and its translates are explicitly understood.}
\GSfield{Assessment}{literal statement: correct; publication match: exact; verification: read the result; development: none; confidence: high}
\end{GSitem}
\begin{GSitem}{1976-417.15}{E}{A marked character packet determines the curve through Torelli}
\GSfield{Kind; source}{theorem; \GSpage{P000703}}
\GSfield{The notes}{The notes claim that the surface can be recovered uniquely from a packet, by a theorem written as Terelli.}
\GSfield{Statement}{Regard a nonempty packet as its embedded analytic subset in the character torus of the fixed oriented marked surface. Translating it to the identity gives the trivial packet; its tangent space identifies \(H^{1,0}(X)\subset H^1(\Sigma_g,\mathbb C)\). Together with the integral lattice and intersection pairing this gives the principally polarized Jacobian, which determines \(X\) up to isomorphism by Torelli.}
\GSfield{Reconstruction}{The embedding, integral marking, and polarization are indispensable. An abstract affine space of dimension \(g\), or an unpolarized torus without the retained data, is not the claimed invariant.}
\GSfield{Fate}{This is the content of \citet[Ch.~II, \S2, Exercises~12--13, pp.~204--205]{tyurin1978periods}.}
\GSfield{Present status}{As of November 2026: the Torelli conclusion is valid with the retained marking and polarization.}
\GSfield{Assessment}{literal statement: correct with missing hypotheses; publication match: exact; verification: read the result; development: none; confidence: high; problem \hyperref[prob:GP-1976-26]{GP-1976-26}}
\end{GSitem}
\begin{GSitem}{1976-417.16}{E}{Unitary characters select a unique point of each line packet}
\GSfield{Kind; source}{construction; \GSpage{P000703}}
\GSfield{The notes}{A small diagram marks unitary characters beside the Jacobian family; its circular central symbol is uncertain.}
\GSfield{Statement}{The map \(\operatorname{Hom}(\pi_1(X),U(1))\to\operatorname{Pic}^0(X)\) is a bijection of real tori. Thus each degree-zero line packet contains exactly one unitary character. This is a real analytic selection, not a holomorphic section of the complex character torus.}
\GSfield{Reconstruction}{The precise assertion reconstructs the diagram using rank-one unitary flat-bundle theory; the uncertain central glyph is not identified as a theta divisor.}
\GSfield{Fate}{Uniqueness is \citet[Ch.~II, \S2, Exercise~10]{tyurin1978periods}; the adjacent Jacobian construction in Exercise~12 supplies the line-bundle setting.}
\GSfield{Present status}{As of November 2026: the unitary character parametrization is standard.}
\GSfield{Assessment}{literal statement: correct with missing hypotheses; publication match: partial; verification: read the result; development: none; confidence: high}
\end{GSitem}
\begin{GSitem}{1976-417.17}{E}{Schwarzian derivative and its chain rule}
\GSfield{Kind; source}{definition; \GSpage{P000703}}
\GSfield{The notes}{The notes define \(S(f)=(f''/f')'-\frac12(f''/f')^2\) and give its composition rule, with zero Schwarzian for fractional-linear coordinate changes.}
\GSfield{Statement}{For a locally univalent holomorphic map,
 \[
 S_z(f)=\frac{f^{(3)}}{f'}-\frac32\left(\frac{f''}{f'}\right)^2,\qquad
 S_t(f\circ g)=(S_z(f)\circ g)(g')^2+S_t(g).
 \]
 A Möbius map has zero Schwarzian; postcomposition by one leaves \(S_z(f)\) unchanged.}
\GSfield{Reconstruction}{Local univalence is needed away from the meromorphic singularities considered next.}
\GSfield{Fate}{See \citet[\S1, properties~1--2]{tyurin1977periods} and \citet[Ch.~I, \S2]{tyurin1978periods}.}
\GSfield{Present status}{As of November 2026: the classical chain rule is exact.}
\GSfield{Assessment}{literal statement: correct; publication match: exact; verification: read the result; development: none; confidence: high}
\end{GSitem}
\begin{GSitem}{1976-417.18}{E}{Schwarzians require a projective atlas to become quadratic differentials}
\GSfield{Kind; source}{claim; \GSpage{P000703}}
\GSfield{The notes}{After the chain rule, the notes say that general quadratic differentials result.}
\GSfield{Statement}{Local Schwarzians transform as a projective connection under arbitrary holomorphic coordinate changes. After choosing a projective atlas, whose transition functions are Möbius, the extra Schwarzian term vanishes and the resulting expressions can be compared as quadratic differentials. The difference of two projective connections is intrinsically a quadratic differential.}
\GSfield{Reconstruction}{Choosing the projective structure is essential; the notebook suppresses that choice.}
\GSfield{Fate}{The projective-atlas assumption is explicit in \citet[Introduction and Definition~1.2]{tyurin1977periods}.}
\GSfield{Present status}{As of November 2026: the claim is correct with the projective-coordinate hypothesis.}
\GSfield{Assessment}{literal statement: correct with missing hypotheses; publication match: exact; verification: read the result; development: annotated; confidence: high}
\end{GSitem}
\begin{GSitem}{1976-417.19}{E}{Double poles and the leading coefficient called a residue}
\GSfield{Kind; source}{claim; \GSpage{P000703-P000704}}
\GSfield{The notes}{The margin allows poles of order at most two and calls the coefficient of \((dz/z)^2\) a residue.}
\GSfield{Statement}{If a meromorphic map has local degree \(m\geq1\), its Schwarzian is
 \[
 S_z(f)=\frac{1-m^2}{2z^2}+O(z^{-1}).
 \]
 For a quadratic differential with a double pole, the coefficient \(a\) of \(a(dz/z)^2\) is invariant under a local holomorphic change of parameter. This is not the residue of a one-form.}
\GSfield{Reconstruction}{This is a necessary local property of a meromorphic Schwarzian, not a sufficient criterion for an arbitrary quadratic differential to have a meromorphic Schwarzian primitive. Resonant logarithmic terms require additional conditions.}
\GSfield{Fate}{The terminology and necessary coefficient are in \citet[\S1, property~4 and equations~(1.3)--(1.4)]{tyurin1977periods}. The stronger sufficiency claim in its following property~5 is not used here; it does not match anything explicitly asserted in the notebook.}
\GSfield{Present status}{As of November 2026: the leading coefficient statement is exact; no converse is inferred.}
\GSfield{Assessment}{literal statement: correct with missing hypotheses; publication match: exact; verification: read the result; development: annotated; confidence: high}
\end{GSitem}
\begin{GSitem}{1976-417.20}{E}{Schwarzian integration and projective monodromy}
\GSfield{Kind; source}{construction; \GSpage{P000704}}
\GSfield{The notes}{The scan, in a diagram omitted from the plain transcription, reads \(f(g(u))=\rho(g)(f(u))\) and says that this reconstructs \(\rho\).}
\GSfield{Statement}{On a Riemann surface with a fixed projective atlas, a holomorphic quadratic differential locally gives \(S(f)=q\), equivalently \(f=y_1/y_2\) for independent solutions of \(y''+\tfrac12q\,y=0\). Analytic continuation on the universal cover produces a locally univalent developing map and a conjugacy class \(\rho:\pi_1(X)\to PGL_2(\mathbb C)\) satisfying the displayed equivariance. For meromorphic data, continuation first gives monodromy of the punctured surface; it descends across poles only under the requisite local meromorphic integrability conditions.}
\GSfield{Reconstruction}{The second-order equation and the distinction between the closed and punctured surface make the notebook's abbreviated reconstruction precise. The stronger uniqueness theorem of the 1977 paper is not imported as an occurrence absent from these pages.}
\GSfield{Fate}{See \citet[\S1, Definitions~1.1--1.3 and equation~(1.6)]{tyurin1977periods} and \citet[Ch.~II, equation~(2.2.7)]{tyurin1978periods}.}
\GSfield{Present status}{As of November 2026: the projective monodromy construction is established with the indicated local hypotheses.}
\GSfield{Assessment}{literal statement: correct with missing hypotheses; publication match: exact; verification: read the result; development: none; confidence: high}
\end{GSitem}

\GStalk{418}{}{}{P000705-P000706}{Catalogue heading for 1977–78 and adjacent next-year context}
This catalogue record marks the 1977--78 academic-year boundary; it is retained as a heading, with no fabricated 1976--77 mathematical occurrences. On P000705 the source starts the seminar's thirty-fourth year and proposes noncommutative algebraic geometry in connection with Bernstein, I.~M.~Gelfand and S.~I.~Gelfand. A distinct heading then names V.~Drinfeld and Atiyah's work on Yang--Mills equations. P000705--P000706 contain the beginning of that adjacent talk: self-dual curvature on \(S^4\), conformal invariance of the Hodge star on two-forms, a heuristic quantization/determinant discussion, the twistor map \(\mathbb{CP}^3\to\mathbb{HP}^1=S^4\) with fibers \(\mathbb{CP}^1\), and a correspondence with bundles trivial on twistor fibers, followed by the quaternionic real structure. These adjacent contents are inventoried for coverage and must be catalogued with the next academic year. The visible date is transcribed as 5/IV-77, in tension with the boundary; the exact month is left unresolved rather than silently changed to September.


\GStalk{437}{1978-02-13}{S. I. Gelfand}{P000753-P000756}{Exterior-algebra modules and algebraic bundles (manuscript: 13 February 1978)}
\textbf{Provenance correction.} The supplied catalogue date, 13 February 1977, is corrected in the structured \texttt{date} field to the manuscript date, 13 February 1978; the catalogue value is preserved here and in the provenance ledger. The scan of P000753 unambiguously reads \ru{13/II--78 г.}; P000757 begins a new lecture, \ru{27/II--78 г.}. Thus this is a 1977--78 talk accidentally included in the requested package, not evidence for a 1976--77 announcement. The occurrence prefix \texttt{1977-437} follows the format rule that YEAR is the first year of the actual academic year. The requested record is retained in this package as an explicitly identified chronological appendix. P000752 is an adjacent physics tail; P000757--P000758 belong to the next, separately dated, S.~I.~Gelfand lecture. They are accounted for in the source ledger and corrections, without assigning their new content to record~437.

The talk explains how finite graded modules over an exterior algebra produce complexes on projective space. A fibrewise exactness condition singles out vector bundles; graded free summands disappear, and every bundle is represented uniquely up to stable module isomorphism. Twisting is expressed through a cosyzygy and a grading shift, followed by an outline of the inverse construction from large-twist sections. I.~M.~Gelfand calls for coherent sheaves and derived categories, and Serre's projective module--sheaf correspondence closes the surviving lecture. The exact publication is \citet[Theorem~1 and its remarks]{bernsteingelfandgelfand1978bundles}: received 10 January 1978, before the manuscript's lecture date, published in the 1978 Russian volume; the English translation bears a 1979 copyright imprint. A submission date is not, by itself, a verified date of public circulation. Class A here denotes subsequent publication of the current work; it does not assert discovery after the lecture. The modern derived equivalence is given there by Theorem~2, but is not silently attributed to the surviving lecture as a fully stated theorem.

\begin{GSitem}{1977-437.1}{E}{Topoi and derived categories as organizing concepts}
\GSfield{Kind; source}{remark; \GSpage{P000753}; I. M. Gelfand}
\GSfield{The notes}{I.~M.~Gelfand identifies \ru{Топос} and \ru{Производная категория} as two primary concepts of category theory.}
\GSfield{Statement}{This is a pedagogical judgement, not a theorem that category theory has precisely two basic notions. A Grothendieck topos is a category equivalent to sheaves of sets on a site; for an abelian category \(\mathcal A\), its derived category inverts quasi-isomorphisms of complexes.}
\GSfield{Reconstruction}{The definitions clarify the two names; the evaluative word primary remains the speaker's judgement.}
\GSfield{Fate}{The frameworks predate this lecture: \citet[Expos\'e~IV]{sga41972topos} and \citet[Chapter~I]{hartshorne1966residues}.}
\GSfield{Present status}{As of 10 November 2026: Both frameworks remain established; the notebook supplies no testable exclusivity claim.}
\GSfield{Assessment}{literal statement: not determinable; publication match: thematic only; verification: publisher metadata checked; development: annotated; confidence: high}
\end{GSitem}
\begin{GSitem}{1977-437.2}{A}{Describe algebraic vector bundles on projective space}
\GSfield{Kind; source}{problem; \GSpage{P000753-P000755}}
\GSfield{The notes}{\ru{Задача: описать все алгебраические расслоения над} \(\mathbb P^n\), for example over \(\mathbb C\). The subsequent theorem specifies the intended description by exterior-algebra modules.}
\GSfield{Statement}{For each integer \(n\geq0\), describe the isomorphism classes of finite-rank algebraic vector bundles on \(\mathbb P^n_{\mathbb C}\). In the sense actually developed in the notes, the answer sought is a bijection with stable isomorphism classes of proper finite-dimensional graded \(\Lambda(\mathbb C^{n+1})\)-modules.}
\GSfield{Reconstruction}{The meaning of describe is fixed by the theorem that follows, rather than by an unrecorded demand for a finite list of invariants.}
\GSfield{Fate}{Answered in this representation-theoretic sense by \citet[Theorem~1]{bernsteingelfandgelfand1978bundles}.}
\GSfield{Present status}{As of 10 November 2026: The correspondence is established. It is not a numerical enumeration of all moduli spaces, and does not imply a solution of every later vector-bundle classification question.}
\GSfield{Assessment}{literal statement: correct; publication match: exact; verification: read the result; development: annotated; confidence: high; problem \hyperref[prob:GP-1977-01]{GP-1977-01}}
\end{GSitem}
\begin{GSitem}{1977-437.3}{E}{Projective space is the scalar quotient of a punctured vector space}
\GSfield{Kind; source}{definition; \GSpage{P000753}}
\GSfield{The notes}{\(\Xi\) is called affine space and \(\xi\sim\alpha\xi\), \(\alpha\in\mathbb C^*\), \(\xi\ne0\); the quotient is \(\mathbb P^n\).}
\GSfield{Statement}{Let \(\Xi\) be a complex vector space of dimension \(n+1\). Then \(\mathbb P(\Xi)=(\Xi\setminus\{0\})/\mathbb C^*\) parametrizes its one-dimensional vector subspaces. The chosen origin and linear structure are essential; a bare affine space does not carry the displayed scalar action canonically.}
\GSfield{Reconstruction}{Replace the informal affine-space wording by a vector space and retain the explicit exclusion of zero.}
\GSfield{Fate}{Classical projective geometry; compare \citet[\S27.8]{stacks2026projconstruction}.}
\GSfield{Present status}{As of 10 November 2026: Standard, with the puncture and dimension made explicit.}
\GSfield{Assessment}{literal statement: correct with missing hypotheses; publication match: exact; verification: read the result; development: annotated; confidence: high}
\end{GSitem}
\begin{GSitem}{1977-437.4}{E}{The negatively graded exterior algebra}
\GSfield{Kind; source}{definition; \GSpage{P000753}}
\GSfield{The notes}{\(\Lambda=\{\xi_0,\ldots,\xi_n\mid\xi_i\xi_j+\xi_j\xi_i=0\}\), \(\dim\Lambda=2^{n+1}\), and \(\Lambda=\bigoplus_{k=0}^{n+1}\Lambda_{-k}\).}
\GSfield{Statement}{Over \(\mathbb C\), put \(\Lambda=\bigwedge\Xi\) with \(\deg\xi=-1\) for \(\xi\in\Xi\). Then \(\Lambda_{-k}=\bigwedge^k\Xi\), \(\dim\Lambda_{-k}=\binom{n+1}{k}\), and the total dimension is \(2^{n+1}\). In characteristic two one must impose \(\xi_i^2=0\) explicitly; the notes use characteristic zero.}
\GSfield{Reconstruction}{The binomial dimensions and characteristic qualification spell out the displayed presentation.}
\GSfield{Fate}{The algebra is classical; these conventions are also fixed in \citet[\S2]{bernsteingelfandgelfand1978bundles}.}
\GSfield{Present status}{As of 10 November 2026: An established definition. The sign of the grading must be retained in subsequent shifts.}
\GSfield{Assessment}{literal statement: correct; publication match: exact; verification: read the result; development: annotated; confidence: high}
\end{GSitem}
\begin{GSitem}{1977-437.5}{E}{Finite graded representations of the exterior algebra}
\GSfield{Kind; source}{definition; \GSpage{P000753}}
\GSfield{The notes}{Objects are \(V=\bigoplus_jV_j\), \(\dim V<\infty\), with \(\Lambda_{-k}V_j\subset V_{j-k}\).}
\GSfield{Statement}{A finite-dimensional graded left \(\Lambda\)-module has a finite-support \(\mathbb Z\)-grading and degree-zero module morphisms. For a generator \(\xi\), multiplication gives \(V_j\to V_{j-1}\). Both the module and every free summand used later are graded.}
\GSfield{Fate}{Standard graded-module language, used in \citet[\S2]{bernsteingelfandgelfand1978bundles}.}
\GSfield{Present status}{As of 10 November 2026: The grading is essential; ungraded stable equivalence would be a different assertion.}
\GSfield{Assessment}{literal statement: correct; publication match: exact; verification: read the result; development: none; confidence: high}
\end{GSitem}
\begin{GSitem}{1977-437.6}{A}{A complex for every nonzero exterior generator}
\GSfield{Kind; source}{construction; \GSpage{P000753}}
\GSfield{The notes}{\(L_\xi(V)=\{\cdots\to V_1\xrightarrow{\xi}V_0\xrightarrow{\xi}V_{-1}\to\cdots\}\), because \(\xi^2=0\).}
\GSfield{Statement}{For \(\xi\in\Xi\setminus\{0\}\), let \(d_j(v)=\xi v\) on \(V_j\). Then \(d_{j-1}d_j=0\), and \(H_j(L_\xi(V))=\ker(d_j)/\operatorname{im}(d_{j+1})\). This is a bounded chain complex of finite-dimensional vector spaces.}
\GSfield{Fate}{The family is part of \citet[\S2]{bernsteingelfandgelfand1978bundles}.}
\GSfield{Present status}{As of 10 November 2026: A standard component of the BGG construction; the displayed differential is literally correct.}
\GSfield{Assessment}{literal statement: correct; publication match: exact; verification: read the result; development: none; confidence: high}
\end{GSitem}
\begin{GSitem}{1977-437.7}{A}{Proper modules: fibre homology concentrated in degree zero}
\GSfield{Kind; source}{definition; \GSpage{P000753}}
\GSfield{The notes}{\(V\) is \ru{правильный} if \(H_i(L_\xi(V))=0\) for every \(\xi\ne0\) and every \(i\ne0\).}
\GSfield{Statement}{A finite-dimensional graded \(\Lambda\)-module is proper when each nonzero \(\xi\in\Xi\) gives homology only in degree zero. No requirement that \(H_0\) be nonzero is imposed: free modules are included.}
\GSfield{Reconstruction}{The terminology must not be confused with a faithful representation in ordinary ring theory.}
\GSfield{Fate}{This is the condition called faithful in the English translation of \citet[\S2]{bernsteingelfandgelfand1978bundles}.}
\GSfield{Present status}{As of 10 November 2026: Established. Here faithful is a translation of the special definition, not injectivity of the action homomorphism.}
\GSfield{Assessment}{literal statement: correct; publication match: exact; verification: read the result; development: annotated; confidence: high}
\end{GSitem}
\begin{GSitem}{1977-437.8}{E}{The Euler characteristic fixes the fibre rank}
\GSfield{Kind; source}{claim; \GSpage{P000754}}
\GSfield{The notes}{For proper \(V\), \(\dim H_0(L_\xi(V))\) is constant and equals \(\sum_i(-1)^i\dim V_i\).}
\GSfield{Statement}{For a bounded finite-dimensional complex, \[\sum_i(-1)^i\dim H_i=\sum_i(-1)^i\dim V_i.\] Thus properness implies \(\dim H_0(L_\xi(V))=\sum_i(-1)^i\dim V_i\), independently of nonzero \(\xi\).}
\GSfield{Reconstruction}{For an independent check, write each dimension as boundary rank plus homology dimension plus the preceding boundary rank; the alternating boundary terms cancel.}
\GSfield{Fate}{An elementary Euler-characteristic identity, used in the construction of \citet[\S2]{bernsteingelfandgelfand1978bundles}.}
\GSfield{Present status}{As of 10 November 2026: Correct. Constant fibre dimension is one part of the vector-bundle argument, not a substitute for its algebraicity.}
\GSfield{Assessment}{literal statement: correct; publication match: exact; verification: read the result; development: annotated; confidence: high}
\end{GSitem}
\begin{GSitem}{1977-437.9}{A}{The homology spaces form an algebraic vector bundle}
\GSfield{Kind; source}{construction; \GSpage{P000754-P000755}}
\GSfield{The notes}{The assignment \(\xi\mapsto H_0(L_\xi(V))\) is the same for \(\xi\) and \(\alpha\xi\); later \(\Phi(V)=H^0(\mathcal L(V))\) is defined.}
\GSfield{Statement}{For a proper module over \(\Lambda(\Xi)\), the complex \(\mathcal L^j(V)=V_{-j}\otimes\mathcal O(j)\) on \(\mathbb P(\Xi)\), with differential \(\sum x_i\otimes\xi_i\), has zero cohomology sheaves outside degree zero, and \(\Phi(V)=\mathcal H^0(\mathcal L(V))\) is locally free. Its fibre at \([\xi]\) identifies with \(H_0(L_\xi(V))\).}
\GSfield{Reconstruction}{Algebraic dependence and local splitting are supplied by the displayed bundle complex: fibre exactness permits elimination of invertible minors at both ends until only a locally free degree-zero term remains.}
\GSfield{Fate}{The precise assertion precedes Theorem~1 of \citet{bernsteingelfandgelfand1978bundles}.}
\GSfield{Present status}{As of 10 November 2026: Established with sheaf cohomology understood as cohomology of this complex, not global cohomology of its terms.}
\GSfield{Assessment}{literal statement: correct; publication match: exact; verification: read the result; development: annotated; confidence: high}
\end{GSitem}
\begin{GSitem}{1977-437.10}{E}{Free exterior-algebra modules give the zero bundle}
\GSfield{Kind; source}{claim; \GSpage{P000754}}
\GSfield{The notes}{For \(V=\Lambda\), all \(H_i(L_\xi(V))\) vanish. Direct sums and grading shifts give \ru{трив. расслоение}, after the more precise phrase \ru{нулевое расслоение}.}
\GSfield{Statement}{Every finite graded free \(\Lambda\)-module \(P\) satisfies \(H_i(L_\xi(P))=0\) for all \(i\) and \(\xi\ne0\), hence \(\Phi(P)=0\), the rank-zero bundle. It does not yield every positive-rank trivial bundle.}
\GSfield{Reconstruction}{Choose \(a\in\Xi^*\) with \(a(\xi)=1\). Interior multiplication \(\iota_a\) satisfies \(\iota_a(\xi\wedge-)+\xi\wedge\iota_a=1\), proving acyclicity in all dimensions. Direct sums and shifts preserve it.}
\GSfield{Fate}{This is the acyclic exterior Koszul complex; see also \citet[\S2]{bernsteingelfandgelfand1978bundles}.}
\GSfield{Present status}{As of 10 November 2026: Correct after disambiguating trivial as zero. The contraction identity is checked in the critical-verification appendix.}
\GSfield{Assessment}{literal statement: correct with missing hypotheses; publication match: exact; verification: read the result; development: annotated; confidence: high}
\end{GSitem}
\begin{GSitem}{1977-437.11}{A}{Every algebraic vector bundle has a proper-module model}
\GSfield{Kind; source}{theorem; \GSpage{P000754-P000755}}
\GSfield{The notes}{Theorem~1, repeated as Theorem~\(1^{\prime}\)(1), says every bundle is \(\Phi(V)\) for a proper \(V\).}
\GSfield{Statement}{For every finite-rank algebraic vector bundle \(\mathcal E\) on \(\mathbb P^n_{\mathbb C}\), there is a finite-dimensional graded proper module \(V\) over \(\bigwedge\mathbb C^{n+1}\) with \(\Phi(V)\cong\mathcal E\).}
\GSfield{Fate}{Exactly \citet[Theorem~1]{bernsteingelfandgelfand1978bundles}.}
\GSfield{Present status}{As of 10 November 2026: Established; the public-talk evidence in this package is February 1978, not February 1977.}
\GSfield{Assessment}{literal statement: correct; publication match: exact; verification: read the result; development: none; confidence: high; problem \hyperref[prob:GP-1977-01]{GP-1977-01}}
\end{GSitem}
\begin{GSitem}{1977-437.12}{A}{Uniqueness is stable graded isomorphism}
\GSfield{Kind; source}{theorem; \GSpage{P000754-P000755}}
\GSfield{The notes}{Theorem~\(1^{\prime}\)(2) states \(\Phi(V)\cong\Phi(V^{\prime})\) iff \(V\oplus P\cong V^{\prime}\oplus P^{\prime}\) for free \(P,P^{\prime}\).}
\GSfield{Statement}{For two finite-dimensional graded proper \(\Lambda\)-modules, their associated bundles are isomorphic exactly when there exist finite graded free modules \(P,P^{\prime}\) with a degree-zero \(\Lambda\)-module isomorphism \(V\oplus P\simeq V^{\prime}\oplus P^{\prime}\). The two added free modules need not coincide.}
\GSfield{Reconstruction}{The grading and the properness of both modules are carried over from the preceding definition.}
\GSfield{Fate}{Exactly the uniqueness clause of \citet[Theorem~1]{bernsteingelfandgelfand1978bundles}.}
\GSfield{Present status}{As of 10 November 2026: Established. The stronger categorical form quotients morphisms factoring through graded free modules.}
\GSfield{Assessment}{literal statement: correct; publication match: exact; verification: read the result; development: annotated; confidence: high; problem \hyperref[prob:GP-1977-01]{GP-1977-01}}
\end{GSitem}
\begin{GSitem}{1977-437.13}{A}{An independently obtained related theorem of Beilinson}
\GSfield{Kind; source}{remark; \GSpage{P000754}}
\GSfield{The notes}{The theorem is attributed to Bernstein, S.~Gelfand and I.~Gelfand, followed by \ru{независимо близкий рез. Бейлинсон}.}
\GSfield{Statement}{The supported historical assertion is a related independent construction, not identity of all formulations. Beilinson describes the derived category of coherent sheaves on \(\mathbb P^n\) through finite complexes formed from \(\mathcal O(-i)\) or \(\Omega^i(i)\), \(0\leq i\leq n\).}
\GSfield{Reconstruction}{The derived-category formulation is taken from the checked publication; the notebook only reports a similar result. The student-year description is not independently verified.}
\GSfield{Fate}{Compare \citet[\S4, Theorem]{beilinson1978coherent}; received 10 January 1978, published 1978.}
\GSfield{Present status}{As of 10 November 2026: Both correspondences are established; their relationship has constructive formulations in \citet[Theorem~6.1]{eisenbudfloystadschreyer2003exterior}.}
\GSfield{Assessment}{literal statement: correct; publication match: partial; verification: read the result; development: annotated; confidence: high}
\end{GSitem}
\begin{GSitem}{1977-437.14}{E}{Twisting line bundles and homogeneous regular functions}
\GSfield{Kind; source}{definition; \GSpage{P000754-P000755; adjacent check P000758}}
\GSfield{The notes}{\(O(k)\) has sections described as homogeneous rational functions of degree \(k\); \(O(1)\) is called the bundle with a section, and then the Hopf bundle, with \(O(-1)\) offered as an alternative.}
\GSfield{Statement}{For \(n\geq1\), let \(\mathcal O(-1)\) be the tautological line bundle and \(\mathcal O(1)\) its dual. A section of \(\mathcal O(k)\) over a Zariski open \(U\subset\mathbb P(\Xi)\) is a homogeneous regular function of degree \(k\) on its preimage in \(\Xi\setminus\{0\}\). On \(x_i\ne0\), division by \(x_i^k\) gives a regular scalar function and proves local rank one. In particular \(H^0(\mathcal O(1))=\Xi^*\). The case \(n=0\) is a point, on which all twists are isomorphic.}
\GSfield{Reconstruction}{Rational sections must be regular on the specified open set. The separately dated P000758 provides corroborating anti-Hopf wording and asks the elementary rank-one question; it is not a second occurrence assigned to record 437.}
\GSfield{Fate}{Classical twists; \citet[\S54]{serre1955coherent} and \citet[Lemma~30.8.1]{stacks2026projective}.}
\GSfield{Present status}{As of 10 November 2026: The convention is settled by the homogeneous-degree definition. Having a nonzero section does not uniquely characterize O(1).}
\GSfield{Assessment}{literal statement: correct with missing hypotheses; publication match: exact; verification: read the result; development: annotated; confidence: high}
\end{GSitem}
\begin{GSitem}{1977-437.15}{A}{The global linear complex on projective space}
\GSfield{Kind; source}{construction; \GSpage{P000754-P000755}}
\GSfield{The notes}{\(\mathcal L_j=V_{-j}\otimes O(j)\); on sections \((df)(\xi)=\xi f(\xi)\), giving \(\mathcal L(V)\).}
\GSfield{Statement}{For any finite graded \(\Lambda\)-module \(V\), define \(\mathcal L^j(V)=V_{-j}\otimes\mathcal O(j)\) and \(d=\sum_{i=0}^n x_i\otimes\xi_i\). The tensor \(\sum x_i\otimes\xi_i\) is independent of dual bases. Since the \(x_i\) commute and the \(\xi_i\) anticommute and square to zero, \(d^2=0\). The sheaf morphisms have the indicated twists, and the complex is bounded.}
\GSfield{Reconstruction}{Superscripts distinguish cohomological degree from the negative internal grading; the differential formula itself is the notebook formula.}
\GSfield{Fate}{This is the functor in \citet[\S3]{bernsteingelfandgelfand1978bundles}.}
\GSfield{Present status}{As of 10 November 2026: Established. An arbitrary module gives a complex; properness is needed for a single vector bundle.}
\GSfield{Assessment}{literal statement: correct; publication match: exact; verification: read the result; development: annotated; confidence: high}
\end{GSitem}
\begin{GSitem}{1977-437.16}{A}{Twisting by O(1) is a shifted cosyzygy}
\GSfield{Kind; source}{question; \GSpage{P000755}}
\GSfield{The notes}{What does \(\mathcal E\mapsto\mathcal E\otimes O(1)\) mean for modules? For \(0\to V\to P\to V^{\prime}\to0\), \(P\) free, set \(W_j=V^{\prime}_{j+1}\); then \(\Phi(W)=\Phi(V)\otimes O(1)\).}
\GSfield{Statement}{Let \(V\) be proper and \(0\to V\to P\to V^{\prime}\to0\) exact in finite graded \(\Lambda\)-modules, with \(P\) graded free. Then the grading shift \(W_j=V^{\prime}_{j+1}\) is proper and \(\Phi(W)\cong\Phi(V)(1)\). Such a free embedding exists because the finite-dimensional exterior algebra is Frobenius and its graded projectives are injective.}
\GSfield{Reconstruction}{The long exact fibre-homology sequence gives \(H_i(V^{\prime})\cong H_{i-1}(V)\). Hence \(W\) is proper. Globally \(\mathcal L^j(W)=\mathcal L^{j-1}(V^{\prime})(1)\); the exact sequence gives the asserted degree-zero cohomology. This also checks the sign of the shift.}
\GSfield{Fate}{Exactly \citet[remark~2 after Theorem~1]{bernsteingelfandgelfand1978bundles}.}
\GSfield{Present status}{As of 10 November 2026: Established. Both the cosyzygy and internal grading shift are needed.}
\GSfield{Assessment}{literal statement: correct; publication match: exact; verification: read the result; development: annotated; confidence: high; problem \hyperref[prob:GP-1977-02]{GP-1977-02}}
\end{GSitem}
\begin{GSitem}{1977-437.17}{E}{Large positive twists are globally generated}
\GSfield{Kind; source}{theorem; \GSpage{P000755}}
\GSfield{The notes}{For sufficiently large \(N\), Serre’s theorem gives enough sections of \(\mathcal E\otimes O(N)\) through every vector in every fibre, called \ru{обильное расслоение}.}
\GSfield{Statement}{If \(\mathcal E\) is coherent on \(\mathbb P^n_{\mathbb C}\), there is \(N_0\) such that for every \(N\geq N_0\) the evaluation map \(H^0(\mathcal E(N))\otimes\mathcal O\twoheadrightarrow\mathcal E(N)\) is surjective. The section space is finite-dimensional. The property used here is global generation, not the definition of an ample vector bundle.}
\GSfield{Reconstruction}{The evaluation formulation specifies the phrase many sections. Ampleness and generation are different properties.}
\GSfield{Fate}{\citet[Chapter~III, \S65]{serre1955coherent}; modern hypotheses in \citet[Lemma~30.14.1]{stacks2026projective}.}
\GSfield{Present status}{As of 10 November 2026: Established; in an inverse BGG algorithm the twist must also be sufficiently regular for exactness.}
\GSfield{Assessment}{literal statement: correct with missing hypotheses; publication match: exact; verification: read the result; development: annotated; confidence: high}
\end{GSitem}
\begin{GSitem}{1977-437.18}{E}{The graded section module and coordinate multiplications}
\GSfield{Kind; source}{construction; \GSpage{P000755}}
\GSfield{The notes}{\(C_N=\Gamma(\mathbb P^n,\mathcal E(N))\), \(\dim\Gamma(O(1))=n+1\), and \(X_i:C_N\to C_{N+1}\) multiplies by the dual coordinate \(x_i\).}
\GSfield{Statement}{Put \(S=\operatorname{Sym}(\Xi^*)\) and \(C_{\geq N}=\bigoplus_{m\geq N}H^0(\mathcal E(m))\), for a vector bundle \(\mathcal E\) on \(\mathbb P(\Xi)\). Multiplication by \(x_i\in\Xi^*\) defines commuting maps between successive components; for a sufficiently large truncation this is a finite graded \(S\)-module.}
\GSfield{Reconstruction}{Finite generation refers to the graded module of all sufficiently large twists, not to a single vector space with a polynomial action.}
\GSfield{Fate}{The classical construction is in \citet[Chapter~III]{serre1955coherent}. Finite generation is stated precisely in \citet[Lemma~30.14.1(5)]{stacks2026projective}.}
\GSfield{Present status}{As of 10 November 2026: Standard. The commutation relations imply the square-zero identity in the ensuing Koszul-dual complex.}
\GSfield{Assessment}{literal statement: correct; publication match: exact; verification: read the result; development: annotated; confidence: high}
\end{GSitem}
\begin{GSitem}{1977-437.19}{A}{Recovering a module from two consecutive section spaces}
\GSfield{Kind; source}{method; \GSpage{P000756}}
\GSfield{The notes}{The notes put \(V=C_N\otimes\Lambda\), \(V^{\prime}=C_{N+1}\otimes\Lambda\), \(\partial(c\otimes\lambda)=\sum_iX_i(c)\otimes\xi_i\lambda\), and \(W=\ker\partial\), followed by \(N\) inverse shifts using free covers.}
\GSfield{Statement}{Let \(\mathcal E\) be an algebraic vector bundle on \(\mathbb P(\Xi)\), \(S=\operatorname{Sym}(\Xi^*)\), and \(M=\bigoplus_{m\geq0}H^0(\mathcal E(m))\). Choose an integer \(N>\max\{0,\operatorname{reg}M\}\). In cohomological degree \(m\), put \(R(M)^m=\operatorname{Hom}_{\mathbb C}(\Lambda,M_m)\), with internal component \[R(M)^m_t=\operatorname{Hom}_{\mathbb C}(\Lambda_{m-t},M_m).\] The differential is \((d_Rf)(\lambda)=\sum_i x_i f(\xi_i\lambda)\); it preserves internal degree. The tail with \(m\geq N\) is exact in degrees strictly greater than \(N\). Adjoin a minimal graded free resolution of its initial kernel, retaining these cohomological indices, to obtain the everywhere exact Tate complex \(T(\mathcal E)\). Then \[V=\ker\bigl(T(\mathcal E)^0\longrightarrow T(\mathcal E)^1\bigr)\] is a finite-dimensional graded proper module and \(\Phi(V)\cong\mathcal E\).}
\GSfield{Reconstruction}{Both gradings, the regularity bound and the degree-zero kernel are supplied here; the notebook gives an abbreviated two-section-space construction. In the modern dual-module convention, \((a\cdot f)(\lambda)=f(\lambda a)\), so \(d_R\) is a degree-zero \(\Lambda\)-module map. By contrast the notebook's raw tensor formula has internal degree \(-1\) and anticommutes with the generators, although its kernel is still a graded submodule. For \(n=0\), \(\mathcal E=\mathcal O\), its raw kernel is \(\mathbb C\xi_0\) in degree \(-1\), hence is not proper. The dual free module \(\operatorname{Hom}_{\mathbb C}(\Lambda,\mathbb C)\) has degrees \(0,\ldots,n+1\), and is noncanonically \(\Lambda(-(n+1))\) under \(Q(a)_j=Q_{a+j}\); omitting that shift changes the answer.}
\GSfield{Fate}{The notebook sketch belongs to the inverse BGG construction; compare \citet[\S~4]{bernsteingelfandgelfand1978bundles}. The explicit conventions above are \citet[\S~2, Corollary~2.4, Theorem~3.7 and the construction in \S~4]{eisenbudfloystadschreyer2003exterior}.}
\GSfield{Present status}{As of 10 November 2026: The construction is established, but the notebook omits the grading of the two free modules; its unqualified kernel cannot be certified as already proper.}
\GSfield{Assessment}{literal statement: correct with missing hypotheses; publication match: partial; verification: read the result; development: annotated; confidence: medium}
\end{GSitem}
\begin{GSitem}{1977-437.20}{E}{Vector bundles must be enlarged to coherent sheaves}
\GSfield{Kind; source}{remark; \GSpage{P000756}; I. M. Gelfand}
\GSfield{The notes}{I.~M.~Gelfand says classification of bundles is not quite the right problem because they do not form an abelian category; one must allow fibre dimensions to vary and introduce sheaves.}
\GSfield{Statement}{For \(n\geq1\), finite-rank locally free sheaves on \(\mathbb P^n\) are not an abelian category with the inherited sheaf morphisms: multiplication by a nonzero linear form gives \(\mathcal O(-1)\to\mathcal O\) whose sheaf cokernel is the structure sheaf of a hyperplane and is not locally free. Coherent sheaves form an abelian category. For \(n=0\), vector bundles are finite-dimensional vector spaces and already form an abelian category.}
\GSfield{Reconstruction}{More intrinsically, that multiplication map is monic and epic within the category of locally free sheaves, since a map from its torsion cokernel to a locally free sheaf is zero, but is not an isomorphism on the hyperplane. An abelian category is balanced, so it cannot have such a map. The point case is exceptional.}
\GSfield{Fate}{Coherent-sheaf formalism predates the lecture: \citet[Chapter~I, \S2]{serre1955coherent}.}
\GSfield{Present status}{As of 10 November 2026: Correct with the positive-dimensional qualification; variable fibre rank alone is not a definition of a coherent sheaf.}
\GSfield{Assessment}{literal statement: correct with missing hypotheses; publication match: exact; verification: read the result; development: annotated; confidence: high}
\end{GSitem}
\begin{GSitem}{1977-437.21}{E}{Serre correspondence is a quotient by finite-dimensional graded modules}
\GSfield{Kind; source}{theorem; \GSpage{P000756}}
\GSfield{The notes}{The notes equate sheaves on \(\mathbb CP^n\) with finitely generated graded \(\mathbb C[x_0,\ldots,x_n]\)-modules up to finitely many homogeneous components.}
\GSfield{Statement}{For \(S=\mathbb C[x_0,\ldots,x_n]\) with standard grading, sheafification induces an equivalence \(\operatorname{grmod}(S)/\operatorname{fdim}(S)\simeq\operatorname{Coh}(\mathbb P^n_{\mathbb C})\). Here the numerator has finitely generated graded modules and degree-zero maps, and the denominator is its Serre subcategory of finite-dimensional modules. A quasi-inverse is the nonnegative graded section module followed by passage to the quotient.}
\GSfield{Reconstruction}{A quotient category replaces the informal phrase ignoring finitely many components; coherence is an essential omitted hypothesis.}
\GSfield{Fate}{\citet[Chapter~III, \S\S57--65]{serre1955coherent}; \citet[Proposition~30.14.4]{stacks2026projective}.}
\GSfield{Present status}{As of 10 November 2026: Established. It concerns algebraic coherent sheaves; arbitrary sheaves are not covered.}
\GSfield{Assessment}{literal statement: correct with missing hypotheses; publication match: exact; verification: read the result; development: annotated; confidence: high}
\end{GSitem}
\begin{GSitem}{1977-437.22}{E}{Sheafification through degree-zero localization}
\GSfield{Kind; source}{construction; \GSpage{P000756}}
\GSfield{The notes}{On \(E_i=\{x_i\ne0\}\), \(\Gamma(E_i,L)=\{m/x_i^s:\deg m=s,\ m\in M\}\).}
\GSfield{Statement}{For a graded \(S\)-module \(M\), set \(\widetilde M(D_+(x_i))=(M_{x_i})_0\). Homogeneous fractions have \(\deg m=s\); equality is equality in the localization, so elements killed by powers of \(x_i\) vanish. The localization maps on overlaps glue these affine modules into a quasi-coherent sheaf, coherent if \(M\) is finitely generated.}
\GSfield{Reconstruction}{The equivalence relation on fractions and the gluing maps are implicit in the notebook.}
\GSfield{Fate}{\citet[Lemma~27.8.4]{stacks2026projconstruction}.}
\GSfield{Present status}{As of 10 November 2026: Established and literally correct when the braces mean localized fractions.}
\GSfield{Assessment}{literal statement: correct; publication match: exact; verification: read the result; development: annotated; confidence: high}
\end{GSitem}
\begin{GSitem}{1977-437.23}{E}{Shifting the polynomial module yields O(N)}
\GSfield{Kind; source}{construction; \GSpage{P000756}}
\GSfield{The notes}{\(O(N)\) arises from \(M=\mathbb C[x_0,\ldots,x_n]\) with \(M_i=S_{i+N}\).}
\GSfield{Statement}{Define \(S(N)_i=S_{i+N}\). Then \(\widetilde{S(N)}\cong\mathcal O_{\mathbb P^n}(N)\). Locally a homogeneous numerator of degree \(s+N\), divided by \(x_i^s\), has homogeneous degree \(N\), exactly the earlier section convention.}
\GSfield{Reconstruction}{The local fraction check compares the final formula with the earlier homogeneous-function definition.}
\GSfield{Fate}{The standard twisting convention of \citet[\S54]{serre1955coherent}.}
\GSfield{Present status}{As of 10 November 2026: Correct; this independently fixes the sign of the polynomial-module shift.}
\GSfield{Assessment}{literal statement: correct; publication match: exact; verification: read the result; development: annotated; confidence: high}
\end{GSitem}

\section{Developments}\label{dev:1976}
No C or D occurrence qualified for a bounded new research contribution.
The subsections below are critical verifications of statements, calculations
and omissions in the notebook, not newly solved open problems. All results
proved or computed in this task are marked with the supplied authorship
macros; their exact scope and independent checks are stated in the text.
Historical classifications remain those in the occurrence index.

\subsection{Critical verification of the early records}

The following checks audit mathematical statements faithfully present in the manuscript. They make no novelty claim and are not developments of an invented open problem. Every calculation was re-derived independently or checked by a second algebraic method; the accompanying \texttt{early\_checks.py} and its output preserve the exact computations. A separate reviewer independently checked the kernel formula as well.

\begin{proposition}[The missing factor in the projector kernel]\ours{}
Let \(|\lambda|>1\), and let \(t=\lambda+\lambda^{-1}\). The character-series expression occurring in the holomorphic discrete-series projector, with its common Haar normalization suppressed, is
\[
 K(\lambda)=\frac1{\lambda-\lambda^{-1}}
 \sum_{n=1}^{\infty}n\lambda^{-n}
 =\frac{\lambda^2}{(\lambda-1)^3(\lambda+1)}.
\]
Consequently
\[
 K(\lambda)^2(t-2)^3(t+2)=1.
\]
Thus the invariant trace expression is proportional, with a consistently chosen holomorphic branch, to \([(t-2)^3(t+2)]^{-1/2}\), as in \citet[Theorem~10, equations~18 and~22]{gelfandgindikin1977complex}. The notebook's \((2-t)^{-3/2}\) lacks a nonconstant factor.
\end{proposition}
\begin{proof}
Put \(r=\lambda^{-1}\). The geometric series converges normally on compact subsets of \(|r|<1\); differentiating it gives
\(\sum_{n\ge1}nr^n=r/(1-r)^2\). Dividing by \(\lambda-\lambda^{-1}\) and simplifying gives the rational formula. Independently,
\[
 t-2=\frac{(\lambda-1)^2}{\lambda},\qquad
 t+2=\frac{(\lambda+1)^2}{\lambda},
\]
so direct substitution proves the squared identity without relying on a branch choice. The phase difference between \(t-2\) and \(2-t\) may be included in the normalization constant on a branch, but \((t+2)^{-1/2}\) cannot be. This is a check of the published character-series kernel, not a claim that arbitrary scalar square roots define a global projector without their boundary-value prescriptions.

As an independent finite check, for \(0<r<1\),
\[
 \sum_{n>N}nr^n=\frac{r^{N+1}((N+1)-Nr)}{(1-r)^2}.
\]
Exact rational arithmetic, using no symbolic summation, verified the finite sum plus this tail at \(\lambda=3/2,2,3,5\), each with \(N=1,2,7,23\). A separate reviewer independently obtained the same rational and trace identities. These checks confirm the algebra; the publication supplies the representation-theoretic projector interpretation.
\end{proof}

\begin{proposition}[Quadratic normalization near a solution]\ours{}
Suppose, near \(y=p=0\),
\[
 y''=a(x)y^2+b(x)yp+c(x)p^2+O((|y|+|p|)^3),\qquad p=y'.
\]
Then a local point transformation whose Jacobian is one on the reference solution removes the two mixed quadratic terms and leaves
\[
 Y_{XX}=A(X)Y^2+O((|Y|+|Y_X|)^3),\qquad
 A=a+\frac16c''-\frac13b'.
\]
\end{proposition}
\begin{proof}
Use \(X=x+f(x)y\), \(Y=y+g(x)y^2\). Along the equation put \(D=\partial_x+p\partial_y+F\partial_p\). The exact transformation formula is
\[
 Y_{XX}=\frac{D^2Y\,DX-DY\,D^2X}{(DX)^3}.
\]
Its quadratic part is
\[
 (a+g'')y^2+(b+4g'-f'')yp+(c+2g-2f')p^2.
\]
Choose \(g=f'-c/2\) and integrate \(f''=(2c'-b)/3\) locally. The latter equation has smooth local solutions; the coordinate Jacobian is one at \(y=0\), so the inverse function theorem applies. The surviving coefficient is the displayed \(A\). Expressing this quadratic polynomial in the new variables changes it only by terms of order at least three.

The symbolic audit expanded the full chain-rule expression with \((y,p)\mapsto(\varepsilon y,\varepsilon p)\). A second check used the concrete polynomials \(a=1+x^2\), \(b=x+2\), \(c=x^3-x\), \(f=x^4/6-x^3/18-2x^2/3\); the mixed terms vanished and the remaining coefficient was \(x^2+x+2/3\). The classical Tresse expression in occurrence~1976-403.1, restricted to the reference solution, independently gives \(H=12A\).
\end{proof}

\begin{example}[A limitation of the density criterion]\ours{}
The analytic equation \(y''=y^3\) has the solution \(y=0\), zero constant and linear transverse terms, and zero quadratic transverse coefficient. Thus the density constructed along this particular solution vanishes. Nevertheless, the classical relative invariant is \(H=6F_{yy}=36y\), which is a nonzero formal series. The equation cannot be formally point-flat. This checks the literal one-solution reading of P000666; it does not rule out a stronger intended statement with additional quantifiers and hypotheses. Direct substitution into the complete \(H\)-formula and the shorter \(F_p=0\) calculation agree.
\end{example}

\begin{proposition}[The projective vector-field invariant]\ours{}
Let \(v(x)\partial_x\) be a specified smooth vector field and let \(X=\phi(x)\), \(\phi'\ne0\). Write \(V(X)=\phi'(x)v(x)\). Then
\[
 2VV_{XX}-V_X^2
 =2vv''-(v')^2+2v^2S(\phi),\qquad
 S(\phi)=\frac{\phi^{(3)}}{\phi'}-
 \frac32\left(\frac{\phi''}{\phi'}\right)^2.
\]
Hence \(I_2=2vv''-(v')^2\) and \(I_n=(v\partial_x)^{n-2}I_2\) are scalar invariants under projective coordinate changes. Under \(v\mapsto cv\), they scale as \(I_n\mapsto c^nI_n\).
\end{proposition}
\begin{proof}
Since \(\partial_X=(\phi')^{-1}\partial_x\),
\[
 V_X=v'+\frac{\phi''}{\phi'}v,
\]
and one more differentiation proves the identity. The Schwarzian is zero for a Möbius map. Also \(V\partial_X=v\partial_x\), giving invariant differentiation and the scaling law by induction. A separate symbolic check used each of \(\phi=(2x+1)/(x+1)\) and \(\phi=(x-2)/(3x+4)\), with \(v=1+x+x^2+x^4\), and obtained exact zero differences. This confirms the notebook formulas for a specified field; it does not eliminate the constant ambiguity inherited from its relative density.
\end{proof}

\begin{proposition}[Auxiliary-field elimination in the written action]\ours{}
With precisely the auxiliary terms on P000669, let
\[
 P=mA+g(A^2-B^2),\qquad Q=mB+2gAB.
\]
Then
\[
 \mathcal L_{\rm aux}=\tfrac12(F^2+G^2)+FP+GQ
\]
has unique stationary values \(F=-P\), \(G=-Q\), and its eliminated value is
\[
 -\frac12m^2(A^2+B^2)-mgA(A^2+B^2)
 -\frac12g^2(A^2+B^2)^2.
\]
\end{proposition}
\begin{proof}
Differentiate with respect to the algebraic variables \(F,G\), then substitute. Independently, completing squares gives
\[
 \mathcal L_{\rm aux}=\tfrac12(F+P)^2+\tfrac12(G+Q)^2
 -\tfrac12(P^2+Q^2).
\]
The identity
\((A^2-B^2)^2+4A^2B^2=(A^2+B^2)^2\)
gives the quartic term, and the two cross terms give the negative cubic coefficient. The code checked both polynomial identities and three independent integer parameter sets. This proves that the two notebook scalar formulas disagree under their printed conventions. Since \(P,Q\) contain no fermions, the operation cannot repair the separate factor-of-two discrepancy in the fermion mass term. The intended pre-existing construction is verified in \citet[\S~4(b)]{ogievetskymezincescu1975superfields}.
\end{proof}

\begin{remark}[What the audit does not settle]
The five unidentified early occurrences concern an unspecified Radon range, an undefined minimality criterion, a damaged global projective question, a cohomological connection without a map or degree, and an uncertainty-principle slogan without a Hamiltonian. None supplies the precise remaining scope required for class C or D development. The completed audits concern literal correctness and finite calculations; they are not substitutes for a literature-supported open question.
\end{remark}


\subsection{Critical checks on the Berezinian and flatness signs}

These calculations audit faithful mathematical errors in the source. They are
elementary critical annotations, not a claim to new research and not a development
of a class C or D problem.

\paragraph{The reciprocal Berezinian on P000673.}\oursComputed{}
For an even block-diagonal matrix \(J=\operatorname{diag}(a,d)\) of size
\(1|1\), with nonzero even scalar entries, the usual Berezinian is
\(\operatorname{Ber}J=a/d\). The first expression on P000673 gives \(d/a\),
whereas the lower expression on the same page gives \(a/d\).
Taking \(a=2,d=3\) gives \(3/2\) versus \(2/3\), so the two expressions
cannot be simultaneous definitions in the stated convention.
The usual normalization is explicitly printed in
\citet[\S1]{bernsteinleites1977integral} and in the earlier density formula
\citep[formula~(3)]{berezinleites1975supermanifolds}.
The correct general formula is
\[
 \operatorname{Ber}\begin{pmatrix}A&B\\C&D\end{pmatrix}
 =\frac{\det(A-BD^{-1}C)}{\det D}
 =\frac{\det A}{\det(D-CA^{-1}B)}.
\]
This test isolates a reciprocal error without depending on conventions for
left versus right odd derivatives.

\paragraph{Flatness on P000684 and in the printed Manin theorem.}\ours{}
Use the notebook's convention
\(\nabla_1\mathbf1=P\mathbf1\) and
\(\nabla_2\mathbf1=L\mathbf1\).
Because both connections commute with \(\nabla\), the coefficientwise
Leibniz rule gives
\[
 [\nabla_1,\nabla_2]\mathbf1
 = (d_1L+LP-d_2P-PL)\mathbf1.
\]
The cyclic filtered-module theorem makes the action on \(\mathbf1\)
faithful for differential operators, hence flatness implies
\[
 d_1L-d_2P=[P,L].
\]
Both the difference sign and the order of the commutator follow from the
declared convention. An unspecified \(\pm\) in front of a commutator cannot
repair a sum of the two derivatives.

For a concrete counterexample to that sum, take
\(\mathcal B=\mathbb Q[t,z]\), \(d=0\),
\(\mathcal M=\mathcal B[D]\), and
\[
 \nabla Q=DQ,\qquad
 \nabla_1Q=\partial_tQ+Qz,\qquad
 \nabla_2Q=\partial_zQ+Qt.
\]
The standard order filtration satisfies the module axioms; the three
connections commute. Their values on \(1\) are \(P=z\) and \(L=t\).
Thus \(d_1L-d_2P=1-1=0=[P,L]\), whereas
\(d_1L+d_2P=2\ne0\). The accompanying exact symbolic check verifies this
on a general test polynomial in \(t,z,D\).

The error is also present in the inspected printed English statement
\citep[Chapter~III, Theorem~2.6, p.~83]{manin1978algebraic}.
That statement uses the opposite names,
\(\nabla_1\mathbf1=L\mathbf1\), \(\nabla_2\mathbf1=P\mathbf1\),
and prints \(d_1P+d_2L=[P,L]\). Its proof requires instead
\(d_1P-d_2L=[L,P]\). This identifies a shared mathematical sign error,
not a failure to transcribe the notebook.

\paragraph{Independent fourth-order check.}\oursComputed{}
For P000681, the companion's separate \texttt{critical\_checks.py} computes
the commutator by two methods. With \(L=D^4+vD^2+wD+z\),
\(P=D^2+AD+v/2\), and \(A_x=0\), its coefficients are
\[
\begin{aligned}
[L,P]_{D^2}&=2v_{xx}-Av_x-2w_x,\\
[L,P]_{D}&=2v_{xxx}+vv_x-w_{xx}-Aw_x-2z_x,\\
[L,P]_{1}&=\tfrac12v_{xxxx}+\tfrac12vv_{xx}+\tfrac12wv_x-z_{xx}-Az_x.
\end{aligned}
\]
The notebook's first two relative coefficients are therefore wrong, while
the third equation agrees. Manin's later equation has the corrected relative
coefficients, with its own time and commutator conventions
\citep[Chapter~III, \S6.3, equation~(17), p.~96]{manin1978algebraic}.


\subsection{Audit of the assertion that changing the complex structure merges any two packets}

The notebook, P000702, third line from the bottom before the final diagram, says
\ru{Любые два неприв. предст. эквивал. относит. нек-рой рим. пов. (вх. в какой-то пакет).}
The scan confirms the two universal quantifiers; the word abbreviated as
\ru{неприв.} is overwritten but is the diplomatic transcription's reading.
The definition earlier on the same page makes equivalence mean
\(E_{\rho_1}\cong E_{\rho_2}\) as holomorphic bundles.  The assertion therefore
includes the following rank-one case, in which it is false.  The author's own
survey explicitly gives uniqueness of the unitary representation in a packet
as an exercise \citep[Ch.~II, \S2, Exercise~10, p.~204]{tyurin1978periods}.
This is an audit
of a literal assertion, not the development of a C or D problem.

\begin{proposition}[Distinct unitary characters remain in distinct packets]\ours{}
Let \(\Sigma\) be a connected, closed, oriented surface of genus \(g\geq1\),
and let \(\chi_1,\chi_2:\pi_1(\Sigma)\to U(1)\) be two distinct characters.
For every complex structure \(J\) compatible with the orientation, the
holomorphic line bundles obtained from the two unitary flat bundles on
\(X=(\Sigma,J)\) are nonisomorphic.  In particular, the notebook's assertion
about arbitrary pairs of irreducible representations is false.
\end{proposition}

\begin{proof}
Write \(L_i\) for the flat line bundle with monodromy \(\chi_i\), with its
parallel Hermitian metric and holomorphic structure given by the
\((0,1)\)-part of the flat connection.  Suppose that
\(\Phi:L_1\to L_2\) is a holomorphic bundle isomorphism.  It is a nowhere-zero
holomorphic section \(s\) of the unitary flat line bundle
\(L=L_2\otimes L_1^{-1}\), whose monodromy is
\(\delta=\chi_2\chi_1^{-1}\).

On any simply connected coordinate neighbourhood choose a parallel frame
\(e\) of norm one for \(L\).  Because the \((0,1)\)-connection defines the
holomorphic structure, \(e\) is holomorphic.  Thus \(s=f e\) with \(f\)
holomorphic and nowhere zero.  The real function
\(u=\log\lVert s\rVert=\log|f|\) is harmonic on that neighbourhood.  Unitary
transition functions make \(u\) a globally defined harmonic function on the
compact connected Riemann surface \(X\); hence \(u\) is constant.  Each
holomorphic function \(f\) therefore has constant modulus and is constant.
Consequently \(s\) is a nonzero parallel section of \(L\).  Parallel
transport around every loop fixes its nonzero value, so \(\delta=1\) and
\(\chi_1=\chi_2\), a contradiction.

For an explicit pair, use the presentation in the notebook.  Take
\(\chi_1=1\), let \(\chi_2(a_1)=e^{i\theta}\) for
\(\theta\notin2\pi\mathbb Z\), and send every other standard generator to
\(1\).  The surface relation holds because \(U(1)\) is abelian.  Both are
one-dimensional complex representations and therefore irreducible.  The
argument applies to every compatible \(J\), including genus \(g\geq2\)
if that was the intended range.  Even allowing a change of marking cannot
turn the nontrivial character into the trivial character.
\end{proof}

\begin{remark}[Independent check]
The proof was checked again by working on the universal cover.  A
holomorphic isomorphism there is multiplication by a nowhere-zero
holomorphic function \(F\) with
\(F(\gamma u)=\delta(\gamma)F(u)\).  Since \(|\delta(\gamma)|=1\),
\(\log|F|\) descends to a harmonic function on \(X\), hence is constant;
the open mapping theorem then makes \(F\) constant and forces
\(\delta=1\).  This check uses neither the higher-rank stability theorem
nor an assumption that all flat holomorphic bundles are stable.
\end{remark}


\subsection{Critical checks for the miscatalogued exterior-algebra lecture}

These checks concern record~437, whose manuscript date is 13 February 1978.
They audit the supplied material without turning it into evidence of a
1976--77 research problem. The results are elementary verifications of the
construction recorded in the notebook; no claim of novelty is intended.

\begin{proposition}[The regular exterior module is acyclic]\ours{}
Let \(k\) be a field, let \(\Xi\) have finite dimension \(m\geq1\), and
let \(0\ne\xi\in\Xi\). Give \(\Lambda=\bigwedge\Xi\) its negative
grading and differential \(d(v)=\xi\wedge v\). Then the complex is
contractible. Consequently every finite direct sum of grading shifts of
\(\Lambda\) gives the zero bundle under the construction of record~437.
\end{proposition}
\begin{proof}
Choose \(a\in\Xi^*\) with \(a(\xi)=1\). Define contraction on a simple
exterior tensor by
\[
 \iota_a(v_1\wedge\cdots\wedge v_r)
 =\sum_{j=1}^r(-1)^{j-1}a(v_j)
 v_1\wedge\cdots\wedge\widehat{v_j}\wedge\cdots\wedge v_r.
\]
Applying this formula to \(\xi\wedge v\), the term contracting the first
factor is \(v\); each remaining term is the negative of the corresponding
term of \(\xi\wedge\iota_a(v)\). Therefore
\(\iota_a d+d\iota_a=1\). If \(dv=0\), this identity gives
\(v=d\iota_a(v)\), so every cycle is a boundary. The homotopy commutes
with taking direct sums and does not depend on the choice of grading
origin. In particular the output has rank zero, which resolves the
notebook's ambiguous use of ``trivial'' in this example.
\end{proof}

\begin{remark}[Independent checks]\oursComputed{}
The identity was independently re-derived by writing
\(\Xi=k\xi\oplus U\), hence
\(\Lambda=\bigwedge U\oplus\xi\wedge\bigwedge U\). The differential
maps the first summand isomorphically to the second and kills the second.
The exact-arithmetic program \texttt{code/critical\_checks.py} additionally
constructs exterior multiplication and contraction by two separate
bit-mask formulas. For \(m=1,\ldots,6\) and three nonzero vectors in each
dimension, all 18 cases satisfy \(d^2=0\), \(\iota_a d+d\iota_a=1\),
and zero homology in every degree. The ranks agree with
\(\binom{m-1}{r}\) in exterior degree \(r\). These finite checks support
the implementation; the preceding proof establishes the assertion for
every \(m\), not an extrapolation from the computed cases.
\end{remark}

\begin{example}[The inverse construction requires a grading convention]\ours{}
Take \(n=0\), so the base is a point, and let
\(\mathcal E=\mathcal O\). For every \(N\geq0\), the two consecutive
section spaces are one-dimensional and \(X_0\) is the identity. If the
two tensor factors in P000756 are given their raw exterior grading, then
\[
 \partial:\Lambda\longrightarrow\Lambda,\qquad
 \partial(1)=\xi_0,\qquad \partial(\xi_0)=0.
\]
Its kernel is \(k\xi_0\), concentrated in internal degree \(-1\).
The action of \(\xi_0\) on that kernel is zero, so its fibre homology
is nonzero in degree \(-1\); the raw kernel is not proper. This is an
explicit check of an omitted grading convention, not a counterexample to
the BGG correspondence. The dual free module
\(\operatorname{Hom}_k(\Lambda,k)\) differs from \(\Lambda\) by a
grading shift, which is essential in the precise inverse construction
\citep[\S4]{bernsteingelfandgelfand1978bundles}.
\end{example}

\begin{proposition}[Sign of the twist operation]\ours{}
Use the negative grading of record~437. If \(V\) is proper and
\(0\to V\to P\to V'\to0\) is an exact sequence of finite graded
exterior-algebra modules with \(P\) free, set \(W_j=V'_{j+1}\).
Then \(W\) is proper and \(\Phi(W)=\Phi(V)(1)\).
\end{proposition}
\begin{proof}
For each nonzero \(\xi\), applying \(L_\xi\) gives a short exact
sequence of complexes. The middle complex is acyclic by the proposition
above, so its homology sequence gives
\(H_j(L_\xi(V'))\cong H_{j-1}(L_\xi(V))\). Consequently the only
possibly nonzero homology of \(V'\) is in chain degree one, and the
displayed grading shift moves it to degree zero. On projective space,
\[
 \mathcal L^j(W)
 =V'_{1-j}\otimes\mathcal O(j)
 =\mathcal L^{j-1}(V')(1).
\]
The short exact sequence of sheaf complexes gives
\(\mathcal H^{-1}(\mathcal L(V'))\cong\mathcal H^0(\mathcal L(V))\).
Reindexing by the displayed formula proves the claimed twist. Depending
on the convention for a shifted complex, a minus sign in its differential
is removed by multiplying degree \(j\) by \((-1)^j\); cohomology and
the bundle are unaffected.
\end{proof}

\begin{remark}[Second check of indexing]
There are two independent index tests: the chain homology sequence
places the cokernel's fibre homology in degree one, whereas the sheaf
cohomology sequence places it in degree minus one. The substitution
\(j\mapsto-j\) in \(\mathcal L^j(V)=V_{-j}\otimes\mathcal O(j)\)
reconciles them and gives \(\mathcal O(1)\), not \(\mathcal O(-1)\).
This is exactly the convention in remark~2 after Theorem~1 of
\citet{bernsteingelfandgelfand1978bundles}.
\end{remark}


\begin{remark}[Independent check of the inverse-module indexing]\ours{}
The conventions of occurrence~1977-437.19 were checked against
\citet[\S2, Corollary~2.4, Theorem~3.7 and \S4]{eisenbudfloystadschreyer2003exterior}.
Put \(Z_m=\ker(T^m\to T^{m+1})\), use internal shifts
\(Q(a)_j=Q_{a+j}\), and use cochain shifts \(K[a]^j=K^{a+j}\).
At a sufficiently large cutoff \(N\), the BGG reciprocity theorem,
applied to the regraded kernel, gives
\(\mathcal L(Z_N(N))\simeq\mathcal E(N)\).
Direct substitution in the definition of the linear complex gives
\[
 \mathcal L(Q(N))=\mathcal L(Q)[-N](N).
\]
Therefore \(\mathcal L(Z_N)\simeq\mathcal E[N]\).
For every integer \(m\), exactness of the Tate complex gives
\(0\to Z_m\to T^m\to Z_{m+1}\to0\).
The middle term is graded free, and its linear sheaf complex is acyclic
by the contraction above. Its associated triangle therefore gives
\(\mathcal L(Z_{m+1})\simeq\mathcal L(Z_m)[1]\).
Iteration yields \(\mathcal L(Z_N)\simeq\mathcal L(Z_0)[N]\), hence
\(\mathcal L(Z_0)\simeq\mathcal E\) with no additional twist.
The bounded complex of locally free sheaves may be split locally away
from its locally free degree-zero cohomology, so its fibres have no
other homology: \(Z_0\) is proper. The point case gives the trivial
exterior module in internal degree zero. On \(\mathbb P^1\), the case
\(\mathcal E=\mathcal O(-1)\) gives the Euler kernel of
\(\mathcal O^2\to\mathcal O(1)\), again checking the grading sign.
These are checks of conventions in the cited construction, not a new
proof of the BGG reciprocity theorem.
\end{remark}

\section{Problem index}\label{problems:1976}
Earlier-year occurrence identifiers are retained as text; current-year identifiers link to their entries.

\subsection{GP-1976-01: Analytic separation of the regular representation of SL(2, R) into series}\label{prob:GP-1976-01}
Construct invariant projectors separating the discrete and continuous parts of \(L^2(SL(2,\mathbb R))\), and realize the two discrete parts as holomorphic Hardy boundary values on explicitly specified invariant domains. This is narrower than the Plancherel decomposition into all individual irreducibles.
\par\noindent\textit{Occurrences:} \hyperref[item:1976-398.5]{1976-398.5}, \hyperref[item:1976-398.14]{1976-398.14}, \hyperref[item:1976-398.16]{1976-398.16}, \hyperref[item:1976-398.17]{1976-398.17}.
\par\noindent\textit{Status (2026-11):} As of 10 November 2026, resolved for this group. Spectral separation predates 1976; the specific invariant-domain and Hardy implementation was published by Gelfand and Gindikin in 1977, received 6 June 1977. The notebook projector formula needs the trace-factor correction recorded in its occurrence.

\subsection{GP-1976-02: Arnold’s Möbius characterization of the Gelfand--Gindikin domains}\label{prob:GP-1976-02}
After Cayley conjugation, identify the two invariant open complex domains with strict Möbius contractions of the closed unit disc or its closed exterior into the corresponding open region in \(\mathbb{CP}^1\).
\par\noindent\textit{Occurrences:} \hyperref[item:1976-398.19]{1976-398.19}.
\par\noindent\textit{Status (2026-11):} As of 10 November 2026, resolved by \citet[Proposition~9]{gelfandgindikin1977complex}. Published 1977 after the 6 September 1976 conjecture; the strictness and real-form conventions are essential.

\subsection{GP-1976-03: Local point linearization of scalar second-order ODEs}\label{prob:GP-1976-03}
For a regular smooth scalar equation \(y''=F(x,y,y')\), determine whether a local nonsingular point transformation takes the equation to \(Y''=0\). The condition concerns the equation on an open set, not the quadratic coefficient along only one solution.
\par\noindent\textit{Occurrences:} \hyperref[item:1976-403.1]{1976-403.1}, \hyperref[item:1976-403.10]{1976-403.10}.
\par\noindent\textit{Status (2026-11):} As of 10 November 2026, classically resolved by simultaneous vanishing of the Lie--Tresse invariants \(I,H\). Arnold explicitly presented nineteenth-century work; Cartan's 1924 connection theory and the 2008/2009 modern reconstruction verify the criterion. The notebook’s weaker one-solution implication is false.

\subsection{GP-1976-04: Extension of four-dimensional supertranslations to the Poincaré superalgebra}\label{prob:GP-1976-04}
Adjoin Lorentz generators acting in the vector representation on translations and the Majorana spin representation on odd generators, with the odd bracket equivariant into translations. Construct the resulting \(N=1\) Lie superalgebra without central charges.
\par\noindent\textit{Occurrences:} \hyperref[item:1976-404.5]{1976-404.5}.
\par\noindent\textit{Status (2026-11):} As of 10 November 2026, resolved before the 1976 question: the Gol'fand--Likhtman algebra circulated and was published in 1971. Likhtman's reproduced preprint carries a 12 April 1971 printing date. The semidirect-product construction verifies the mixed Jacobi identity through equivariance.

\subsection{GP-1976-10: Stokes formula for closed superdomains with even boundary}\label{prob:GP-1976-10}
Construct the appropriate integral forms and boundary restriction for an oriented smooth \(p|q\) supermanifold with a specified embedded even boundary, and establish \(\int_{\mathcal D}d\omega=\int_{\partial\mathcal D}\iota^*\omega\) for compactly supported \((p-1)\)-degree integral forms; the notebook tests the boundary \(u_1^2+u_2^2+c\xi_1\xi_2-1=0\).
\par\noindent\textit{Occurrences:} \hyperref[item:1976-406.16]{1976-406.16}.
\par\noindent\textit{Status (2026-11):} Solved by \citet[\S12]{bernsteinleites1977integral}, submitted 6 October 1976 and published in 1977. The submission predates the dated seminar; public circulation before that date has not been established. As of November 2026 this exact smooth-boundary statement is settled.

\subsection{GP-1975-26: Reconstructing higher-rank commuting ordinary differential operators}\label{prob:GP-1975-26}
Recover a commuting algebra of scalar ordinary differential operators of generic spectral rank \(r>1\) from correctly normalized pointed-curve, spectral-bundle and local functional data, at least for smooth spectral curves and data in general position. A curve and an unframed bundle alone are not sufficient input.
\par\noindent\textit{Occurrences:} 1975-391.3, 1975-391.4, \hyperref[item:1976-410.7]{1976-410.7}, \hyperref[item:1976-412.3]{1976-412.3}.
\par\noindent\textit{Status (2026-11):} As of November 2026, the generic smooth-curve inverse problem is solved by the abstract Drinfeld correspondence and Krichever's effective reconstruction. \citet[Theorem~2.3]{krichever1978commutative} corrects the earlier 1977 implementation; its received date is 15 February 1977. No claim of a complete classification of all singular or globally regular real cases is attached to this occurrence.

\subsection{GP-1976-20: Boundary-compatible hypersimplicial chains for the first Pontryagin class}\label{prob:GP-1976-20}
Extend the double-zero and one-section-change loci to rational configuration-space chains assigned to hypersimplices, compatible with faces and satisfying \(\partial C(H)=C(\partial H)\), in the first-Pontryagin construction.
\par\noindent\textit{Occurrences:} \hyperref[item:1976-413.11]{1976-413.11}.
\par\noindent\textit{Status (2026-11):} As of November 2026: solved before the 1977 exposition by the 1975 GGL construction. The notebook does not specify a separate stronger variant. Compare \citet[\S18 and \S\S27--28]{gabrielovgelfandlosik1975calculus}.

\subsection{GP-1976-21: Range and inversion of the line transform in three-space}\label{prob:GP-1976-21}
Characterize globally smooth rapidly decreasing line data that are integrals of a Schwartz function on \(\mathbb C^3\) or \(\mathbb R^3\), and reconstruct that function.
\par\noindent\textit{Occurrences:} \hyperref[item:1976-414.2]{1976-414.2}, \hyperref[item:1976-414.6]{1976-414.6}.
\par\noindent\textit{Status (2026-11):} As of November 2026: solved before the talk by \citet[\S\S7--8, Theorems 1--2]{gelfand1967planes}. The notebook omits the global data class and real/complex cycle distinction.

\subsection{GP-1976-22: Admissible algebraic complexes of lines}\label{prob:GP-1976-22}
For a closed irreducible algebraic three-dimensional line complex covering a dense open subset of \(\mathbb C^3\), determine when the differential inversion form descends from the scalar line data restricted to the complex, and classify such complexes near generic regular lines.
\par\noindent\textit{Occurrences:} \hyperref[item:1976-414.9]{1976-414.9}, \hyperref[item:1976-414.10]{1976-414.10}.
\par\noindent\textit{Status (2026-11):} As of November 2026: the specified admissibility and generic classification were established in \citet[\S\S2--5]{gelfandgraev1968complexes}. The notebook's wording about already differentiated data is a partial match.

\subsection{GP-1976-23: Poncelet closure for two conics}\label{prob:GP-1976-23}
For two nonsingular conics with smooth incidence elliptic curve, determine whether the consistent tangent iteration has a closed polygon, and characterize its period.
\par\noindent\textit{Occurrences:} \hyperref[item:1976-416.1]{1976-416.1}, \hyperref[item:1976-416.2]{1976-416.2}.
\par\noindent\textit{Status (2026-11):} As of November 2026: classical; the translation point on the incidence elliptic curve is torsion precisely in the closing case. The exact source cited in the notes is \citet[\S I(d)]{griffiths1976abel}.

\subsection{GP-1976-24: Counting Young-graph paths between prescribed endpoints}\label{prob:GP-1976-24}
Given \(\mu\subseteq\lambda\), count saturated Young-graph paths from \(\mu\) to \(\lambda\), equivalently compute the multiplicity of \(V^\mu\) in \(V^\lambda|_{S_{|\mu|}}\).
\par\noindent\textit{Occurrences:} \hyperref[item:1976-415.6]{1976-415.6}.
\par\noindent\textit{Status (2026-11):} As of November 2026: the count is the skew-tableau number \(f^{\lambda/\mu}\), given exactly by the classical Jacobi--Trudi/Aitken determinant. The notebook states no further asymptotic or positivity demand. The checked later derivation is \citet[equation~(1.2)]{stanley2011skew}.

\subsection{GP-1976-25: Ulam asymptotics and the Plancherel first row}\label{prob:GP-1976-25}
For a uniform permutation, or the ranks of an independent continuous sample, determine the leading asymptotic longest increasing-subsequence length.
\par\noindent\textit{Occurrences:} \hyperref[item:1976-415.14]{1976-415.14}, \hyperref[item:1976-415.18]{1976-415.18}.
\par\noindent\textit{Status (2026-11):} As of November 2026: \(L_n/(2\sqrt n)\to1\) in probability. This is the 1977 first-order result of \citet[Corollaries~1--2]{vershikkerov1977asymptotics}; the separate fluctuation discussion is not folded into this solved question.

\subsection{GP-1976-26: Description and Torelli information of rank-one representation packets}\label{prob:GP-1976-26}
For a fixed marked compact Riemann surface \(X\) and degree-zero line bundle \(L\), describe all characters with underlying holomorphic bundle \(L\), and determine the information about \(X\) retained by that embedded packet.
\par\noindent\textit{Occurrences:} \hyperref[item:1976-417.9]{1976-417.9}, \hyperref[item:1976-417.13]{1976-417.13}, \hyperref[item:1976-417.15]{1976-417.15}.
\par\noindent\textit{Status (2026-11):} As of November 2026: the packet is a coset of the exponential of holomorphic periods; its embedded tangent space, integral marking, and intersection pairing recover the principally polarized Jacobian and hence the curve. These are the classical constructions and exercises in \citet[Ch.~II, \S\S2--3]{tyurin1978periods}.

\subsection{GP-1977-01: Algebraic description of vector bundles on projective space}\label{prob:GP-1977-01}
For each \(n\geq0\), describe the finite-rank algebraic vector bundles on \(\mathbb P^n_{\mathbb C}\) up to isomorphism. The precise requested form in the notes is a classification by proper graded exterior-algebra modules up to adding graded free summands. This problem belongs to the manuscript lecture of 13 February 1978; its corrected GP-1977 identifier records the actual academic year, while this package retains the requested catalogue record.
\par\noindent\textit{Occurrences:} \hyperref[item:1977-437.2]{1977-437.2}, \hyperref[item:1977-437.11]{1977-437.11}, \hyperref[item:1977-437.12]{1977-437.12}.
\par\noindent\textit{Status (2026-11):} Solved in the stated representation-theoretic sense by \citet[Theorem~1]{bernsteingelfandgelfand1978bundles}, received 10 January 1978 and published in the 1978 volume. Theorem~2 gives the derived categorical equivalence. No earlier problem ID in the supplied index states this exact projective-bundle problem; affine freeness (the Quillen--Suslin theorem) has different hypotheses.

\subsection{GP-1977-02: Represent twisting by O(1) on the exterior-module side}\label{prob:GP-1977-02}
Given a proper graded exterior-algebra module \(V\), determine an operation whose associated bundle is \(\Phi(V)(1)\). The notes ask and answer this question within the same, miscatalogued 1978 lecture.
\par\noindent\textit{Occurrences:} \hyperref[item:1977-437.16]{1977-437.16}.
\par\noindent\textit{Status (2026-11):} Solved: embed \(V\) into a finite graded free module, take the cokernel \(V^{\prime}\), and use \(W_j=V^{\prime}_{j+1}\). See \citet[remark~2 after Theorem~1]{bernsteingelfandgelfand1978bundles}. The fibre long exact sequence independently checks the shift.

\section{Corrections to the transcription and translation}\label{corr:1976}
The evidence distinguishes a transcription correction from a mathematical emendation of wording faithfully preserved in the notebook.
\subsection*{P000659-P000662 (both; provenance): Record 398 catalogue page range}
\noindent\textit{Transcription reads.} Catalogue: P000659-P000660\par
\noindent\textit{Critical reading.} Substantive lecture continues through P000662\par
\noindent\textit{Evidence.} P661 continues the same notation and main theorem; P662 ends with Arnold’s conjecture; the next dated session begins P663.\par
\subsection*{P000659 (both; faithful error): All linear symmetries of the quadric}
\noindent\textit{Transcription reads.} Any linear transformation preserving the hyperboloid is g 1\textasciicircum{}\{-1\}gg 2\par
\noindent\textit{Critical reading.} Restrict to the identity component SO\_0(2,2)\par
\noindent\textit{Evidence.} The scan really says any. Transposition preserves determinant and reverses orientation of M 2(R), so is outside the left-right image.\par
\subsection*{P000660-P000661 (both; ambiguity): Degenerate hyperplane sections}
\noindent\textit{Transcription reads.} Delta=0: paraboloid/cylinder,2 parameters\par
\noindent\textit{Critical reading.} For Delta=0, zeta nonzero: paraboloid,3 parameters; zeta=0: cylinder,2 parameters\par
\noindent\textit{Evidence.} The scan conflates labels; P661 independently assigns E three parameters. Setting alpha=zeta gives the two distinct normal forms.\par
\subsection*{P000662 (both; faithful error): Hardy theorem}
\noindent\textit{Transcription reads.} Analytic continuation into D\_+ and D\_-\par
\noindent\textit{Critical reading.} Hardy holomorphic continuation with uniform L2 bounds and specified boundary values\par
\noindent\textit{Evidence.} The published 1977 theorem supplies the missing functional hypotheses; this is an emendation of shorthand, not a transcription change.\par
\subsection*{P000662 (both; faithful error): Definition of distinguished boundary}
\noindent\textit{Transcription reads.} Smallest boundary set giving uniqueness of holomorphic continuation\par
\noindent\textit{Critical reading.} Smallest closed norming boundary for a specified function algebra; use an appropriate noncompact formulation here\par
\noindent\textit{Evidence.} In the disc algebra, two disjoint nontrivial closed boundary arcs are uniqueness sets, but every closed norming set contains the whole circle.\par
\subsection*{P000662 (both; faithful error): Displayed projection kernel}
\noindent\textit{Transcription reads.} c(2-tr(g 1 g 2\textasciicircum{}\{-1\}))\textasciicircum{}\{-3/2\}\par
\noindent\textit{Critical reading.} c[(tr(g 1 g 2\textasciicircum{}\{-1\})-2)\textasciicircum{}3(tr(g 1 g 2\textasciicircum{}\{-1\})+2)]\textasciicircum{}\{-1/2\}, with consistent holomorphic branch\par
\noindent\textit{Evidence.} The scan agrees with both transcriptions. Gelfand-Gindikin 1977 equation 22 and an independent character-series sum supply the missing nonconstant factor.\par
\subsection*{P000662 (both; faithful error): Arnold’s conjecture and real circle}
\noindent\textit{Transcription reads.} SL(2, R) preserves the circle; maps disc into itself\par
\noindent\textit{Critical reading.} Use the Cayley-conjugate real form and strict contraction to describe the open domain\par
\noindent\textit{Evidence.} Standard real Möbius matrices preserve the real projective line. Published 1977 Proposition 9 identifies the corresponding strict semigroups.\par
\subsection*{P000663 (both; provenance): Record 400 boundary}
\noindent\textit{Transcription reads.} Catalogue assigns P000662-P000664 to Kirillov\par
\noindent\textit{Critical reading.} Only Kirillov’s announced title precedes the explicitly named Fuchs lecture\par
\noindent\textit{Evidence.} The scan begins 13/IX-76, lists Kirillov’s heading, and then marks D.B.Fuchs. No separate Kirillov exposition intervenes.\par
\subsection*{P000664 (both; faithful error): Fuchs acyclicity theorem}
\noindent\textit{Transcription reads.} All homologies are zero\par
\noindent\textit{Critical reading.} H\textasciicircum{}0(W)=ground field; H\textasciicircum{}j(W)=0 for j>0, or use the augmented complex\par
\noindent\textit{Evidence.} The notebook faithfully omits the augmentation/positive-degree restriction. Cartan 1950 II, section 2 Theorem 1 and the explicit contraction retain the constant 1.\par
\subsection*{P000664 (both; faithful error): Kirillov counting rule}
\noindent\textit{Transcription reads.} As many invariant alternating k-forms as symmetric forms of degree(k+1)/2\par
\noindent\textit{Critical reading.} Primitive odd cohomology generators correspond to indecomposable invariant symmetric polynomials\par
\noindent\textit{Evidence.} SU(2), k=7 is a counterexample to the unrestricted rule. Cartan 1950 II, section 4 Theorem 2 gives the corrected correspondence.\par
\subsection*{P000664 (both; faithful error): Kirillov transgression direction}
\noindent\textit{Transcription reads.} The odd form is the differential of a symmetric expression\par
\noindent\textit{Critical reading.} The invariant polynomial of Weil degree 2 m is d of a transgression element of degree 2 m-1\par
\noindent\textit{Evidence.} The degree-one differential cannot lower degree; its primitive exterior component has degree 2 m-1.\par
\subsection*{P000664 (both; ambiguity): Doubtful phrase after the straightening question}
\noindent\textit{Transcription reads.} Linear menu\par
\noindent\textit{Critical reading.} Unresolved reading; interpret only the subsequent displayed linear equation\par
\noindent\textit{Evidence.} The scan does not justify a confident lexical correction. The next explicit formula is a linear scalar second-order equation.\par
\subsection*{P000665 (both; faithful error): Two line fields equivalence}
\noindent\textit{Transcription reads.} Arbitrary two direction fields in R 3\par
\noindent\textit{Critical reading.} Transverse ordered line fields whose span is contact\par
\noindent\textit{Evidence.} The scan later mentions the contact structure. Without nonintegrability the claimed ODE equivalence does not hold.\par
\subsection*{P000666 (both; faithful error): Vanishing density implies straightening}
\noindent\textit{Transcription reads.} If the density vanishes, the equation is formally straightenable\par
\noindent\textit{Critical reading.} One-solution vanishing is insufficient; require the complete formal curvature criterion, or suitable stated hypotheses\par
\noindent\textit{Evidence.} The analytic equation y-double-prime=y\textasciicircum{}3 has zero quadratic density on y=0 but nonzero Tresse H=36 y. The intended extra quantifiers are not supplied.\par
\subsection*{P000666 (both; faithful error): Vector field from the density}
\noindent\textit{Transcription reads.} v=a\textasciicircum{}\{-2/5\}partial\_x\par
\noindent\textit{Critical reading.} Restrict to a nonvanishing density coefficient and retain its constant-scale ambiguity\par
\noindent\textit{Evidence.} At a=0 the expression is singular; In rescales as c\textasciicircum{}n when v rescales as c.\par
\subsection*{P000664-P000667 (both; provenance): Records 403/404 boundary}
\noindent\textit{Transcription reads.} Catalogue overlaps Arnold through P000668 and Novikov from P000666\par
\noindent\textit{Critical reading.} Arnold ends P000666; Novikov starts P000667\par
\noindent\textit{Evidence.} The scanP 667 has 20/IX-76 and an explicit S.P.Novikov heading.\par
\subsection*{P000667 (both; faithful error): Definition of da}
\noindent\textit{Transcription reads.} da=dim a\par
\noindent\textit{Critical reading.} da=parity(a), for homogeneous a\par
\noindent\textit{Evidence.} The scan says dim a; P668 explicitly changes to degree. Supercommutator signs depend on parity, not dimension.\par
\subsection*{P000668 (both; faithful error): Historical attribution}
\noindent\textit{Transcription reads.} Gelfand and Likhtman\par
\noindent\textit{Critical reading.} Yu. A. Gol’fand and E. P. Likhtman\par
\noindent\textit{Evidence.} Likhtman’s author history and the 1971 reprinted preprint identify Gol’fand. The diplomatic spelling is preserved as a notebook attribution error, not silently attributed to I. M. Gelfand.\par
\subsection*{P000668 (both; faithful error): Supermultiplet field table}
\noindent\textit{Transcription reads.} psi: scalar field\par
\noindent\textit{Critical reading.} psi: Majorana spinor field\par
\noindent\textit{Evidence.} The scan clearly repeats scalar. The gamma-matrix kinetic term and displayed supersymmetry variations require a spinor.\par
\subsection*{P000669 (both; faithful error): Fermion kinetic term}
\noindent\textit{Transcription reads.} barPsi gamma\textasciicircum{}beta partial\_mu Psi\par
\noindent\textit{Critical reading.} Contract the same index: barPsi gamma\textasciicircum{}mu partial\_mu Psi\par
\noindent\textit{Evidence.} Both distinct letters are visible in the scan. The written term has two free Lorentz indices and is not a scalar.\par
\subsection*{P000668-P000669 (both; faithful error): Cubic scalar term after auxiliary elimination}
\noindent\textit{Transcription reads.} +mg A(A\textasciicircum{}2+B\textasciicircum{}2)\par
\noindent\textit{Critical reading.} -mg A(A\textasciicircum{}2+B\textasciicircum{}2), using the P669 auxiliary terms\par
\noindent\textit{Evidence.} Substitution F=-mA-g(A\textasciicircum{}2-B\textasciicircum{}2), G=-mB-2 gAB and independent completion of squares agree.\par
\subsection*{P000668 (both; ambiguity): Quartic scalar coefficient}
\noindent\textit{Transcription reads.} g\_2\textasciicircum{}2\par
\noindent\textit{Critical reading.} g\textasciicircum{}2/2 under the P669 conventions; the paleographic fraction remains uncertain\par
\noindent\textit{Evidence.} The scan has two 2 s but no clear bar. Exact auxiliary elimination forces the stated coefficient.\par
\subsection*{P000668-P000669 (both; faithful error): Fermion mass mismatch}
\noindent\textit{Transcription reads.} P668:-im/2 barpsi psi; P669:-im barPsi Psi\par
\noindent\textit{Critical reading.} Reconcile the missing factor 1/2 in P669 if the same m and fermion normalization are intended\par
\noindent\textit{Evidence.} Auxiliary elimination affects only scalar terms and cannot change the displayed fermion mass coefficient. Both transcriptions preserve the mismatch.\par
\subsection*{P000669 (both; faithful error): Invariance of the Lagrangian}
\noindent\textit{Transcription reads.} The Lagrangian is invariant\par
\noindent\textit{Critical reading.} The action is invariant; the Lagrangian density may change by a total derivative\par
\noindent\textit{Evidence.} This is the standard component supersymmetry statement, verified in the 1975 review.\par
\subsection*{P000669 (both; faithful error): Confining forces}
\noindent\textit{Transcription reads.} Forces must increase with distance\par
\noindent\textit{Critical reading.} A growing separation energy suffices; linear potential gives asymptotically constant force\par
\noindent\textit{Evidence.} The scan is explicit. Differentiating V(R)=sigma R disproves the stated necessity.\par
\subsection*{P000669 (both; provenance): Records 404/405/406 boundary}
\noindent\textit{Transcription reads.} Catalogue windows overlap all three records\par
\noindent\textit{Critical reading.} Novikov ends before the Migdal heading on P669; Bernstein begins P670\par
\noindent\textit{Evidence.} The scanP 669 names Migdal after the five variations; P670 has 18/X-76 andBernstein.\par
\subsection*{P000670 (both; faithful error): last exponent in superfunction expansion}
\noindent\textit{Transcription reads.} \(\xi_q^{\alpha_k}\)\par
\noindent\textit{Critical reading.} \(\xi_q^{\alpha_q}\)\par
\noindent\textit{Evidence.} Scan and Russian transcription agree on k; the indexed expansion requires q.\par
\subsection*{P000670 (en; translation): caption for excluded even term in image of odd generator}
\noindent\textit{Transcription reads.} disallowed, since it preserves parity\par
\noindent\textit{Critical reading.} disallowed because the morphism must preserve parity\par
\noindent\textit{Evidence.} Russian wording explains the parity requirement; the quadratic term itself is even.\par
\subsection*{P000671 (both; faithful error): second superfunction ring}
\noindent\textit{Transcription reads.} \(g(v,\xi_1,\xi_2)\)\par
\noindent\textit{Critical reading.} \(g(v,\eta_1,\eta_2)\)\par
\noindent\textit{Evidence.} The target coordinates were named eta; the scan itself reuses xi. This is a harmless faithful notation slip.\par
\subsection*{P000672 (both; faithful error): circle-normalizing substitution}
\noindent\textit{Transcription reads.} \((u_1-\tfrac12\xi_1\xi_2)^2+u_2^2-1\)\par
\noindent\textit{Critical reading.} \((u_1-\xi_1\xi_2/(2u_1))^2+u_2^2-1\quad(u_1\ne0)\)\par
\noindent\textit{Evidence.} Direct expansion; scan confirms no u1 denominator. The corrected formula is local, with the u2 chart covering the remaining boundary.\par
\subsection*{P000672 (both; faithful error): half-line invariance calculation}
\noindent\textit{Transcription reads.} Expansion of \(g_{11}\) alone\par
\noindent\textit{Critical reading.} Top coefficient \(g_{12}+(u g_0)^{\prime}\), including \(\operatorname{Ber}D\Phi=1+\xi_1\xi_2\)\par
\noindent\textit{Evidence.} Full pullback calculation, not a diplomatic emendation; the notebook check is incomplete.\par
\subsection*{P000673 (both; faithful error): first block determinant definition}
\noindent\textit{Transcription reads.} \(\det(A_{11})^{-1}\det(A_{22}-A_{21}A_{11}^{-1}A_{12})\)\par
\noindent\textit{Critical reading.} \(\det(A_{11})/\det(A_{22}-A_{21}A_{11}^{-1}A_{12})\)\par
\noindent\textit{Evidence.} Scan confirms reciprocal; lower Berezinian and exp(str log) on the same page require the corrected normalization. Diagonal 1|1 check.\par
\subsection*{P000673-P000674 (both; provenance): record provenance}
\noindent\textit{Transcription reads.} Continuation of the 18 October talk without qualification\par
\noindent\textit{Critical reading.} Retrospective summary of earlier Bernstein talks from I. Dorfman's notebook\par
\noindent\textit{Evidence.} Explicit handwritten heading on P000673; retained within record406 without manufacturing a new session.\par
\subsection*{P000675 (both; faithful error): claim about quarks}
\noindent\textit{Transcription reads.} The Yang--Mills equations describe quarks\par
\noindent\textit{Critical reading.} The displayed source-free equations describe gauge fields; quark matter requires a coupled covariant Dirac equation and a current\par
\noindent\textit{Evidence.} Faithful wording in both versions; the displayed unknown is the connection, not a spinor matter field.\par
\subsection*{P000675 (both; faithful error): Maxwell curvature indices}
\noindent\textit{Transcription reads.} \(F^{\mu\nu}=\partial_\mu A_\nu-\partial_\nu A_\mu\)\par
\noindent\textit{Critical reading.} \(F_{\mu\nu}=\partial_\mu A_\nu-\partial_\nu A_\mu;\quad F^{\mu\nu}=g^{\mu\alpha}g^{\nu\beta}F_{\alpha\beta}\)\par
\noindent\textit{Evidence.} Scan confirms the mixed index placement; Minkowski raising is not the identity.\par
\subsection*{P000676 (both; faithful error): conjugated potential in gauge transformation}
\noindent\textit{Transcription reads.} \(B_\mu\mapsto S^{-1}B^\mu S+S^{-1}\partial_\mu S\)\par
\noindent\textit{Critical reading.} \(B_\mu\mapsto S^{-1}B_\mu S+S^{-1}\partial_\mu S\)\par
\noindent\textit{Evidence.} The superscript is present in the scan; all terms need the same covariant index.\par
\subsection*{P000677 (both; ambiguity): curvature two-form normalization}
\noindent\textit{Transcription reads.} \(T_aF^a_{\mu\nu}dx^\mu\wedge dx^\nu\)\par
\noindent\textit{Critical reading.} \(\tfrac12T_aF^a_{\mu\nu}dx^\mu\wedge dx^\nu\quad\hbox{or sum only over }\mu<\nu\)\par
\noindent\textit{Evidence.} The notes omit a summation convention; if Einstein summation is intended, the factor1/2 is necessary.\par
\subsection*{P000675-P000677 (both; provenance): speaker initials in catalogue}
\noindent\textit{Transcription reads.} V. N. Popov\par
\noindent\textit{Critical reading.} Retain catalogue initials with explicit conflict: the exact papers are by D. A. Popov (with L. I. Daikhin)\par
\noindent\textit{Evidence.} Handwritten heading says only Popov; exact 1975 primary papers and author metadata do not justify silently replacing the catalogued speaker.\par
\subsection*{P000678 (both; faithful error): partition terminology}
\noindent\textit{Transcription reads.} partitions into nonnegative summands\par
\noindent\textit{Critical reading.} partitions into positive parts, or identify all zero padding\par
\noindent\textit{Evidence.} Unlimited zeros otherwise give infinitely many decompositions; handwriting faithfully uses nonnegative.\par
\subsection*{P000678 (both; faithful error): Dedekind function name}
\noindent\textit{Transcription reads.} Dedekind \(\Theta\)-function\par
\noindent\textit{Critical reading.} Dedekind \(\eta\)-function\par
\noindent\textit{Evidence.} The displayed product \(q^{1/24}\) times the Euler product defines eta; the preceding Jacobi theta reference is separate.\par
\subsection*{P000679 (both; faithful error): raw zeta reflection}
\noindent\textit{Transcription reads.} \(\zeta(s)=\zeta(\dim\mathfrak g/2-s)\)\par
\noindent\textit{Critical reading.} \(\Lambda(s)=\Lambda(\dim\mathfrak g/2-s),\quad\Lambda(s)=(2\pi)^{-s}\Gamma(s)\zeta(s)\)\par
\noindent\textit{Evidence.} Actual scan contains no completion; the Mellin transform of \(\eta(it)^d\) gives the corrected relation. The raw values at \(0\) and \(d/2\) differ.\par
\subsection*{P000679 (both; faithful error): Cartan off-diagonal inequality}
\noindent\textit{Transcription reads.} \(a_{ij}<0\quad(i\ne j)\)\par
\noindent\textit{Critical reading.} \(a_{ij}\le0\quad(i\ne j)\)\par
\noindent\textit{Evidence.} The scan has a strict sign; adjacent symmetric-zero remark and standard Cartan examples require allowing zeros.\par
\subsection*{P000680 (both; faithful error): number of generators}
\noindent\textit{Transcription reads.} \(n\) generators \(h_i,e_i,f_i\)\par
\noindent\textit{Critical reading.} \(3n\) generators, in three families indexed by \(i=1,\ldots,n\)\par
\noindent\textit{Evidence.} The page lists three families; n is the size of the Cartan matrix.\par
\subsection*{P000680 (both; faithful error): general Chevalley brackets}
\noindent\textit{Transcription reads.} \([h_i,e_j]=-a_{ij}e_j,\ [h_i,f_j]=a_{ij}f_j,\ [e_i,f_i]=\delta_{ij}h_i\)\par
\noindent\textit{Critical reading.} \([h_i,e_j]=a_{ij}e_j,\ [h_i,f_j]=-a_{ij}f_j,\ [e_i,f_j]=\delta_{ij}h_i\)\par
\noindent\textit{Evidence.} The scan confirms the signs and repeated i; its later sl2 example fixes the intended convention. An internally consistent opposite convention would also require changing [e,f].\par
\subsection*{P000680 (both; ambiguity): Serre relations and quotient}
\noindent\textit{Transcription reads.} Only positive-generator relation and unrestricted maximal normal subalgebra\par
\noindent\textit{Critical reading.} Include negative-generator Serre relations for i≠j; in the contragredient construction require the ideal to meet the Cartan subalgebra trivially\par
\noindent\textit{Evidence.} These conditions are absent or abbreviated in the notebook, not mistranscribed.\par
\subsection*{P000680 (both; faithful error): loop-algebra notation}
\noindent\textit{Transcription reads.} \(SU(3)\otimes\mathbb C[x,x^{-1}]\)\par
\noindent\textit{Critical reading.} \(\mathfrak{sl}_3(\mathbb C)\otimes\mathbb C[x,x^{-1}]\)\par
\noindent\textit{Evidence.} The source is constructing complex Lie algebras; a compact Lie group is not a vector-space tensor factor.\par
\subsection*{P000678-P000680 (both; provenance): session date}
\noindent\textit{Transcription reads.} 22 November 1976 throughout the packet\par
\noindent\textit{Critical reading.} Date supplied by catalogue; the handwritten start is undated and the closing greeting says Happy New Year\par
\noindent\textit{Evidence.} Preserve date with provenance rather than infer additional sessions.\par
\subsection*{P000681 (both; faithful error): fourth-order Lax system}
\noindent\textit{Transcription reads.} \(v_t=v_{xx}-Av_x-2w_x;\quad w_t=2v_{xxx}+2vv_x-w_{xx}-Aw_x-2z_x\)\par
\noindent\textit{Critical reading.} \(v_t=2v_{xx}-Av_x-2w_x;\quad w_t=2v_{xxx}+vv_x-w_{xx}-Aw_x-2z_x\)\par
\noindent\textit{Evidence.} Both coefficients are confirmed by the handwritten scan. Independent Ore-product and test-function calculations agree; Manin ChapterIII§6 equation17 has the corrected relative coefficients.\par
\subsection*{P000681 (both; ambiguity): pulse normalization and omitted component}
\noindent\textit{Transcription reads.} \(\sigma w=4\mu_2\mu_{2x}+6\mu_{2xx}+4\mu_{1x};\quad z=\ldots\)\par
\noindent\textit{Critical reading.} Retain the recorded formula; do not certify it as a complete solution without z and a parameter/sign convention\par
\noindent\textit{Evidence.} Manin printedp80 gives a closely related but differently signed/scaled family. No unique emendation is justified from this excerpt.\par
\subsection*{P000681 (both; faithful error): positive recurrence period}
\noindent\textit{Transcription reads.} \(2\pi/\sigma\)\par
\noindent\textit{Critical reading.} \(2\pi/|\sigma|\quad(\sigma\ne0)\)\par
\noindent\textit{Evidence.} The page allows any real sigma; positive temporal period requires absolute value and excludes zero.\par
\subsection*{P000682 (both; ambiguity): two-quadrics ambient and missing data}
\noindent\textit{Transcription reads.} \(\text{hyperelliptic curve}\longrightarrow\text{two quadrics in }\mathbb C^5\)\par
\noindent\textit{Critical reading.} Retain as ambiguous; a plausible candidate is the genus-two rank-two fixed odd-determinant moduli space in projective P5\par
\noindent\textit{Evidence.} Actual scan says C5. The classical moduli theorem has further hypotheses absent from the note; affine-chart interpretation is possible but unrecorded.\par
\subsection*{P000683 (both; faithful error): graded connection isomorphism lower bound}
\noindent\textit{Transcription reads.} \(i\ge-1\)\par
\noindent\textit{Critical reading.} \(i\ge0\text{ for the induced graded isomorphism}\)\par
\noindent\textit{Evidence.} At \(i=-1\) the source graded piece is zero or undefined while the target is \(M_0\), a free rank-one module. The mere inclusion \(\nabla M_i\subset M_{i+1}\) may include \(i=-1\).\par
\subsection*{P000683 (both; faithful error): Theorem2 module order}
\noindent\textit{Transcription reads.} \((\mathcal O,\mathcal B)\text{-bimodule}\)\par
\noindent\textit{Critical reading.} \((\mathcal B,\mathcal O)\text{-bimodule}\)\par
\noindent\textit{Evidence.} The preceding definitions and subsequent formulas use a left B-action and right O-action; the reversed order is visible in the scan.\par
\subsection*{P000683 (both; faithful error): Theorem2 Leibniz rule}
\noindent\textit{Transcription reads.} \(\nabla_1(bm)=\partial_1b+\nabla_1m\)\par
\noindent\textit{Critical reading.} \(\nabla_1(bm)=(\partial_1b)m+b\nabla_1m\)\par
\noindent\textit{Evidence.} The scan visibly omits both factors; the printed Manin theorem assumes a connection, hence the full Leibniz rule.\par
\subsection*{P000683 (both; ambiguity): Theorem2 commutator sign}
\noindent\textit{Transcription reads.} \(\partial_1L_\varphi=\pm[L_\varphi,P]\)\par
\noindent\textit{Critical reading.} \(\partial_1L_\varphi=[P,L_\varphi]\text{ when }\nabla_1\mathbf1=P\mathbf1\)\par
\noindent\textit{Evidence.} Direct application of the connection to \(L_\varphi\mathbf1=\mathbf1\varphi\) fixes the sign; Manin Theorem 2.5, printed p. 83, agrees.\par
\subsection*{P000682; P000684 (both; faithful error): zero-curvature derivative and bracket signs}
\noindent\textit{Transcription reads.} \(\partial_1L+\partial_2P=\pm[L,P]\)\par
\noindent\textit{Critical reading.} \(\partial_1L-\partial_2P=[P,L]\text{ for }\nabla_1\mathbf1=P\mathbf1,\ \nabla_2\mathbf1=L\mathbf1\)\par
\noindent\textit{Evidence.} The plus sign is faithful to the scan. Flatness expansion and the scalar counterexample P=z,L=t disprove it. Manin Theorem2.6 printedp83 repeats a plus-sign error under swapped labels.\par
\subsection*{P000686 (both; ambiguity): graph notation}
\noindent\textit{Transcription reads.} \(\varphi(t)=\operatorname{graph}\varphi(t)\)\par
\noindent\textit{Critical reading.} \(\varphi(t)=\operatorname{graph}\Phi(t)\)\par
\noindent\textit{Evidence.} The scan defines capital Phi as the flow derivative; map and graph need distinct symbols.\par
\subsection*{P000686 (both; faithful error): symmetric-form equivalence}
\noindent\textit{Transcription reads.} \(\sigma=0\text{ on }\Gamma\)\par
\noindent\textit{Critical reading.} \(\sigma=0\text{ on }\operatorname{graph}\Gamma\)\par
\noindent\textit{Evidence.} A two-form restricts to the graph, not to the linear map.\par
\subsection*{P000687 (both; faithful error): Maslov cycle}
\noindent\textit{Transcription reads.} Sigma is embedded and two-sided\par
\noindent\textit{Critical reading.} The one-dimensional-intersection stratum is a smooth cooriented hypersurface.\par
\noindent\textit{Evidence.} Higher-intersection strata are singular; faithful mathematical shorthand.\par
\subsection*{P000687 (both; faithful error): rho}
\noindent\textit{Transcription reads.} tangent to R-perp; fibre of the normal bundle\par
\noindent\textit{Critical reading.} Tangent to the total conormal bundle at the endpoint covector.\par
\noindent\textit{Evidence.} The scan writes T(R-perp); the fibre generally has another dimension.\par
\subsection*{P000688-P000689 (both; faithful error): individual trace terms}
\noindent\textit{Transcription reads.} \(\int_X\varphi(x\gamma x^{-1})\,dx\)\par
\noindent\textit{Critical reading.} Regroup by conjugacy classes or specify a fundamental domain.\par
\noindent\textit{Evidence.} An individual term is not Gamma-invariant; DKV1979 section5 gives the grouped formula.\par
\subsection*{P000688 (both; faithful error): half-sum of roots}
\noindent\textit{Transcription reads.} \(\rho=\frac12\sum\text{positive roots}\)\par
\noindent\textit{Critical reading.} \(\rho=\frac12\sum_{\alpha>0}m_\alpha\alpha\)\par
\noindent\textit{Evidence.} Restricted-root multiplicities are needed in general.\par
\subsection*{P000689 (both; ambiguity): final remainder}
\noindent\textit{Transcription reads.} \(O(R^{n-r-1})\)\par
\noindent\textit{Critical reading.} Retain the uncertain symbol; the matched general theorem has \(O(R^{n-1})\).\par
\noindent\textit{Evidence.} PNG inspected; DKV1979 Theorem8.8 verifies the weaker estimate.\par
\subsection*{P000690 (both; faithful error): Zakharov--Shabat equation}
\noindent\textit{Transcription reads.} \(L_t+P_y=[P,L]\)\par
\noindent\textit{Critical reading.} \(L_t-P_y=[P,L]\)\par
\noindent\textit{Evidence.} The scan has plus. With the displayed commuting connections and commutator AB-BA, expansion gives the corrected sign.\par
\subsection*{P000691 (both; faithful error): node parameters in real case}
\noindent\textit{Transcription reads.} \(\alpha_i,\beta_i\in k\); \(\alpha_i=\overline{\beta_i}\) for \(k=\mathbb R\)\par
\noindent\textit{Critical reading.} Use complexified parameters for real descent and require all \(2N\) points distinct.\par
\noindent\textit{Evidence.} Real distinct conjugate pairs cannot both lie in \(\mathbb R\); Manin III.5.1 explicitly chooses complex node parameters.\par
\subsection*{P000692 (both; faithful error): filtration shift}
\noindent\textit{Transcription reads.} \(M_i=M_{(i-rN)}\)\par
\noindent\textit{Critical reading.} \(M_i=M_{(i+rN)}\quad(M_0=0,\ \mathbf1\in M_1)\)\par
\noindent\textit{Evidence.} Actual PNG has minus. The preceding vanishing lemma forces a positive shift; Manin III.5.5 uses \(+rN+1\) with generator in \(M_0\).\par
\subsection*{P000692 (both; faithful error): generator coefficient equation}
\noindent\textit{Transcription reads.} \(\sum_j(\alpha_i^jE_r-\beta_i^jg_i)m_i\)\par
\noindent\textit{Critical reading.} \(\sum_j(\alpha_i^jE_r-\beta_i^jg_i)m_j\)\par
\noindent\textit{Evidence.} Actual PNG and Russian have \(m_i\); substitution of the displayed polynomial forces \(m_j\).\par
\subsection*{P000692 (both; faithful error): endomorphism ring}
\noindent\textit{Transcription reads.} \(M_r(\mathcal B[\lambda])=\operatorname{End}_{\mathcal B}(N)\)\par
\noindent\textit{Critical reading.} \(M_r(\mathcal B[\lambda])=\operatorname{End}_{\mathcal B[\lambda]}(N)\)\par
\noindent\textit{Evidence.} Polynomial matrices also commute with multiplication by lambda; arbitrary B-linear maps need not.\par
\subsection*{P000692 (both; faithful error): flatness equation}
\noindent\textit{Transcription reads.} \(\partial_xd_t-\partial_td_x=[d_x,d_t]\)\par
\noindent\textit{Critical reading.} \(\partial_xd_t-\partial_td_x=-[d_x,d_t]\)\par
\noindent\textit{Evidence.} Scan has no minus; expansion of \([\partial_x+d_x,\partial_t+d_t]\) fixes the sign.\par
\subsection*{P000693 (both; faithful error): left time exponential}
\noindent\textit{Transcription reads.} \[\begin{aligned}g_i(x)&=e^{-xd_x(\alpha_i)+td_t(\alpha_i)}g_i(0,0)\\&\qquad{}\times e^{xd_x(\beta_i)+td_t(\beta_i)}.\end{aligned}\]\par
\noindent\textit{Critical reading.} \[\begin{aligned}g_i(x,t)&=e^{-xd_x(\alpha_i)-td_t(\alpha_i)}g_i(0,0)\\&\qquad{}\times e^{xd_x(\beta_i)+td_t(\beta_i)}.\end{aligned}\]\par
\noindent\textit{Evidence.} PNG shows both plus time signs; Manin III5.12 equation10 and direct differentiation require the left minus.\par
\subsection*{P000694 (both; faithful error): polyhedral manifold criterion}
\noindent\textit{Transcription reads.} each codimension-one simplex belongs to exactly two top simplices\par
\noindent\textit{Critical reading.} This is an incidence condition; manifold links or flattenability must also be required.\par
\noindent\textit{Evidence.} Russian and scan agree. MacPherson1977/78 pp105–106 supplies the missing geometric hypotheses.\par
\subsection*{P000694 (both; ambiguity): angle-defect motivation}
\noindent\textit{Transcription reads.} sum of external angles is not pi\par
\noindent\textit{Critical reading.} Preserve the ambiguous angle description; for a closed Euclidean triangulated surface the standard defect is 2pi minus the incident face-angle sum.\par
\noindent\textit{Evidence.} The scan does not identify the external angles or justify a numerical replacement.\par
\subsection*{P000695 (both; faithful error): determinantal coorientation}
\noindent\textit{Transcription reads.} \(Z\subset\operatorname{Hom}(\mathbb R^m,\mathbb R^n)\) has canonical coorientation\par
\noindent\textit{Critical reading.} Restrict this assertion to the square corank-two stratum used in the argument.\par
\noindent\textit{Evidence.} Arbitrary real determinantal strata need not be coorientable; here the normal space is \(\operatorname{Hom}(\ker A,\operatorname{coker}A)\), with both factors of dimension two.\par
\subsection*{P000697 (both; faithful error): range theorem}
\noindent\textit{Transcription reads.} closedness is necessary and sufficient without a function class\par
\noindent\textit{Critical reading.} Add the global smooth, rapidly decreasing, density-covariant data class.\par
\noindent\textit{Evidence.} GGS (1967), Theorem 1, states class \(\Phi\) as an independent hypothesis.\par
\subsection*{P000697 (both; ambiguity): inversion cycle dimension}
\noindent\textit{Transcription reads.} one-dimensional cycle\par
\noindent\textit{Critical reading.} A real two-cycle, for example a complex one-dimensional projective line.\par
\noindent\textit{Evidence.} The PNG clearly says one-dimensional; the integrand is a \((1,1)\)-form. GGS (1967), Theorem 2, explicitly says real dimension \(2k\).\par
\subsection*{P000697 (both; ambiguity): admissibility question}
\noindent\textit{Transcription reads.} recover phi from Hphi on K\par
\noindent\textit{Critical reading.} Distinguish this from recovering the inversion form from scalar data phi restricted to K.\par
\noindent\textit{Evidence.} Gelfand--Graev (1968), section 2, Definition 2, specifies descent of the differential inversion operator.\par
\subsection*{P000698 (both; faithful error): all linear complexes isomorphic}
\noindent\textit{Transcription reads.} any other is isomorphic to it\par
\noindent\textit{Critical reading.} Any other nondegenerate linear complex is projectively isomorphic to it.\par
\noindent\textit{Evidence.} The rank of the associated alternating form distinguishes ranks four and two; rank two gives a special complex.\par
\subsection*{P000699 (both; provenance): record415 date}
\noindent\textit{Transcription reads.} Catalogue 4 April 1977\par
\noindent\textit{Critical reading.} 28 March 1977; retain catalogue date separately\par
\noindent\textit{Evidence.} The PNG has a single clear 28/III-77 heading, Gelfand intervention, then Vershik heading on the same page with no new date.\par
\subsection*{P000700 (both; faithful error): hook formula attribution}
\noindent\textit{Transcription reads.} 1955\par
\noindent\textit{Critical reading.} 1954\par
\noindent\textit{Evidence.} The Frame, Robinson and Thrall publisher record dates Canadian Journal of Mathematics, volume 6, to 1954.\par
\subsection*{P000700 (both; ambiguity): hook length}
\noindent\textit{Transcription reads.} cells to the right and below\par
\noindent\textit{Critical reading.} include the distinguished cell itself\par
\noindent\textit{Evidence.} The displayed hook formula requires \(h=\text{arm}+\text{leg}+1\); otherwise a removable cell would have length zero.\par
\subsection*{P000700 (both; faithful error): branching sentence}
\noindent\textit{Transcription reads.} Every irreducible of \(S_{n-1}\) occurs once.\par
\noindent\textit{Critical reading.} only the immediate predecessor diagrams occur once\par
\noindent\textit{Evidence.} The next notebook line specifies predecessor diagrams; Young branching theorem.\par
\subsection*{P000700 (both; transcription): regular representation}
\noindent\textit{Transcription reads.} Functions on \(S^n\)\par
\noindent\textit{Critical reading.} Functions on \(S_n\)\par
\noindent\textit{Evidence.} Context is the finite symmetric group and its regular representation of dimension \(n!\).\par
\subsection*{P000701 (both; faithful error): rotated coordinates}
\noindent\textit{Transcription reads.} \(X=(x+y)/2;\quad Y=(x-y)/2\)\par
\noindent\textit{Critical reading.} \(X=(x-y)/2;\quad Y=(x+y)/2\), for the stated graph normalization.\par
\noindent\textit{Evidence.} PNG confirms the transcribed sum/difference; Vershik--Kerov1977p1024 has the corrected order.\par
\subsection*{P000701 (both; faithful error): limit-shape formula}
\noindent\textit{Transcription reads.} interior arcsine formula alone\par
\noindent\textit{Critical reading.} Add \(\Omega(X)=|X|\) for \(|X|>1\).\par
\noindent\textit{Evidence.} The real square root is defined only on \([-1,1]\); the 1977 paper gives both branches.\par
\subsection*{P000701 (both; faithful error): fluctuation claim}
\noindent\textit{Transcription reads.} white noise; fourth-root scaled uncentered functional\par
\noindent\textit{Critical reading.} not an identified valid CLT; use a separately labelled later Kerov Gaussian-field theorem\par
\noindent\textit{Evidence.} Both radical indices 4 are visible on the PNG. Ivanov--Olshanski, Theorem 7.1, uses centered profile fluctuations and a generalized Gaussian field with nonwhite covariance.\par
\subsection*{P000701 (both; faithful error): subsequence definition}
\noindent\textit{Transcription reads.} \(i_1<\cdots<i_n\) inside a maximum over \(r\)\par
\noindent\textit{Critical reading.} \(i_1<\cdots<i_r\)\par
\noindent\textit{Evidence.} The scan confirms \(n\); \(r\) is the length being maximized.\par
\subsection*{P000701 (both; faithful error): 1961 author}
\noindent\textit{Transcription reads.} Shenstead\par
\noindent\textit{Critical reading.} Schensted\par
\noindent\textit{Evidence.} The cited 1961 Canadian Journal paper is C. Schensted, Longest Increasing and Decreasing Subsequences.\par
\subsection*{P000702 (both; provenance): speaker initials}
\noindent\textit{Transcription reads.} Catalogue A.S.Tyurin\par
\noindent\textit{Critical reading.} Scan A.Tyurin; publication identification A.N.Tyurin\par
\noindent\textit{Evidence.} The scan abbreviates the first name only. Both exactly matching primary papers are signed A.N.Tyurin; the supplied catalogue initials conflict.\par
\subsection*{P000702 (both; faithful error): packet dimension in GL(n)}
\noindent\textit{Transcription reads.} \((n^2-1)(g-1)\)\par
\noindent\textit{Critical reading.} \(n^2(g-1)+1\); the displayed value is for a fixed determinant connection.\par
\noindent\textit{Evidence.} Connections on a stable degree-zero bundle form a torsor for \(H^0(X,\operatorname{End}(E)\otimes K_X)\); Riemann--Roch. Rank one gives \(g\), not zero.\par
\subsection*{P000702 (both; faithful error): any two irreducible representations}
\noindent\textit{Transcription reads.} any two are in one packet for some Riemann surface\par
\noindent\textit{Critical reading.} false: distinct unitary characters cannot be in one packet for any complex structure\par
\noindent\textit{Evidence.} Unitary flat Hom line has no nowhere-zero holomorphic section unless its character is trivial; Tyurin (1978), Chapter II, section 2, Exercise 10. Full proof independently checked.\par
\subsection*{P000702 (both; faithful error): affine bundle diagram}
\noindent\textit{Transcription reads.} whole character variety to moduli of bundles is an affine bundle\par
\noindent\textit{Critical reading.} restrict to stable degree-zero underlying bundles, and specify determinant conventions\par
\noindent\textit{Evidence.} Packet automorphisms and unstable underlying bundles prevent the unrestricted claim.\par
\subsection*{P000703 (both; ambiguity): section notation}
\noindent\textit{Transcription reads.} \(f_i\) on \(U_j\)\par
\noindent\textit{Critical reading.} \(f_i\) on \(U_i\)\par
\noindent\textit{Evidence.} The cocycle equation and local restriction \(df_i/f_i\) on \(U_i\) fix the index; the handwritten subscript is ambiguous.\par
\subsection*{P000703 (both; faithful error): construct a differential with specified divisor}
\noindent\textit{Transcription reads.} arbitrary divisor\par
\noindent\textit{Critical reading.} degree-zero divisor for a compact surface\par
\noindent\textit{Evidence.} Sum of residues of a meromorphic differential on a compact Riemann surface is zero.\par
\subsection*{P000703 (both; faithful error): Torelli name}
\noindent\textit{Transcription reads.} Terelli\par
\noindent\textit{Critical reading.} Torelli\par
\noindent\textit{Evidence.} The notebook claim and matching Tyurin (1978), Chapter II, section 2, Exercise 13, explicitly invoke the Torelli theorem.\par
\subsection*{P000703 (both; faithful error): Schwarzian gluing}
\noindent\textit{Transcription reads.} general quadratic differentials\par
\noindent\textit{Critical reading.} fix a projective atlas; otherwise a projective connection\par
\noindent\textit{Evidence.} The written chain rule includes the extra Schwarzian term; Tyurin (1977), Introduction, assumes Möbius coordinate changes.\par
\subsection*{P000704 (both; transcription): bottom monodromy diagram}
\noindent\textit{Transcription reads.} diagram only; equation absent from plain text\par
\noindent\textit{Critical reading.} \(f(g(u))=\rho(g)(f(u))\)\par
\noindent\textit{Evidence.} The actual PNG contains this equation and a phrase saying that rho is thereby recovered. It belongs to the Tyurin continuation.\par
\subsection*{P000705-P000706 (both; provenance): year and talk boundary}
\noindent\textit{Transcription reads.} included pages after 1976--77 Tyurin record\par
\noindent\textit{Critical reading.} next-academic-year heading with a separately named adjacent Drinfeld talk\par
\noindent\textit{Evidence.} Catalogue record 418 is the 1977--78 heading; P000705 explicitly marks the thirty-fourth seminar year and names V. Drinfeld in a new heading.\par
\subsection*{P000705 (both; ambiguity): visible meeting date}
\noindent\textit{Transcription reads.} 5/IV-77\par
\noindent\textit{Critical reading.} preserve the visible reading; exact month unresolved\par
\noindent\textit{Evidence.} The transcribed IV is compatible with visible strokes but conflicts with the new-year boundary; replacing it by IX would be an inference, not a secure scan reading.\par
\subsection*{P000753 (both; provenance): Date and academic-year assignment}
\noindent\textit{Transcription reads.} Catalogue record~437 and supplied talk list: 13 February 1977.\par
\noindent\textit{Critical reading.} Manuscript: \ru{13/II--78 г.}, 13 February 1978; academic year 1977--78.\par
\noindent\textit{Evidence.} The date is plainly visible in the scan and is correctly transcribed in both TeX files. P000757 explicitly begins 27 February 1978. This corrects catalogue provenance, not either diplomatic date.\par
\subsection*{P000752 (both; provenance): Left boundary of the record-437 envelope}
\noindent\textit{Transcription reads.} The supplied record range includes this physics page.\par
\noindent\textit{Critical reading.} Do not attribute the sine-Gordon S-matrix and chiral-field tail to S. I. Gelfand's vector-bundle lecture.\par
\noindent\textit{Evidence.} The scan contains equations (46)--(51), soliton scattering and a chiral field, with no S.~I.~Gelfand heading. P000753 begins a new dated lecture and resets numbering to (1).\par
\subsection*{P000757-P000758 (both; provenance): Right boundary of the record-437 envelope}
\noindent\textit{Transcription reads.} The supplied record range includes the next lecture.\par
\noindent\textit{Critical reading.} Treat 27 February 1978 as a separate S. I. Gelfand lecture; its visible repetition and line-bundle exercise are adjacent evidence only.\par
\noindent\textit{Evidence.} The scan of P000757 has \ru{27/II--78 г.}; P000758 continues its constructions. The archive list places a 27-Feb-78 record immediately after the erroneous 13-Feb-77 entry.\par
\subsection*{P000753 (both; ambiguity): Projective quotient}
\noindent\textit{Transcription reads.} \(\Xi\) is an affine space; \(\Xi/{\sim}=\mathbb P^n\).\par
\noindent\textit{Critical reading.} Use a vector space \(\Xi\) of dimension \(n+1\) and \((\Xi\setminus\{0\})/\mathbb C^*\).\par
\noindent\textit{Evidence.} The scan explicitly adds \(\xi\ne0\), so the puncture is recoverable. Scalar multiplication and the following exterior algebra require the linear structure; the clarification does not overwrite the notebook.\par
\subsection*{P000754 (both; ambiguity): Free-module example}
\noindent\textit{Transcription reads.} Free modules give trivial bundles.\par
\noindent\textit{Critical reading.} Free modules give the zero bundle, of rank zero.\par
\noindent\textit{Evidence.} The scan first says \ru{нулевое расслоение} and then abbreviates \ru{трив.}. Its preceding formula has all homology groups zero. The contraction proof in the critical checks verifies the interpretation.\par
\subsection*{P000754-P000755 (both; ambiguity): Hopf convention and sections}
\noindent\textit{Transcription reads.} \(O(1)\) is the one with a section; \(O(1)\) is the Hopf bundle (or \(O(-1)\)).\par
\noindent\textit{Critical reading.} For the displayed homogeneous-degree convention, \(O(-1)\) is tautological and \(O(1)\) its dual; sections are regular on the specified open set.\par
\noindent\textit{Evidence.} The scan preserves this uncertainty. P000758 explicitly says anti-Hopf. Many positive twists have nonzero sections, so that property cannot identify degree one; the chart calculation resolves the sign.\par
\subsection*{P000755 (both; faithful error): Serre theorem terminology}
\noindent\textit{Transcription reads.} \ru{обильное расслоение}, translated as ample bundle, for having sections through every fibre vector.\par
\noindent\textit{Critical reading.} The property asserted here is globally generated.\par
\noindent\textit{Evidence.} The explanatory clause in the scan is the surjectivity of the evaluation map. This is the actual hypothesis used by the construction; it is distinct from the definition of ampleness.\par
\subsection*{P000756 (both; ambiguity): Inverse construction, equations (18)--(20)}
\noindent\textit{Transcription reads.} Two free modules and their kernel, without an explicit total grading.\par
\noindent\textit{Critical reading.} Specify the grading shifts, the dual exterior module or its shifted identification, and a sufficiently regular section-module tail.\par
\noindent\textit{Evidence.} The scan leaves these data unstated. With the raw grading the case \(n=0\), \(X_0=1\), gives kernel \(\mathbb C\xi_0\) in degree \(-1\), which is not proper. This demonstrates a missing convention, not failure of the published BGG theorem.\par
\subsection*{P000756 (both; faithful error): Serre citation date}
\noindent\textit{Transcription reads.} Article in the collection \ru{Расслоенные пр-ва}; 1953.\par
\noindent\textit{Critical reading.} The original FAC article is 1955; the Russian collection cited in the BGG paper is 1957.\par
\noindent\textit{Evidence.} The scan indeed writes 1953, so the transcription is faithful. \citet[title page]{serre1955coherent} gives 1955; \citet[reference~9]{bernsteingelfandgelfand1978bundles} gives the Russian collection, Moscow 1957, pp.~372--453.\par
\subsection*{P000756 (both; faithful error): Scope of Serre correspondence}
\noindent\textit{Transcription reads.} Sheaves on complex projective space.\par
\noindent\textit{Critical reading.} Algebraic coherent sheaves, with a quotient of the finitely generated graded-module category by finite-dimensional modules.\par
\noindent\textit{Evidence.} The scan omits coherent. \citet[Chapter~III]{serre1955coherent} and \citet[Proposition~30.14.4]{stacks2026projective} give the exact scope; the omission cannot be extended to arbitrary sheaves.\par
